% !TEX program = xelatex
% !TEX encoding = UTF-8
%% ============================================================
%%  Analiță și Algebră pentru Olimpiadă — VOLUM COMPLET
%%  Teorie + Probleme de Analiză + Probleme de Algebră
%%  Compilare: xelatex (x3) pentru cuprins complet
%% ============================================================
\documentclass[11pt,a4paper]{report}

%% ---- encoding (compatibil pdflatex si xelatex) ----
\usepackage{iftex}
\ifXeTeX
  \usepackage{fontspec}
  \setmainfont{Latin Modern Roman}
\else
  \usepackage[utf8]{inputenc}
  \usepackage[T1]{fontenc}
\fi

%% ---- limbă ----
\usepackage[romanian]{babel}

%% ---- matematică ----
\usepackage{amsmath,amssymb,amsthm}
\usepackage{mathtools}
\usepackage{mathrsfs}

%% ---- layout ----
\usepackage{multicol}
\usepackage{graphicx}
\usepackage{caption}
%% KDP: lets TeX loosen a line instead of pushing inline math into the margin
\setlength{\emergencystretch}{3em}
%% KDP: a heading never stays alone at the bottom of a page. Otherwise a breakable
%% tcolorbox that cannot start under it is pushed to the next page and drawn over
%% the running head (tcolorbox: "Using nobreak failed"). Same trick as needspace.sty.
\makeatletter
\newcommand{\needlines}[1]{\par\begingroup\dimen@=#1\baselineskip
  \vskip\z@\@plus\dimen@ \penalty-100 \vskip\z@\@plus-\dimen@
  \vskip\dimen@ \penalty9999 \vskip-\dimen@ \vskip\z@skip\endgroup}
\makeatother
\AtBeginDocument{%
  \pretocmd{\section}{\needlines{9}}{}{\typeout{needlines: section not patched}}%
  \pretocmd{\subsection}{\needlines{8}}{}{\typeout{needlines: subsection not patched}}%
  \pretocmd{\subsubsection}{\needlines{6}}{}{\typeout{needlines: subsubsection not patched}}}
\usepackage{geometry}
%% Geometry identical to the printed books (VRAI LIVRES / interieurs_corriges).
%% Values measured from the shipped PDFs:
%%   - text block: 356.96pt = 4.9578in wide  => left = right = 2.2cm
%%   - header: top margin 34.64pt + headheight 22pt => head rule at 56.64pt
%%   - footer: page-number baseline 32.77pt above the bottom edge
%% Do NOT change these without regenerating the covers: the spine width depends
%% on the page count, and the page count depends on the text block.
\geometry{paperwidth=6.69in,paperheight=9.45in,
          left=2.2cm,right=2.2cm,top=34.64pt,bottom=32.77pt,
          includehead,includefoot,headsep=8pt,footskip=16pt}
\usepackage{enumitem}
\usepackage{xcolor}
\usepackage{tikz}
\usepackage{tcolorbox}
\tcbuselibrary{skins,breakable}
\usepackage{titlesec}
\usepackage{fancyhdr}
\usepackage[colorlinks=true,linkcolor=blue!60!black,urlcolor=blue!60!black]{hyperref}

%% ---- paleta cromatică unificată ----
\definecolor{ink}{HTML}{1C2230}
\definecolor{grey}{HTML}{6B7280}
\definecolor{olm}{HTML}{2A9D8F}
\definecolor{ojm}{HTML}{E76F51}
\definecolor{onm}{HTML}{6A4C93}
\definecolor{solcol}{HTML}{0B7A52}
\definecolor{teoblue}{RGB}{25,55,105}
\definecolor{obsgray}{RGB}{90,90,90}
\definecolor{acc}{RGB}{80,40,10}

%% ---- box-uri tehnice ----
\newtcolorbox{teobox}[1]{%
  enhanced,breakable,sharp corners,
  colback=onm!6,colframe=onm,
  colbacktitle=onm,coltitle=white,
  borderline west={4pt}{0pt}{onm},
  boxrule=0.4pt,left=10pt,top=4pt,bottom=4pt,
  title={#1},fonttitle=\bfseries}
\newtcolorbox{tehbox}[1]{%
  enhanced,breakable,sharp corners,
  colback=ojm!6,colframe=ojm,
  colbacktitle=ojm,coltitle=white,
  borderline west={4pt}{0pt}{ojm},
  boxrule=0.4pt,left=10pt,top=4pt,bottom=4pt,
  title={#1},fonttitle=\bfseries}

%% ---- teoreme ----
\theoremstyle{plain}
\newtheorem{teorema}{Teoremă}[chapter]
\newtheorem{propozitie}[teorema]{Propoziție}
\newtheorem{lema}[teorema]{Lemă}
\newtheorem{corolar}[teorema]{Corolar}
\theoremstyle{definition}
\newtheorem{definitie}[teorema]{Definiție}
\newtheorem{exemplu}[teorema]{Exemplu}
\newtheorem{aplicatie}[teorema]{Aplicație}
\theoremstyle{remark}
\newtheorem{observatie}[teorema]{Observație}
% obs redefinit mai jos ca newenvironment

%% ---- bare colorate pe teoreme ----
\tcolorboxenvironment{teorema}{%
  enhanced,breakable,sharp corners,boxrule=0pt,frame hidden,
  borderline west={4pt}{0pt}{onm},
  colback=onm!5,left=10pt,right=4pt,top=3pt,bottom=3pt}
\tcolorboxenvironment{propozitie}{%
  enhanced,breakable,sharp corners,boxrule=0pt,frame hidden,
  borderline west={3.5pt}{0pt}{onm!75},
  colback=onm!3,left=10pt,right=4pt,top=3pt,bottom=3pt}
\tcolorboxenvironment{lema}{%
  enhanced,breakable,sharp corners,boxrule=0pt,frame hidden,
  borderline west={3pt}{0pt}{onm!55},
  colback=onm!2,left=10pt,right=4pt,top=3pt,bottom=3pt}
\tcolorboxenvironment{corolar}{%
  enhanced,breakable,sharp corners,boxrule=0pt,frame hidden,
  borderline west={2.5pt}{0pt}{onm!40},
  colback=white,left=10pt,right=4pt,top=2pt,bottom=2pt}
\tcolorboxenvironment{definitie}{%
  enhanced,breakable,sharp corners,boxrule=0pt,frame hidden,
  borderline west={4pt}{0pt}{olm},
  colback=olm!5,left=10pt,right=4pt,top=3pt,bottom=3pt}
\tcolorboxenvironment{exemplu}{%
  enhanced,breakable,sharp corners,boxrule=0pt,frame hidden,
  borderline west={3.5pt}{0pt}{ojm},
  colback=ojm!4,left=10pt,right=4pt,top=3pt,bottom=3pt}
\tcolorboxenvironment{aplicatie}{%
  enhanced,breakable,sharp corners,boxrule=0pt,frame hidden,
  borderline west={3pt}{0pt}{ojm!80},
  colback=ojm!3,left=10pt,right=4pt,top=3pt,bottom=3pt}
\tcolorboxenvironment{observatie}{%
  enhanced,breakable,sharp corners,boxrule=0pt,frame hidden,
  borderline west={2.5pt}{0pt}{grey},
  colback=white,left=10pt,right=4pt,top=2pt,bottom=2pt}

%% ---- macro-uri algebră ----
\newcommand{\NN}{\mathbb{N}}
\newcommand{\ZZ}{\mathbb{Z}}
\newcommand{\QQ}{\mathbb{Q}}
\newcommand{\RR}{\mathbb{R}}
\newcommand{\CC}{\mathbb{C}}
\DeclareMathOperator{\ord}{ord}
\DeclareMathOperator{\Ker}{Ker}
\DeclareMathOperator{\Ima}{Im}
\DeclareMathOperator{\car}{car}
\DeclareMathOperator{\Hom}{Hom}
\DeclareMathOperator{\End}{End}
\DeclareMathOperator{\Aut}{Aut}
\DeclareMathOperator{\Inn}{Inn}
\DeclareMathOperator{\Orb}{Orb}
\DeclareMathOperator{\Stab}{Stab}
\newcommand{\cl}[1]{\widehat{#1}}
\DeclareMathOperator{\Tr}{Tr}
\DeclareMathOperator{\rang}{rang}
\DeclareMathOperator{\Span}{Span}
\DeclareMathOperator{\adj}{adj}
\DeclareMathOperator{\sgn}{sgn}
\DeclareMathOperator{\diag}{diag}
\DeclareMathOperator{\Spec}{Spec}
\newcommand{\GL}{\mathrm{GL}}

%% ---- macro-uri culegeri (titluri colorate) ----
\newcommand{\capitolcurent}{Introducere}
\newcommand{\sectiunealg}[2]{%
  \clearpage\vspace*{6mm}%
  \addcontentsline{toc}{section}{\textcolor{blue!60!black}{#1}}%
  {\noindent\color{ink}\rule{\linewidth}{3pt}}\par\vspace{2mm}%
  {\noindent\fontsize{34}{38}\selectfont\bfseries\color{ink}#1}\par\vspace{2mm}%
  {\noindent\large\itshape\textcolor{grey}{#2}}\par\vspace{2mm}%
  {\noindent\color{ink}\rule{\linewidth}{3pt}}\par\vspace{7mm}}
\newcommand{\parteamare}[2]{%
  \clearpage\vspace*{4mm}%
  \addcontentsline{toc}{subsection}{#1}%
  {\noindent\color{ink}\rule{\linewidth}{2.4pt}}\par\vspace{2mm}%
  {\noindent\fontsize{30}{34}\selectfont\bfseries\color{ink}#1}\par\vspace{1.5mm}%
  {\noindent\large\itshape\textcolor{grey}{#2}}\par\vspace{2mm}%
  {\noindent\color{ink}\rule{\linewidth}{2.4pt}}\par\vspace{6mm}}
\newcommand{\parteamareenunt}[1]{%
  \vspace*{4mm}%
  \addcontentsline{toc}{subsection}{#1}%
  {\noindent\color{ink}\rule{\linewidth}{2.4pt}}\par\vspace{2mm}%
  {\noindent\fontsize{30}{34}\selectfont\bfseries\color{ink}#1}\par\vspace{1.5mm}%
  {\noindent\color{ink}\rule{\linewidth}{2.4pt}}\par\vspace{6mm}}
\newcommand{\subsectiuneclasa}[3]{%
  \clearpage\renewcommand{\capitolcurent}{#3}\vspace*{3mm}%
  \addcontentsline{toc}{subsection}{\quad #1}%
  {\noindent\color{ink}\rule{\linewidth}{1.6pt}}\par\vspace{2mm}%
  {\noindent\Huge\bfseries\color{ink}#1}\par\vspace{1.5mm}%
  {\noindent\large\itshape\textcolor{grey}{#2}}\par\vspace{2mm}%
  {\noindent\color{ink}\rule{\linewidth}{1.6pt}}\par\vspace{5mm}}
\newcommand{\sectiunemare}[2]{%
  \vspace{4mm}%
  {\noindent\color{ink}\rule{\linewidth}{1.2pt}}\par\vspace{1.5mm}%
  {\noindent\huge\bfseries\color{ink}#1}\par\vspace{1mm}%
  {\noindent\small\itshape\textcolor{grey}{#2}}\par\vspace{1.5mm}%
  {\noindent\color{ink}\rule{\linewidth}{1.2pt}}\par\vspace{3mm}}
\newcommand{\nivel}[3]{%
  \vspace{5mm}%
  \colorlet{lvl}{#1}%
  {\noindent\color{#1}\rule{\linewidth}{2.2pt}}\par\vspace{1.5mm}%
  {\noindent\LARGE\bfseries\color{#1}#2}\par\vspace{1mm}%
  {\noindent\small\itshape\textcolor{grey}{#3}}\par\vspace{3mm}}
\newcommand{\nivelE}[2]{%
  \colorlet{lvl}{#1}%
  \vspace{4mm}{\noindent\color{#1}\rule{\linewidth}{1.5pt}}\par\vspace{1mm}%
  {\noindent\Large\bfseries\color{#1}#2}\par\vspace{2mm}}
\newcommand{\prob}[3]{%
  \vspace{3mm}%
  {\noindent\color{#1}\rule{3pt}{\baselineskip}\hspace{4pt}%
   \textbf{Problema #2}%
   \ifx&#3&\else\ {\small\itshape\textcolor{grey}{--- #3}}\fi}\par\vspace{1mm}}
\newcommand{\sol}{\par\vspace{2mm}{\noindent\small\textcolor{solcol}{\textbf{Soluție.}}\ }}
\newenvironment{parti}{\begin{enumerate}[label=\textbf{(\alph*)},leftmargin=*,itemsep=2pt]}{\end{enumerate}}

%% ---- titluri capitole armonizate ----
\titleformat{\chapter}[block]
  {\normalfont\huge\bfseries\color{ink}}{}
  {0pt}{\vspace*{2mm}\noindent\textcolor{ink}{\rule{\linewidth}{2.4pt}}\vspace{2mm}\newline}
  [\vspace{2mm}\noindent\textcolor{ink}{\rule{\linewidth}{2.4pt}}\vspace{5mm}]
\titleformat{\section}[block]
  {\normalfont\Large\bfseries\color{ink}}{}
  {0pt}{\vspace*{1mm}\noindent\textcolor{ink}{\rule{\linewidth}{1.2pt}}\vspace{1.5mm}\newline}
  [\vspace{1.5mm}\noindent\textcolor{ink}{\rule{\linewidth}{1.2pt}}\vspace{3mm}]
\titleformat{\subsection}[hang]
  {\normalfont\large\bfseries\color{ink}}{}
  {0pt}{}[\vspace{1mm}\noindent\textcolor{grey}{\rule{\linewidth}{0.5pt}}\vspace{2mm}]

%% ---- antet/subsol ----
\setlength{\headheight}{22pt}
\pagestyle{fancy}\fancyhf{}
\lhead{\small\textcolor{grey}{\titluserie}}
\rhead{\small\textcolor{grey}{\capitolcurent}}
\cfoot{\small\textcolor{grey}{\thepage}}
\renewcommand{\headrulewidth}{0.2pt}
\renewcommand{\headrule}{\hbox to\headwidth{\color{grey}\leaders\hrule height \headrulewidth\hfill}}


%% ---- macro-uri suplimentare algebra (Stil B: clasa a XII-a) ----
\colorlet{lvl}{olm}
\newcommand{\Mn}{M_n(\mathbb{C})}
\newcommand{\On}{O_n}
\newcommand{\demo}{\par\vspace{1.6mm}\noindent\textit{\textcolor{solcol}{Demonstrație.}}\ \ignorespaces}
\newcounter{cat}
\newtheoremstyle{probstyle}
  {8pt}{3pt}{\normalfont}{0pt}{\bfseries\large}{}{0.55em}
  {{\color{lvl}\rule[-2.5pt]{3.2pt}{15pt}\hspace{0.55em}}\color{ink}\thmname{#1}\thmnumber{ #2}\thmnote{ \normalfont\small\textcolor{grey}{--- #3}}}
\theoremstyle{probstyle}
\newtheorem{problemaX}{Problema}
\renewcommand{\theproblemaX}{\ifcase\value{cat}\or OLM\or OJM\or ONM\fi.\arabic{problemaX}}
\newenvironment{problema}[1][]
  {\par\vspace{2.5mm}{\color{grey!40}\hrule height 0.4pt}\vspace{0.5mm}\ifx\relax#1\relax\begin{problemaX}\else\begin{problemaX}[#1]\fi}
  {\end{problemaX}}
\newcommand{\theproblema}{\theproblemaX}
\newtheorem*{obsX}{Observație}
\newenvironment{obs}{\begin{obsX}\small\color{ink}}{\end{obsX}}
\newtheorem*{lemat}{Lemă}
\renewcommand{\proofname}{\textcolor{solcol}{Demonstrație}}
\newenvironment{solutie}
  {\par\vspace{1.6mm}\noindent\textit{\textcolor{solcol}{\textbf{Soluție.}}}\enspace\ignorespaces}
  {\hfill\textcolor{solcol}{$\square$}\par\medskip}
\setcounter{tocdepth}{2}
\AtBeginDocument{\renewcommand{\proofname}{Demonstrație}}
\renewcommand{\qedsymbol}{$\blacksquare$}
\setlength{\parskip}{2pt}
\color{ink}

%% ---- parametri de volum (se redefinesc in fisierul fiecarui volum) ----
\newcommand{\titluserie}{Analiză pentru Olimpiadă}
\newcommand{\numarvolum}{1}
\newcommand{\titluvolum}{Analiză, clasa a XI-a}
\newcommand{\subtitluvolum}{}
\newcommand{\isbnvolum}{}
\newcommand{\despreacestvolum}{}
\newcommand{\celelaltevolume}{}

%% ---- compatibilitate tex4ht (build-ul EPUB) ----
%% tex4ht nu poate traduce \halign-ul din mediile *smallmatrix; la conversia in
%% MathML le inlocuim cu variantele normale, care produc <mtable> corect.
\ifdefined\HCode
  \renewenvironment{psmallmatrix}{\begin{pmatrix}}{\end{pmatrix}}
  \renewenvironment{bsmallmatrix}{\begin{bmatrix}}{\end{bmatrix}}
  \renewenvironment{vsmallmatrix}{\begin{vmatrix}}{\end{vmatrix}}
  \renewenvironment{smallmatrix}{\begin{matrix}}{\end{matrix}}
\fi

%% ---- trimiteri catre alt volum (capitolul Constructii e comun volumelor 3 si 4) ----
%% \refv{eticheta}{numar}{volum}: daca eticheta e in volumul curent, \ref obisnuit; altfel numarul ei
%% din volumul in care se afla, urmat de "din volumul N".
\makeatletter
\newcommand{\invol}[1]{\ifnum\numarvolum=#1\relax\else\ din volumul~#1\fi}
\newcommand{\refv}[3]{\@ifundefined{r@#1}{#2\invol{#3}}{\ref{#1}}}
\makeatother


\renewcommand{\titluserie}{Algebră pentru Olimpiadă}
\renewcommand{\numarvolum}{4}
\renewcommand{\titluvolum}{clasa a XII-a}
\renewcommand{\subtitluvolum}{Grupuri, inele, corpuri și polinoame peste un corp}
\renewcommand{\despreacestvolum}{Cuprinde structurile algebrice de clasa a XII-a: grupuri, inele, corpuri și polinoame peste un corp.}
\renewcommand{\celelaltevolume}{analizei de clasa a XI-a, analizei de clasa a XII-a și algebrei liniare}

\begin{document}
%% ============================================================
%%  PAGINA DE TITLU, COPYRIGHT, CUPRINS, CUVÂNT ÎNAINTE
%%  (comun tuturor volumelor; textul variabil vine din macro-uri)
%% ============================================================
\thispagestyle{empty}
\vspace*{2.2cm}
\begin{center}
  {\color{ink}\rule{\linewidth}{3pt}}\vspace{6mm}\\
  {\fontsize{30}{36}\selectfont\bfseries\color{ink}\titluserie}\\[4mm]
  {\color{onm}\rule{0.4\linewidth}{1.5pt}}\\[6mm]
  {\LARGE \titluvolum}\\[3mm]
  {\large\itshape\textcolor{grey}{\subtitluvolum}}\\[8mm]
  {\normalsize\textcolor{grey}{Teoria completă și probleme}}\\[10mm]
  {\color{ink}\rule{\linewidth}{3pt}}\\[8mm]
  {\large\itshape Vlad-Ștefan Constantinescu}
\end{center}
\vfill
\clearpage

%% ---------------- pagina de copyright ----------------
\thispagestyle{empty}
\vspace*{1cm}
\noindent{\small
\titluserie, \titluvolum\\
Volumul \numarvolum{} dintr-o serie de patru volume\\[2mm]
Autor: Vlad-Ștefan Constantinescu\\
Ediția a doua, București, 2026\\[3mm]
\copyright{} 2026 Vlad-Ștefan Constantinescu. Toate drepturile rezervate.\\[2mm]
Nicio parte a acestei lucrări nu poate fi reprodusă sau transmisă sub nicio formă și prin niciun mijloc,
electronic sau mecanic, fără acordul prealabil scris al autorului, cu excepția unor scurte citate
în recenzii și în lucrări de referință.\\[3mm]
Tehnoredactat de autor în \LaTeX{} (Xe\LaTeX).\\[3mm]
\ifx\isbnvolum\empty\else ISBN: \isbnvolum\fi

}
\vfill
\clearpage

\tableofcontents
\clearpage

%% ============================================================
%%  CUVÂNT ÎNAINTE
%% ============================================================
\thispagestyle{empty}
\addcontentsline{toc}{chapter}{Cuvânt înainte}
\vspace*{1.5cm}
{\noindent\color{onm}\rule{\linewidth}{3pt}}\par\vspace{3mm}
{\noindent\fontsize{28}{32}\selectfont\bfseries\color{ink} Cuvânt înainte}\par\vspace{3mm}
{\noindent\color{onm}\rule{\linewidth}{3pt}}\par\vspace{8mm}

\noindent \textbf{Despre acest volum.} Acesta este volumul \numarvolum{} dintr-o serie de patru volume.
\despreacestvolum{} Volumele care îl însoțesc sunt dedicate \celelaltevolume.

\medskip
\noindent Volumul are două mari componente: \textbf{teoria}, cu demonstrații complete și exemple, și \textbf{problemele}, grupate pe trei niveluri de dificultate (OLM $\cdot$ OJM $\cdot$ ONM). În partea de teorie am inclus două secțiuni pe care le consider esențiale și care lipsesc adesea din cărțile de olimpiadă: una de \textbf{redactare} (cum se scrie o demonstrație riguroasă și ce greșeli de formă penalizează juriul) și una de \textbf{construcții} (ce obiecte să încerci când enunțul nu îți dă o direcție).

\medskip
\noindent \textbf{Cele trei niveluri.} Pe tot parcursul cărții, problemele sunt clasificate după cele trei
etape ale Olimpiadei Naționale de Matematică: \textbf{OLM} (olimpiada locală, adică etapa pe școală și pe oraș),
\textbf{OJM} (olimpiada județeană) și \textbf{ONM} (olimpiada națională, etapa finală). În linii mari,
problemele OLM sunt probleme de concurs accesibile, OJM corespunde unei olimpiade regionale solide, iar
problemele ONM sunt de nivelul unei finale naționale: comparabile, pentru materia acestei cărți, cu cele mai
grele probleme din concursurile studențești precum Putnam.

\medskip
\noindent \textbf{Analiză vs.\ algebră.} Analiza matematică este un subiect continuu: materia clasei a XII-a nu este altceva decât o extensie firească a celei din clasa a XI-a, iar limitele, integralele și seriile trăiesc în aceeași lume. De aceea am construit o punte explicită între cele două clase. Algebra stă altfel: structurile algebrice din clasa a XII-a (grupuri, inele, corpuri, polinoame) formează un bloc coerent în sine, cu fire clare între ele, dar nu au nicio legătură cu algebra liniară din clasa a XI-a. Cele două jumătăți de algebră sunt lumi separate, și am ales să le tratez ca atare.

\medskip
\noindent \textbf{Despre metode.} Am îmbinat, deliberat, două perspective. Pe de o parte, tehnicile care se văd la olimpiadele din România, de la faza locală până la cea națională. Pe de altă parte, unele metode pe care le-am întâlnit în primii doi ani de facultate din sistemul \emph{CPGE} (\textit{classes préparatoires aux grandes écoles}) din Franța (funcții generatoare, comparări serie-integrală, asimptotică prin singularități) și pe care le-am găsit perfect utilizabile și la nivel de olimpiadă, adesea mai clare și mai puternice decât abordările ad-hoc. Puntea dintre clasele a XI-a și a XII-a este locul unde aceste două perspective se întâlnesc cel mai firesc.

\medskip
\noindent \textbf{Despre programă.} Nivelurile OLM $\cdot$ OJM $\cdot$ ONM indică \emph{dificultatea}, nu strict
programa etapei corespunzătoare. Am inclus deliberat, la un anumit nivel, probleme care folosesc noțiuni ce nu
figurează oficial în programa acelei etape: de exemplu probleme cu derivate la OLM și la OJM, deși, nefiind în
programă, ele nu pot apărea efectiv acolo. Le-am păstrat ca exercițiu de tipul „și dacă ar pica totuși la OLM
sau la OJM'': ideea de bază este accesibilă și la acel nivel, iar antrenamentul pe ele nu poate strica.

\medskip
\noindent \textbf{Despre redactare.} Am încercat să redactez așa cum se cere la olimpiadă: enunț și demonstrez
fiecare lemă pe care o folosesc, împart soluțiile în pași numiți explicit, verific ipotezele înainte de a aplica
o teoremă și marchez clar unde se încheie fiecare argument. Redactarea nu este un moft: la concurs, ea face
diferența dintre o idee bună și punctajul maxim.

\medskip
\noindent \textbf{Despre probleme.} Aproape toate sunt inventate de mine. Excepțiile (problemele clasice și cele preluate de la concursuri) sunt indicate explicit. Am încercat, de asemenea, să includ cel puțin o problemă pentru fiecare rezultat enunțat în partea de teorie: o lemă care nu „lucrează'' nicăieri nu merită ținută minte.

\vspace{1cm}
\begin{flushright}
{\itshape Vlad-Ștefan Constantinescu}\\
{\small\textcolor{grey}{București, 2026}}
\end{flushright}
\clearpage

%% ============================================================
%%  MULȚUMIRI
%% ============================================================
\thispagestyle{empty}
\addcontentsline{toc}{chapter}{Mulțumiri}
\vspace*{1.5cm}
{\noindent\color{onm}\rule{\linewidth}{3pt}}\par\vspace{3mm}
{\noindent\fontsize{28}{32}\selectfont\bfseries\color{ink} Mulțumiri}\par\vspace{3mm}
{\noindent\color{onm}\rule{\linewidth}{3pt}}\par\vspace{8mm}

\noindent Înainte de orice, aș vrea să îi mulțumesc unui vechi prieten, \textbf{Deluen Eliaz}, care mi-a fost profesor în clasele 8--9 la Liceul Francez „Anna de Noailles''. A fost primul om care a crezut în talentul meu matematic, și asta contează mai mult decât orice lecție.

\medskip
\noindent Îi mulțumesc \textbf{domnului Cristian Teodor Mangra}, profesorul meu de matematică la Colegiul Național „Tudor Vianu'' în clasele 10--12. A fost mereu aliatul nostru și ne-a sprijinit în pregătirea noastră.

\medskip
\noindent Îi mulțumesc \textbf{domnului Flavian Georgescu} și \textbf{domnului Traian Preda}, care m-au pregătit pentru olimpiadă și mi-au deschis accesul la tehnici și perspective pe care manualele nu le ating.

\medskip
\noindent Îi mulțumesc \textbf{domnului Cătălin Gherghe}, președintele Societății de Matematică din România, organizatorul Olimpiadei Naționale de Matematică și al Lotului Olimpic de Seniori. Îi mulțumesc nu doar pentru efortul organizatoric pe care îl depune, ci și pentru pregătirea pe care ne-o oferă la Cercul de Excelență.



\clearpage\color{ink}
\part{ALGEBRĂ}
\chapter{Construcții: ce obiecte să încerci}\label{ch:obiecte}

\emph{Notă:} acesta este pandantul algebric al capitolului „Redactare și probleme de construcții'' din partea de analiză --- o \emph{atitudine}, nu un capitol de teorie, și de aceea stă la începutul întregii părți de algebră: se folosește peste tot, și la clasa a XI-a (matrice), și la clasa a XII-a (structuri algebrice). La analiză, subpunctele de tip „dați un exemplu / rezultă că...?'' se rezolvau dintr-un catalog scurt de funcții ($|x|$, Dirichlet, $\sqrt{x^2+\varepsilon}$...); la algebră, ele se rezolvă dintr-un catalog la fel de scurt de \emph{obiecte}. Exemplele de mai jos invocă noțiuni tratate riguros abia în capitolele următoare (blocuri, Vandermonde, Cayley, teorema de structură a abelienelor); la prima citire se rețin \emph{reflexele}, iar la capitolele respective se revine aici.

\section{În general}

La ONM, un subpunct b) de tip „dați un exemplu'' sau „rămâne adevărat dacă...?'' nu pică din cer: \textbf{subpunctul a) este busola lui}. Punctul a) fixează \emph{invariantul} pe care problema îl testează (la a XI-a: rangul, urma, inversabilitatea, comutarea $AB=BA$; la a XII-a: ordinul elementelor, comutativitatea, normalitatea, cardinalul unei mulțimi de soluții), iar b) cere aproape întotdeauna un obiect în care \emph{exact acel invariant} se comportă altfel. Înainte de a răsfoi catalogul, gândește logic ce ți-a arătat a).

\begin{tehbox}{Rețeta în trei pași (comună celor două clase)}
\begin{enumerate}[leftmargin=*, itemsep=2pt]
\item \textbf{Identifică proprietatea-cheie din a).} Ce anume s-a demonstrat și \emph{ce ipoteză a fost esențială}? La a XI-a: inversabilitatea, rangul, urma, dimensiunea, faptul că matricele comută. La a XII-a: comutativitatea, primalitatea ordinului, normalitatea, faptul că un anumit număr divide altul.
\item \textbf{Întreabă-te unde poate eșua.} Negarea concluziei lui b) cere un obiect în care ipoteza esențială din a) pică --- de obicei \emph{cel mai mic} astfel de obiect. Demonstrează întâi mici \emph{restricții obligatorii} (ce dimensiune minimă trebuie să aibă matricea? ce ordine sunt posibile? poate fi obiectul comutativ?): fiecare restricție taie din spațiul de căutare, exact ca la analiză.
\item \textbf{Caută în catalogul clasei tale, de la mic la mare.} Nu inventa obiecte exotice. La a XI-a: rezolvă cazul $2\times2$ prin calcul direct și \emph{ridică-l pe blocuri} la orice dimensiune (secțiunea următoare); ține la îndemână rotația $P$ și blocul Jordan $N$. La a XII-a: $\ZZ_n$, produsele $\ZZ_{n_1}\times\dots\times\ZZ_{n_l}$, permutările $S_n$, grupurile de matrice; la inele: $\ZZ_n$ cu $n$ compus, $M_2(K)$, produsele de inele. Și nu uita: de multe ori b) \emph{refolosește} exact obiectul sau morfismul construit la a) --- nu reinventa ce ți s-a dat deja.
\end{enumerate}
\end{tehbox}

\section{Clasa a XI-a: matrice}

La clasa a XI-a obiectele sunt matricele, iar două tehnici acoperă majoritatea problemelor de existență: \emph{trucul blocurilor} (orice fenomen $2\times2$ se ridică la orice dimensiune) și \emph{recunoașterea circulantelor} (rotația deghizată în matrice).

\subsection{Cazuri mici și trucul blocurilor}

La problemele de existență cu matrice --- „găsiți $A,B\in M_n(\CC)$ cu proprietatea...'' --- dimensiunea $n$ este adesea o sperietoare falsă: fenomenul se produce deja la $2\times2$, iar de acolo se \emph{ridică} la orice $n$ printr-un truc de matrice pe blocuri (Capitolul~\refv{ch:blocuri}{5}{3}). Exemplul mic se găsește prin calcul; blocurile îl fac să funcționeze la orice dimensiune.

\begin{tehbox}{De la $2\times2$ la $n\times n$: rețeta}
\begin{enumerate}[leftmargin=*, itemsep=2pt]
\item \textbf{Rezolvă cazul mic prin bash, impunând o formă simplă.} O matrice din $M_2$ are $4$ necunoscute --- două matrice, $8$: prea multe pentru un calcul orb. Impune o formă care respectă simetria problemei: \emph{ambele diagonale}, ambele triunghiulare, sau matrice-unitate $E_{ij}$. Sistemul rămas devine banal.
\item \textbf{Ridică exemplul cu blocuri de zerouri.} Scufundarea $A\mapsto\begin{psmallmatrix}A&0\\0&0\end{psmallmatrix}$ păstrează adunarea, înmulțirea și puterile:
\[
\begin{pmatrix}A&0\\0&0\end{pmatrix}\begin{pmatrix}B&0\\0&0\end{pmatrix}=\begin{pmatrix}AB&0\\0&0\end{pmatrix},
\]
deci orice identitate „fără $I_n$'' verificată în colțul $2\times2$ rămâne verificată în $M_n$: dimensiunea mare doar \emph{găzduiește} colțul.
\item \textbf{Dacă enunțul implică $I_n$ sau inversabilitatea, umple cu $I$, nu cu $0$.} Scufundarea $A\mapsto\begin{psmallmatrix}A&0\\0&I_{n-2}\end{psmallmatrix}$ păstrează produsele și inversele și duce $I_2$ în $I_n$ --- potrivită pentru ecuații de tip $X^k=I_n$ sau pentru exemple în $\GL_n$. (Atenție: aceasta \textbf{nu} păstrează sumele --- e un truc pur multiplicativ.)
\end{enumerate}
\end{tehbox}

\textbf{Exemplu (divizori ai lui zero de orice dimensiune).} \emph{Găsiți două matrice nenule $A,B\in M_n(\CC)$ cu $AB=O_n$.}

\emph{Pasul 1 (bash la $2\times2$, impunând forma diagonală).} Căutăm $A=\diag(a_1,a_2)$, $B=\diag(b_1,b_2)$; produsul e $\diag(a_1b_1,\,a_2b_2)$, deci sistemul este $a_1b_1=a_2b_2=0$, cu $A,B\neq O_2$. Alegem zerourile \emph{încrucișat}:
\[
A_0=\begin{pmatrix}1&0\\0&0\end{pmatrix},\qquad
B_0=\begin{pmatrix}0&0\\0&1\end{pmatrix},\qquad
A_0B_0=O_2 .
\]
\emph{Pasul 2 (ridicarea prin blocuri).} Pentru orice $n\ge2$ punem
\[
\begin{gathered}
A=\begin{pmatrix}A_0&0\\0&0\end{pmatrix},\quad
B=\begin{pmatrix}B_0&0\\0&0\end{pmatrix}\in M_n(\CC)\\
\Longrightarrow\quad
AB=\begin{pmatrix}A_0B_0&0\\0&0\end{pmatrix}=O_n,
\end{gathered}
\]
cu $A,B\neq O_n$. $\blacksquare$

Aceeași scufundare produce, gata calibrate, exemplele-tip ale clasei a XI-a: \emph{nilpotent de indice exact $k\le n$} --- blocul Jordan $J_k$ umplut cu zerouri (teorema lui Jordan pentru matrice nilpotente, din capitolul „Triangularizare, diagonalizare și forma Jordan''\invol{3}); \emph{idempotent de rang $r$} --- $\diag(I_r,0)$, cu $\Tr=\rang$ (Lema~\refv{lema:idemptrrang}{7.24}{3}); \emph{pereche necomutativă în $M_n$}, $n\ge2$ --- $E_{11}$ și $E_{12}$ umplute cu zerouri ($E_{11}E_{12}=E_{12}$, dar $E_{12}E_{11}=O$); \emph{element de ordin $k$ în $\GL_n(\RR)$} --- rotația de unghi $\frac{2\pi}k$ umplută cu $I_{n-2}$ (a doua scufundare!). Morala: pentru problemele de existență, „cazul $2\times2$ \emph{este} cazul general''.

\subsection{Matrice circulantă}

Un obiect care apare obsesiv la determinanți și sisteme „cu simetrie ciclică'' --- și care merită recunoscut din prima privire.

\begin{tehbox}{Circulanta $=$ polinom în rotația $P$}
\emph{Matricea circulantă} asociată numerelor $c_0,c_1,\dots,c_{n-1}$ este matricea ale cărei linii se obțin una din alta prin rotire cu un pas spre dreapta:
\[
C=\mathrm{circ}(c_0,c_1,\dots,c_{n-1})=\begin{pmatrix}
c_0&c_1&\cdots&c_{n-1}\\
c_{n-1}&c_0&\cdots&c_{n-2}\\
\vdots&&\ddots&\vdots\\
c_1&c_2&\cdots&c_0
\end{pmatrix}.
\]
Piesa fundamentală este cea mai simplă dintre ele, \textbf{matricea de rotație} $P=\mathrm{circ}(0,1,0,\dots,0)$, care \emph{rotește coordonatele}:
\[
P\begin{pmatrix}x_1\\x_2\\\vdots\\x_n\end{pmatrix}=\begin{pmatrix}x_2\\\vdots\\x_n\\x_1\end{pmatrix},
\qquad P^n=I_n .
\]
Observația-cheie: $P^k$ rotește cu $k$ pași, iar
\[
\begin{gathered}
C=c_0I_n+c_1P+c_2P^2+\dots+c_{n-1}P^{n-1}=f(P),\\
f(t)=c_0+c_1t+\dots+c_{n-1}t^{n-1}:
\end{gathered}
\]
\textbf{orice circulantă este un polinom în rotația $P$}. De aici curg toate proprietățile:
\begin{enumerate}[leftmargin=*, itemsep=2pt]
\item \textbf{Cum acționează la înmulțire.} $Cx=c_0x+c_1Px+\dots+c_{n-1}P^{n-1}x$ --- o \emph{combinație ponderată a rotitelor} vectorului $x$; pe componente, $(Cx)_i=\sum_k c_k\,x_{i+k}$, cu indicii citiți ciclic (modulo $n$). Înmulțirea cu o circulantă este deci o „mediere ciclică'' a coordonatelor --- de aceea circulantele apar natural la sisteme și recurențe ciclice.
\item \textbf{Produsul a două circulante e circulantă, iar circulantele comută:} $f(P)\,g(P)=(fg)(P)=g(P)\,f(P)$. Calculul cu circulante este calcul cu \emph{polinoame}, redus modulo $X^n-1$ (căci $P^n=I_n$) --- în limbajul clasei a XII-a, circulantele formează un inel comutativ.
\item \textbf{Vectori proprii gratuiți.} Fie $\omega=\cos\frac{2\pi}n+i\sin\frac{2\pi}n$ și, pentru $j=0,1,\dots,n-1$, vectorul $v_j=\bigl(1,\omega^j,\omega^{2j},\dots,\omega^{(n-1)j}\bigr)^t$. Rotirea unei progresii geometrice o înmulțește cu rația: $Pv_j=\omega^j v_j$; aplicând polinomul $f$,
\[
C\,v_j=f(\omega^j)\,v_j .
\]
Cei $n$ vectori sunt liniar independenți (matricea lor e Vandermonde cu noduri distincte $\omega^0,\dots,\omega^{n-1}$, vezi determinantul Vandermonde din Capitolul~\refv{ch:detspeciali}{10}{3}), deci am diagonalizat dintr-un foc \emph{orice} circulantă, și
\[
\boxed{\;\det C=\prod_{j=0}^{n-1}f(\omega^j)=\prod_{j=0}^{n-1}\bigl(c_0+c_1\omega^j+c_2\omega^{2j}+\dots+c_{n-1}\omega^{(n-1)j}\bigr).\;}
\]
\item \textbf{Bonusul liniei constante:} $v_0=(1,1,\dots,1)^t$ dă valoarea proprie $f(1)=c_0+c_1+\dots+c_{n-1}$ --- suma pe linie. De aceea, la determinanții circulanți, factorul $c_0+\dots+c_{n-1}$ „se vede'' imediat (adună toate coloanele la prima).
\end{enumerate}
\end{tehbox}

\textbf{Problemă-tip (acțiunea rotației asupra unei matrice).} \emph{Fie $A=(a_{ij})\in M_n(\CC)$ oarecare. Găsiți matrice $X,Y\in M_n(\CC)$, independente de $A$, astfel încât: \textbf{(a)} $AX$ să fie $A$ cu coloanele rotite ciclic cu un pas spre dreapta; \textbf{(b)} $YA$ să fie $A$ cu prima linie \emph{ștearsă}: liniile $2,\dots,n$ urcă toate cu un pas, iar ultima linie se completează cu zerouri.}

\emph{Soluție.} \textbf{(a)} Chiar rotația: $X=P$. Coloana $j$ a lui $P$ are un singur $1$, pe linia $j-1$, deci $(AP)_{ij}=a_{i,j-1}$ (indici modulo $n$): fiecare coloană a lui $AP$ „vine'' de la vecina din stânga --- rotire cu un pas spre dreapta. (Regula generală: înmulțirea \emph{la dreapta} acționează pe coloane, cea \emph{la stânga} pe linii --- simetric, $PA$ rotește liniile cu un pas în sus.)

\textbf{(b)} Rotația pură nu ajunge: în $PA$ prima linie a lui $A$ nu dispare, ci se \emph{întoarce} pe ultima poziție. Ștergere $=$ rotire \emph{fără} întoarcere: anulăm exact „$1$-ul de întoarcere'' al lui $P$ --- cel din colțul $(n,1)$, singurul din \textbf{prima coloană} a circulantei --- și obținem
\[
Y=N=\begin{pmatrix}0&1&&\\&0&\ddots&\\&&\ddots&1\\&&&0\end{pmatrix},
\qquad
(NA)_{ij}=\begin{cases}a_{i+1,j},&i<n,\\[1mm] 0,&i=n,\end{cases}
\]
exact cerința. $\blacksquare$

\emph{Morala:} $N$ este blocul Jordan nilpotent $J_n$ (teorema lui Jordan pentru matrice nilpotente, din capitolul „Triangularizare, diagonalizare și forma Jordan''\invol{3}) --- circulanta $P$ „cu memoria tăiată''. Comparația spune totul: $P$ rotește la nesfârșit ($P^n=I_n$), pe când $N$ tot împinge în sus și șterge --- după $n$ pași nu mai rămâne nimic ($N^n=O_n$). Nilpotența lui $N$ „se vede'' din poveste, fără niciun calcul.

\textbf{Exemplu (shortlist ONM).} \emph{Găsiți o matrice $X\in M_n(\ZZ)$ astfel încât $X^n=I_n$, dar $X^k\neq I_n$ pentru $1\le k\le n-1$.}

\emph{Soluție.} Citim întâi busola: cerința $X^n=I_n$, $X^k\neq I_n$ spune „element de ordin \emph{exact} $n$'', iar ipoteza $M_n(\ZZ)$ --- intrări \emph{întregi} --- interzice candidații „leneși'': $\diag(\omega,1,\dots,1)$ cere numere complexe, iar rotația plană de unghi $\frac{2\pi}n$ cere sinusuri și cosinusuri. Ne trebuie un obiect de ordin $n$ cu intrări întregi --- și el stă deja pe masă: $X=P$, rotația, care are doar intrări $0$ și $1$. Fiecare coordonată revine pe locul ei după exact $n$ rotiri, deci $P^n=I_n$. Rămâne să vedem că $P^k\neq I_n$ pentru $1\le k\le n-1$ --- două argumente, fiecare instructiv:
\begin{itemize}[leftmargin=*, itemsep=1pt]
\item \emph{Prin invariant (urma).} $P^k$ este matricea rotirii cu $k$ pași: $(P^k)_{ij}=1$ exact când $j\equiv i+k\pmod n$. O intrare diagonală ar cere $j=i$, adică $n\mid k$ --- imposibil pentru $1\le k\le n-1$. Deci $\Tr(P^k)=0\neq n=\Tr(I_n)$: urma este invariantul care le desparte (busola de la începutul capitolului, în acțiune).
\item \emph{Prin vectori proprii.} Privim pe $P$ ca matrice complexă (avem voie: $M_n(\ZZ)\subset M_n(\CC)$). $Pv_1=\omega v_1$, cu $\omega=\cos\frac{2\pi}n+i\sin\frac{2\pi}n$ rădăcină \emph{primitivă} a unității; dacă am avea $P^k=I_n$, aplicând pe $v_1$ ar rezulta $\omega^k=1$, deci $n\mid k$. Așadar ordinul lui $P$ este \emph{exact} $n$. $\blacksquare$
\end{itemize}

\emph{Bonus:} mulțimea $\{I_n,P,P^2,\dots,P^{n-1}\}$ este un grup ciclic de ordin $n$ format din matrice cu intrări întregi: rotația $P$ este încarnarea matriceală a lui $\ZZ_n$ --- clasa a XI-a și clasa a XII-a, în aceeași matrice. \emph{Semnalul la concurs:} un determinant sau un sistem în care fiecare linie este precedenta „rotită'' $\to$ scrie matricea ca $f(P)$ și factorizează prin rădăcinile unității.

\section{Clasa a XII-a: structuri algebrice}

La clasa a XII-a obiectele sunt grupurile și inelele; catalogul de mai jos acoperă aproape toate subpunctele „dați un exemplu'' de la ONM.

\subsection{Catalogul de grupuri}

\begin{tehbox}{Gândește simplu: cele cinci familii standard}
\begin{enumerate}[leftmargin=*, itemsep=2pt]
\item \textbf{Ciclicele $\ZZ_n$} --- prototipul abelian; prin teorema de clasificare (Teorema~\refv{teo:clasif}{2.71}{4}), $\ZZ_n$ este \emph{singurul} grup ciclic de ordin $n$, deci orice enunț despre „un grup ciclic'' se testează direct aici. Are element de ordin $d$ exact pentru $d\mid n$, iar ecuația $x^k=e$ are exact $(k,n)$ soluții (Propoziția~\refv{prop:solutii}{2.75}{4}).
\item \textbf{Produsele $\ZZ_{n_1}\times\ZZ_{n_2}\times\dots\times\ZZ_{n_l}$} --- abelienele cu ordine \emph{controlate}: $\ord(x_1,\dots,x_l)=[\ord x_1,\dots,\ord x_l]$, deci alegi componentele ca să obții exact spectrul de ordine dorit. Piesele-vedetă: grupul lui Klein $\ZZ_2\times\ZZ_2$ (cel mai mic abelian \emph{neciclic}) și $\ZZ_2^{\,k}$ (toate elementele de ordin $\le2$). Esențial: prin teorema de structură (Corolarul~\refv{cor:abelianfinit}{2.153}{4}), \emph{orice} grup abelian finit este un astfel de produs --- căutarea printre exemplele abeliene finite \textbf{se încheie} în această familie; nu există altceva.
\item \textbf{Permutările $S_n$ (și $A_n$)} --- sursa universală de neabelian: prin Cayley (Teorema~\refv{teo:cayley}{2.93}{4}), \emph{orice} grup finit trăiește într-un $S_n$. $S_3$ este cel mai mic grup neabelian (ordin $6$). Ordinul unei permutări $=$ c.m.m.m.c.-ul lungimilor ciclilor, deci construiești ieftin elemente de orice ordin: $(1\,2\,3)(4\,5)\in S_5$ are ordinul $6$. Iar $A_4$ (ordin $12$, fără subgrup de ordin $6$) este contraexemplul-standard la reciproca teoremei lui Lagrange (observația de după Corolarul~\refv{cor:lagrange}{2.85}{4}).
\item \textbf{Matricele} --- terenul infinit și neabelian: în $\GL_2(\RR)$, rotația de unghi $\frac{2\pi}{n}$ e un element de ordin $n$ într-un grup \emph{infinit}; matricele superior triunghiulare cu $1$ pe diagonală, peste $\ZZ_p$, dau $p$-grupuri neabeliene de ordin $p^3$; iar cuaternionii $Q_8=\{\pm1,\pm i,\pm j,\pm k\}$ (cu $i^2=j^2=k^2=ijk=-1$) --- neabelian de ordin $8$ în care \emph{toate} subgrupurile sunt normale și care are un \emph{singur} element de ordin $2$.
\item \textbf{Diedralele și „abelienele multiplicative''.} Grupul diedral $D_n$ (simetriile poligonului regulat cu $n$ laturi: $n$ rotații $+$ $n$ reflexii, ordin $2n$, neabelian pentru $n\ge3$) are multe elemente de ordin $2$ fără a fi abelian --- reflexiile. Iar în structuri multiplicative, abelienele „gata făcute'': $U_n\subset\CC^*$ (rădăcinile unității), $U(\ZZ_n)$ --- de pildă $U(\ZZ_8)=\{\cl1,\cl3,\cl5,\cl7\}\simeq$ Klein, abelian neciclic; și reuniunea $\bigcup_{n\ge1}U_n$: grup \emph{infinit} în care orice element are ordin \emph{finit}.
\end{enumerate}
\end{tehbox}

\subsection{Dicționarul: condiția din enunț $\to$ obiectul de încercat}

\begin{tehbox}{Condiție cerută $\to$ primul candidat}
\begin{enumerate}[leftmargin=*, itemsep=1pt]
\item „abelian, dar neciclic'' $\to$ $\ZZ_2\times\ZZ_2$ (general: $\ZZ_p\times\ZZ_p$).
\item „neabelian cât mai mic'' $\to$ $S_3$ (ordin $6$); la ordin $8$: $D_4$ sau $Q_8$.
\item „element de ordin $d$ prescris'' $\to$ $\ZZ_d$; sau cicli disjuncți cu $[\,\text{lungimi}\,]=d$ în $S_n$.
\item „$x^2=e$ pentru toți $x$'' (exponent $2$) $\to$ $\ZZ_2^{\,k}$.
\item „infinit, dar cu elemente de ordin finit'' $\to$ rotații în $\GL_2(\RR)$; $\bigcup_nU_n\subset\CC^*$.
\item „$x^k=e$ cu mai mult de $k$ soluții'' $\to$ Klein pentru $k=2$ (patru soluții). Morala contrastului: în grupuri ciclice numărul e exact $(k,n)$ (Propoziția~\refv{prop:solutii}{2.75}{4}), iar într-un \emph{corp} nu se poate depăși $k$ (Teorema~\refv{teo:numarradacini}{4.16}{4}) --- deci exemplul trebuie căutat departe de ciclice.
\item „subgrup nenormal'' $\to$ $\langle(1\,2)\rangle$ în $S_3$.
\item „contraexemplu la reciproca lui Lagrange'' $\to$ $A_4$, fără subgrup de ordin $6$.
\item „același ordin, neizomorfe'' $\to$ $\ZZ_4$ vs.\ Klein (ordin $4$); $\ZZ_6$ vs.\ $S_3$ (ordin $6$).
\item „exact un element de ordin $2$'' $\to$ $\ZZ_4$ (abelian) sau $Q_8$ (neabelian).
\item „$HK$ nu e subgrup'' $\to$ $H=\langle(1\,2)\rangle$, $K=\langle(1\,3)\rangle$ în $S_3$: $|HK|=\frac{|H|\,|K|}{|H\cap K|}=4\nmid6$ (formula produsului, Lema~\refv{lema:produsHK}{2.159}{4}, plus Lagrange).
\end{enumerate}
\end{tehbox}

\begin{obs}[Checklist-ul ordinelor mici]
Contraexemplele mici trăiesc la ordine mici, iar la multe ordine „nu ai unde căuta'': la ordin prim $p$ grupul e ciclic (Corolarul~\refv{cor:lagrange}{2.85}{4}, punctul (iii)), la ordin $p^2$ e abelian (Corolarul~\refv{cor:p2}{2.131}{4}). Acțiunea e concentrată la $4,6,8,12$: ordin $4$: $\ZZ_4$, Klein; ordin $6$: $\ZZ_6$, $S_3$; ordin $8$: $\ZZ_8$, $\ZZ_4\times\ZZ_2$, $\ZZ_2^3$, $D_4$, $Q_8$; ordin $12$: $\ZZ_{12}$, $\ZZ_2\times\ZZ_6$, $D_6$, $A_4$ (și încă unul, diciclicul). Când cauți un contraexemplu, parcurge lista în această ordine --- dacă nu e aici, abia apoi urcă.
\end{obs}

\subsection{Catalogul de inele}

Aceeași filozofie funcționează la inele și corpuri --- doar catalogul se schimbă:

\begin{tehbox}{Piesele standard la inele}
\begin{enumerate}[leftmargin=*, itemsep=1pt]
\item $\ZZ_n$ cu $n$ \emph{compus} --- divizori ai lui zero ($\cl2\cdot\cl3=\cl0$ în $\ZZ_6$) și idempotenți nebanali ($\cl3^2=\cl3$, $\cl4^2=\cl4$ în $\ZZ_6$).
\item $M_2(K)$ --- necomutativ, cu nilpotenți ($E_{12}^2=O$) și divizori ai lui zero; grupul unităților e $\GL_2(K)$.
\item $\ZZ_2^{\,k}$ (sau $(\mathcal P(M),\Delta,\cap)$) --- inelele booleene.
\item Produsul $A\times B$ --- mașina de idempotenți: $(1,0)^2=(1,0)$; tot el strică integritatea: $(1,0)(0,1)=(0,0)$.
\item $\ZZ[i]$, $\ZZ[\sqrt d\,]$ --- inele integre cu unități controlate prin normă.
\item $K[X]$ și $\ZZ_p[X]/(f)$ --- polinoamele și corpurile finite cu $p^{\deg f}$ elemente.
\item Matricele superior triunghiulare --- exemplu-tip de inel necomutativ cu nilpotenți „vizibili''.
\end{enumerate}
\end{tehbox}

Și aici busola rămâne subpunctul a): dacă a) a demonstrat, de pildă, că într-un inel \emph{fără nilpotenți} idempotenții sunt centrali, atunci b) („rămâne adevărat în general?'') te trimite exact la un inel \emph{cu} nilpotenți (adică, din catalog, la matrice).



\part{Clasa a XII-a: Structuri algebrice}
\chapter{Grupuri}

\section{Definiții preliminare}

\subsection{Legi de compoziție internă}

\begin{definitie}
O \emph{lege de compoziție internă} pe mulțimea nevidă $M$ este o funcție $*:M\times M\to M$, $(x,y)\mapsto x*y$. O submulțime $H\subseteq M$ este \emph{parte stabilă} dacă $x,y\in H\Rightarrow x*y\in H$; în acest caz $*$ induce o lege de compoziție pe $H$.
\end{definitie}

\begin{definitie}
Legea $*$ este \emph{asociativă} dacă $(x*y)*z=x*(y*z)$ pentru orice $x,y,z$, respectiv \emph{comutativă} dacă $x*y=y*x$ pentru orice $x,y$. Un element $e\in M$ este \emph{element neutru} dacă $x*e=e*x=x$ pentru orice $x\in M$. Un element $x$ este \emph{simetrizabil} dacă există $x'\in M$ cu $x*x'=x'*x=e$; $x'$ se numește \emph{simetricul} lui $x$. O structură $(M,*)$ asociativă și cu element neutru se numește \emph{monoid}.
\end{definitie}

\begin{definitie}[semigrup]
O pereche $(S,\cdot)$ formată dintr-o mulțime nevidă și o lege de compoziție \emph{asociativă} se numește \emph{semigrup}; dacă legea este și comutativă, semigrupul se numește \emph{comutativ} (abelian). Un semigrup cu element neutru este exact un \emph{monoid} (definit mai sus). Astfel orice grup este monoid și orice monoid este semigrup (grupuri $\subset$ monoizi $\subset$ semigrupuri, ca familii de structuri): trecerea de la semigrup la monoid adaugă axioma elementului neutru, iar cea de la monoid la grup adaugă simetrizabilitatea tuturor elementelor.
\end{definitie}

\begin{propozitie}[unicitate]\label{prop:unicitate}
\emph{(i)} Dacă legea are element neutru, acesta este unic. \emph{(ii)} Într-un monoid, simetricul unui element simetrizabil este unic. \emph{(iii)} Dacă $x,y$ sunt simetrizabile într-un monoid, atunci $x*y$ este simetrizabil și $(x*y)'=y'*x'$ (atenție la ordine!), iar $(x')'=x$.
\end{propozitie}

\begin{proof}
(i) Dacă $e$ și $e'$ sunt neutre, $e=e*e'=e'$. (ii) Dacă $x'$ și $x''$ sunt simetrice ale lui $x$, atunci
$x''=x''*e=x''*(x*x')=(x''*x)*x'=e*x'=x'$. (iii) Calcul direct:
$(x*y)*(y'*x')=x*(y*y')*x'=x*x'=e$ și analog în cealaltă ordine; egalitatea $(x')'=x$ rezultă din simetria definiției.
\end{proof}

\begin{observatie}[notații și puteri]
În notație \emph{multiplicativă} scriem $xy$, $1$ sau $e$ pentru neutru și $x^{-1}$ pentru simetric (numit \emph{invers}); în notație \emph{aditivă} scriem $x+y$, $0$ și $-x$ (numit \emph{opus}). Definim $x^0=e$, $x^{n+1}=x^n\cdot x$ și, pentru elemente simetrizabile, $x^{-n}=(x^{-1})^n$. Într-un monoid au loc $x^{m+n}=x^m x^n$ și $(x^m)^n=x^{mn}$ pentru toți exponenții pentru care puterile au sens. \textbf{Capcană:} în general $(xy)^n\neq x^n y^n$; egalitatea are loc pentru orice $n$ dacă $xy=yx$ (inducție).
\end{observatie}

\begin{exemplu}
$\NN$ este parte stabilă a lui $(\ZZ,+)$, dar în $(\NN,+)$ singurul element simetrizabil este $0$. În monoidul $(\ZZ,\cdot)$ elementele simetrizabile sunt $\pm 1$.
\end{exemplu}

\subsection{Grupuri: definiție și exemple fundamentale}

\begin{definitie}
$(G,\cdot)$ este \emph{grup} dacă legea este asociativă, are element neutru și orice element este simetrizabil. Grupul este \emph{abelian} (comutativ) dacă legea este comutativă. \emph{Ordinul grupului} este cardinalul $|G|$.
\end{definitie}

\begin{exemplu}[exemplele fundamentale]\label{ex:grupuri}
\begin{enumerate}[label=(\alph*),nosep]
\item Grupuri aditive: $(\ZZ,+)$, $(\QQ,+)$, $(\RR,+)$, $(\CC,+)$ --- abeliene, infinite.
\item Grupuri multiplicative: $(\QQ^*,\cdot)$, $(\RR^*,\cdot)$, $(\CC^*,\cdot)$, $((0,\infty),\cdot)$.
\item $(\ZZ_n,+)$ --- grupul claselor de resturi modulo $n$, abelian, $|\ZZ_n|=n$.
\item $U_n=\{z\in\CC\mid z^n=1\}=\{1,\varepsilon,\varepsilon^2,\dots,\varepsilon^{n-1}\}$, unde $\varepsilon=\cos\frac{2\pi}{n}+i\sin\frac{2\pi}{n}$: \emph{grupul rădăcinilor de ordinul $n$ ale unității}, abelian, cu $n$ elemente, față de înmulțire.
\item $(S_n,\circ)$ --- grupul permutărilor mulțimii $\{1,\dots,n\}$, $|S_n|=n!$, \emph{neabelian} pentru $n\ge 3$.
\item $GL_2(\RR)$ --- grupul matricelor $2\times 2$ inversabile, cu înmulțirea; neabelian.
\item $G=\RR\setminus\{1\}$ cu $x\circ y=x+y-xy$: neutrul este $0$, iar simetricul lui $x$ este $\frac{x}{x-1}$ (verificare directă). Vom vedea în §\ref{sec:morfisme} că $(G,\circ)\simeq(\RR^*,\cdot)$.
\end{enumerate}
\end{exemplu}

\begin{propozitie}[reguli de calcul în grup]\label{prop:simplificare}
Într-un grup $G$: \emph{(i)} au loc \emph{simplificările} $ax=ay\Rightarrow x=y$ și $xa=ya\Rightarrow x=y$; \emph{(ii)} ecuațiile $ax=b$ și $xa=b$ au soluții unice, $x=a^{-1}b$, respectiv $x=ba^{-1}$.
\end{propozitie}

\begin{proof}
(i) Înmulțim la stânga (respectiv la dreapta) cu $a^{-1}$. (ii) Verificare directă; unicitatea rezultă din (i).
\end{proof}

\begin{observatie}
Din (ii): în tabla unui grup finit, fiecare element apare \textbf{exact o dată} pe fiecare linie și pe fiecare coloană (proprietatea de ,,sudoku``). Este primul test de verificat când se cere să se decidă dacă o tablă dată definește un grup.
\end{observatie}

\begin{propozitie}\label{prop:x2e}
Dacă $x^2=e$ pentru orice $x\in G$, atunci $G$ este abelian.
\end{propozitie}

\begin{proof}
Ipoteza spune că fiecare element este propriul invers. Atunci $xy=(xy)^{-1}=y^{-1}x^{-1}=yx$.
\end{proof}

\subsection{Subgrupuri. Ordinul unui element}

\begin{definitie}
$H\subseteq G$ este \emph{subgrup} al grupului $G$ (notăm $H\le G$) dacă $H$ este parte stabilă și $(H,\cdot)$ este grup cu legea indusă. Subgrupurile $\{e\}$ și $G$ se numesc \emph{improprii}.
\end{definitie}

\begin{propozitie}[criteriul de subgrup]\label{prop:criteriu}
Fie $H\subseteq G$ nevidă. Atunci $H\le G$ dacă și numai dacă $xy^{-1}\in H$ pentru orice $x,y\in H$.
\end{propozitie}

\begin{proof}
Necesitatea este clară. Reciproc: luând $y=x$ obținem $e=xx^{-1}\in H$; apoi $x^{-1}=ex^{-1}\in H$ pentru orice $x\in H$; în fine $xy=x(y^{-1})^{-1}\in H$, deci $H$ este parte stabilă, iar asociativitatea se moștenește.
\end{proof}

\begin{propozitie}[cazul finit]\label{prop:finitstabil}
Dacă $H\subseteq G$ este \textbf{finită}, nevidă și stabilă la operație, atunci $H\le G$.
\end{propozitie}

\begin{proof}
Fie $x\in H$. Puterile $x,x^2,x^3,\dots$ rămân în $H$, care este finită, deci există $i<j$ cu $x^i=x^j$. Simplificând în $G$, $x^{j-i}=e\in H$. Dacă $j-i=1$, atunci $x=e$ și $x^{-1}=e\in H$; altfel $x^{-1}=x^{j-i-1}\in H$. Concluzia rezultă din Propoziția~\ref{prop:criteriu}.
\end{proof}

\begin{lema}[un grup nu e reuniunea a două subgrupuri proprii]\label{lema:reuniune}
Dacă $H,K\le G$ și $H\cup K=G$, atunci $H=G$ sau $K=G$.
\end{lema}

\begin{proof}
Presupunem, prin absurd, că $H\neq G$ și $K\neq G$. Cum $H\cup K=G$ dar $K\neq G$, există $h\in H\setminus K$; simetric, există $k\in K\setminus H$. Elementul $hk$ aparține lui $G=H\cup K$, deci $hk\in H$ sau $hk\in K$.
\begin{itemize}[nosep,leftmargin=1.4em]
  \item Dacă $hk\in H$, atunci $k=h^{-1}(hk)\in H$, contrazicând $k\notin H$.
  \item Dacă $hk\in K$, atunci $h=(hk)k^{-1}\in K$, contrazicând $h\notin K$.
\end{itemize}
Ambele cazuri sunt imposibile, deci presupunerea a fost falsă: $H=G$ sau $K=G$.
\end{proof}

\begin{observatie}
Rezultatul este optim: un grup \emph{poate} fi reuniunea a \textbf{trei} subgrupuri proprii --- de pildă grupul lui Klein $\ZZ_2\times\ZZ_2$ este reuniunea celor trei subgrupuri de ordin $2$ ale sale. Lema arată doar că două nu ajung niciodată.
\end{observatie}

\begin{definitie}
Pentru $x\in G$, mulțimea $\langle x\rangle=\{x^k\mid k\in\ZZ\}$ este un subgrup, numit \emph{subgrupul ciclic generat de $x$}. \emph{Ordinul} lui $x$, notat $\ord(x)$, este cel mai mic $n\ge 1$ cu $x^n=e$, dacă există; altfel $x$ are ordin infinit.
\end{definitie}

\begin{teorema}[proprietățile ordinului]\label{teo:ordin}
Fie $x\in G$ cu $\ord(x)=n$ finit. Atunci:
\emph{(i)} $x^m=e\iff n\mid m$;\quad
\emph{(ii)} $x^i=x^j\iff i\equiv j \pmod n$;\quad
\emph{(iii)} $\langle x\rangle=\{e,x,\dots,x^{n-1}\}$ are exact $n$ elemente, adică $|\langle x\rangle|=\ord(x)$.
Dacă $x$ are ordin infinit, toate puterile $x^k$, $k\in\ZZ$, sunt distincte.
\end{teorema}

\begin{proof}
(i) Dacă $n\mid m$, evident $x^m=e$. Reciproc, scriem $m=nq+r$ cu $0\le r<n$; atunci $e=x^m=(x^n)^q x^r=x^r$, iar minimalitatea lui $n$ forțează $r=0$. (ii) $x^i=x^j\iff x^{i-j}=e$, apoi (i). (iii) Din (ii), elementele $e,x,\dots,x^{n-1}$ sunt distincte, iar orice putere $x^k$ coincide cu $x^{k \bmod n}$. Cazul infinit: $x^i=x^j$ cu $i\neq j$ ar da $x^{|i-j|}=e$, contradicție.
\end{proof}

\begin{propozitie}[ordinul unei puteri]\label{prop:ordputere}
Dacă $\ord(x)=n$ și $k\in\ZZ$, atunci $\ord(x^k)=\dfrac{n}{(n,k)}$, unde $(n,k)$ este c.m.m.d.c.
\end{propozitie}

\begin{proof}
Fie $d=(n,k)$. Avem $(x^k)^m=e\iff n\mid km\iff \frac{n}{d}\mid \frac{k}{d}\,m\iff \frac{n}{d}\mid m$, ultima echivalență deoarece $\bigl(\frac{n}{d},\frac{k}{d}\bigr)=1$ (lema lui Gauss din aritmetică). Cel mai mic astfel de $m\ge1$ este $\frac{n}{d}$.
\end{proof}

\begin{propozitie}[ordinul: invers și conjugare]\label{prop:ordinv}
În orice grup $G$: \emph{(i)} $\ord(x^{-1})=\ord(x)$; \emph{(ii)} $\ord(g^{-1}xg)=\ord(x)$ pentru orice $g\in G$; \emph{(iii)} $\ord(xy)=\ord(yx)$ pentru orice $x,y\in G$.
\end{propozitie}

\begin{proof}
(i) $(x^{-1})^k=(x^k)^{-1}$, deci $x^k=e\iff(x^{-1})^k=e$: cele două elemente au exact aceleași puteri egale cu $e$. (ii) $(g^{-1}xg)^k=g^{-1}x^kg$ (produsul se telescopează), deci $(g^{-1}xg)^k=e\iff x^k=e$. (iii) $yx=x^{-1}(xy)x$, apoi aplicăm (ii). Toate afirmațiile rămân valabile și pentru ordin infinit.
\end{proof}

\begin{teorema}[ordinul unui produs de elemente care comută]\label{teo:ordprodus}
Fie $a,b\in G$ cu $ab=ba$, $\ord a=m$, $\ord b=n$, și fie $d=(m,n)$, $\ell=[m,n]$ (c.m.m.m.c.). Atunci
\[
\ord(ab)\ \Big|\ \ell
\qquad\text{și}\qquad
\frac{\ell}{d}\ \Big|\ \ord(ab).
\]
În particular, dacă $(m,n)=1$, atunci $\ord(ab)=mn$.
\end{teorema}

\begin{proof}
\emph{Prima divizibilitate}: $(ab)^{\ell}=a^{\ell}b^{\ell}=e$, deoarece $m\mid\ell$ și $n\mid\ell$ (aici este esențial $ab=ba$); apoi Teorema~\ref{teo:ordin}(i).

\emph{A doua}: fie $k=\ord(ab)$ și $c=a^k$. Din $(ab)^k=a^kb^k=e$ rezultă $c=a^k=b^{-k}$, deci $c\in\langle a\rangle\cap\langle b\rangle$. Atunci $c^m=(a^m)^k=e$ și $c^n=(b^n)^{-k}=e$, deci $\ord(c)$ divide și $m$, și $n$, prin urmare divide $d=(m,n)$; în particular $c^d=e$. Aceasta înseamnă $a^{kd}=e$ și $b^{kd}=e$, de unde $m\mid kd$ și $n\mid kd$, adică (simplificând cu $d$, care divide $m$ și $n$)
\[
\frac md\ \Big|\ k\qquad\text{și}\qquad\frac nd\ \Big|\ k.
\]
Observăm acum că $\bigl(\frac md,\frac nd\bigr)=1$: dacă un prim $r$ le-ar divide pe amândouă, atunci $rd\mid m$ și $rd\mid n$, deci $rd\mid(m,n)=d$ --- absurd. Un număr divizibil cu doi factori primi între ei este divizibil cu produsul lor, deci
\[
\frac md\cdot\frac nd=\frac{mn}{d^2}=\frac{\ell}{d}\ \Big|\ k.
\]
\emph{Cazul coprim}: pentru $d=1$ avem $\ell=mn$, iar cele două divizibilități dau $mn\mid\ord(ab)\mid mn$.
\end{proof}

\begin{corolar}[{element de ordin $[m,n]$}]\label{cor:lcm}
Dacă un grup \textbf{abelian} conține elemente de ordine $m$ și $n$, atunci conține un element de ordin $[m,n]$.
\end{corolar}

\begin{proof}
Scriem $\ell=[m,n]=p_1^{\alpha_1}\cdots p_r^{\alpha_r}$. Pentru fiecare $i$, puterea $p_i^{\alpha_i}$ divide $m$ sau $n$ (exponentul lui $p_i$ din $\ell$ este atins în unul dintre ele); dacă, de pildă, $p_i^{\alpha_i}\mid m$, atunci elementul $c_i=a^{m/p_i^{\alpha_i}}$ are ordinul $\frac{m}{(m,\,m/p_i^{\alpha_i})}=p_i^{\alpha_i}$ (Propoziția~\ref{prop:ordputere}). Obținem $c_1,\dots,c_r$ cu ordine $p_1^{\alpha_1},\dots,p_r^{\alpha_r}$, două câte două prime între ele; aplicând repetat cazul coprim din Teorema~\ref{teo:ordprodus}, produsul $c_1c_2\cdots c_r$ are ordinul $p_1^{\alpha_1}\cdots p_r^{\alpha_r}=\ell$.
\end{proof}

\begin{observatie}
(a) Comutativitatea este esențială: în $S_3$, transpozițiile $\sigma=(1\,2)$ și $\tau=(2\,3)$ au ordinul $2$, dar $\sigma\tau$ este un ciclu de lungime $3$, deci $\ord(\sigma\tau)=3\nmid[2,2]$. (b) Ambele margini din Teorema~\ref{teo:ordprodus} se ating: pentru $b=a^{-1}$ avem $\ord(ab)=\ord(e)=1=\ell/d$, iar în cazul coprim $\ord(ab)=\ell$. Între aceste margini, $\ord(ab)$ nu este determinat doar de $m$ și $n$.
\end{observatie}

\begin{propozitie}[intersecția de subgrupuri]\label{prop:intersectie}
Intersecția oricărei familii de subgrupuri ale lui $G$ este un subgrup al lui $G$.
\end{propozitie}

\begin{proof}
$e$ aparține fiecărui subgrup, deci și intersecției; dacă $x,y$ aparțin fiecărui subgrup din familie, atunci $xy^{-1}$ la fel, și aplicăm criteriul (Propoziția~\ref{prop:criteriu}).
\end{proof}

% --- grupuri mici + exponent, plasate la Ordinul unui element ---
\subsubsection*{Exponentul unui grup}

\begin{definitie}[exponentul grupului]
\emph{Exponentul} unui grup $G$, notat $\exp(G)$, este cel mai mic $n\ge1$ cu proprietatea $x^n=e$ pentru \textbf{orice} $x\in G$, dacă un astfel de $n$ există. Pentru $G$ finit, exponentul există (de exemplu $n=|G|$ convine, prin Corolarul~\ref{cor:lagrange}).
\end{definitie}

\begin{propozitie}[proprietățile exponentului]\label{prop:exponent}
Fie $G$ un grup finit. Atunci:
\emph{(i)} $\exp(G)$ este c.m.m.m.c.-ul ordinelor elementelor lui $G$ și, pentru $m\in\NN^*$: $x^m=e$ pentru orice $x$ $\iff\exp(G)\mid m$;
\emph{(ii)} $\exp(G)$ divide $|G|$;
\emph{(iii)} $\exp(G)$ și $|G|$ au \textbf{aceiași divizori primi};
\emph{(iv)} dacă $G$ este \textbf{abelian}, există un element de ordin $\exp(G)$; în consecință, $\exp(G)$ este cel mai mare ordin de element din $G$.
\end{propozitie}

\begin{proof}
(i) $x^m=e$ pentru toți $x$ $\iff\ord(x)\mid m$ pentru toți $x$ (Teorema~\ref{teo:ordin}(i)) $\iff L\mid m$, unde $L$ este c.m.m.m.c.-ul ordinelor; cel mai mic asemenea $m$ este chiar $L$.
(ii) $x^{|G|}=e$ pentru orice $x$ (Corolarul~\ref{cor:lagrange}), apoi (i).
(iii) Dacă un prim $q$ divide $\exp(G)$, atunci $q\mid|G|$ din (ii). Reciproc, dacă $p\mid|G|$, teorema lui Cauchy (Teorema~\ref{teo:cauchy}) dă un element de ordin $p$, deci $p\mid\exp(G)$ prin (i).
(iv) Enumerăm $G=\{x_1,\dots,x_k\}$ și aplicăm repetat Corolarul~\ref{cor:lcm}: obținem succesiv elemente de ordin $[\ord x_1,\ord x_2]$, apoi $[\ord x_1,\ord x_2,\ord x_3]$ etc., iar la final un element de ordin $L=\exp(G)$.
\end{proof}

\begin{corolar}[criteriu de ciclicitate]\label{cor:expciclic}
Un grup abelian finit $G$ este ciclic dacă și numai dacă $\exp(G)=|G|$.
\end{corolar}

\begin{proof}
Dacă $G=\langle a\rangle$, atunci $\exp(G)\ge\ord(a)=|G|$, iar $\exp(G)\mid|G|$ dă egalitatea. Reciproc, prin Propoziția~\ref{prop:exponent}(iv) există un element de ordin $\exp(G)=|G|$, care generează $G$.
\end{proof}

\begin{observatie}
Ipoteza de comutativitate la (iv) și la corolar este esențială: $\exp(S_3)=[1,2,3]=6=|S_3|$, dar $S_3$ nu are niciun element de ordin $6$ (nu este ciclic). Invers, chiar și în cazul abelian exponentul poate fi strict mai mic decât ordinul: grupul lui Klein are $\exp=2$ și $4$ elemente. Punctul (iv) este exact ideea din a doua demonstrație a Teoremei~\ref{teo:ciclicitate}: într-un subgrup finit al lui $(K^*,\cdot)$, ordinul maxim coincide cu exponentul.
\end{observatie}

\begin{definitie}[subgrupul generat de o mulțime]
Fie $(G,\cdot)$ un grup și $X\subseteq G$. \emph{Subgrupul generat de $X$} este
\[
\langle X\rangle=\bigcap\{H\mid X\subseteq H,\ H\le G\},
\]
intersecția tuturor subgrupurilor lui $G$ care conțin $X$.
\end{definitie}

\begin{propozitie}[descrierea subgrupului generat]
$\langle X\rangle$ este cel mai mic subgrup al lui $G$ (față de incluziune) care conține $X$, iar pentru $X\neq\varnothing$
\[
\langle X\rangle=\{x_1^{k_1}x_2^{k_2}\cdots x_n^{k_n}\mid n\in\NN^*,\ x_i\in X,\ k_i\in\ZZ\},
\]
mulțimea tuturor produselor finite de puteri întregi ale elementelor lui $X$. În particular $\langle x\rangle=\{x^k\mid k\in\ZZ\}$.
\end{propozitie}
\begin{proof}
Intersecția unei familii de subgrupuri este subgrup (conține $e$; e închisă la $xy^{-1}$), deci $\langle X\rangle\le G$, conține $X$ și este inclus în orice subgrup ce conține $X$ --- adică e cel mai mic. Fie $M$ mulțimea produselor finite din enunț. Orice subgrup ce conține $X$ conține și $M$ (închidere la produse și inverse), deci $M\subseteq\langle X\rangle$. Reciproc, $M$ este el însuși subgrup (produsul a două produse finite e produs finit; inversul lui $x_1^{k_1}\cdots x_n^{k_n}$ este $x_n^{-k_n}\cdots x_1^{-k_1}$) și conține $X$, deci $\langle X\rangle\subseteq M$.
\end{proof}

\section{Subgrupuri remarcabile}

Grupăm aici subgrupurile atașate canonic unui element sau unei submulțimi --- centralizator, centru, normalizator, subgrup normal --- împreună cu proprietățile lor de bază, folosite constant în restul fișei.

\begin{definitie}[centralizator; centru]
Pentru $a\in G$, \emph{centralizatorul} lui $a$ este
\[
C_G(a)=\{x\in G\mid xa=ax\};
\]
pentru o submulțime $M\subseteq G$, \emph{centralizatorul} lui $M$ este
$C_G(M)=\{x\in G\mid xm=mx,\ \forall m\in M\}=\bigcap_{m\in M}C_G(m)$.
\emph{Centrul} grupului este $Z(G)=C_G(G)=\{a\in G\mid ax=xa,\ \forall x\in G\}$.
\end{definitie}

\begin{propozitie}\label{prop:centralizator}
\emph{(i)} $C_G(a)\le G$ pentru orice $a$; în consecință $C_G(M)\le G$ și $Z(G)=\bigcap_{a\in G}C_G(a)\le G$.
\emph{(ii)} $\langle a\rangle\le C_G(a)$, iar $a\in Z(G)\iff C_G(a)=G$.
\emph{(iii)} $G$ este abelian $\iff Z(G)=G$.
\end{propozitie}

\begin{proof}
(i) $e\in C_G(a)$. Fie $x,y\in C_G(a)$. Din $ya=ay$, înmulțind la stânga și la dreapta cu $y^{-1}$, obținem $ay^{-1}=y^{-1}a$, deci
\[
(xy^{-1})a=x(y^{-1}a)=x(ay^{-1})=(xa)y^{-1}=(ax)y^{-1}=a(xy^{-1}),
\]
adică $xy^{-1}\in C_G(a)$; aplicăm criteriul. Restul rezultă din Propoziția~\ref{prop:intersectie}. (ii) și (iii): direct din definiții ($a$ comută cu propriile puteri).
\end{proof}

\begin{definitie}[normalizator]
Pentru $x\in G$ și $M\subseteq G$ notăm $xMx^{-1}=\{xmx^{-1}\mid m\in M\}$ (\emph{conjugata} lui $M$). \emph{Normalizatorul} lui $M$ în $G$ este
\[
N_G(M)=\{x\in G\mid xMx^{-1}=M\},
\]
echivalent, mulțimea acelor $x$ cu $xM=Mx$ (egalitate de mulțimi).
\end{definitie}

\begin{propozitie}\label{prop:normalizator}
\emph{(i)} $N_G(M)\le G$ pentru orice submulțime $M$.
\emph{(ii)} Dacă $H\le G$, atunci $H\le N_G(H)$ și $C_G(H)\le N_G(H)$.
\end{propozitie}

\begin{proof}
(i) $e\in N_G(M)$. Dacă $xMx^{-1}=M$, aplicând $x^{-1}\,(\cdot)\,x$ ambilor membri obținem $M=x^{-1}Mx$, deci $x^{-1}\in N_G(M)$. Pentru închidere:
\[
(xy)M(xy)^{-1}=x\,(yMy^{-1})\,x^{-1}=xMx^{-1}=M.
\]
(ii) Pentru $h\in H$: $hHh^{-1}\subseteq H$ (închiderea lui $H$); aplicând același lucru lui $h^{-1}$ rezultă $h^{-1}Hh\subseteq H$, echivalent $H\subseteq hHh^{-1}$, deci egalitate și $h\in N_G(H)$. Dacă $x\in C_G(H)$, atunci $xhx^{-1}=h$ pentru orice $h\in H$, deci $xHx^{-1}=H$.
\end{proof}

\begin{observatie}
Dacă $M$ este \textbf{finită}, incluziunea $xMx^{-1}\subseteq M$ implică deja egalitatea (funcția $m\mapsto xmx^{-1}$ este injectivă de la $M$ în $M$, deci surjectivă); pentru mulțimi infinite o singură incluziune nu este suficientă.
\end{observatie}

\begin{definitie}[subgrup normal]\label{def:normal}
Un subgrup $H\le G$ se numește \emph{normal} (notăm $H\trianglelefteq G$) dacă $xHx^{-1}=H$ pentru orice $x\in G$; echivalent, $N_G(H)=G$; echivalent, $xH=Hx$ pentru orice $x\in G$. Astfel, $N_G(H)$ este cel mai mare subgrup al lui $G$ în care $H$ este normal.
\end{definitie}

\begin{propozitie}[criteriul de normalitate]\label{prop:normalcrit}
$H\trianglelefteq G$ dacă și numai dacă $xHx^{-1}\subseteq H$ pentru \textbf{orice} $x\in G$ (este suficientă o singură incluziune, cerută însă pentru toți $x$).
\end{propozitie}

\begin{proof}
Aplicând incluziunea și lui $x^{-1}$ obținem $x^{-1}Hx\subseteq H$, echivalent $H\subseteq xHx^{-1}$; deci egalitate.
\end{proof}

\begin{exemplu}[subgrupuri normale și nenormale]\label{ex:normal}
\begin{enumerate}[label=(\alph*),nosep]
\item Într-un grup abelian orice subgrup este normal ($xhx^{-1}=h$).
\item $Z(G)\trianglelefteq G$ și, mai general, orice subgrup al centrului este normal.
\item Orice subgrup de \textbf{indice 2} este normal: pentru $x\notin H$, clasa la stânga $xH$ și clasa la dreapta $Hx$ sunt amândouă egale cu $G\setminus H$ (clasele la dreapta partiționează $G$ din motive simetrice cu Lema~\ref{lema:clase}), deci $xH=Hx$; pentru $x\in H$ egalitatea este evidentă.
\item În $S_3$, subgrupul $H=\{e,(1\,2)\}$ \textbf{nu} este normal: $(1\,3)(1\,2)(1\,3)^{-1}=(2\,3)\notin H$.
\item Nucleul oricărui morfism de grupuri este subgrup normal (Teorema~\ref{teo:nucleu}).
\end{enumerate}
\end{exemplu}

\begin{lema}[Bézout pentru exponenți și subgrupuri]\label{lema:bezoutexp}
Fie $H\le G$, $x\in G$ și $n,p\in\NN^*$. Dacă $x^n\in H$ și $x^p\in H$, atunci $x^{(n,p)}\in H$. În particular, dacă $(n,p)=1$, atunci $x\in H$.
\end{lema}

\begin{proof}
Fie $d=(n,p)$; relația lui Bézout dă $u,v\in\ZZ$ cu $un+vp=d$. Atunci
\[
x^{d}=x^{un+vp}=(x^n)^u\,(x^p)^v\in H,
\]
deoarece $H$ este închis la puteri întregi (inclusiv negative) și la produse.
\end{proof}

\begin{observatie}
Lema se aplică de obicei cu $H=Z(G)$ (sau cu un centralizator): dacă $x^n\in Z(G)$ și $x^p\in Z(G)$, atunci $x^{(n,p)}\in Z(G)$; în particular, pentru $(n,p)=1$ rezultă $x\in Z(G)$. O aplicație tipică este Corolarul~\ref{cor:nn1}.
\end{observatie}

\begin{exemplu}
În $(\ZZ_n,+)$, ordinul lui $\cl{k}$ este $\dfrac{n}{(n,k)}$ (versiunea aditivă a Propoziției~\ref{prop:ordputere}, cu $\cl k=k\cdot\cl 1$).
\end{exemplu}

\section{Morfisme și izomorfisme de grupuri}\label{sec:morfisme}

\begin{definitie}
Fie $(G,\cdot)$ și $(G',*)$ grupuri. O funcție $f:G\to G'$ este \emph{morfism} dacă
$f(xy)=f(x)*f(y)$ pentru orice $x,y\in G$. Un morfism bijectiv este \emph{izomorfism} (scriem $G\simeq G'$). Un morfism $G\to G$ este \emph{endomorfism}; un izomorfism $G\to G$ este \emph{automorfism}.
\end{definitie}

\begin{propozitie}\label{prop:morfism}
Dacă $f:G\to G'$ este morfism, atunci $f(e)=e'$, $f(x^{-1})=f(x)^{-1}$ și $f(x^n)=f(x)^n$ pentru orice $n\in\ZZ$.
\end{propozitie}

\begin{proof}
$f(e)=f(e\cdot e)=f(e)*f(e)$ și simplificăm în $G'$. Apoi $f(x)*f(x^{-1})=f(xx^{-1})=e'$, deci $f(x^{-1})=f(x)^{-1}$ prin unicitatea inversului. Ultima afirmație: inducție pentru $n\ge0$, iar pentru $n<0$ se combină cu precedenta.
\end{proof}

\begin{propozitie}\label{prop:compunere}
Compunerea a două morfisme este morfism, iar inversa unui izomorfism este izomorfism (deci $\simeq$ este relație de echivalență).
\end{propozitie}

\begin{proof}
Prima parte, calcul direct. Pentru a doua: fie $u=f^{-1}(a)$, $v=f^{-1}(b)$; din $f(uv)=f(u)*f(v)=a*b$ rezultă $f^{-1}(a*b)=uv=f^{-1}(a)f^{-1}(b)$.
\end{proof}

\begin{definitie}
\emph{Nucleul} și \emph{imaginea} morfismului $f:G\to G'$ sunt
$\Ker f=\{x\in G\mid f(x)=e'\}$ și $\Ima f=f(G)$.
\end{definitie}

\begin{propozitie}[imagini directe și inverse]\label{prop:imagini}
Fie $f:G\to G'$ un morfism. \emph{(i)} Pentru orice subgrup $K\le G$, imaginea $f(K)=\{f(x)\mid x\in K\}$ este subgrup al lui $G'$; în particular $\Ima f=f(G)\le G'$. \emph{(ii)} Pentru orice subgrup $K'\le G'$, preimaginea $f^{-1}(K')=\{x\in G\mid f(x)\in K'\}$ este subgrup al lui $G$.
\end{propozitie}

\begin{proof}
(i) $e'=f(e)\in f(K)$; pentru $x,y\in K$, $f(x)f(y)^{-1}=f(xy^{-1})\in f(K)$ (Propoziția~\ref{prop:morfism}), și aplicăm criteriul de subgrup. (ii) Analog: $e\in f^{-1}(K')$, iar din $f(x),f(y)\in K'$ rezultă $f(xy^{-1})=f(x)f(y)^{-1}\in K'$.
\end{proof}

\begin{teorema}[nucleul]\label{teo:nucleu}
Fie $f:G\to G'$ un morfism. Atunci:
\emph{(i)} $\Ker f$ este subgrup \textbf{normal} al lui $G$;\quad
\emph{(ii)} $f$ este injectiv $\iff\Ker f=\{e\}$.
\end{teorema}

\begin{proof}
(i) $\Ker f=f^{-1}(\{e'\})$ este subgrup (Propoziția~\ref{prop:imagini}(ii)). Normalitatea: pentru $x\in G$ și $k\in\Ker f$,
\[
f(xkx^{-1})=f(x)\,f(k)\,f(x)^{-1}=f(x)\,e'\,f(x)^{-1}=e',
\]
deci $x(\Ker f)x^{-1}\subseteq\Ker f$ pentru orice $x$, iar Propoziția~\ref{prop:normalcrit} încheie.
(ii) Dacă $f$ este injectiv și $f(x)=e'=f(e)$, atunci $x=e$. Reciproc, dacă $\Ker f=\{e\}$ și $f(x)=f(y)$, atunci $f(xy^{-1})=e'$, deci $xy^{-1}=e$, adică $x=y$.
\end{proof}

\begin{propozitie}[morfismele și ordinul]\label{prop:morford}
Dacă $f:G\to G'$ este morfism și $x$ are ordin finit, atunci $\ord f(x)\mid\ord x$. Dacă $f$ este injectiv, $\ord f(x)=\ord x$ (inclusiv în cazul infinit).
\end{propozitie}

\begin{proof}
$f(x)^{\ord x}=f(x^{\ord x})=e'$, deci $\ord f(x)\mid \ord x$ (Teorema~\ref{teo:ordin}(i)). Dacă $f$ este injectiv și $f(x)^m=e'=f(e)$, atunci $x^m=e$, deci $\ord x\mid m$; luând $m=\ord f(x)$ obținem egalitatea. În cazul infinit, dacă $f(x)$ ar avea ordin finit $m$, ar rezulta $x^m=e$, contradicție.
\end{proof}

\begin{observatie}[invarianți la izomorfism]\label{obs:invarianti}
Un izomorfism păstrează: cardinalul, comutativitatea, ciclicitatea, ordinul fiecărui element (Propoziția~\ref{prop:morford}), numărul soluțiilor ecuației $x^k=e$ și numărul elementelor de un ordin dat. \textbf{Metoda standard} de a arăta că $G\not\simeq G'$: găsiți un astfel de invariant care diferă.
\end{observatie}

\begin{tehbox}{De ce ,,izomorf`` înseamnă ,,aceeași structură``}
Un izomorfism $f:G\to G'$ este, în esență, o \emph{redenumire} a elementelor: fiecare $x\in G$ primește numele nou $f(x)\in G'$, iar condiția de morfism $f(xy)=f(x)\,f(y)$ spune că tabla operației se copiază întocmai --- compui întâi și traduci, sau traduci întâi și compui: obții același lucru. De aceea \textbf{orice proprietate formulată exclusiv în limbajul operației de grup se transportă prin $f$}; orice teoremă despre operația lui $G$ devine, cuvânt cu cuvânt, o teoremă despre $G'$. Concret, dicționarul:
\begin{enumerate}[leftmargin=*, itemsep=1pt]
\item \textbf{Elementul neutru merge în elementul neutru.} Din $f(e)=f(e\cdot e)=f(e)f(e)$, simplificând în $G'$, rezultă $f(e_G)=e_{G'}$. ,,Rolul`` de neutru nu ține de numele elementului, ci de comportarea lui la operație --- iar comportarea se traduce.
\item \textbf{Inversele rămân perechi.} $f(x)\,f(x^{-1})=f(xx^{-1})=e_{G'}$, deci $f(x^{-1})=f(x)^{-1}$: perechile element--invers se corespund exact.
\item \textbf{Puterile și ordinele se păstrează.} Prin inducție $f(x^n)=f(x)^n$ pentru orice $n\in\ZZ$; cum $f$ e injectivă, $x^n=e\iff f(x)^n=e'$, deci $\ord(f(x))=\ord(x)$ (Propoziția~\ref{prop:morford}). Fiind și bijecție, $G$ și $G'$ au \emph{exact același număr de elemente de fiecare ordin} --- același ,,spectru`` de ordine.
\item \textbf{Ecuațiile au același număr de soluții.} $x$ verifică $x^k=a$ $\iff$ $f(x)$ verifică $y^k=f(a)$: bijecția $x\mapsto f(x)$ duce mulțimea soluțiilor unei ecuații exact în mulțimea soluțiilor ecuației traduse. În particular $\#\{x\in G:x^k=e\}=\#\{y\in G':y^k=e'\}$ pentru orice $k$.
\item \textbf{Subgrupurile se corespund unu-la-unu.} Imaginea unui subgrup e subgrup (Propoziția~\ref{prop:imagini}), iar $H\mapsto f(H)$ este o \emph{bijecție} între subgrupurile lui $G$ și cele ale lui $G'$, cu inversa $K\mapsto f^{-1}(K)$. Ea păstrează: incluziunile, cardinalul ($|f(H)|=|H|$), indicele ($f$ duce clasele $xH$ în clasele $f(x)f(H)$), \emph{normalitatea} ($H\trianglelefteq G\iff f(H)\trianglelefteq G'$, căci $f(gxg^{-1})=f(g)f(x)f(g)^{-1}$) și ciclicitatea, cu $f(\langle x\rangle)=\langle f(x)\rangle$. Întreaga ,,hartă`` a subgrupurilor este deci aceeași.
\item \textbf{Obiectele canonice se corespund.} $f(Z(G))=Z(G')$ și $f\big(C_G(a)\big)=C_{G'}\big(f(a)\big)$ (din $ax=xa\iff f(a)f(x)=f(x)f(a)$, surjectivitatea acoperind tot $G'$); conjugarea se transportă prin $f(gxg^{-1})=f(g)\,f(x)\,f(g)^{-1}$, deci clasele de conjugare merg în clase de conjugare de aceleași cardinale; comutatorii merg în comutatori.
\item \textbf{Proprietățile globale se moștenesc.} $G$ abelian $\iff G'$ abelian; $G$ ciclic $\iff G'$ ciclic, generatorii corespunzându-se ($G=\langle x\rangle\iff G'=\langle f(x)\rangle$) --- în particular cele două grupuri au același număr de generatori.
\item \textbf{Pentru grupuri finite: aceeași tablă.} Scriind tabla lui Cayley a lui $G$ și înlocuind fiecare element prin numele lui nou $f(x)$, obținem exact tabla lui $G'$: două grupuri finite sunt izomorfe dacă și numai dacă tablele lor coincid după o renumerotare a elementelor.
\end{enumerate}
\textbf{Ce NU vede izomorfismul:} natura concretă a elementelor (numere, matrice, permutări, clase de resturi) și notația operației (aditivă sau multiplicativă). $(\RR,+)$ și $((0,\infty),\cdot)$ sunt ,,același grup``, deși elementele lor arată complet diferit --- logaritmul e doar dicționarul dintre ele. \textbf{Atenție:} o simplă \emph{bijecție} nu transportă structura --- $\ZZ_4$ și grupul lui Klein au ambele $4$ elemente, dar nu sunt izomorfe ($\ZZ_4$ are un element de ordin $4$, Klein nu); toată forța stă în condiția de morfism. \textbf{Morala de concurs:} pentru $G\simeq G'$ construiești explicit dicționarul $f$ și verifici morfism $+$ bijecție (sau transport de structură); pentru $G\not\simeq G'$ cauți în lista de mai sus un invariant care diferă.
\end{tehbox}

\begin{exemplu}[izomorfisme și non-izomorfisme]\label{ex:noniso}
\begin{enumerate}[label=(\alph*),nosep]
\item $\exp:(\RR,+)\to((0,\infty),\cdot)$ este izomorfism: $e^{x+y}=e^xe^y$, bijectivă cu inversa $\ln$.
\item Pentru $G=\RR\setminus\{1\}$ cu $x\circ y=x+y-xy$ (Exemplul~\ref{ex:grupuri}(g)), funcția $f:G\to\RR^*$, $f(x)=1-x$, este izomorfism: $f(x\circ y)=1-x-y+xy=(1-x)(1-y)=f(x)f(y)$, bijectivă. (\emph{Transport de structură} --- tipar frecvent în manuale și la olimpiadă.)
\item $(\RR^*,\cdot)\not\simeq(\CC^*,\cdot)$: ecuația $x^4=1$ are $2$ soluții în $\RR^*$ și $4$ în $\CC^*$.
\item $(\RR,+)\not\simeq(\RR^*,\cdot)$: în $\RR^*$, elementul $-1$ are ordinul $2$, dar în $(\RR,+)$ singurul element de ordin finit este $0$.
\item $(\QQ,+)\not\simeq(\ZZ,+)$: $\ZZ$ este ciclic, $\QQ$ nu. Într-adevăr, dacă $\QQ=\langle a/b\rangle$ cu $a\neq0$, atunci $\frac{a}{2b}=k\cdot\frac ab$ ar da $k=\frac12\notin\ZZ$.
\end{enumerate}
\end{exemplu}

\begin{propozitie}[transport de structură]\label{prop:transport}
Fie $(G,\cdot)$ un grup, $M$ o mulțime nevidă și $\varphi:M\to G$ o funcție \textbf{bijectivă}. Definim pe $M$ legea
\[
x*y=\varphi^{-1}\bigl(\varphi(x)\cdot\varphi(y)\bigr).
\]
Atunci $(M,*)$ este grup, iar $\varphi:(M,*)\to(G,\cdot)$ este izomorfism. Elementul neutru este $\varphi^{-1}(e)$, simetricul lui $x$ este $\varphi^{-1}\bigl(\varphi(x)^{-1}\bigr)$, iar dacă $G$ este abelian, $(M,*)$ este abelian.
\end{propozitie}

\begin{proof}
Prin construcție, $\varphi(x*y)=\varphi(x)\varphi(y)$. Asociativitatea:
\[
\varphi\bigl(x*(y*z)\bigr)=\varphi(x)\bigl(\varphi(y)\varphi(z)\bigr)=\bigl(\varphi(x)\varphi(y)\bigr)\varphi(z)=\varphi\bigl((x*y)*z\bigr),
\]
și $\varphi$ este injectivă. Neutrul: $x*\varphi^{-1}(e)=\varphi^{-1}\bigl(\varphi(x)\,e\bigr)=x$ și simetric la stânga; simetricul: $x*\varphi^{-1}\bigl(\varphi(x)^{-1}\bigr)=\varphi^{-1}(e)$. Așadar $(M,*)$ este grup, iar $\varphi$ este morfism bijectiv.
\end{proof}

\begin{exemplu}[legile ,,de manual``]\label{ex:transport}
Legile de pe $\RR$ de forma $x*y=xy-m(x+y)+m^2+m$ sunt transporturi prin funcții afine: pentru $\varphi(x)=x-m$ avem
\[
\begin{gathered}
\varphi(x)\varphi(y)=(x-m)(y-m)=xy-m(x+y)+m^2,\\
\text{deci}\qquad \varphi^{-1}\bigl(\varphi(x)\varphi(y)\bigr)=x*y.
\end{gathered}
\]
Prin urmare $\bigl(\RR\setminus\{m\},*\bigr)\simeq(\RR^*,\cdot)$, cu neutrul $\varphi^{-1}(1)=m+1$ și simetricul lui $x$ egal cu $m+\frac{1}{x-m}$; toate proprietățile (asociativitate, comutativitate, elemente simetrizabile) vin \emph{gratuit} din Propoziția~\ref{prop:transport}. Exemplul~\ref{ex:noniso}(b) este tot un transport, prin $\varphi(x)=1-x$. \textbf{Strategia de concurs}: dată o lege ,,inventată``, căutați $\varphi$ afină, $\varphi(x)=\alpha x+\beta$, cu $\varphi(x*y)=\varphi(x)\varphi(y)$ (identificare de coeficienți); analog pentru transport din $(\RR,+)$.
\end{exemplu}

\begin{observatie}[endomorfismele lui $(\ZZ_n,+)$ și $(\QQ,+)$]\label{obs:endo}
(a) Orice endomorfism $f$ al lui $(\ZZ_n,+)$ este de forma $f(\cl k)=\cl k\cdot\cl a$, unde $\cl a=f(\cl 1)$ (din aditivitate, $f(\cl k)=k\,f(\cl 1)$); $f$ este automorfism dacă și numai dacă $\cl a$ este inversabil în inelul $\ZZ_n$, adică $(a,n)=1$.
(b) Orice endomorfism $f$ al lui $(\QQ,+)$ este $f(x)=cx$ cu $c=f(1)$: din aditivitate $f(nx)=nf(x)$ pentru $n\in\ZZ$, apoi $q\,f\!\left(\frac pq\right)=f(p)=p\,f(1)$, deci $f\!\left(\frac pq\right)=\frac pq\,f(1)$. (Aceeași schemă rezolvă ecuația funcțională Cauchy pe $\QQ$.)
\end{observatie}

\begin{definitie}
Pentru grupurile $G$ și $H$ notăm $\Hom(G,H)$ mulțimea morfismelor de la $G$ la $H$, $\End(G)=\Hom(G,G)$ mulțimea \emph{endomorfismelor} și $\Aut(G)$ mulțimea \emph{automorfismelor} lui $G$.
\end{definitie}

\begin{propozitie}[{structura lui $\Hom(G,H)$}]\label{prop:homgrup}
Dacă $H$ este \textbf{abelian}, atunci $\Hom(G,H)$ este grup abelian față de \emph{produsul punctual}
\[
(f\cdot g)(x)=f(x)\,g(x)
\]
(atenție: nu compunerea!). Elementul neutru este morfismul trivial $x\mapsto e_H$, iar simetricul lui $f$ este $x\mapsto f(x)^{-1}$.
\end{propozitie}

\begin{proof}
Închiderea: $(f\cdot g)(xy)=f(x)f(y)\,g(x)g(y)=f(x)g(x)\,f(y)g(y)$, unde comutativitatea lui $H$ a permis mutarea lui $f(y)$ peste $g(x)$; deci $f\cdot g$ este morfism. Funcția $\bar f(x)=f(x)^{-1}$ este morfism: $\bar f(xy)=\bigl(f(x)f(y)\bigr)^{-1}=f(y)^{-1}f(x)^{-1}=f(x)^{-1}f(y)^{-1}$, din nou prin comutativitate. Asociativitatea, comutativitatea și proprietatea neutrului se verifică punctual, fiind moștenite din $H$.
\end{proof}

\begin{observatie}
Fără comutativitatea lui $H$, produsul punctual a două morfisme poate să nu mai fie morfism: pentru $f=g=1_{S_3}$, produsul punctual este $x\mapsto x^2$, care nu este endomorfism al lui $S_3$ (Propoziția~\ref{prop:patrat}, $S_3$ nefiind abelian).
\end{observatie}

\begin{propozitie}[{$\Hom(\ZZ,H)$}]\label{prop:homZ}
Pentru orice grup $H$, funcția $f\mapsto f(1)$ este o \textbf{bijecție} de la $\Hom(\ZZ,H)$ la $H$: orice morfism de la $(\ZZ,+)$ la $H$ este de forma $f_a(k)=a^k$, pentru un unic $a\in H$. Dacă $H$ este abelian, aceasta este izomorfism de grupuri:
\[
\Hom(\ZZ,H)\simeq H.
\]
\end{propozitie}

\begin{proof}
Orice morfism satisface $f(k)=f(k\cdot 1)=f(1)^k$ (Propoziția~\ref{prop:morfism}), deci este determinat de $a=f(1)$; reciproc, pentru orice $a\in H$, $f_a(k)=a^k$ este morfism, din $a^{k+l}=a^ka^l$. Așadar $f\mapsto f(1)$ este bijecție. Dacă $H$ este abelian, $(f\cdot g)(1)=f(1)g(1)$, deci bijecția respectă operațiile.
\end{proof}

\begin{teorema}[{$\Hom(\ZZ_m,\ZZ_n)$}]\label{teo:homZmZn}
Fie $d=(m,n)$. Atunci
\[
\Hom(\ZZ_m,\ZZ_n)\simeq\ZZ_{d}.
\]
Mai precis: morfismele sunt exact $f_{\cl a}(\cl x)=x\cdot\cl a$, unde $\cl a$ parcurge subgrupul
$K=\{\cl a\in\ZZ_n\mid m\cdot\cl a=\cl 0\}=\bigl\langle\cl{n/d}\bigr\rangle$, ciclic de ordin $d$, iar $f\mapsto f(\cl 1)$ este izomorfism de la $\Hom(\ZZ_m,\ZZ_n)$ (cu adunarea punctuală) la $K$.
\end{teorema}

\begin{proof}
\emph{Determinare}: orice morfism satisface $f(\cl x)=x\cdot f(\cl 1)$ (aditivitate). \emph{Buna definire}: valoarea $\cl a=f(\cl 1)$ trebuie să satisfacă $m\cl a=f(\cl m)=f(\cl 0)=\cl 0$; reciproc, dacă $m\cl a=\cl 0$, formula $f_{\cl a}(\cl x)=x\cl a$ este bine definită (pentru $x\equiv x'\pmod m$, diferența $(x-x')\cl a$ este multiplu de $m\cl a=\cl 0$) și este evident morfism. Așadar $f\mapsto f(\cl 1)$ este bijecție de la $\Hom(\ZZ_m,\ZZ_n)$ la $K$, și respectă adunarea punctuală: $(f+g)(\cl 1)=f(\cl 1)+g(\cl 1)$.

\emph{Structura lui $K$}: condiția $m\cl a=\cl 0$ este ecuația ,,$y^m=e$`` în notație aditivă, care în grupul ciclic $(\ZZ_n,+)$ are exact $(m,n)=d$ soluții (Propoziția~\ref{prop:solutii}). $K$ este subgrup (verificare directă sau Propoziția~\ref{prop:imagini}(ii)) al grupului ciclic $\ZZ_n$, deci ciclic (Teorema~\ref{teo:subgrupciclic}); având ordinul $d$, este unicul subgrup de ordin $d$, anume $\langle\cl{n/d}\rangle$, și $K\simeq\ZZ_d$ (Teorema~\ref{teo:clasif}).
\end{proof}

\begin{corolar}
$\Hom(\ZZ_m,\ZZ_n)$ conține doar morfismul nul $\iff(m,n)=1$. Pentru $m=n$ regăsim $\End(\ZZ_n)$ cu $n$ elemente, în acord cu Observația~\ref{obs:endo}(a).
\end{corolar}

\begin{propozitie}[{$\End(G)$ ca monoid; $\Aut(G)=U(\End(G))$}]\label{prop:endmonoid}
$(\End(G),\circ)$ este monoid (în general necomutativ), cu elementul neutru $1_G$, iar
\[
\Aut(G)=U\bigl(\End(G)\bigr):
\]
automorfismele sunt exact elementele inversabile ale acestui monoid. În particular, $(\Aut(G),\circ)$ este grup.
\end{propozitie}

\begin{proof}
Compunerea a două endomorfisme este endomorfism (Propoziția~\ref{prop:compunere}), este asociativă și are neutrul $1_G$: monoid. Dacă $f\in\Aut(G)$, inversa $f^{-1}$ este morfism (tot Propoziția~\ref{prop:compunere}), deci $f^{-1}\in\End(G)$ și $f$ este inversabil în monoid. Reciproc, dacă $f\circ g=g\circ f=1_G$, atunci $f$ este bijectivă, deci automorfism. În fine, elementele inversabile ale oricărui monoid formează grup față de operația indusă --- aceeași verificare ca la grupul unităților unui inel: $(f\circ g)^{-1}=g^{-1}\circ f^{-1}$.
\end{proof}

\begin{definitie}[automorfisme interioare]
Pentru $g\in G$, funcția $i_g:G\to G$, $i_g(x)=gxg^{-1}$, se numește \emph{automorfismul interior} asociat lui $g$. Notăm $\Inn(G)=\{i_g\mid g\in G\}$.
\end{definitie}

\begin{teorema}[{$\Inn(G)\trianglelefteq\Aut(G)$}]\label{teo:inn}
\emph{(i)} Fiecare $i_g$ este automorfism al lui $G$, iar $\Phi:G\to\Aut(G)$, $\Phi(g)=i_g$, este morfism de grupuri cu
\[
\Ima\Phi=\Inn(G)\qquad\text{și}\qquad\Ker\Phi=Z(G);
\]
în particular $\Inn(G)\le\Aut(G)$ și $i_g=i_h\iff h^{-1}g\in Z(G)$.
\emph{(ii)} Pentru orice $f\in\Aut(G)$ și orice $g\in G$,
\[
f\circ i_g\circ f^{-1}=i_{f(g)}.
\]
În consecință, $\Inn(G)$ este subgrup \textbf{normal} al lui $\Aut(G)$.
\end{teorema}

\begin{proof}
(i) $i_g(xy)=gxyg^{-1}=(gxg^{-1})(gyg^{-1})=i_g(x)i_g(y)$, deci $i_g$ este endomorfism; calcul direct: $(i_g\circ i_h)(x)=g(hxh^{-1})g^{-1}=(gh)x(gh)^{-1}$, adică $i_g\circ i_h=i_{gh}$, deci $\Phi$ este morfism, iar $i_g\circ i_{g^{-1}}=i_e=1_G$ arată că $i_g$ este bijectiv (automorfism). Imaginea unui morfism este subgrup (Propoziția~\ref{prop:imagini}), deci $\Inn(G)\le\Aut(G)$. Nucleul: $i_g=1_G\iff gxg^{-1}=x$ pentru orice $x\iff g\in Z(G)$; apoi $i_g=i_h\iff i_{h^{-1}g}=1_G\iff h^{-1}g\in Z(G)$.

(ii) Pentru orice $x\in G$:
\[
\bigl(f\circ i_g\circ f^{-1}\bigr)(x)=f\bigl(g\,f^{-1}(x)\,g^{-1}\bigr)=f(g)\,x\,f(g)^{-1}=i_{f(g)}(x).
\]
Deci $f\,\Inn(G)\,f^{-1}\subseteq\Inn(G)$ pentru orice $f\in\Aut(G)$, iar criteriul de normalitate (Propoziția~\ref{prop:normalcrit}) încheie.
\end{proof}

\begin{lema}[centru trivial]\label{lema:centrutrivial}
Dacă $Z(G)=\{e\}$, atunci $Z\bigl(\Aut(G)\bigr)=\{1_G\}$.
\end{lema}

\begin{proof}
Fie $f\in Z(\Aut(G))$. În particular, $f$ comută cu fiecare automorfism interior: $f\circ i_g=i_g\circ f$, adică $f\circ i_g\circ f^{-1}=i_g$. Dar din Teorema~\ref{teo:inn}(ii), $f\circ i_g\circ f^{-1}=i_{f(g)}$; deci $i_{f(g)}=i_g$, ceea ce înseamnă $g^{-1}f(g)\in Z(G)=\{e\}$, adică $f(g)=g$, pentru \textbf{orice} $g\in G$. Așadar $f=1_G$.
\end{proof}

\begin{observatie}
(a) Din $\Ker\Phi=Z(G)$ și Teorema~\ref{teo:nucleu}: dacă $Z(G)=\{e\}$, morfismul $g\mapsto i_g$ este injectiv, deci $G\simeq\Inn(G)\le\Aut(G)$ --- grupurile cu centru trivial se ,,scufundă`` în propriul grup de automorfisme. (b) Pentru grupuri abeliene $\Inn(G)=\{1_G\}$, iar automorfismele interesante sunt cele ,,exterioare``: de pildă $\Aut(\ZZ_n,+)\simeq\bigl(U(\ZZ_n),\cdot\bigr)$, prin $\cl a\mapsto f_{\cl a}$, cu $f_{\cl a}\circ f_{\cl b}=f_{\cl a\cl b}$ (Observația~\ref{obs:endo}(a)); iar din Propoziția~\ref{prop:homZ}, $\End(\ZZ,+)$ corespunde lui $\ZZ$ ($f_a(k)=ak$) și $\Aut(\ZZ,+)=\{\pm1_\ZZ\}\simeq\ZZ_2$.
\end{observatie}

\subsection*{Comutativitate prin aplicații de tip putere}

\begin{propozitie}\label{prop:patrat}
\emph{(i)} $f:G\to G$, $f(x)=x^2$, este endomorfism dacă și numai dacă $G$ este abelian. \emph{(ii)} La fel pentru $f(x)=x^{-1}$.
\end{propozitie}

\begin{proof}
(i) $f(xy)=f(x)f(y)\iff xyxy=x^2y^2\iff yx=xy$ (simplificări la stânga cu $x$ și la dreapta cu $y$). (ii) $(xy)^{-1}=x^{-1}y^{-1}\iff y^{-1}x^{-1}=x^{-1}y^{-1}$, adică inversele comută, echivalent cu comutativitatea.
\end{proof}

\begin{teorema}[exponenți consecutivi]\label{teo:consec}
Dacă există $n\in\NN^*$ astfel încât $(xy)^k=x^k y^k$ pentru $k\in\{n,n+1,n+2\}$ și pentru orice $x,y\in G$, atunci $G$ este abelian.
\end{teorema}

\begin{proof}
Din $(xy)^{n+1}=(xy)^n(xy)$ obținem $x^{n+1}y^{n+1}=x^ny^nxy$; înmulțind la stânga cu $x^{-n}$ și la dreapta cu $y^{-1}$ rezultă
\[
xy^n=y^nx\qquad\text{pentru orice }x,y\in G,
\]
adică $y^n\in Z(G)$ pentru orice $y$. Aplicând același raționament perechii $(n+1,n+2)$, obținem $y^{n+1}\in Z(G)$ pentru orice $y$. Atunci
\[
(xy)\,y^n=x\,y^{n+1}=y^{n+1}x=y\,(y^n x)=y\,(x y^n)=(yx)\,y^n,
\]
și simplificând la dreapta cu $y^n$ rezultă $xy=yx$.
\end{proof}

\begin{teorema}[endomorfismul $x\mapsto x^n$]\label{teo:endon}
Fie $n\ge2$ și presupunem că $f:G\to G$, $f(x)=x^n$, este endomorfism, adică $(xy)^n=x^ny^n$ pentru orice $x,y\in G$. Atunci:
\emph{(i)} $(yx)^{n-1}=x^{n-1}y^{n-1}$ pentru orice $x,y\in G$;
\emph{(ii)} dacă, în plus, $f$ este \textbf{surjectiv}, atunci $x^{n-1}\in Z(G)$ pentru orice $x\in G$.
\end{teorema}

\begin{proof}
(i) Scriind produsul alternat și grupând altfel factorii,
\[
(xy)^n=\underbrace{(xy)(xy)\cdots(xy)}_{n}=x\,\underbrace{(yx)(yx)\cdots(yx)}_{n-1}\,y=x\,(yx)^{n-1}y.
\]
Egalând cu $x^ny^n$ și simplificând la stânga cu $x$, la dreapta cu $y$: $(yx)^{n-1}=x^{n-1}y^{n-1}$.

(ii) Schimbând $x$ cu $y$ în (i): $(xy)^{n-1}=y^{n-1}x^{n-1}$. Atunci
\[
x^ny^n=(xy)^n=(xy)\,(xy)^{n-1}=xy\cdot y^{n-1}x^{n-1}=x\,y^n\,x^{n-1},
\]
și simplificând la stânga cu $x$ obținem
\[
x^{n-1}y^n=y^nx^{n-1}\qquad\text{pentru orice }x,y\in G:
\]
elementul $x^{n-1}$ comută cu orice putere a $n$-a. Dar $f$ fiind surjectiv, \emph{orice} $g\in G$ se scrie $g=y^n$, deci $x^{n-1}g=gx^{n-1}$ pentru toți $g$, adică $x^{n-1}\in Z(G)$.
\end{proof}

\begin{corolar}\label{cor:nn1}
Dacă $x\mapsto x^n$ și $x\mapsto x^{n+1}$ sunt ambele endomorfisme \textbf{surjective} ale lui $G$ (cu $n\ge2$), atunci $G$ este abelian.
\end{corolar}

\begin{proof}
Din Teorema~\ref{teo:endon}(ii), $x^{n-1}\in Z(G)$ și $x^{n}\in Z(G)$ pentru orice $x$. Cum $(n-1,n)=1$, Lema~\ref{lema:bezoutexp} aplicată subgrupului $H=Z(G)$ dă $x\in Z(G)$ pentru orice $x$, adică $Z(G)=G$.
\end{proof}

\begin{observatie}
Comparați cu Teorema~\ref{teo:consec}: acolo se cer \emph{trei} exponenți consecutivi, dar fără surjectivitate; aici doar doi, în schimbul surjectivității. Ambele scheme (plus identitatea (i), valabilă fără nicio ipoteză suplimentară) sunt punctul de plecare standard pentru problemele de tip ,,$(xy)^n=x^ny^n$``.
\end{observatie}

\section{Grupuri ciclice}

\begin{definitie}
$G$ este \emph{ciclic} dacă există $x\in G$ (numit \emph{generator}) cu $G=\langle x\rangle$. Exemple: $(\ZZ,+)=\langle 1\rangle$, $(\ZZ_n,+)=\langle\cl 1\rangle$, $U_n=\langle\varepsilon\rangle$. Orice grup ciclic este abelian.
\end{definitie}

\begin{teorema}[clasificarea grupurilor ciclice]\label{teo:clasif}
Orice grup ciclic infinit este izomorf cu $(\ZZ,+)$; orice grup ciclic de ordin $n$ este izomorf cu $(\ZZ_n,+)$.
\end{teorema}

\begin{proof}
Fie $G=\langle x\rangle$. Cazul infinit: $f:\ZZ\to G$, $f(k)=x^k$, este morfism ($x^{k+l}=x^kx^l$), surjectiv prin definiție și injectiv deoarece puterile sunt distincte. Cazul finit: $f:\ZZ_n\to G$, $f(\cl k)=x^k$, este \emph{bine definită} și injectivă datorită echivalenței $x^i=x^j\iff i\equiv j\pmod n$ (Teorema~\ref{teo:ordin}(ii)), surjectivă și morfism ca mai sus.
\end{proof}

\begin{teorema}[subgrupurile lui $\ZZ$]\label{teo:subgrupZ}
Subgrupurile lui $(\ZZ,+)$ sunt exact mulțimile $n\ZZ=\{nk\mid k\in\ZZ\}$, cu $n\in\NN$.
\end{teorema}

\begin{proof}
Fiecare $n\ZZ$ este subgrup (criteriul: $nk-nl=n(k-l)$). Reciproc, fie $H\le\ZZ$. Dacă $H=\{0\}$, luăm $n=0$. Altfel $H$ conține un element pozitiv (odată cu $x$ conține $-x$); fie $n$ cel mai mic element pozitiv din $H$. Evident $n\ZZ\subseteq H$. Pentru $h\in H$, împărțirea cu rest $h=nq+r$, $0\le r<n$, dă $r=h-nq\in H$, deci $r=0$ prin minimalitate. Așadar $H=n\ZZ$.
\end{proof}

\begin{teorema}[subgrupurile grupurilor ciclice]\label{teo:subgrupciclic}
Fie $G=\langle x\rangle$ ciclic de ordin $n$. Atunci: \emph{(i)} orice subgrup al lui $G$ este ciclic, de forma $\langle x^m\rangle$ cu $m\mid n$; \emph{(ii)} pentru fiecare divizor $d$ al lui $n$ există \textbf{exact un} subgrup de ordin $d$, anume $\langle x^{n/d}\rangle$.
\end{teorema}

\begin{proof}
(i) Fie $H\le G$, $H\neq\{e\}$, și $m=\min\{k\ge 1\mid x^k\in H\}$. Ca în demonstrația Teoremei~\ref{teo:subgrupZ} (împărțire cu rest la exponent), orice $x^k\in H$ are $m\mid k$; în particular, din $x^n=e\in H$ rezultă $m\mid n$ și $H=\langle x^m\rangle$. (ii) $\langle x^{n/d}\rangle$ are ordinul $\frac{n}{(n,\,n/d)}=d$ (Propoziția~\ref{prop:ordputere}). Unicitatea: un subgrup de ordin $d$ este $\langle x^m\rangle$ cu $m\mid n$, de ordin $\frac{n}{m}=d$, deci $m=\frac{n}{d}$ este determinat.
\end{proof}

\begin{propozitie}[generatorii]\label{prop:generatori}
Dacă $G=\langle x\rangle$ are ordinul $n$, atunci $x^k$ este generator dacă și numai dacă $(k,n)=1$. În particular, $G$ are exact $\varphi(n)$ generatori, unde $\varphi$ este indicatorul lui Euler.
\end{propozitie}

\begin{proof}
$x^k$ generează $G$ $\iff\ord(x^k)=n\iff\frac{n}{(n,k)}=n\iff(n,k)=1$.
\end{proof}

\begin{propozitie}[numărul soluțiilor lui $y^d=e$]\label{prop:solutii}
Într-un grup ciclic de ordin $n$, ecuația $y^d=e$ are exact $(d,n)$ soluții.
\end{propozitie}

\begin{proof}
Scriem $y=x^t$, $0\le t\le n-1$. Atunci $y^d=e\iff n\mid td\iff\frac{n}{(n,d)}\mid t$, ceea ce dă exact $(n,d)$ valori ale lui $t$ în intervalul considerat.
\end{proof}

\begin{propozitie}[grupuri fără subgrupuri proprii]\label{prop:farasub}
Dacă $G\neq\{e\}$ are ca subgrupuri doar $\{e\}$ și $G$, atunci $G$ este finit, de ordin prim.
\end{propozitie}

\begin{proof}
Fie $x\neq e$; atunci $\langle x\rangle=G$, deci $G$ este ciclic. Dacă $G$ ar fi infinit, $G\simeq\ZZ$ (Teorema~\ref{teo:clasif}) ar avea subgrupul propriu $2\ZZ$, contradicție. Deci $|G|=n$; dacă $n=ab$ cu $1<a<n$, subgrupul $\langle x^a\rangle$ are ordinul $b\notin\{1,n\}$, contradicție. Așadar $n$ este prim.
\end{proof}

\section{Grupul permutărilor: cicluri și signatură}\label{sec:permutari}

Fixăm $n\ge2$. În $S_n$ produsul $\sigma\tau=\sigma\circ\tau$ se calculează de la dreapta la stânga: întâi $\tau$, apoi $\sigma$.

\begin{definitie}[cicluri; transpoziții]\label{def:ciclu}
Fie $k\ge2$ și $a_1,\dots,a_k\in\{1,\dots,n\}$ distincte. \emph{Ciclul} $(a_1\,a_2\,\dots\,a_k)\in S_n$ este permutarea care duce $a_1\mapsto a_2\mapsto\dots\mapsto a_k\mapsto a_1$ și fixează celelalte elemente; $k$ este \emph{lungimea} ciclului, iar $\{a_1,\dots,a_k\}$ este \emph{suportul} lui. Ciclurile de lungime $2$ se numesc \emph{transpoziții}. Două cicluri sunt \emph{disjuncte} dacă suporturile lor sunt disjuncte.
\end{definitie}

\begin{teorema}[descompunerea în cicluri disjuncte]\label{teo:cicluri}
Orice permutare $\sigma\in S_n$, $\sigma\neq e$, se scrie ca produs de cicluri disjuncte două câte două, în mod unic până la ordinea factorilor, iar ciclurile disjuncte comută între ele. Dacă lungimile acestor cicluri sunt $k_1,\dots,k_r$, atunci
\[
\ord(\sigma)=[k_1,\dots,k_r].
\]
\end{teorema}

\begin{proof}
Relația $a\sim b\iff b=\sigma^j(a)$ pentru un $j\in\ZZ$ este de echivalență pe $\{1,\dots,n\}$; clasele ei se numesc \emph{orbitele} lui $\sigma$. Fie $a$ fixat și $k\ge1$ cel mai mic cu $\sigma^k(a)=a$ (există, căci $\sigma$ are ordin finit). Elementele $a,\sigma(a),\dots,\sigma^{k-1}(a)$ sunt distincte (din $\sigma^i(a)=\sigma^j(a)$ cu $0\le i<j<k$ ar rezulta $\sigma^{j-i}(a)=a$), iar orice $\sigma^j(a)$ este unul dintre ele (împărțim $j$ cu rest la $k$); deci ele formează orbita lui $a$, pe care $\sigma$ acționează exact ca ciclul $(a\ \sigma(a)\ \dots\ \sigma^{k-1}(a))$. Produsul ciclurilor asociate orbitelor cu cel puțin două elemente coincide cu $\sigma$ pe fiecare orbită, deci peste tot. Unicitatea: în orice astfel de descompunere, suporturile ciclurilor sunt exact orbitele cu cel puțin două elemente, iar pe fiecare dintre ele ciclul este determinat de $\sigma$.

Dacă $\gamma,\delta$ sunt cicluri disjuncte, atunci pe suportul lui $\gamma$ ambele produse $\gamma\delta$ și $\delta\gamma$ acționează ca $\gamma$ (căci $\delta$ fixează acele elemente, iar $\gamma$ le trimite tot în suportul său), simetric pe suportul lui $\delta$, iar în rest ambele sunt identitatea; deci $\gamma\delta=\delta\gamma$.

Un ciclu $\gamma=(a_1\,\dots\,a_k)$ are ordinul $k$: $\gamma^j(a_1)=a_{1+j}$ (indici modulo $k$), deci $\gamma^j\neq e$ pentru $1\le j<k$, iar $\gamma^k=e$. Dacă $\sigma=\gamma_1\cdots\gamma_r$ cu cicluri disjuncte, acestea comută, deci $\sigma^m=\gamma_1^m\cdots\gamma_r^m$; fiecare $\gamma_i^m$ mișcă doar elemente din suportul lui $\gamma_i$, iar suporturile sunt disjuncte, deci $\sigma^m=e\iff\gamma_i^m=e$ pentru orice $i\iff k_i\mid m$ pentru orice $i\iff[k_1,\dots,k_r]\mid m$.
\end{proof}

\begin{definitie}[signatura]\label{def:sgn}
O \emph{inversiune} a permutării $\sigma\in S_n$ este o pereche $(i,j)$ cu $i<j$ și $\sigma(i)>\sigma(j)$. Notând $\mathrm{inv}(\sigma)$ numărul inversiunilor, \emph{signatura} (semnul) lui $\sigma$ este $\sgn(\sigma)=(-1)^{\mathrm{inv}(\sigma)}$. Permutarea $\sigma$ este \emph{pară} dacă $\sgn(\sigma)=1$ și \emph{impară} dacă $\sgn(\sigma)=-1$.
\end{definitie}

\begin{teorema}[proprietățile signaturii]\label{teo:sgn}
\emph{(i)} $\sgn(\sigma)=\displaystyle\prod_{1\le i<j\le n}\frac{\sigma(j)-\sigma(i)}{j-i}$.
\emph{(ii)} $\sgn(\sigma\tau)=\sgn(\sigma)\sgn(\tau)$; deci $\sgn:S_n\to(\{\pm1\},\cdot)$ este morfism de grupuri, iar nucleul său $A_n$ (\emph{grupul altern}, format din permutările pare) este subgrup normal al lui $S_n$.
\emph{(iii)} Orice transpoziție are signatura $-1$, iar un ciclu de lungime $k$ are signatura $(-1)^{k-1}$. În consecință, dacă descompunerea lui $\sigma$ în cicluri disjuncte are ciclurile de lungimi $k_1,\dots,k_r$, atunci $\sgn(\sigma)=(-1)^{(k_1-1)+\dots+(k_r-1)}$.
\end{teorema}

\begin{proof}
(i) Când $\{i,j\}$ parcurge perechile neordonate de elemente distincte, $\{\sigma(i),\sigma(j)\}$ le parcurge pe aceleași, bijectiv; deci produsul modulelor numărătorilor este egal cu produsul numitorilor, iar produsul are câte un factor negativ pentru fiecare inversiune.

(ii) Din (i),
\[
\sgn(\sigma\tau)=\prod_{i<j}\frac{\sigma(\tau(j))-\sigma(\tau(i))}{\tau(j)-\tau(i)}\cdot\prod_{i<j}\frac{\tau(j)-\tau(i)}{j-i}.
\]
Al doilea produs este $\sgn(\tau)$. În primul, fiecare factor depinde doar de perechea neordonată $\{\tau(i),\tau(j)\}$ (schimbând ordinea, numărătorul și numitorul își schimbă amândouă semnul), iar aceasta parcurge toate perechile $\{a,b\}$ cu $a<b$; deci primul produs este $\prod_{a<b}\frac{\sigma(b)-\sigma(a)}{b-a}=\sgn(\sigma)$. Nucleul unui morfism este subgrup normal (Teorema~\ref{teo:nucleu}).

(iii) Transpoziția $(1\,2)$ are o singură inversiune, deci signatura $-1$. O transpoziție oarecare $(a\,b)$ este egală cu $\pi(1\,2)\pi^{-1}$, pentru orice $\pi\in S_n$ cu $\pi(1)=a$, $\pi(2)=b$ (verificare directă pe elemente), deci, prin (ii), $\sgn(a\,b)=\sgn(\pi)\cdot(-1)\cdot\sgn(\pi)^{-1}=-1$. Pentru un ciclu, verificăm direct pe fiecare $a_i$ că
\[
(a_1\,a_2\,\dots\,a_k)=(a_1\,a_k)(a_1\,a_{k-1})\cdots(a_1\,a_2),
\]
un produs de $k-1$ transpoziții, deci signatura lui este $(-1)^{k-1}$. Ultima afirmație rezultă din (ii) și Teorema~\ref{teo:cicluri}.
\end{proof}

\begin{exemplu}
În $S_3$: $e$ și ciclurile de lungime $3$, $(1\,2\,3)$ și $(1\,3\,2)$, sunt pare (ordin $1$, respectiv $3$), iar cele trei transpoziții sunt impare (ordin $2$); așadar $A_3=\{e,(1\,2\,3),(1\,3\,2)\}$. În $S_4$, permutarea $(1\,2)(3\,4)$ are ordinul $[2,2]=2$ și este pară, iar $(1\,2\,3\,4)$ are ordinul $4$ și este impară.
\end{exemplu}

\section{Marile teoreme din teoria grupurilor}\label{sec:mari}

\emph{Notă}: Lagrange și Cauchy fac parte din nucleul programei; acțiunile de grupuri, teoremele lui Sylow și teorema lui Frobenius o depășesc, dar sunt reunite aici --- alături de Cayley --- ca \emph{marile teoreme cu nume} din teoria grupurilor.

\subsection*{Teorema lui Lagrange}

\begin{definitie}
Pentru $H\le G$ și $x\in G$, mulțimea $xH=\{xh\mid h\in H\}$ se numește \emph{clasă la stânga} a lui $x$ în raport cu $H$.
\end{definitie}

\begin{lema}\label{lema:clase}
\emph{(i)} $x\in xH$;\quad \emph{(ii)} $xH=yH\iff x^{-1}y\in H$;\quad \emph{(iii)} două clase la stânga sunt fie egale, fie disjuncte;\quad \emph{(iv)} $|xH|=|H|$ pentru orice $x$.
\end{lema}

\begin{proof}
(i) $x=xe$. (ii) Dacă $x^{-1}y=h\in H$, atunci $y=xh$, deci $yH=x(hH)=xH$. Reciproc, $y\in yH=xH$ dă $y=xh$, deci $x^{-1}y=h\in H$. (iii) Dacă $z\in xH\cap yH$, atunci $x^{-1}z,y^{-1}z\in H$, deci $x^{-1}y=(x^{-1}z)(y^{-1}z)^{-1}\in H$ și aplicăm (ii). (iv) Funcția $h\mapsto xh$ este o bijecție de la $H$ la $xH$ (injectivă prin simplificare, surjectivă prin definiție).
\end{proof}

\begin{teorema}[Lagrange]\label{teo:lagrange}
Dacă $G$ este un grup finit și $H\le G$, atunci $|H|$ divide $|G|$. Numărul claselor la stânga, notat $[G:H]$ (\emph{indicele} lui $H$), satisface $|G|=[G:H]\cdot|H|$.
\end{teorema}

\begin{proof}
Din Lema~\ref{lema:clase}, clasele la stânga formează o partiție a lui $G$ în mulțimi de câte $|H|$ elemente.
\end{proof}

\begin{corolar}\label{cor:lagrange}
Fie $G$ finit cu $|G|=n$. Atunci: \emph{(i)} $\ord(x)\mid n$ pentru orice $x\in G$; \emph{(ii)} $x^{n}=e$ pentru orice $x\in G$; \emph{(iii)} dacă $n$ este prim, $G$ este ciclic și orice element $x\neq e$ îl generează.
\end{corolar}

\begin{proof}
(i) $\ord(x)=|\langle x\rangle|$ divide $n$ (Teorema~\ref{teo:ordin}(iii) și Lagrange). (ii) $x^n=(x^{\ord x})^{n/\ord x}=e$. (iii) Pentru $x\neq e$, $\ord(x)\in\{1,n\}$ și $\ord(x)\neq 1$, deci $\langle x\rangle$ are $n$ elemente, adică este $G$.
\end{proof}

\begin{observatie}
\textbf{Reciproca teoremei lui Lagrange este falsă}: nu pentru orice divizor $d$ al lui $|G|$ există un subgrup de ordin $d$. Exemplul clasic: grupul $A_4$ al permutărilor pare din $S_4$ ($12$ elemente) nu are subgrup de ordin $6$. Într-adevăr, un subgrup de ordin $6$ ar avea indice $2$, deci ar fi normal și ar conține toate pătratele elementelor lui $A_4$; dar pătratul oricărui $3$-ciclu este tot un $3$-ciclu, iar $A_4$ are $8$ ciclii de lungime $3$ --- prea mulți pentru un subgrup cu $6$ elemente. Reciproca este însă adevărată pentru grupurile ciclice (Teorema~\ref{teo:subgrupciclic}).
\end{observatie}

\begin{propozitie}[subgrupurile finite ale lui $\CC^*$]\label{prop:subgrupC}
Orice subgrup finit cu $n$ elemente al grupului $(\CC^*,\cdot)$ este egal cu $U_n$.
\end{propozitie}

\begin{proof}
Fie $G\le\CC^*$ cu $|G|=n$. Din Corolarul~\ref{cor:lagrange}(ii), $z^n=1$ pentru orice $z\in G$, deci $G\subseteq U_n$. Cum $|G|=n=|U_n|$, avem egalitate.
\end{proof}

\begin{observatie}
Tehnica de \emph{împerechere} $x\leftrightarrow x^{-1}$ (dintr-o demonstrație clasică a faptului că un grup de ordin par conține un element de ordin $2$) este un instrument frecvent la ONM; ea reapare la Teorema~\ref{teo:produs} (produsul tuturor elementelor unui grup abelian finit).
\end{observatie}

\subsection*{Teorema lui Cauchy}

\begin{teorema}[Cauchy]\label{teo:cauchy}
Fie $G$ un grup finit și $p$ un număr prim cu $p\mid|G|$. Atunci $G$ conține un element de ordin $p$.
\end{teorema}

\begin{proof}[Demonstrație (McKay)]
Considerăm mulțimea $p$-uplurilor
\[
S=\bigl\{(x_1,x_2,\dots,x_p)\in G^p\ \big|\ x_1x_2\cdots x_p=e\bigr\}.
\]
Primele $p-1$ componente se aleg liber, iar ultima este determinată, $x_p=(x_1\cdots x_{p-1})^{-1}$; deci $|S|=|G|^{p-1}$, număr divizibil cu $p$ (căci $p\mid|G|$ și $p-1\ge1$).

Fie $\sigma:S\to S$ ,,rotația`` $\sigma(x_1,x_2,\dots,x_p)=(x_2,\dots,x_p,x_1)$. Ea este bine definită: dacă $x_1x_2\cdots x_p=e$, atunci
\[
x_2\cdots x_p\,x_1=x_1^{-1}(x_1x_2\cdots x_p)\,x_1=e
\]
(același truc de conjugare ca în Propoziția~\ref{prop:ordinv}(iii)); este evident bijectivă și $\sigma^p=\mathrm{id}$.

Grupăm elementele lui $S$ în \emph{orbite} $O_t=\{t,\sigma(t),\dots,\sigma^{p-1}(t)\}$; relația $s\in O_t$ este de echivalență (simetria: $s=\sigma^k(t)\Rightarrow t=\sigma^{p-k}(s)$), deci orbitele partiționează $S$. Susținem că fiecare orbită are exact $p$ elemente sau exact $1$ element, ultimul caz însemnând tupluri \emph{constante} $(x,x,\dots,x)$. Într-adevăr, dacă $\sigma^k(t)=t$ pentru un $1\le k\le p-1$, alegem $u,v\in\ZZ$ cu $uk+vp=1$ (posibil, căci $(k,p)=1$) și obținem
\[
\sigma(t)=\sigma^{uk+vp}(t)=(\sigma^k)^u(t)=t,
\]
deci toate componentele lui $t$ coincid.

Fie $N$ numărul tuplurilor constante din $S$, adică numărul elementelor $x\in G$ cu $x^p=e$. Numărând pe orbite:
\[
|G|^{p-1}=|S|=N+p\cdot(\text{numărul orbitelor cu }p\text{ elemente}),
\]
deci $p\mid N$. Cum $(e,\dots,e)\in S$, avem $N\ge1$, deci $N\ge p\ge2$: există $x\neq e$ cu $x^p=e$, iar $p$ fiind prim, $\ord(x)\in\{1,p\}$, adică $\ord(x)=p$.
\end{proof}

\begin{observatie}
Demonstrația este un exemplu-tipar de \emph{numărare cu simetrie de rotație}: aceeași mulțime se numără în două moduri, folosind că orbitele nebanale au exact $p$ elemente --- aici intervine esențial primalitatea lui $p$. Cazul $p=2$ dă din nou existența unui element de ordin $2$ într-un grup de ordin par. Împreună cu Lagrange, Cauchy dă echivalența: $G$ finit are un element de ordin $p$ prim $\iff p\mid|G|$.
\end{observatie}

\begin{corolar}[ridicarea la puterea $p$]\label{lema:puterep}
Fie $G$ un grup finit cu $n$ elemente, $p\ge 2$ un întreg, și $H=\{x^p\mid x\in G\}$ mulțimea puterilor de ordin $p$. Atunci
\[
(p,n)=1 \iff H=G.
\]
(Aici $H=G$ înseamnă că aplicația $x\mapsto x^p$ este surjectivă; pe un grup finit aceasta echivalează cu injectivitatea.)
\end{corolar}

\begin{proof}
$(\Rightarrow)$ Fie $(p,n)=1$; există $a,b\in\ZZ$ cu $ap+bn=1$ (Bézout). Pentru orice $g\in G$, ordinul lui $g$ divide $n$ (Lagrange), deci $g^{n}=e$ și astfel $g^{bn}=e$. Atunci
\[
g=g^{ap+bn}=\bigl(g^{a}\bigr)^{p}\,g^{bn}=\bigl(g^{a}\bigr)^{p}\in H,
\]
deci $H=G$.

$(\Leftarrow)$ Prin contrapoziție: presupunem $(p,n)=d>1$ și fie $q$ un divizor prim al lui $d$. Cum $q\mid n$, teorema lui Cauchy (Teorema~\ref{teo:cauchy}) dă un element $x$ de ordin $q$. Cum $q\mid p$, avem $x^{p}=e=e^{p}$; deci $x\neq e$ și $e$ au aceeași imagine, aplicația $x\mapsto x^{p}$ nu e injectivă, deci (pe mulțime finită) nici surjectivă: $H\neq G$.
\end{proof}

\begin{observatie}
Enunțul nu cere ca $G$ să fie abelian și nu afirmă că $H$ este subgrup --- în grupuri neabeliene, $x\mapsto x^p$ nu e neapărat morfism, iar $H$ poate să nu fie stabil. Când $(p,n)=1$, ridicarea la puterea $p$ este chiar o \emph{bijecție} a lui $G$ (permutarea inversă e $x\mapsto x^{a}$ cu $ap\equiv 1\pmod n$), fapt des folosit la ONM pentru a extrage rădăcini de ordin $p$ în grupuri finite.
\end{observatie}

\subsection*{Teorema lui Cayley}

\begin{teorema}[Cayley]\label{teo:cayley}
Orice grup $G$ este izomorf cu un subgrup al grupului $\bigl(S(G),\circ\bigr)$ al bijecțiilor mulțimii $G$. În particular, orice grup finit cu $n$ elemente este izomorf cu un subgrup al lui $S_n$.
\end{teorema}

\begin{proof}
Pentru $g\in G$, \emph{translația la stânga} $\lambda_g:G\to G$, $\lambda_g(x)=gx$, este bijectivă (ecuația $gx=b$ are soluție unică --- Propoziția~\ref{prop:simplificare}), deci $\lambda_g\in S(G)$. Funcția $\lambda:G\to S(G)$, $\lambda(g)=\lambda_g$, este morfism:
\[
\lambda_{gh}(x)=(gh)x=g(hx)=\bigl(\lambda_g\circ\lambda_h\bigr)(x),
\]
și este injectivă: $\lambda_g=1_G$ înseamnă $gx=x$ pentru toți $x$, deci $g=e$; prin urmare $\lambda_g=\lambda_h\Rightarrow g=h$. Imaginea $\lambda(G)=\{\lambda_g\mid g\in G\}$ este parte stabilă a lui $S(G)$ (căci $\lambda_g\circ\lambda_h=\lambda_{gh}$), conține $1_G=\lambda_e$ și inversele ($\lambda_g^{-1}=\lambda_{g^{-1}}$), deci este subgrup; iar $\lambda:G\to\lambda(G)$ este morfism bijectiv, deci izomorfism. Așadar $G\simeq\lambda(G)\le S(G)$. Pentru $|G|=n$, numerotarea elementelor lui $G$ dă $S(G)\simeq S_n$.
\end{proof}

\begin{observatie}
Teorema explică rolul special al grupurilor de permutări: \emph{orice} grup ,,trăiește`` într-un $S(X)$. Scufundarea nu este economică ($S_n$ are $n!$ elemente), dar este conceptual fundamentală.
\end{observatie}

\subsection*{Teoremele lui Sylow}

Fie $|G|=p^a m$ cu $p$ prim, $a\ge1$ și $p\nmid m$. Un subgrup de ordin $p^a$ (puterea maximă a lui $p$ permisă de Lagrange) se numește \emph{$p$-subgrup Sylow}; notăm $n_p$ numărul lor.

\begin{teorema}[Sylow]\label{teo:sylow}
Fie $G$ un grup finit cu $|G|=p^am$, $p\nmid m$. Atunci:
\begin{enumerate}[label=\emph{(\Roman*)},nosep,leftmargin=2.2em]
\item \emph{(existența)} $G$ are cel puțin un $p$-subgrup Sylow;
\item \emph{(conjugarea)} orice subgrup al lui $G$ de ordin o putere a lui $p$ este conținut într-un conjugat al unui $p$-subgrup Sylow; în particular, oricare două $p$-subgrupuri Sylow sunt conjugate;
\item \emph{(numărarea)} $n_p=[G:N_G(P)]$, $\;n_p\mid m\;$ și $\;n_p\equiv1\pmod p$.
\end{enumerate}
\end{teorema}

Demonstrația --- prin acțiuni de grupuri --- se află în secțiunea \emph{Teorie avansată} (pagina~\pageref{dem:sylow}). Reținem consecința folosită cel mai des:

\begin{corolar}\label{cor:sylownormal}
$n_p=1$ dacă și numai dacă $p$-subgrupul Sylow este \textbf{normal} în $G$.
\end{corolar}

\begin{proof}
Conjugatele unui $p$-subgrup Sylow $P$ epuizează mulțimea celor $n_p$ subgrupuri Sylow (Sylow~(II)), iar numărul lor este $[G:N_G(P)]$; $P$ este normal $\iff$ nu are alte conjugate $\iff n_p=1$ (vezi și Teorema~\ref{teo:conjsub}).
\end{proof}

\subsection*{Teorema lui Frobenius}

\begin{teorema}[Frobenius, 1895]\label{teo:frobgrup}
Fie $G$ un grup finit și $n$ un divizor al lui $|G|$. Atunci numărul soluțiilor ecuației $x^n=e$ în $G$ este un \textbf{multiplu de $n$}:
\[
n\ \Big|\ \#\{x\in G\mid x^n=e\}.
\]
\end{teorema}

Pentru grupuri ciclice regăsim numărul exact $(n,|G|)=n$ (Propoziția~\ref{prop:solutii}); teorema afirmă că în orice grup acest număr rămâne divizibil cu $n$. Cazul $n=p$ prim redă concluzia demonstrației McKay a teoremei lui Cauchy. Demonstrația (numărare cu acțiuni, lema lui Burnside și coliere aperiodice) se află în secțiunea \emph{Teorie avansată} (pagina~\pageref{dem:frobenius}).

\begin{corolar}[Frobenius, forma generală cu c.m.m.d.c.]\label{cor:frobgcd}
Fie $G$ un grup finit și $n\in\NN^*$ \textbf{arbitrar} (nu neapărat divizor al lui $|G|$). Atunci
\[
\bigl(n,\,|G|\bigr)\ \Big|\ \#\{x\in G\mid x^n=e\}.
\]
\end{corolar}

\begin{proof}
Fie $d=(n,|G|)$. Susținem că \emph{cele două ecuații au aceleași soluții}: $\{x\mid x^n=e\}=\{x\mid x^d=e\}$. Într-adevăr, dacă $x^d=e$, atunci $x^n=(x^d)^{n/d}=e$, căci $d\mid n$. Reciproc, dacă $x^n=e$, atunci $\ord(x)\mid n$; dar $\ord(x)\mid|G|$ (Lagrange), deci $\ord(x)$ divide $(n,|G|)=d$, adică $x^d=e$. Cum $d\mid|G|$, Teorema~\ref{teo:frobgrup} aplicată lui $d$ dă $d\mid\#\{x\mid x^d=e\}=\#\{x\mid x^n=e\}$.
\end{proof}

\begin{obs}
Cele două formulări sunt deci echivalente (pentru $n\mid|G|$ avem $d=n$ și corolarul redă teorema), iar ipoteza $n\mid|G|$ din enunțul clasic nu se poate pur și simplu omite: în $\ZZ_2$, ecuația $x^4=e$ are $2$ soluții --- nu e multiplu de $4$, dar e multiplu de $(4,2)=2$. Forma cu c.m.m.d.c.\ este cea ,,corectă`` pentru $n$ arbitrar.
\end{obs}

\section{Teorie avansată}\label{sec:avansat}

\emph{Notă}: întreaga secțiune \textbf{depășește programa clasei a XII-a}; este însă standard în pregătirea pentru baraje, loturi și olimpiade studențești și se sprijină exclusiv pe rezultatele deja demonstrate în fișă. Cele trei subsecțiuni tratează, pe rând, structurile factor, acțiunile de grupuri (cu demonstrațiile teoremelor lui Sylow și Frobenius) și clasificările structurale (grupuri abeliene, simple, rezolubile).

\subsection{Relații de echivalență, congruențe și grupuri factor}\label{sec:avansat1}

\begin{definitie}[relații binare; relații de echivalență]\label{def:echiv}
O \emph{relație binară} pe mulțimea nevidă $A$ este o submulțime $\rho\subseteq A\times A$; scriem $x\,\rho\,y$ pentru $(x,y)\in\rho$. Relația $\rho$ este: \emph{reflexivă} dacă $x\,\rho\,x$ pentru orice $x$; \emph{simetrică} dacă $x\,\rho\,y\Rightarrow y\,\rho\,x$; \emph{tranzitivă} dacă $x\,\rho\,y$ și $y\,\rho\,z\Rightarrow x\,\rho\,z$. O relație reflexivă, simetrică și tranzitivă se numește \emph{relație de echivalență}.
\end{definitie}

\begin{definitie}[clase, mulțimea cât, sistem de reprezentanți]
Fie $\rho$ o relație de echivalență pe $A$. \emph{Clasa de echivalență} a lui $x$ este $\widehat x=\{y\in A\mid x\,\rho\,y\}$. Mulțimea claselor,
\[
A/\rho=\{\widehat x\mid x\in A\},
\]
se numește \emph{mulțimea cât} (mulțimea factor) a lui $A$ prin $\rho$. Un \emph{sistem complet de reprezentanți} este o submulțime $S\subseteq A$ care conține \textbf{exact câte un} element din fiecare clasă.
\end{definitie}

\begin{exemplu}
Congruența modulo $n$ pe $\ZZ$: $x\,\rho\,y\iff n\mid x-y$ este relație de echivalență; clasele sunt $\cl0,\cl1,\dots,\cl{n-1}$, deci $\ZZ/\rho=\ZZ_n$, iar $S=\{0,1,\dots,n-1\}$ este un sistem complet de reprezentanți (la fel $\{1,2,\dots,n\}$ sau orice alegere de câte un element din fiecare clasă).
\end{exemplu}

\begin{definitie}[partiție]
O \emph{partiție} a lui $A$ este o familie de submulțimi nevide ale lui $A$, două câte două disjuncte, a căror reuniune este $A$.
\end{definitie}

\begin{teorema}[echivalențe $\leftrightarrow$ partiții]\label{teo:partitie}
\emph{(i)} Clasele unei relații de echivalență pe $A$ formează o partiție a lui $A$ (notată $A/\rho$). \emph{(ii)} Reciproc, orice partiție a lui $A$ provine dintr-o unică relație de echivalență: $x\,\rho\,y\iff x$ și $y$ aparțin aceluiași bloc al partiției.
\end{teorema}

\begin{proof}
(i) $x\in\widehat x$ (reflexivitate), deci clasele sunt nevide și acoperă $A$. Dacă $z\in\widehat x\cap\widehat y$, atunci $x\,\rho\,z$ și $y\,\rho\,z$, deci (simetrie și tranzitivitate) $x\,\rho\,y$; de aici $\widehat x=\widehat y$, prin dublă incluziune cu tranzitivitatea. Așadar două clase sunt egale sau disjuncte. (ii) Verificare directă a celor trei proprietăți; clasele sunt exact blocurile.
\end{proof}

\begin{definitie}[relațiile canonice pe un grup]
Fie $G$ un grup și $H\subseteq G$ o submulțime. Definim pe $G$ relațiile
\[
x\,\rho_H\,y\iff x^{-1}y\in H\qquad\text{și}\qquad x\,\rho'_H\,y\iff xy^{-1}\in H .
\]
\end{definitie}

\begin{propozitie}[relațiile induse de un subgrup]\label{prop:rhoH}
$\rho_H$ este relație de echivalență pe $G$ \textbf{dacă și numai dacă} $H$ este subgrup al lui $G$ (analog pentru $\rho'_H$). În acest caz, clasa lui $x$ față de $\rho_H$ este $xH$, iar față de $\rho'_H$ este $Hx$:
\[
G/\rho_H=\{xH\mid x\in G\},\qquad G/\rho'_H=\{Hx\mid x\in G\}.
\]
$\rho_H$ și $\rho'_H$ se numesc \emph{relațiile de echivalență induse de subgrupul $H$} (la stânga, respectiv la dreapta).
\end{propozitie}

\begin{proof}
Presupunem că $\rho_H$ este echivalență. Reflexivitatea dă $e=x^{-1}x\in H$. Pentru $h\in H$ avem $e\,\rho_H\,h$ (căci $e^{-1}h=h\in H$), deci prin simetrie $h\,\rho_H\,e$, adică $h^{-1}=h^{-1}e\in H$. Pentru $h,k\in H$: $e\,\rho_H\,h$ și $h\,\rho_H\,hk$ (căci $h^{-1}(hk)=k\in H$), deci prin tranzitivitate $e\,\rho_H\,hk$, adică $hk\in H$. Așadar $H\le G$.
Reciproc, pentru $H\le G$, verificările sunt cele din Lema~\ref{lema:clase}: reflexivitate din $e\in H$; simetrie din $y^{-1}x=(x^{-1}y)^{-1}\in H$; tranzitivitate din $x^{-1}z=(x^{-1}y)(y^{-1}z)\in H$. Clasa lui $x$: $\{y\mid x^{-1}y\in H\}=xH$. Pentru $\rho'_H$, totul este simetric, cu clasele $Hx$.
\end{proof}

\begin{observatie}
În acest limbaj, teorema lui Lagrange (Teorema~\ref{teo:lagrange}) spune că partiția $G/\rho_H$ are toate blocurile de cardinal $|H|$. Numărul claselor la stânga este egal cu numărul claselor la dreapta: $xH\mapsto Hx^{-1}$ este o bijecție bine definită între $G/\rho_H$ și $G/\rho'_H$ (din $xH=yH\iff x^{-1}y\in H\iff Hx^{-1}=Hy^{-1}$), deci indicele $[G:H]$ nu depinde de parte.
\end{observatie}

\begin{teorema}[normalitatea în limbajul claselor]\label{teo:normalrho}
Pentru $H\le G$ sunt echivalente: \emph{(i)} $H\trianglelefteq G$; \emph{(ii)} $\rho_H=\rho'_H$; \emph{(iii)} partițiile $G/\rho_H$ și $G/\rho'_H$ coincid.
\end{teorema}

\begin{proof}
(i)$\iff$(ii): relațiile coincid $\iff$ clasa fiecărui $x$ este aceeași în ambele, adică $xH=Hx$ pentru orice $x$ --- exact Definiția~\ref{def:normal}. (ii)$\Rightarrow$(iii) evident. (iii)$\Rightarrow$(ii): blocul care îl conține pe $x$ este unic în fiecare partiție, anume $xH$, respectiv $Hx$; dacă familiile de blocuri coincid, atunci $xH=Hx$ pentru fiecare $x$.
\end{proof}

\begin{definitie}[congruență]\label{def:congruenta}
O relație de echivalență $\rho$ pe grupul $G$ se numește \emph{congruență} dacă este \emph{compatibilă cu operația}:
\[
x\,\rho\,y\ \text{ și }\ x'\,\rho\,y'\ \Longrightarrow\ xx'\,\rho\,yy' .
\]
\end{definitie}

\begin{observatie}[compatibilitatea cu inversarea]\label{obs:invcompat}
Într-un \textbf{grup}, dintr-o congruență rezultă automat: $x\,\rho\,y\Rightarrow x^{-1}\,\rho\,y^{-1}$. Într-adevăr, din $x^{-1}\,\rho\,x^{-1}$ și $x\,\rho\,y$ obținem $e=x^{-1}x\,\rho\,x^{-1}y$; apoi din aceasta și $y^{-1}\,\rho\,y^{-1}$ obținem $y^{-1}\,\rho\,x^{-1}yy^{-1}=x^{-1}$, iar simetria încheie. (Într-un monoid, condiția de compatibilitate cu inversarea s-ar impune separat, acolo unde are sens.)
\end{observatie}

\begin{teorema}[congruențele provin din subgrupuri normale]\label{teo:congruenta}
Fie $G$ un grup și $\rho\subseteq G\times G$ o relație de echivalență. Atunci $\rho$ este \textbf{congruență} dacă și numai dacă există un subgrup normal $N\trianglelefteq G$ astfel încât $\rho=\rho_N$. În acest caz $N=\widehat e$ (clasa elementului neutru) și $\rho_N=\rho'_N$.
\end{teorema}

\begin{proof}
($\Leftarrow$) Fie $N\trianglelefteq G$; $\rho_N$ este echivalență (Propoziția~\ref{prop:rhoH}). Compatibilitatea: dacă $x^{-1}y=n\in N$ și $x'^{-1}y'=n'\in N$, atunci
\[
(xx')^{-1}(yy')=x'^{-1}(x^{-1}y)\,y'=\bigl(x'^{-1}n\,x'\bigr)\bigl(x'^{-1}y'\bigr)=\underbrace{x'^{-1}nx'}_{\in N}\cdot\,n'\in N,
\]
unde $x'^{-1}nx'\in N$ prin normalitate. Egalitatea $\rho_N=\rho'_N$: Teorema~\ref{teo:normalrho}.

($\Rightarrow$) Fie $\rho$ congruență și $N=\widehat e=\{x\in G\mid x\,\rho\,e\}$. \emph{$N$ este subgrup}: $e\in N$; pentru $x,y\in N$, compatibilitatea dă $xy\,\rho\,e\cdot e=e$; iar $x\,\rho\,e$ implică $x^{-1}\,\rho\,e^{-1}=e$ (Observația~\ref{obs:invcompat}). \emph{$N$ este normal}: pentru $x\in N$ și $g\in G$, din $g\,\rho\,g$, $x\,\rho\,e$ și $g^{-1}\,\rho\,g^{-1}$ rezultă, aplicând compatibilitatea de două ori, $gxg^{-1}\,\rho\,geg^{-1}=e$, deci $gxg^{-1}\in N$; Propoziția~\ref{prop:normalcrit} încheie. \emph{$\rho=\rho_N$}: din $x\,\rho\,y$, înmulțind cu $x^{-1}\,\rho\,x^{-1}$, obținem $e\,\rho\,x^{-1}y$, adică $x^{-1}y\in N$; reciproc, din $e\,\rho\,x^{-1}y$, înmulțind cu $x\,\rho\,x$, obținem $x\,\rho\,y$.
\end{proof}

\begin{propozitie}[congruența asociată unui morfism]\label{prop:congmorf}
Pentru orice morfism $f:G\to G'$, relația
\[
x\sim y\iff f(x)=f(y)
\]
este o congruență pe $G$ și coincide cu $\rho_{\Ker f}=\rho'_{\Ker f}$.
\end{propozitie}

\begin{proof}
$f(x)=f(y)\iff f(x^{-1}y)=e'\iff x^{-1}y\in\Ker f$ și, la fel, $\iff f(xy^{-1})=e'\iff xy^{-1}\in\Ker f$. Deci $\sim$ este simultan $\rho_{\Ker f}$ și $\rho'_{\Ker f}$; compatibilitatea cu operația este imediată ($f(xx')=f(x)f(x')=f(y)f(y')=f(yy')$) --- în acord cu Teorema~\ref{teo:congruenta}, întrucât $\Ker f\trianglelefteq G$ (Teorema~\ref{teo:nucleu}).
\end{proof}

\begin{teorema}[grupul factor / grupul cât]\label{teo:factor}
Fie $N\trianglelefteq G$. Pe mulțimea claselor
\[
G/N:=G/\rho_N=\{xN\mid x\in G\}
\]
operația $(xN)\cdot(yN)=(xy)N$ este \textbf{bine definită} (independentă de reprezentanți) și dă o structură de grup, numit \emph{grupul factor} (sinonim: \emph{grupul cât}) al lui $G$ prin $N$: neutrul este $N=eN$, iar $(xN)^{-1}=x^{-1}N$. Funcția
\[
\pi:G\to G/N,\qquad\pi(x)=xN\qquad(\text{\emph{proiecția canonică}})
\]
este morfism surjectiv cu $\Ker\pi=N$; pentru $G$ finit, $|G/N|=[G:N]=|G|/|N|$.
\textbf{Reciproc}: dacă pentru un subgrup $H\le G$ formula $(xH)(yH)=(xy)H$ este bine definită pe clasele la stânga, atunci $H\trianglelefteq G$.
\end{teorema}

\begin{proof}
\emph{Buna definire}: dacă $x'=xn$ și $y'=ym$ cu $n,m\in N$, atunci $x'y'=xnym=xy\,(y^{-1}ny)\,m\in xyN$, folosind $y^{-1}ny\in N$ (normalitate). Asociativitatea, neutrul și inversele se verifică pe reprezentanți, moștenite din $G$. $\pi$ este morfism chiar prin definiția operației, evident surjectiv, iar $\pi(x)=eN\iff xN=N\iff x\in N$.
\emph{Reciproca}: fie $x\in G$ și $h\in H$. Perechile de reprezentanți $(x,\,x^{-1})$ și $(xh,\,x^{-1})$ definesc aceleași clase, deci buna definire cere $(x\cdot x^{-1})H=(xh\cdot x^{-1})H$, adică $xhx^{-1}\in H$ pentru orice $x,h$; criteriul de normalitate (Propoziția~\ref{prop:normalcrit}) încheie.
\end{proof}

\begin{exemplu}
$\ZZ/n\ZZ$: pe $(\ZZ,+)$, subgrupul $n\ZZ$ este normal (grup abelian), clasele lui $\rho_{n\ZZ}$ sunt clasele de resturi, iar grupul factor este exact $(\ZZ_n,+)$. Construcția lui $\ZZ_n$ din clasa a X-a este deci cazul fondator al întregii teorii.
\end{exemplu}

\begin{teorema}[teorema fundamentală de izomorfism]\label{teo:fundamental}
Fie $f:G\to G'$ un morfism de grupuri. Atunci funcția
\[
\bar f:G/\Ker f\longrightarrow\Ima f,\qquad\bar f\bigl(x\,\Ker f\bigr)=f(x),
\]
este bine definită și este \textbf{izomorfism} de grupuri:
\[
G/\Ker f\ \simeq\ \Ima f .
\]
\textbf{Consecință}: dacă $f:G\to G'$ este morfism \textbf{surjectiv}, atunci $G/\Ker f\simeq G'$.
\end{teorema}

\begin{proof}
Notăm $K=\Ker f$. Buna definire și injectivitatea, dintr-o singură echivalență: $xK=yK\iff x^{-1}y\in K\iff f(x)=f(y)$ (Propoziția~\ref{prop:congmorf}). Morfism: $\bar f\bigl((xK)(yK)\bigr)=\bar f(xyK)=f(xy)=f(x)f(y)=\bar f(xK)\,\bar f(yK)$. Surjectivitatea pe $\Ima f$ este evidentă. Consecința: $\Ima f=G'$.
\end{proof}

\begin{exemplu}[aplicații ale teoremei fundamentale]\label{ex:fundamental}
\begin{enumerate}[label=(\alph*),nosep]
\item $f:\ZZ\to\ZZ_n$, $f(k)=\cl k$: surjectiv, $\Ker f=n\ZZ$, deci $\ZZ/n\ZZ\simeq\ZZ_n$.
\item $f:(\RR,+)\to(\CC^*,\cdot)$, $f(t)=\cos2\pi t+i\sin2\pi t$: morfism (formulele de adunare), cu imaginea $U=\{z\in\CC\mid|z|=1\}$ (cercul unitate) și nucleul $\ZZ$; deci $\RR/\ZZ\simeq U$.
\item $\Phi:G\to\Aut(G)$, $\Phi(g)=i_g$, are $\Ima\Phi=\Inn(G)$ și $\Ker\Phi=Z(G)$ (Teorema~\ref{teo:inn}); deci $G/Z(G)\simeq\Inn(G)$.
\item $\mathrm{sgn}:S_n\to(\{\pm1\},\cdot)$ ($n\ge2$; Teorema~\ref{teo:sgn}): surjectiv (o transpoziție are semnul $-1$), cu nucleul $A_n$ (grupul altern al permutărilor pare); deci $S_n/A_n\simeq\ZZ_2$ și $|A_n|=n!/2$.
\end{enumerate}
\end{exemplu}

\begin{teorema}[teorema de corespondență]\label{teo:corespondenta}
Fie $N\trianglelefteq G$ și $\pi:G\to G/N$ proiecția canonică. Atunci
\[
H\ \longmapsto\ \pi(H)=H/N=\{hN\mid h\in H\}
\]
este o \textbf{bijecție} între subgrupurile lui $G$ care conțin $N$ și subgrupurile lui $G/N$, cu inversa $\overline H\mapsto\pi^{-1}(\overline H)$. Bijecția păstrează incluziunile și normalitatea: pentru $N\le H\le G$,
\[
H\trianglelefteq G\iff H/N\trianglelefteq G/N .
\]
Pentru $G$ finit, $\bigl|\pi^{-1}(\overline H)\bigr|=|\overline H|\cdot|N|$.
\end{teorema}

\begin{proof}
$\pi(H)$ și $\pi^{-1}(\overline H)$ sunt subgrupuri (Propoziția~\ref{prop:imagini}), iar $\pi^{-1}(\overline H)\supseteq\pi^{-1}(\{eN\})=N$. Compunerile: $\pi\bigl(\pi^{-1}(\overline H)\bigr)=\overline H$ din surjectivitatea lui $\pi$; iar pentru $N\le H\le G$: $x\in\pi^{-1}(\pi(H))\iff xN=hN$ pentru un $h\in H\iff x\in hN\subseteq H$ (căci $N\subseteq H$), deci $\pi^{-1}(\pi(H))=H$. Păstrarea incluziunilor este evidentă. Normalitatea: dacă $H\trianglelefteq G$, atunci $(gN)(hN)(gN)^{-1}=(ghg^{-1})N\in H/N$; reciproc, din $(ghg^{-1})N\in H/N$ rezultă $ghg^{-1}\in HN=H$. Cardinalul: $\pi^{-1}(\overline H)$ este reuniunea disjunctă a celor $|\overline H|$ clase, fiecare de cardinal $|N|$.
\end{proof}

\begin{lema}\label{lema:HNsubgrup}
Fie $G$ un grup, $H\le G$ și $N\trianglelefteq G$. Atunci $HN=\{hn\mid h\in H,\ n\in N\}$ este subgrup al lui $G$, $N\trianglelefteq HN$ și $H\cap N\trianglelefteq H$.
\end{lema}
\begin{proof}
$e=ee\in HN$. Pentru $h_1n_1,h_2n_2\in HN$, folosind normalitatea ($Nh=hN$):
$(h_1n_1)(h_2n_2)^{-1}=h_1n_1n_2^{-1}h_2^{-1}=h_1h_2^{-1}\,(h_2 n_1n_2^{-1}h_2^{-1})\in HN$,
căci $h_2n_1n_2^{-1}h_2^{-1}\in N$. Deci $HN\le G$. $N\trianglelefteq HN$ pentru că $N\trianglelefteq G$; iar $H\cap N\trianglelefteq H$ deoarece pentru $h\in H$, $h(H\cap N)h^{-1}\subseteq H$ ($H$ subgrup) și $\subseteq N$ ($N$ normal), deci $\subseteq H\cap N$.
\end{proof}

\begin{teorema}[a doua teoremă de izomorfism]\label{teo:izo2}
Fie $G$ un grup, $H\le G$ și $N\trianglelefteq G$. Atunci
\[
\frac{HN}{N}\ \cong\ \frac{H}{H\cap N}.
\]
\end{teorema}
\begin{proof}
Definim $\varphi:H\to HN/N$ prin $\varphi(h)=hN$; este morfism, ca restricție la $H$ a proiecției canonice $HN\to HN/N$. Este \emph{surjectiv}: orice element al lui $HN/N$ este $hnN=hN=\varphi(h)$ (căci $nN=N$). Nucleul: $\varphi(h)=N\iff hN=N\iff h\in N$, deci $\Ker\varphi=H\cap N$. Teorema fundamentală de izomorfism dă $H/(H\cap N)\cong\Ima\varphi=HN/N$.
\end{proof}

\begin{teorema}[a treia teoremă de izomorfism]\label{teo:izo3}
Fie $G$ un grup și $M,N\trianglelefteq G$ cu $M\subseteq N$. Atunci $M\trianglelefteq N$, $N/M\trianglelefteq G/M$ și
\[
\frac{G/M}{N/M}\ \cong\ \frac{G}{N}.
\]
\end{teorema}
\begin{proof}
$M\trianglelefteq N$ e clar ($M$ normal în $G$, deci și în subgrupul $N$). Definim $\psi:G/M\to G/N$ prin $\psi(gM)=gN$; este \emph{bine definită}, căci $M\subseteq N$ (dacă $g_1M=g_2M$ atunci $g_1^{-1}g_2\in M\subseteq N$, deci $g_1N=g_2N$). Este morfism surjectiv, iar
\[
\Ker\psi=\{gM\mid gN=N\}=\{gM\mid g\in N\}=N/M.
\]
Teorema fundamentală de izomorfism dă $(G/M)\big/(N/M)\cong G/N$ (în particular $N/M\trianglelefteq G/M$, fiind nucleu).
\end{proof}

\subsection{Acțiuni de grupuri; teoremele lui Sylow și Frobenius}\label{sec:avansat2}

Această subsecțiune dezvoltă limbajul acțiunilor de grupuri și îl folosește pentru a demonstra complet teoremele lui Sylow (Teorema~\ref{teo:sylow}) și a lui Frobenius (Teorema~\ref{teo:frobgrup}), enunțate în §\ref{sec:mari}.

\subsubsection*{Acțiuni de grupuri și ecuația claselor}

Acțiunile sunt instrumentul unificator: ele traduc probleme de grupuri în numărări de orbite, iar teoremele lui Sylow se demonstrează integral prin acest limbaj.

\begin{definitie}[acțiune de grup]\label{def:actiune}
O \emph{acțiune} a grupului $G$ pe mulțimea nevidă $X$ este o funcție $G\times X\to X$, $(g,x)\mapsto g\cdot x$, cu proprietățile
\[
e\cdot x=x\qquad\text{și}\qquad g\cdot(h\cdot x)=(gh)\cdot x\qquad(\forall\,g,h\in G,\ \forall\,x\in X).
\]
\end{definitie}

\begin{propozitie}[acțiuni $=$ morfisme în $S(X)$]\label{prop:actmorf}
A da o acțiune a lui $G$ pe $X$ revine la a da un morfism $\varphi:G\to S(X)$, prin corespondența $\varphi(g)(x)=g\cdot x$. Acțiunea se numește \emph{fidelă} dacă $\varphi$ este injectiv.
\end{propozitie}

\begin{proof}
Dintr-o acțiune: $\varphi(g)\circ\varphi(g^{-1})=\varphi(g^{-1})\circ\varphi(g)=\varphi(e)=1_X$ (axiomele), deci fiecare $\varphi(g)$ este bijecție, iar a doua axiomă spune exact $\varphi(gh)=\varphi(g)\circ\varphi(h)$. Reciproc, orice morfism $\varphi:G\to S(X)$ definește o acțiune prin $g\cdot x=\varphi(g)(x)$.
\end{proof}

\begin{exemplu}[acțiuni fundamentale]\label{ex:actiuni}
\begin{enumerate}[label=(\alph*),nosep]
\item \emph{Translațiile}: $G$ acționează pe el însuși prin $g\cdot x=gx$; acțiunea este fidelă (demonstrația teoremei lui Cayley).
\item \emph{Conjugarea pe elemente}: $G$ acționează pe $G$ prin $g\cdot x=gxg^{-1}$ (axiomele: $e\cdot x=x$, iar $g\cdot(h\cdot x)=g(hxh^{-1})g^{-1}=(gh)x(gh)^{-1}=(gh)\cdot x$).
\item \emph{Conjugarea pe submulțimi}: $G$ acționează pe mulțimea submulțimilor lui $G$ prin $g\cdot M=gMg^{-1}$; în particular pe mulțimea \emph{subgrupurilor} (conjugatul $gHg^{-1}$ al unui subgrup este tot subgrup, căci $x\mapsto gxg^{-1}$ este bijecție ce păstrează operația).
\item \emph{Acțiunea pe clase}: pentru $H\le G$, grupul $G$ acționează pe mulțimea claselor la stânga $\{xH\mid x\in G\}$ prin $g\cdot(xH)=(gx)H$; formula este bine definită ($xH=yH\Rightarrow(gx)H=(gy)H$) și \textbf{nu cere normalitatea} lui $H$.
\end{enumerate}
\end{exemplu}

\begin{definitie}[orbită; stabilizator]
Pentru o acțiune a lui $G$ pe $X$ și $x\in X$ definim
\[
\Orb(x)=G\cdot x=\{g\cdot x\mid g\in G\}\subseteq X,
\qquad
\Stab(x)=\{g\in G\mid g\cdot x=x\}\subseteq G .
\]
\end{definitie}

\begin{teorema}[orbită--stabilizator]\label{teo:orbstab}
\emph{(i)} Relația $x\sim y\iff\exists\,g\in G:\ y=g\cdot x$ este o relație de echivalență pe $X$, ale cărei clase sunt orbitele; în particular, \textbf{orbitele partiționează} $X$.
\emph{(ii)} $\Stab(x)\le G$.
\emph{(iii)} Corespondența $g\,\Stab(x)\mapsto g\cdot x$ este o bijecție bine definită de la mulțimea claselor la stânga $\{g\,\Stab(x)\mid g\in G\}$ la $\Orb(x)$; în particular, pentru $G$ finit,
\[
|\Orb(x)|=[\,G:\Stab(x)\,]\qquad\text{și}\qquad|G|=|\Orb(x)|\cdot|\Stab(x)| .
\]
\end{teorema}

\begin{proof}
(i) Reflexivă ($e\cdot x=x$), simetrică ($y=g\cdot x\Rightarrow x=g^{-1}\cdot y$), tranzitivă (compunerea acțiunilor); clasa de echivalență a lui $x$ este chiar $\Orb(x)$, iar clasele unei relații de echivalență formează o partiție. (ii) $e\in\Stab(x)$; dacă $g,h\in\Stab(x)$, atunci $h^{-1}\cdot x=h^{-1}\cdot(h\cdot x)=x$, deci $(gh^{-1})\cdot x=g\cdot x=x$; criteriul de subgrup. (iii)
\[
\begin{aligned}
g\cdot x=h\cdot x&\iff(h^{-1}g)\cdot x=x\iff h^{-1}g\in\Stab(x)\\
&\iff g\,\Stab(x)=h\,\Stab(x),
\end{aligned}
\]
deci corespondența este bine definită și injectivă; surjectivitatea pe orbită este definiția. Egalitățile finale: teorema lui Lagrange.
\end{proof}

\begin{teorema}[ecuația claselor pentru o acțiune]\label{teo:ecclase}
Fie $G$ un grup finit care acționează pe o mulțime finită $X$, și fie $x_1,\dots,x_r$ un sistem de reprezentanți ai orbitelor. Atunci
\[
|X|=\sum_{i=1}^{r}|\Orb(x_i)|=\sum_{i=1}^{r}[\,G:\Stab(x_i)\,].
\]
Separând orbitele cu un singur element (adică punctele fixe $X^{G}=\{x\mid g\cdot x=x\ \forall g\in G\}$) de cele cu cel puțin două, se obține forma
\[
|X|=|X^{G}|+\sum_{\substack{i\\ |\Orb(x_i)|\ge 2}}[\,G:\Stab(x_i)\,],
\]
în care fiecare termen al sumei este un divizor $>1$ al lui $|G|$.
\end{teorema}

\begin{proof}
Orbitele formează o partiție a lui $X$ (Teorema~\ref{teo:orbstab}(i)), deci $|X|$ este suma cardinalelor orbitelor distincte; prin orbită--stabilizator, $|\Orb(x_i)|=[\,G:\Stab(x_i)\,]$, divizor al lui $|G|$ prin Lagrange. O orbită are un singur element exact când $x_i\in X^{G}$, ceea ce dă a doua formă.
\end{proof}

\begin{teorema}[conjugarea pe elemente; ecuația claselor]\label{teo:clase}
Pentru acțiunea prin conjugare a lui $G$ pe el însuși:
\emph{(i)} orbita lui $x$ este \emph{clasa de conjugare} $\mathcal C(x)=\{gxg^{-1}\mid g\in G\}$, iar $\Stab(x)=C_G(x)$; deci $|\mathcal C(x)|=[G:C_G(x)]$ pentru $G$ finit;
\emph{(ii)} $\mathcal C(x)=\{x\}\iff x\in Z(G)$;
\emph{(iii)} \emph{(ecuația claselor)} dacă $G$ este finit și $x_1,\dots,x_r$ sunt reprezentanți ai claselor de conjugare cu cel puțin două elemente, atunci
\[
|G|=|Z(G)|+\sum_{i=1}^{r}\,[\,G:C_G(x_i)\,],
\]
fiecare termen al sumei fiind $>1$ și divizor al lui $|G|$.
\end{teorema}

\begin{proof}
(i) $g\cdot x=x\iff gxg^{-1}=x\iff gx=xg\iff g\in C_G(x)$; apoi orbită--stabilizator. (ii) Direct din (i). (iii) Orbitele partiționează $G$; clasele cu un singur element sunt, prin (ii), exact elementele centrului, iar celelalte au cardinalele $[G:C_G(x_i)]>1$, divizori ai lui $|G|$.
\end{proof}

\begin{teorema}[submulțimi conjugate]\label{teo:conjsub}
Pentru acțiunea prin conjugare pe submulțimile lui $G$: $\Stab(M)=N_G(M)$ (normalizatorul!), deci, pentru $G$ finit, numărul conjugatelor distincte ale lui $M$ este
\[
\bigl|\{gMg^{-1}\mid g\in G\}\bigr|=[\,G:N_G(M)\,].
\]
În particular, pentru un subgrup $H\le G$: numărul conjugatelor lui $H$ este $[G:N_G(H)]$, iar $H\trianglelefteq G$ dacă și numai dacă $H$ este singurul său conjugat.
\end{teorema}

\begin{proof}
$g\cdot M=M\iff gMg^{-1}=M\iff g\in N_G(M)$; apoi Teorema~\ref{teo:orbstab} și Definiția~\ref{def:normal}.
\end{proof}

\begin{propozitie}[$p^k$ și centrul]\label{prop:pkcentru}
Fie $G$ un grup finit și $p$ prim. Presupunem că $p^k\mid|G|$ pentru un $k\ge1$. Atunci fie $p^k$ divide ordinul unui subgrup \emph{propriu} al lui $G$, fie $p\mid|Z(G)|$.
\end{propozitie}
\begin{proof}
Folosim ecuația claselor $|G|=|Z(G)|+\sum_{i=1}^r[G:C_G(g_i)]$, suma fiind luată peste clasele de conjugare cu mai mult de un element (deci $g_i\notin Z(G)$, adică $C_G(g_i)\neq G$: centralizatoarele sunt subgrupuri \emph{proprii}). Dacă pentru vreun $i$ avem $p^k\mid|C_G(g_i)|$, concluzia e atinsă (subgrup propriu). Altfel, pentru fiecare $i$, din $p^k\mid|G|=[G:C_G(g_i)]\cdot|C_G(g_i)|$ și $p^k\nmid|C_G(g_i)|$ rezultă $p\mid[G:C_G(g_i)]$. Atunci $p$ divide suma $\sum_i[G:C_G(g_i)]$ și divide $|G|$, deci divide diferența $|Z(G)|=|G|-\sum_i[G:C_G(g_i)]$.
\end{proof}

\begin{corolar}[centrul $p$-grupurilor]\label{cor:pcentru}
Dacă $|G|=p^k$ cu $p$ prim și $k\ge1$ ($G$ este un \emph{$p$-grup}), atunci $Z(G)\neq\{e\}$.
\end{corolar}

\begin{proof}
În ecuația claselor, fiecare termen $[G:C_G(x_i)]$ este un divizor $>1$ al lui $p^k$, deci multiplu de $p$. Rezultă $|Z(G)|\equiv|G|\equiv0\pmod p$, deci $|Z(G)|\ge p$.
\end{proof}

\begin{lema}\label{lema:GZ}
Dacă $G/Z(G)$ este ciclic, atunci $G$ este abelian.
\end{lema}

\begin{proof}
Fie $G/Z(G)=\langle g\,Z(G)\rangle$. Orice element al lui $G$ aparține unei clase $g^i\,Z(G)$, deci se scrie $g^iz$ cu $z\in Z(G)$. Atunci, pentru $x=g^iz$ și $y=g^jw$ ($z,w\in Z(G)$):
\[
xy=g^iz\,g^jw=g^{i+j}zw=g^jw\,g^iz=yx,
\]
folosind centralitatea lui $z,w$ și comutarea puterilor lui $g$ între ele.
\end{proof}

\begin{corolar}\label{cor:p2}
Orice grup de ordin $p^2$ ($p$ prim) este abelian.
\end{corolar}

\begin{proof}
Din Corolarul~\ref{cor:pcentru} și Lagrange, $|Z(G)|\in\{p,\,p^2\}$. Dacă $|Z(G)|=p$, atunci $|G/Z(G)|=p$, deci $G/Z(G)$ este ciclic (Corolarul~\ref{cor:lagrange}(iii)), iar Lema~\ref{lema:GZ} dă $G$ abelian, adică $Z(G)=G$ --- contradicție. Deci $Z(G)=G$.
\end{proof}

\begin{propozitie}[grupuri de ordin $p^2$]\label{prop:ordp2struct}
Orice grup de ordin $p^2$ ($p$ prim) este izomorf fie cu $\ZZ_{p^2}$, fie cu $\ZZ_p\times\ZZ_p$ (caz în care orice element $\neq e$ are ordinul $p$).
\end{propozitie}

\begin{proof}
$G$ este abelian (Corolarul~\ref{cor:p2}). Dacă are un element de ordin $p^2$, este ciclic, $\simeq\ZZ_{p^2}$. Altfel orice $x\neq e$ are ordinul $p$; alegem $a\neq e$ și $b\notin\langle a\rangle$. Atunci $\langle a\rangle\cap\langle b\rangle=\{e\}$ (ordinele fiind prime), iar $(\cl i,\cl j)\mapsto a^ib^j$ este morfism (căci $G$ e abelian) de la $\ZZ_p\times\ZZ_p$ la $G$, injectiv ($a^ib^j=e\Rightarrow a^i=b^{-j}\in\langle a\rangle\cap\langle b\rangle$) și deci bijectiv, $|G|=p^2$. Așadar $G\simeq\ZZ_p\times\ZZ_p$.
\end{proof}

\begin{lema}[normalizatorul crește]\label{lema:normcreste}
Fie $G$ un $p$-grup finit și $H<G$ un subgrup propriu. Atunci $H\subsetneq N_G(H)$.
\end{lema}
\begin{proof}
Inducție după $|G|$. Centrul $Z=Z(G)$ este netrivial (Corolarul~\ref{cor:pcentru}). Dacă $Z\not\subseteq H$, alegem $z\in Z\setminus H$: fiind central, $z$ normalizează $H$, deci $z\in N_G(H)\setminus H$. Dacă $Z\subseteq H$, trecem la $G/Z$, un $p$-grup mai mic, în care $H/Z<G/Z$ este propriu; prin ipoteza de inducție $H/Z\subsetneq N_{G/Z}(H/Z)$. Prin teorema de corespondență $N_{G/Z}(H/Z)=N_G(H)/Z$, deci $H\subsetneq N_G(H)$.
\end{proof}

\begin{propozitie}[lanț normal în $p$-grupuri]\label{prop:pchain}
Fie $G$ un $p$-grup și $H<G$ un subgrup propriu. Atunci există un subgrup $K$ al lui $G$ cu $H\trianglelefteq K$ și $K/H$ ciclic de ordin $p$.
\end{propozitie}
\begin{proof}
Prin Lema~\ref{lema:normcreste}, $H\subsetneq N_G(H)$, iar $H\trianglelefteq N_G(H)$ ($H$ e normal în normalizatorul său). Grupul factor $N_G(H)/H$ este un $p$-grup netrivial, deci (teorema lui Cauchy) conține un element de ordin $p$, care generează un subgrup de ordin $p$ al lui $N_G(H)/H$. Prin teorema de corespondență, acesta este $K/H$ pentru un subgrup $K$ cu $H\trianglelefteq K\le N_G(H)$, iar $K/H$ este ciclic de ordin $p$.
\end{proof}

\begin{lema}[punctele fixe ale $p$-grupurilor]\label{lema:fix}
Fie $P$ un $p$-grup finit care acționează pe mulțimea finită $X$ și
\[
X^{P}=\{x\in X\mid g\cdot x=x\ \text{pentru orice }g\in P\}
\]
mulțimea punctelor fixe. Atunci $|X|\equiv|X^{P}|\pmod p$.
\end{lema}

\begin{proof}
Orbitele partiționează $X$ și au cardinalele $[P:\Stab(x)]$, divizori ai lui $|P|$, deci puteri ale lui $p$. Orbitele cu un singur element sunt exact punctele fixe; toate celelalte au cardinal divizibil cu $p$. Sumând pe orbite, $|X|\equiv|X^P|\pmod p$.
\end{proof}

\subsubsection*{Lema lui Burnside}

\begin{lema}[Burnside; numărarea orbitelor]\label{lema:burnside}
Dacă grupul finit $G$ acționează pe mulțimea finită $X$, atunci
\[
\sum_{g\in G}\bigl|\mathrm{Fix}(g)\bigr|=|G|\cdot(\text{numărul de orbite}),\qquad \mathrm{Fix}(g)=\{x\in X\mid g\cdot x=x\}.
\]
În particular, $\sum_{g\in G}|\mathrm{Fix}(g)|\equiv0\pmod{|G|}$.
\end{lema}

\begin{proof}
Numărăm în două moduri perechile $(g,x)$ cu $g\cdot x=x$. Sumând după $g$ obținem $\sum_g|\mathrm{Fix}(g)|$. Sumând după $x$ obținem $\sum_x|\Stab(x)|=\sum_x\frac{|G|}{|\Orb(x)|}=|G|\sum_x\frac1{|\Orb(x)|}$, iar pe fiecare orbită $O$ suma $\sum_{x\in O}\frac1{|O|}=1$, deci totalul este $|G|\cdot(\text{numărul de orbite})$.
\end{proof}

\subsubsection*{Demonstrația teoremelor lui Sylow}\label{dem:sylow}

Demonstrăm Teorema~\ref{teo:sylow} punct cu punct, folosind exclusiv acțiuni de grupuri. Fie $|G|=p^am$, $p\nmid m$.

\paragraph{(I) Existența (Wielandt).}
Fie $\Omega$ mulțimea submulțimilor lui $G$ cu exact $p^a$ elemente; atunci $|\Omega|=\binom{p^am}{p^a}$. \emph{Arătăm întâi $p\nmid\binom{p^am}{p^a}$}: scriem $\binom{p^am}{p^a}=\prod_{j=0}^{p^a-1}\frac{p^am-j}{p^a-j}$; factorul $j=0$ este $m$, nedivizibil cu $p$, iar pentru $1\le j\le p^a-1$, cu $c$ exponentul lui $p$ în $j$ (deci $c<a$), atât $p^am-j$ cât și $p^a-j$ au exponentul $p$-adic $c$, deci fiecare factor are exponentul $0$. Așadar $p\nmid|\Omega|$. Grupul $G$ acționează pe $\Omega$ prin translație la stânga, $g\cdot X=gX$; orbitele partiționează $\Omega$, deci există o orbită $O$ cu $p\nmid|O|$. Fie $X\in O$ și $H=\Stab(X)=\{g\mid gX=X\}$; prin orbită--stabilizator (Teorema~\ref{teo:orbstab}), $|O|=[G:H]$, deci $p^a\mid|H|$, adică $|H|\ge p^a$. Pe de altă parte, fixând $x_0\in X$, aplicația $h\mapsto hx_0$ trimite $H$ injectiv în $X$ (căci $hX=X$), deci $|H|\le|X|=p^a$. Prin urmare $|H|=p^a$: $H$ este $p$-subgrup Sylow. \emph{(Aceasta regenerează teorema lui Cauchy}: un element $x\neq e$ al lui $H$ are ordinul $p^j$, iar $x^{p^{j-1}}$ are ordinul $p$.\emph{)}

\paragraph{(II) Conjugarea.}
Fie $P$ un $p$-subgrup Sylow și $Q\le G$ cu $|Q|$ o putere a lui $p$. $Q$ acționează pe mulțimea claselor $X=\{xP\mid x\in G\}$ prin $q\cdot(xP)=(qx)P$ (Exemplul~\ref{ex:actiuni}(d)); $|X|=[G:P]=m$, nedivizibil cu $p$. Prin lema punctelor fixe (Lema~\ref{lema:fix}), $|X^Q|\equiv m\not\equiv0\pmod p$, deci există o clasă fixă $gP$: $(qg)P=gP$ pentru orice $q\in Q$, adică $g^{-1}Qg\subseteq P$, deci $Q\subseteq gPg^{-1}$. Dacă $Q$ este el însuși Sylow, egalitatea $|Q|=|gPg^{-1}|=p^a$ forțează $Q=gPg^{-1}$: oricare două Sylow-uri sunt conjugate.

\paragraph{(III) Numărarea.}
Fie $\Sigma$ mulțimea $p$-subgrupurilor Sylow; $G$ acționează pe $\Sigma$ prin conjugare, tranzitiv (punctul (II)), cu $\Stab(P)=N_G(P)$ (Teorema~\ref{teo:conjsub}), deci $n_p=|\Sigma|=[G:N_G(P)]$. Din $P\le N_G(P)\le G$ rezultă $[G:P]=[G:N_G(P)][N_G(P):P]$, deci $n_p\mid[G:P]=m$. Pentru congruență, restrângem acțiunea la $P$: dacă $Q\in\Sigma$ este fixat de tot $P$, atunci $P\subseteq N_G(Q)$, deci $P,Q$ sunt Sylow-uri ale lui $N_G(Q)$; prin (II) aplicat în $N_G(Q)$ ele sunt conjugate acolo, dar $Q\trianglelefteq N_G(Q)$, deci $P=Q$. Astfel $\Sigma^P=\{P\}$, iar Lema~\ref{lema:fix} dă $n_p=|\Sigma|\equiv|\Sigma^P|=1\pmod p$.

\subsubsection*{Demonstrația teoremei lui Frobenius}\label{dem:frobenius}

\begin{definitie}[funcția lui Möbius]\label{def:mobius}
Pentru $n\in\NN^*$ definim $\mu(1)=1$, $\mu(n)=(-1)^k$ dacă $n$ este produsul a $k$ numere prime distincte și $\mu(n)=0$ dacă $n$ se divide cu pătratul unui număr prim.
\end{definitie}

\begin{propozitie}[proprietatea de bază; inversiunea Möbius]\label{prop:mobius}
\emph{(i)} $\displaystyle\sum_{d\mid n}\mu(d)$ este egală cu $1$ pentru $n=1$ și cu $0$ pentru $n>1$.
\emph{(ii)} (\emph{inversiunea Möbius}) Dacă $(M_d)_{d\ge1}$ și $(N_d)_{d\ge1}$ sunt numere (sau, mai general, elemente ale unui grup abelian) cu $N_d=\sum_{e\mid d}M_e$ pentru orice $d$, atunci $M_d=\sum_{e\mid d}\mu(d/e)\,N_e$ pentru orice $d$.
\end{propozitie}

\begin{proof}
(i) Fie $n>1$ cu divizorii primi distincți $q_1,\dots,q_k$, $k\ge1$. Divizorii $d$ cu $\mu(d)\neq0$ sunt produsele a câte $j$ dintre $q_1,\dots,q_k$, pentru care $\mu(d)=(-1)^j$; suma este $\sum_{j=0}^{k}\binom kj(-1)^j=(1-1)^k=0$.
(ii) Înlocuim $N_e$ și schimbăm ordinea sumării: pentru $c\mid d$ fixat, $e$ parcurge multiplii lui $c$ care divid $d$, iar $f=d/e$ parcurge divizorii lui $d/c$. Obținem
\[
\sum_{e\mid d}\mu(d/e)\,N_e=\sum_{e\mid d}\mu(d/e)\sum_{c\mid e}M_c=\sum_{c\mid d}M_c\sum_{f\mid d/c}\mu(f)=M_d,
\]
căci, prin (i), suma interioară este $1$ pentru $c=d$ și $0$ în rest.
\end{proof}

Folosim funcția lui Möbius și inversiunea Möbius (Propoziția~\ref{prop:mobius}(ii)).

\begin{lema}[coliere aperiodice]\label{lema:coliere}
Pentru întregi pozitivi $a\mid b$,\quad $\displaystyle\sum_{f\mid a}\mu(f)\binom{b/f}{a/f}\equiv0\pmod b$.
\end{lema}

\begin{proof}
$\binom{b/f}{a/f}$ numără șirurile binare de lungime $b$ cu $a$ cifre de $1$ care sunt invariante la rotația de perioadă $f$ (un astfel de șir este determinat de primele $b/f$ poziții, care conțin $a/f$ cifre de $1$). Prin inversiune Möbius, suma numără șirurile \emph{aperiodice} (fără nicio simetrie de rotație netrivială). Cele $b$ rotații ale unui șir aperiodic sunt distincte, deci șirurile aperiodice se grupează în orbite de câte $b$, iar numărul lor este multiplu de $b$.
\end{proof}

\emph{Demonstrația Teoremei~\ref{teo:frobgrup}} (D.~Speyer). Fie $n\mid|G|$. Pentru $d\mid|G|$ notăm $M_d$ numărul elementelor de ordin exact $d$ și $N_d=\#\{x:x^d=e\}=\sum_{e\mid d}M_e$. Grupul $G$ acționează prin translație la stânga pe mulțimea $\Omega$ a submulțimilor cu $n$ elemente. Un element $g$ de ordin $d$ fixează $X\in\Omega$ dacă și numai dacă $X$ este reuniune de clase la stânga ale lui $\langle g\rangle$; aceasta cere $d\mid n$, iar atunci se aleg $n/d$ din cele $|G|/d$ clase, deci $|\mathrm{Fix}(g)|=\binom{|G|/d}{n/d}$ (și $0$ dacă $d\nmid n$). Lema lui Burnside (Lema~\ref{lema:burnside}) dă
\[
\sum_{d\mid n}\binom{|G|/d}{n/d}\,M_d\equiv0\pmod{|G|}.\tag{$\dagger$}
\]
Înlocuind $M_d=\sum_{e\mid d}\mu(d/e)N_e$ și schimbând ordinea sumării (cu $f=d/e$):
\[
\sum_{e\mid n}N_e\,S_e\equiv0\pmod{|G|},\qquad S_e:=\sum_{f\mid n/e}\mu(f)\binom{|G|/(ef)}{n/(ef)}.\tag{$\ddagger$}
\]
Prin Lema~\ref{lema:coliere} cu $a=n/e$, $b=|G|/e$ (și $a\mid b$, căci $n\mid|G|$), fiecare $S_e$ este divizibil cu $|G|/e$. Procedăm prin inducție după $n$ (baza $n=1$: $N_1=1$). Pentru orice divizor $e<n$ al lui $n$, ipoteza de inducție dă $e\mid N_e$, deci $N_eS_e$ este divizibil cu $e\cdot\frac{|G|}{e}=|G|$. Din $(\ddagger)$ rezultă că și termenul $e=n$ este divizibil cu $|G|$; dar $S_n=\binom{|G|/n}{1}=|G|/n$, deci
\[
N_n\cdot\frac{|G|}{n}\equiv0\pmod{|G|}\quad\Longrightarrow\quad n\mid N_n.\qquad\blacksquare
\]

\subsubsection*{Aplicații}

\begin{aplicatie}[grupurile de ordin 15 sunt ciclice]\label{apl:15}
Fie $|G|=15=3\cdot5$. Din Sylow~III: $n_3\equiv1\pmod3$ și $n_3\mid5$, deci $n_3=1$; analog $n_5\equiv1\pmod5$ și $n_5\mid3$, deci $n_5=1$. Fie $P\trianglelefteq G$ Sylow de ordin $3$ și $Q\trianglelefteq G$ Sylow de ordin $5$ (Corolarul~\ref{cor:sylownormal}); $P\cap Q=\{e\}$, căci ordinul intersecției divide și $3$, și $5$. Fie $x\in P$ de ordin $3$ și $y\in Q$ de ordin $5$. Comutatorul lor stă în intersecție:
\[
xyx^{-1}y^{-1}=\underbrace{x\,\bigl(yx^{-1}y^{-1}\bigr)}_{\in P,\ \text{căci }P\trianglelefteq G}=\underbrace{\bigl(xyx^{-1}\bigr)\,y^{-1}}_{\in Q,\ \text{căci }Q\trianglelefteq G}\ \in\ P\cap Q=\{e\},
\]
deci $xy=yx$. Atunci $\ord(xy)=3\cdot5=15$ (Teorema~\ref{teo:ordprodus}, cazul coprim), deci $G=\langle xy\rangle$ este ciclic.
\end{aplicatie}

\begin{propozitie}[produs de două numere prime]\label{prop:pqciclic}
Fie $p<q$ prime cu $p\nmid q-1$. Atunci orice grup de ordin $pq$ este ciclic. (Teoremele lui Sylow sunt astfel punctul de intrare în clasificarea grupurilor finite mici.)
\end{propozitie}

\begin{proof}
Din Sylow (Teorema~\ref{teo:sylow}): $n_q\equiv1\pmod q$ și $n_q\mid p$ dau $n_q=1$; iar $n_p\equiv1\pmod p$, $n_p\mid q$ și $p\nmid q-1$ dau $n_p=1$. Deci Sylow-urile $P$ (ordin $p$) și $Q$ (ordin $q$) sunt normale, cu $P\cap Q=\{e\}$. Exact ca în Aplicația~\ref{apl:15}, generatorii lor comută, iar produsul are ordinul $pq$ (Teorema~\ref{teo:ordprodus}), deci $G$ este ciclic.
\end{proof}


\subsubsection*{Numere ciclice}

\begin{definitie}
Un număr natural $n\ge1$ se numește \emph{număr ciclic} dacă \textbf{orice} grup de ordin $n$ este ciclic.
\end{definitie}

\begin{teorema}[caracterizarea numerelor ciclice]\label{teo:numereciclice}
Un număr $n$ este ciclic dacă și numai dacă
\[
\bigl(n,\varphi(n)\bigr)=1 .
\]
\end{teorema}

\begin{proof}
Vom folosi în mod repetat că, pentru $q$ prim, $\Aut(C_q)\simeq(\ZZ_q^*,\cdot)$ este ciclic de ordin $q-1$ (Teorema~\ref{teo:ciclicitate}).

\smallskip
\emph{$(n,\varphi(n))=1\Rightarrow n$ este liber de pătrate.} Dacă $p^2\mid n$ pentru un prim $p$, atunci $p\mid n$ și, cum $\varphi(p^2)=p(p-1)$ divide $\varphi(n)$, avem $p\mid\varphi(n)$; deci $p\mid(n,\varphi(n))$, contradicție. Așadar $n=p_1p_2\cdots p_k$ cu primi distincți, iar condiția $(n,\varphi(n))=1$ (cu $\varphi(n)=\prod_i(p_i-1)$) se rescrie:
\[
p_i\nmid p_j-1\quad\text{pentru orice }i,j.\tag{$\ast$}
\]

\smallskip
\emph{Necesitatea (}$n$ ciclic $\Rightarrow(n,\varphi(n))=1$\emph{).} Prin contrapunere: presupunem $(n,\varphi(n))>1$ și construim un grup \textbf{ne}ciclic de ordin $n$. Fie $p$ prim cu $p\mid n$ și $p\mid\varphi(n)$.
\begin{itemize}[nosep,leftmargin=1.4em]
\item Dacă $p^2\mid n$: grupul abelian $C_p\times C_{n/p}$ are ordinul $n$, dar exponentul $[p,\,n/p]=n/p<n$ (căci $p\mid n/p$), deci nu are element de ordin $n$: neciclic.
\item Dacă $p\,\|\,n$ (exact o dată), din $p\mid\varphi(n)=\prod_{q}\varphi(q^{a_q})$ rezultă $p\mid\varphi(q^{a})=q^{a-1}(q-1)$ pentru un prim $q\neq p$; cum $p\neq q$, obligatoriu $p\mid q-1$, adică $q\equiv1\pmod p$. Fie $H\le\ZZ_q^*$ subgrupul (ciclic) de ordin $p$ (există, căci $p\mid q-1=|\ZZ_q^*|$). Mulțimea aplicațiilor afine
\[
G_0=\bigl\{\,x\mapsto ax+b\ \big|\ a\in H,\ b\in\ZZ_q\,\bigr\}
\]
este grup față de compunere, de ordin $|H|\cdot q=pq$, și este \textbf{neabelian} (compunerea $x\mapsto a(x+b)$ vs.\ $x\mapsto ax+b$ diferă pentru $a\neq1$, $b\neq0$). Atunci $G_0\times C_{n/(pq)}$ este neabelian, deci neciclic, de ordin $n$.
\end{itemize}

\smallskip
\emph{Suficiența (}$(n,\varphi(n))=1\Rightarrow n$ ciclic\emph{).} Fie $G$ de ordin $n$ liber de pătrate, cu condiția $(\ast)$; arătăm prin inducție după $k$ că $G$ este ciclic. Pentru $k=1$, $|G|=p$ este prim, deci ciclic. Pentru pasul de inducție, fie $p$ cel mai mic prim care divide $n$ și $P$ un $p$-subgrup Sylow (de ordin $p$, deci $P\simeq C_p$). Grupul $N_G(P)/C_G(P)$ se scufundă în $\Aut(P)\simeq C_{p-1}$, deci ordinul său divide $p-1$; dar divide și $n$. Orice divizor prim al lui $p-1$ este $<p$, deci (prin minimalitatea lui $p$) nu divide $n$; așadar $\bigl(p-1,n\bigr)=1$ și $N_G(P)=C_G(P)$, adică $P\le Z\bigl(N_G(P)\bigr)$. Prin \emph{teorema complementului normal a lui Burnside} (rezultat clasic, a cărui demonstrație folosește transferul și depășește cadrul fișei), $G$ admite un \emph{$p$-complement normal}: un subgrup normal $M\trianglelefteq G$ de ordin $n/p$, cu $G=M\rtimes P$. Ordinul $n/p$ satisface aceeași condiție, deci prin inducție $M\simeq C_{n/p}$ este ciclic. Acțiunea prin conjugare $P\to\Aut(M)$ are imaginea de ordin divizor al lui $(p,\varphi(n/p))$; dar $\varphi(n/p)\mid\varphi(n)$ și $(p,\varphi(n))=1$, deci acțiunea este trivială: $G=M\times P\simeq C_{n/p}\times C_p$. Cum $(n/p,p)=1$, produsul este ciclic (Teorema~\ref{teo:ordprodus}, cazul coprim), deci $G\simeq C_n$.
\end{proof}

\begin{observatie}[ciclic vs.\ abelian --- o precizare importantă]\label{obs:ciclicabelian}
Condiția $(n,\varphi(n))=1$ caracterizează numerele pentru care orice grup de ordin $n$ este \textbf{ciclic}, \emph{nu} doar abelian. Cele două noțiuni chiar diferă: pentru $n=4$, orice grup de ordin $4$ este abelian ($C_4$ sau grupul lui Klein), dar $(4,\varphi(4))=(4,2)=2\neq1$; într-adevăr, grupul lui Klein nu este ciclic. Mai general, orice $n=p^2$ este ,,număr abelian`` (grupurile de ordin $p^2$ sunt abeliene --- Corolarul~\ref{cor:p2}) fără a fi număr ciclic. Mulțimea numerelor $n$ pentru care orice grup de ordin $n$ este abelian este deci strict mai mare. Caracterizarea ei (teoremă clasică, pe care o cităm fără demonstrație) este: $n=p_1^{a_1}\cdots p_k^{a_k}$ este ,,număr abelian`` dacă și numai dacă toți exponenții satisfac $a_i\le2$ (adică $n$ este liber de cuburi) și $p_i\nmid p_j^{a_j}-1$ pentru orice $i,j$. Numerele abeliene mai mici decât $200$ care nu sunt ciclice sunt $4,9,25,45,49,99,121,153,169,175$. De pildă $55=5\cdot11$ \textbf{nu} este abelian: cum $5\mid11-1$, construcția din demonstrația Teoremei~\ref{teo:numereciclice} dă un grup neabelian de ordin $55$.
\end{observatie}

\begin{corolar}
Orice număr liber de pătrate $n$ cu proprietatea $(\ast)$ --- niciun factor prim nu divide (vreun alt factor prim) $-1$ --- este ciclic. În particular, orice număr prim, orice produs $pq$ de primi cu $p\nmid q-1$ și $q\nmid p-1$ (de exemplu $15=3\cdot5$, $33=3\cdot11$, $35=5\cdot7$) este ciclic: orice grup de acel ordin este $C_n$.
\end{corolar}

\begin{proof}
Direct din Teorema~\ref{teo:numereciclice} și rescrierea $(\ast)$ a condiției $(n,\varphi(n))=1$.
\end{proof}

\subsection{Produse directe de grupuri}\label{sec:produsedirecte}

\begin{definitie}
Fie $G_1,\dots,G_n$ grupuri. \emph{Produsul direct} $G_1\times\cdots\times G_n$ este produsul cartezian al mulțimilor, înzestrat cu operația pe componente
\[
(x_1,\dots,x_n)(y_1,\dots,y_n)=(x_1y_1,\dots,x_ny_n).
\]
\end{definitie}

\begin{propozitie}
$G_1\times\cdots\times G_n$ este grup, cu neutrul $(e_1,\dots,e_n)$ și inversul $(x_1,\dots,x_n)^{-1}=(x_1^{-1},\dots,x_n^{-1})$; este abelian dacă și numai dacă toate $G_i$ sunt abeliene, iar $|G_1\times\cdots\times G_n|=\prod_i|G_i|$. Dacă fiecare $x_i$ are ordin finit, atunci
\[
\ord\big((x_1,\dots,x_n)\big)=\big[\ord(x_1),\dots,\ord(x_n)\big]\quad(\text{c.m.m.m.c.}).
\]
\end{propozitie}
\begin{proof}
Axiomele de grup se verifică pe componente. Pentru ordin: $(x_1,\dots,x_n)^k=(x_1^k,\dots,x_n^k)=(e_1,\dots,e_n)$ dacă și numai dacă $\ord(x_i)\mid k$ pentru fiecare $i$, adică dacă și numai dacă $[\ord(x_1),\dots,\ord(x_n)]\mid k$; cel mai mic astfel de $k$ este c.m.m.m.c.
\end{proof}

\begin{corolar}[criteriul de ciclicitate a produsului]\label{cor:produsciclic}
Pentru grupuri ciclice, $C_m\times C_n\cong C_{mn}$ dacă și numai dacă $(m,n)=1$. În particular $C_2\times C_3\cong C_6$ este ciclic, pe când $C_2\times C_2$ (grupul lui Klein) nu este ciclic.
\end{corolar}
\begin{proof}
Fie $x,y$ generatori ai lui $C_m$, respectiv $C_n$. Ordinul lui $(x,y)$ este $[m,n]$. Dacă $(m,n)=1$, atunci $[m,n]=mn=|C_m\times C_n|$, deci $(x,y)$ generează, iar produsul e ciclic. Dacă $d=(m,n)>1$, atunci $[m,n]=\tfrac{mn}{d}<mn$: orice element are ordinul divizor al lui $[m,n]<mn$, deci nu există generator (Teorema despre ordinul unui produs de elemente care comută).
\end{proof}

\subsection{Grupuri abeliene finit generate; grupuri simple și rezolubile}\label{sec:avansat3}

Această subsecțiune încununează teoria grupurilor cu trei rezultate structurale majore, construite pe Sylow, teorema de corespondență și grupurile factor.

\subsubsection*{Grupuri abeliene finit generate}

Lucrăm în grupul aditiv $\ZZ^n$ al $n$-uplurilor de întregi. O listă $b_1,\dots,b_r\in\ZZ^n$ este \emph{bază integrală} (sau $\ZZ$-bază) a unui subgrup $H\le\ZZ^n$ dacă orice element al lui $H$ se scrie \textbf{unic} sub forma $m_1b_1+\dots+m_rb_r$ cu $m_i\in\ZZ$ (și fiecare astfel de combinație aparține lui $H$).

\begin{observatie}
Elementele $b_1,\dots,b_n$ formează o $\ZZ$-bază a lui $\ZZ^n$ însuși dacă și numai dacă matricea (întreagă) cu liniile $b_i$ are determinantul $\pm1$ (matrice \emph{unimodulară}); inversa unei astfel de matrice este tot întreagă. Schimbările de bază admisibile în $\ZZ^n$ sunt deci exact transformările unimodulare.
\end{observatie}

\begin{teorema}[bază adaptată; ,,stacked basis``]\label{teo:adaptata}
Fie $H\le\ZZ^n$ un subgrup netrivial. Atunci există o bază integrală $(b_1,\dots,b_n)$ a lui $\ZZ^n$, un întreg $s$ cu $1\le s\le n$ și întregi pozitivi $k_1,\dots,k_s$ astfel încât $(k_1b_1,\dots,k_sb_s)$ este bază integrală a lui $H$.
\end{teorema}

\begin{proof}
Inducție după $n$. Pentru $n=1$: orice subgrup netrivial al lui $\ZZ$ este $k\ZZ$ (Teorema~\ref{teo:subgrupZ}), deci $b_1=1$, $k_1=k$.

Fie $n>1$. Printre toate perechile (bază $(u_1,\dots,u_n)$ a lui $\ZZ^n$; element $h\in H$), scriem $h=\sum_i m_iu_i$ și alegem $k_1$ = \emph{cea mai mică valoare strict pozitivă} a unui coeficient $m_i$ care apare. Putem presupune că indicele este $1$: există o bază $(u_1,\dots,u_n)$ și $h_0=k_1u_1+m_2u_2+\dots+m_nu_n\in H$.

\emph{Toți $m_i$ se divid cu $k_1$}: scriem $m_i=q_ik_1+r_i$, $0\le r_i<k_1$, și punem $b_1=u_1+\sum_{i\ge2}q_iu_i$. Atunci $(b_1,u_2,\dots,u_n)$ este tot bază (schimbare unimodulară) și $h_0=k_1b_1+\sum_{i\ge2}r_iu_i$. Minimalitatea lui $k_1$ forțează $r_i=0$, deci $h_0=k_1b_1$.

Fie $\iota:\ZZ^{n-1}\to\ZZ^n$ morfismul injectiv $(m_2,\dots,m_n)\mapsto\sum_{i\ge2}m_iu_i$ și $H'=\iota^{-1}(H)$. Orice $h\in H$ se scrie unic $h=m\,b_1+\iota(h'')$; scriind $m=qk_1+r$ ($0\le r<k_1$), avem $h-qh_0=r\,b_1+\iota(\dots)$, iar minimalitatea lui $k_1$ dă $r=0$. Așadar $k_1\mid m$ și $\iota(h'')=h-qk_1b_1\in H$, deci $h''\in H'$. Prin ipoteza de inducție, $H'$ are o bază adaptată: există o bază $(b'_2,\dots,b'_n)$ a lui $\ZZ^{n-1}$ și $k_2,\dots,k_s>0$ cu $(k_2b'_2,\dots,k_sb'_s)$ bază a lui $H'$. Punând $b_i=\iota(b'_i)$, lista $(b_1,\dots,b_n)$ este bază a lui $\ZZ^n$, iar $(k_1b_1,\dots,k_sb_s)$ este bază a lui $H$.
\end{proof}

\begin{lema}\label{lema:fg}
Un grup abelian netrivial $G$ este \emph{finit generat} dacă și numai dacă există $n\ge1$ și un morfism surjectiv $\pi:\ZZ^n\to G$.
\end{lema}

\begin{proof}
Dacă $\pi:\ZZ^n\to G$ este surjectiv, atunci $G$ este generat de $\pi(e_1),\dots,\pi(e_n)$ (baza canonică). Reciproc, dacă $G=\langle g_1,\dots,g_n\rangle$ (aditiv), morfismul $\pi(m_1,\dots,m_n)=m_1g_1+\dots+m_ng_n$ este bine definit (comutativitate) și surjectiv.
\end{proof}

\begin{teorema}[structura grupurilor abeliene finit generate]\label{teo:structura}
Fie $G$ un grup abelian finit generat, netrivial. Atunci există $n\ge1$ și $s$ cu $0\le s\le n$ astfel încât
\[
G\ \simeq\ C_{k_1}\times C_{k_2}\times\dots\times C_{k_s}\times\ZZ^{\,n-s},
\]
unde $C_{k_i}$ este grup ciclic de ordin $k_i>0$ (dacă $s=0$, $G\simeq\ZZ^n$). Mai mult, se pot alege $k_1,\dots,k_s$ cu $k_1>1$ și $k_1\mid k_2\mid\dots\mid k_s$ (\emph{factorii invarianți}), caz în care descompunerea este unic determinată.
\end{teorema}

\begin{proof}
Prin Lema~\ref{lema:fg}, fie $\pi:\ZZ^n\to G$ surjectiv și $H=\Ker\pi$, deci $G\simeq\ZZ^n/H$ (teorema fundamentală, Teorema~\ref{teo:fundamental}). Dacă $H=\{0\}$, $G\simeq\ZZ^n$. Altfel, alegem o bază adaptată $(b_1,\dots,b_n)$ cu $(k_1b_1,\dots,k_sb_s)$ bază a lui $H$ (Teorema~\ref{teo:adaptata}). Morfismul care duce $\sum m_ib_i$ în \[ (\cl{m_1},\dots,\cl{m_s},m_{s+1},\dots,m_n)\in C_{k_1}\times\dots\times C_{k_s}\times\ZZ^{n-s} \] este surjectiv și are nucleul exact $H$; teorema fundamentală încheie. Reducerea la forma cu $k_1\mid\dots\mid k_s$ și unicitatea sunt un pas suplimentar standard (nu-l detaliem).
\end{proof}

\begin{corolar}[grupuri abeliene finite]\label{cor:abelianfinit}
Orice grup abelian finit netrivial este produs direct de grupuri ciclice:
$G\simeq C_{k_1}\times\dots\times C_{k_n}$.
\end{corolar}

\begin{proof}
În Teorema~\ref{teo:structura}, finitudinea forțează $s=n$ (nu poate apărea factor $\ZZ$).
\end{proof}

\begin{observatie}[descompunerea primară]
Combinând cu lema chineză (Teorema~\ref{teo:crt}), fiecare $C_{k}$ cu $k=p_1^{a_1}\cdots p_r^{a_r}$ se descompune $C_k\simeq C_{p_1^{a_1}}\times\dots\times C_{p_r^{a_r}}$. Astfel orice grup abelian finit este produs de grupuri ciclice de ordin putere de prim (forma cu \emph{divizori elementari}). De exemplu, grupurile abeliene de ordin $8$ sunt $C_8$, $C_4\times C_2$, $C_2\times C_2\times C_2$.
\end{observatie}

\begin{teorema}[produsul tuturor elementelor; generalizarea teoremei lui Wilson]\label{teo:produs}
Fie $G$ un grup \textbf{abelian} finit și $H=\{x\in G\mid x^2=e\}$. Atunci:
\emph{(i)} $H\le G$;\quad
\emph{(ii)} $\displaystyle\prod_{x\in G}x=\prod_{x\in H}x$, adică produsul tuturor elementelor se reduce la produsul elementelor de ordin cel mult $2$;\quad
\emph{(iii)} dacă $G$ are \textbf{exact un} element $t$ de ordin $2$, produsul tuturor elementelor este $t$; în toate celelalte cazuri produsul este $e$.
\end{teorema}

\begin{proof}
(i) $e\in H$, iar pentru $x,y\in H$ avem $(xy^{-1})^2=x^2(y^{-1})^2=e$ (comutativitate), deci se aplică criteriul de subgrup.

(ii) Elementele din $G\setminus H$ satisfac $x\neq x^{-1}$ și se grupează în perechi disjuncte $\{x,x^{-1}\}$ cu produsul $e$; grupul fiind abelian, putem rearanja factorii, iar aceste perechi se simplifică.

(iii) Dacă $H=\{e\}$, produsul este $e$; dacă $H=\{e,t\}$, produsul este $t$. Rămâne cazul în care $H$ conține două elemente distincte $a,b\neq e$. Atunci $K=\{e,a,b,ab\}$ este un subgrup cu $4$ elemente al lui $H$: $ab\notin\{e,a,b\}$ (căci $ab=e\Rightarrow b=a$, $ab=a\Rightarrow b=e$, $ab=b\Rightarrow a=e$), iar închiderea se verifică direct ($a\cdot ab=b$, $b\cdot ab=a$, $(ab)^2=a^2b^2=e$). Clasele $xK$ partiționează $H$ (Lema~\ref{lema:clase}), iar produsul pe fiecare clasă este
\[
x\cdot xa\cdot xb\cdot xab \;=\; x^4\,(ab)^2 \;=\; e,
\]
folosind comutativitatea și $x^4=(x^2)^2=e$. Înmulțind pe clase, $\prod_{x\in H}x=e$.
\end{proof}

\begin{observatie}
Pentru $G=(\ZZ_p^*,\cdot)$ cu $p$ prim impar, singurul element de ordin $2$ este $-\cl1$ (soluțiile ecuației $x^2=\cl1$ într-un corp sunt $\pm\cl1$), deci teorema dă $\cl1\cdot\cl2\cdots\cl{p-1}=-\cl1$ --- exact teorema lui Wilson (Teorema~\ref{teo:wilson}), văzută acum ca fapt de teoria grupurilor.
\end{observatie}

\subsubsection*{Grupuri simple}

\begin{definitie}
Un grup netrivial $G$ este \emph{simplu} dacă singurele sale subgrupuri normale sunt $\{e\}$ și $G$.
\end{definitie}

\begin{lema}\label{lema:abeliansimplu}
Un grup abelian este simplu dacă și numai dacă este ciclic de ordin prim.
\end{lema}

\begin{proof}
Într-un grup abelian orice subgrup este normal (Exemplul~\ref{ex:normal}(a)); deci simplitatea revine la a nu avea subgrupuri proprii netriviale, ceea ce, prin Propoziția~\ref{prop:farasub}, caracterizează grupurile ciclice de ordin prim. Reciproc, $C_p$ este simplu (Corolarul~\ref{cor:lagrange}(iii)).
\end{proof}

\begin{lema}[formula produsului]\label{lema:produsHK}
Fie $H,K$ subgrupuri finite ale unui grup $G$. Atunci mulțimea $HK=\{hk\mid h\in H,\ k\in K\}$ satisface
\[
|HK|\cdot|H\cap K|=|H|\cdot|K|.
\]
\end{lema}

\begin{proof}
Considerăm $\psi:H\times K\to HK$, $\psi(h,k)=hk$ (notația $\psi$, pentru a nu o confunda cu funcția lui Möbius). Avem $h_1k_1=h_2k_2\iff h_2^{-1}h_1=k_2k_1^{-1}\in H\cap K$; notând $x=h_2^{-1}h_1\in H\cap K$, perechile cu aceeași imagine sunt exact $(h_1x^{-1},xk_1)$, în număr de $|H\cap K|$. Deci fiecare element din $HK$ are exact $|H\cap K|$ preimagini, de unde $|H\times K|=|HK|\cdot|H\cap K|$.
\end{proof}

\begin{lema}\label{lema:sylownormalsimplu}
Dacă un grup finit $G$ are un $p$-subgrup Sylow normal propriu netrivial (echivalent, $n_p=1$ cu $p^a<|G|$), atunci $G$ nu este simplu. În particular, dacă din Sylow~III rezultă $n_p=1$ pentru un prim $p\mid|G|$ cu $|G|\neq p^a$, grupul nu e simplu.
\end{lema}

\begin{proof}
Un $p$-subgrup Sylow normal netrivial și diferit de $G$ este un subgrup normal propriu netrivial (Corolarul~\ref{cor:sylownormal}).
\end{proof}

\begin{teorema}[grupuri de ordine mici]\label{teo:simplemici}
Următoarele grupuri finite nu sunt simple (deci singurele grupuri simple de ordin $<60$ sunt cele ciclice de ordin prim):
\emph{(i)} orice grup de ordin $p^k$ cu $k\ge2$ (are subgrup normal de ordin $p$);
\emph{(ii)} orice grup de ordin $pq$ cu $p<q$ prime (are Sylow $q$-subgrup normal);
\emph{(iii)} grupurile de ordin $18,50,54$ (Sylow-ul de la primul impar este unic), $20,28,40,44,45,52$ (Sylow-ul de la cel mai mare prim este unic), $42$ (Sylow $7$ unic), $30$ și $56$ (numărare de elemente), $12,24,36,48$ (intersecții de Sylow-uri).
Orice număr compus mai mic decât $60$ se află într-unul dintre cazurile (i), (ii), (iii) (verificare directă pe lista numerelor compuse), iar un grup de ordin prim este ciclic, deci simplu.
\end{teorema}

\begin{proof}
(i) Un grup de ordin $p^k$ are centru netrivial (Corolarul~\ref{cor:pcentru}); prin Cauchy, $Z(G)$ conține un element de ordin $p$, generând un subgrup normal de ordin $p$ (subgrup al centrului --- Exemplul~\ref{ex:normal}(b)).

(ii) Sylow~III: $n_q\equiv1\pmod q$ și $n_q\mid p$; cum $1<p<q$, singura posibilitate este $n_q=1$, deci Sylow-ul de ordin $q$ este normal (Lema~\ref{lema:sylownormalsimplu}).

(iii) \emph{Ordin $18=2\cdot3^2$}: $n_3\equiv1\pmod3$, $n_3\mid2\Rightarrow n_3=1$: Sylow $3$ (ordin $9$) normal. Analog $50=2\cdot5^2$ (Sylow $25$) și $54=2\cdot3^3$ (Sylow $27$). \emph{Ordin $42=2\cdot3\cdot7$}: $n_7\equiv1\pmod7$, $n_7\mid6\Rightarrow n_7=1$. \emph{Ordin $30$}: dacă $n_3>1$ atunci $n_3=10$, dând $20$ elemente de ordin $3$; dacă $n_5>1$ atunci $n_5=6$, dând $24$ elemente de ordin $5$; nu pot coexista în $30$ elemente, deci un Sylow (de ordin $3$ sau $5$) este unic, deci normal.

\emph{Ordinele $20=2^2\cdot5$, $28=2^2\cdot7$, $40=2^3\cdot5$, $44=2^2\cdot11$, $45=3^2\cdot5$, $52=2^2\cdot13$}: pentru cel mai mare prim $q$ ($5,7,5,11,5,13$), Sylow~III dă $n_q\equiv1\pmod q$ și $n_q\mid4,4,8,4,9,4$; singurul divizor $\equiv1\pmod q$ este $1$ (divizorii sunt $1,2,4$, respectiv $1,2,4,8$, respectiv $1,3,9$). Deci $n_q=1$ și Sylow-ul de ordin $q$ este normal.

\emph{Ordin $56=2^3\cdot7$}: $n_7\equiv1\pmod7$ și $n_7\mid8$ dau $n_7\in\{1,8\}$. Dacă $n_7=8$, cele $8$ subgrupuri de ordin $7$ se intersectează două câte două trivial (ordin prim), deci conțin $8\cdot6=48$ elemente de ordin $7$. Rămân exact $8$ elemente de ordin diferit de $7$; orice $2$-subgrup Sylow (de ordin $8$) nu conține elemente de ordin $7$, deci coincide cu mulțimea acestor $8$ elemente. Așadar $n_2=1$. În ambele cazuri un Sylow este normal.

\emph{Ordinele $12,24,36,48$}: luăm $p=2$ pentru $12,24,48$ și $p=3$ pentru $36$; $p$-subgrupurile Sylow au ordinul $4,8,16$, respectiv $9$. Dacă $n_p=1$, am terminat. Altfel fie $H\neq K$ două $p$-subgrupuri Sylow și $D=H\cap K$. Din formula produsului (Lema~\ref{lema:produsHK}), $|H|^2/|D|=|HK|\le|G|$, deci $|D|\ge|H|^2/|G|$; cum $|D|$ este un divizor propriu al lui $|H|$, obținem $|D|=2$ (ordin $12$), $|D|=4$ (ordin $24$), $|D|=8$ (ordin $48$), respectiv $|D|=3$ (ordin $36$). În fiecare caz $[H:D]=[K:D]=p$, iar un subgrup de indice $p$ al unui $p$-grup este normal în acesta (Lema~\ref{lema:normcreste}: $D\subsetneq N_H(D)$ forțează $N_H(D)=H$). Prin urmare $N=N_G(D)$ conține $H$ și $K$, deci și mulțimea $HK$, de unde $|N|\ge|HK|=|H|^2/|D|$, adică $|N|\ge8,16,32$, respectiv $27$. Cum $|N|$ este multiplu de $|H|$ și divide $|G|$, rezultă $|N|=|G|$ în toate cele patru cazuri: $D$ este subgrup normal, netrivial și propriu al lui $G$.
\end{proof}

\begin{propozitie}[$A_5$ este simplu]\label{prop:A5}
Grupul altern $A_5$ (de ordin $60$) este simplu și neabelian; el este cel mai mic grup simplu neabelian. (Aceasta explică de ce clasificarea de mai sus se oprește la $60$.)
\end{propozitie}

\begin{proof}
Determinăm clasele de conjugare din $|\mathcal C(x)|=[A_5:C_{A_5}(x)]$ (Teorema~\ref{teo:clase}). Un ciclu de lungime $5$ are $C_{S_5}(x)=\langle x\rangle$ (ordin $5$), inclus în $A_5$; deci clasa sa din $S_5$ (de $24$ elemente) se sparge în $A_5$ în două clase de câte $60/5=12$. Analog: dublele transpoziții dau o clasă de $15$, iar ciclii de lungime $3$ o clasă de $20$. Cu identitatea, cardinalele claselor sunt $1,15,20,12,12$. Un subgrup normal $N\trianglelefteq A_5$ este o reuniune de clase care conține identitatea, deci $|N|=1+(\text{sumă de termeni din }\{15,20,12,12\})$, iar $|N|$ trebuie să dividă $60$. Verificând, singurele astfel de sume care divid $60$ sunt $|N|=1$ și $|N|=1+15+20+12+12=60$ (celelalte, $16,21,13,25,36,28,33,40,45,48$, nu divid $60$). Deci $N\in\{\{e\},A_5\}$: $A_5$ este simplu. Neabelian: $(1\,2\,3)(3\,4\,5)\neq(3\,4\,5)(1\,2\,3)$.
\end{proof}

\subsubsection*{Grupuri rezolubile}

\begin{definitie}
Un grup $G$ este \emph{rezolubil} (solubil) dacă există un lanț finit de subgrupuri
\[
\{e\}=G_0\trianglelefteq G_1\trianglelefteq\dots\trianglelefteq G_m=G
\]
în care fiecare $G_{i-1}$ este normal în $G_i$, iar toate câturile $G_i/G_{i-1}$ sunt abeliene.
\end{definitie}

\begin{exemplu}
$S_4$ este rezolubil: cu $V=\{e,(1\,2)(3\,4),(1\,3)(2\,4),(1\,4)(2\,3)\}$ (grupul lui Klein) și $A_4$ (permutările pare), avem
\[
\{e\}\trianglelefteq V\trianglelefteq A_4\trianglelefteq S_4,
\]
cu câturile $V\simeq C_2\times C_2$ (abelian), $A_4/V\simeq C_3$ și $S_4/A_4\simeq C_2$ (Exemplul~\ref{ex:fundamental}(d)) --- toate abeliene.
\end{exemplu}

\begin{propozitie}[proprietăți de închidere]\label{prop:rezolubil}
Fie $G$ un grup și $N\trianglelefteq G$. Atunci:
\emph{(i)} dacă $G$ este rezolubil, orice subgrup $H\le G$ este rezolubil;
\emph{(ii)} dacă $G$ este rezolubil, $G/N$ este rezolubil;
\emph{(iii)} dacă $N$ și $G/N$ sunt rezolubile, atunci $G$ este rezolubil.
\end{propozitie}

\begin{proof}[Schiță]
(i) Intersectăm lanțul lui $G$ cu $H$: $H_i=H\cap G_i$; atunci $H_{i-1}\trianglelefteq H_i$ și $H_i/H_{i-1}$ se scufundă în $G_i/G_{i-1}$ (abelian), deci este abelian. (ii) Proiectăm lanțul prin $\pi:G\to G/N$: $\pi(G_i)$ formează un lanț cu câturi factori ai unor grupuri abeliene, deci abeliene. (iii) Concatenăm un lanț al lui $N$ cu preimaginile unui lanț al lui $G/N$ (teorema de corespondență, Teorema~\ref{teo:corespondenta}, asigură normalitatea și, prin teorema fundamentală, câturile corecte).
\end{proof}

\begin{observatie}
\label{obs:simplurezol}Un grup simplu rezolubil este obligatoriu ciclic de ordin prim: lanțul din definiție trebuie să aibă un cât $G/G_{m-1}$ abelian netrivial, forțând $G_{m-1}=\{e\}$ și $G=G_m$ abelian, deci (fiind simplu) $\simeq C_p$ (Lema~\ref{lema:abeliansimplu}).
\end{observatie}

\begin{teorema}[$S_n$ nu este rezolubil pentru $n\ge5$]\label{teo:Snrezol}
Pentru $n\ge5$, grupurile $A_n$ și $S_n$ nu sunt rezolubile.
\end{teorema}

\begin{proof}
$A_5$ este simplu și neabelian (Propoziția~\ref{prop:A5}), deci nu este $\simeq C_p$; prin Observația~\ref{obs:simplurezol}, un grup simplu rezolubil ar fi ciclic de ordin prim --- contradicție, deci $A_5$ nu este rezolubil. Pentru $n\ge5$, $A_5$ se scufundă în $A_n$ și în $S_n$ (permutările pare ale lui $\{1,\dots,5\}$, fixând $\{6,\dots,n\}$), iar un subgrup al unui grup rezolubil este rezolubil (Propoziția~\ref{prop:rezolubil}(i)). Prin urmare nici $A_n$, nici $S_n$ nu sunt rezolubile pentru $n\ge5$.
\end{proof}

\begin{teorema}[Abel--Ruffini]\label{teo:abelruffini}
Ecuația polinomială generală de grad $\ge5$ nu se rezolvă prin radicali.
\end{teorema}

Legătura cu cele de mai sus, prin teoria lui Galois: grupul Galois al polinomului general de grad $n$ este $S_n$, iar rezolvarea prin radicali a ecuației corespunde exact rezolubilității acestui grup. Nerezolubilitatea lui $S_n$ pentru $n\ge5$ (Teorema~\ref{teo:Snrezol}) dă deci concluzia.

\chapter{Inele și corpuri}

\section{Inele: definiții și reguli de calcul}

\begin{definitie}
$(A,+,\cdot)$ este \emph{inel} dacă $(A,+)$ este grup abelian (cu neutrul $0$), înmulțirea este asociativă și distributivă față de adunare: $a(b+c)=ab+ac$, $(a+b)c=ac+bc$. Inelul este \emph{unitar} dacă înmulțirea are element neutru $1$ și \emph{comutativ} dacă înmulțirea este comutativă.
\end{definitie}

\begin{observatie}[cele două convenții; convenția noastră]\label{obs:conventieinel}
În literatură circulă \textbf{două definiții} ale noțiunii de inel:
\begin{itemize}[nosep,leftmargin=1.4em]
\item \emph{fără unitate}: se cer doar axiomele de mai sus (grup abelian, asociativitate, distributivitate). Un astfel de obiect se numește uneori \emph{pseudo-inel} sau, în engleză, \emph{rng} (de la „ring'' fără „i'' de la „identity''). Exemplu: $2\ZZ$ cu operațiile uzuale --- e închis la $+$ și $\cdot$, dar nu conține $1$.
\item \emph{cu unitate}: se cere în plus existența lui $1$ cu $1\cdot a=a\cdot 1=a$.
\end{itemize}
\textbf{În programa românească --- și în tot acest volum --- „inel'' înseamnă întotdeauna \emph{inel unitar}.} Așadar $2\ZZ$ \emph{nu} este inel pentru noi (este doar un ideal al lui $\ZZ$), iar când spunem „subinel'' cerem ca acesta să conțină același $1$. Este o convenție, nu o teoremă: la olimpiadă enunțul o presupune tacit, dar merită știut că altundeva „inel'' poate însemna altceva.
\end{observatie}

\begin{definitie}[morfism de inele; morfism unitar]\label{def:morfinel}
Fie $A,B$ inele. O aplicație $f:A\to B$ se numește \emph{morfism de inele} dacă păstrează ambele operații:
\[
f(x+y)=f(x)+f(y),\qquad f(xy)=f(x)f(y)\qquad\text{pentru orice }x,y\in A.
\]
Morfismul se numește \emph{unitar} dacă, în plus,
\[
f(1_A)=1_B .
\]
Un morfism bijectiv se numește \emph{izomorfism} (inversul lui este automat morfism); un morfism $f:A\to A$ se numește \emph{endomorfism}.
\end{definitie}

\begin{observatie}[de ce $f(1)=1$ trebuie cerut separat]
Din $f(x+y)=f(x)+f(y)$ rezultă întotdeauna $f(0)=0$ și $f(-x)=-f(x)$ (proprietăți de morfism de grupuri aditive), \textbf{dar} $f(1_A)=1_B$ \emph{nu} rezultă din axiome. Contraexemplu: $f:\ZZ\to\ZZ\times\ZZ$, $f(x)=(x,0)$, păstrează adunarea și înmulțirea, însă $f(1)=(1,0)\neq(1,1)=1_{\ZZ\times\ZZ}$. De aceea, \textbf{în acest volum, prin „morfism de inele'' înțelegem implicit morfism unitar} ($f(1)=1$) --- convenția standard în programa românească; ea este esențială, de pildă, la subinelul prim (Propoziția~\ref{prop:subinelprim}) și la faptul că un morfism duce elementele inversabile în elemente inversabile: $f(u^{-1})=f(u)^{-1}$.
\end{observatie}

\begin{exemplu}
$(\ZZ,+,\cdot)$, $(\QQ,+,\cdot)$, $(\RR,+,\cdot)$, $(\CC,+,\cdot)$, $(\ZZ_n,+,\cdot)$ --- comutative; $(M_2(\RR),+,\cdot)$ --- necomutativ; $\ZZ[i]=\{a+bi\mid a,b\in\ZZ\}$ (\emph{întregii lui Gauss}) și $\ZZ[\sqrt d\,]=\{a+b\sqrt d\mid a,b\in\ZZ\}$ --- comutative; inelele de polinoame $A[X]$ (§\ref{sec:polinoame}).
\end{exemplu}

\begin{propozitie}[reguli de calcul]\label{prop:reguli}
Într-un inel $A$: \emph{(i)} $a\cdot 0=0\cdot a=0$; \emph{(ii)} $a(-b)=(-a)b=-(ab)$ și $(-a)(-b)=ab$; \emph{(iii)} $a(b-c)=ab-ac$; \emph{(iv)} distributivitatea se extinde la sume finite.
\end{propozitie}

\begin{proof}
(i) $a\cdot0=a(0+0)=a\cdot0+a\cdot0$; simplificăm în grupul $(A,+)$. (ii) $ab+a(-b)=a(b+(-b))=a\cdot0=0$, deci $a(-b)$ este opusul lui $ab$; restul analog.
\end{proof}

\begin{observatie}
Dacă într-un inel $1=0$, atunci $x=x\cdot 1=x\cdot 0=0$ pentru orice $x$, deci $A=\{0\}$. De aceea la corpuri se cere explicit $1\neq 0$.
\end{observatie}

\begin{definitie}
$U(A)=\{a\in A\mid a\text{ inversabil față de }\cdot\}$ este mulțimea \emph{unităților} inelului $A$.
\end{definitie}

\begin{propozitie}
$(U(A),\cdot)$ este grup, numit \emph{grupul unităților}.
\end{propozitie}

\begin{proof}
$1\in U(A)$; dacă $a,b\in U(A)$, atunci $(ab)(b^{-1}a^{-1})=1=(b^{-1}a^{-1})(ab)$, deci $ab\in U(A)$, iar $a^{-1}\in U(A)$ evident.
\end{proof}

\begin{exemplu}[{$U(\ZZ[i])$}]\label{ex:gauss}
$U(\ZZ)=\{\pm1\}$. În $\ZZ[i]$ definim \emph{norma} $N(a+bi)=a^2+b^2=|a+bi|^2$; ea este multiplicativă: $N(zw)=|zw|^2=N(z)N(w)$. Dacă $uv=1$ în $\ZZ[i]$, atunci $N(u)N(v)=1$ în $\NN$, deci $N(u)=1$, adică $u\in\{\pm1,\pm i\}$; reciproc, acestea sunt inversabile. Așadar $U(\ZZ[i])=\{\pm1,\pm i\}$. (Norma multiplicativă --- unealtă standard pentru inele de tip $\ZZ[\sqrt d\,]$.)
\end{exemplu}

\begin{definitie}[centrul unui inel]\label{def:centruinel}
\emph{Centrul} inelului $A$ este
\[
Z(A)=\{a\in A\mid ax=xa\ \text{pentru orice }x\in A\}.
\]
$A$ este comutativ dacă și numai dacă $Z(A)=A$.
\end{definitie}

\begin{definitie}[subinel]\label{def:subinel}
O submulțime $B\subseteq A$ se numește \emph{subinel} al lui $A$ dacă este stabilă la operațiile lui $A$ și devine ea însăși inel cu operațiile induse; concret (în convenția noastră, cu unitate):
\[
1_A\in B,\qquad x-y\in B\ \text{ și }\ xy\in B\qquad\text{pentru orice }x,y\in B.
\]
(Mulțimea $B$ este nevidă, căci $1_A\in B$; condiția $x-y\in B$ dă $0=1_A-1_A\in B$ și apoi $-y=0-y\in B$, deci $(B,+)$ este subgrup al lui $(A,+)$, prin criteriul de subgrup.) Exemple: $\ZZ\subseteq\QQ\subseteq\RR\subseteq\CC$; $\ZZ[i]\subseteq\CC$; matricele scalare în $M_n(K)$; iar $2\ZZ$ \textbf{nu} este subinel al lui $\ZZ$ (nu conține $1$).
\end{definitie}

\begin{propozitie}\label{prop:centruinel}
$Z(A)$ este un subinel al lui $A$ care conține $1$; în plus, $(Z(A),+)$ este subgrup al lui $(A,+)$. Orice element inversabil din $Z(A)$ are inversul tot în $Z(A)$.
\end{propozitie}

\begin{proof}
$0,1\in Z(A)$. Dacă $a,b\in Z(A)$, atunci pentru orice $x$: $(a-b)x=ax-bx=xa-xb=x(a-b)$ și $(ab)x=a(bx)=a(xb)=(ax)b=(xa)b=x(ab)$, deci $a-b,ab\in Z(A)$. Dacă $a\in Z(A)$ este inversabil, din $ax=xa$ rezultă (înmulțind cu $a^{-1}$ la stânga și la dreapta) $xa^{-1}=a^{-1}x$.
\end{proof}

\begin{lema}[Bézout aditiv pentru centru]\label{lema:bezoutcentru}
Fie $a\in A$ și $n,p\in\ZZ$. Dacă $na\in Z(A)$ și $pa\in Z(A)$ (multipli întregi), atunci $(n,p)\,a\in Z(A)$. În particular, dacă $(n,p)=1$, atunci $a\in Z(A)$.
\end{lema}

\begin{proof}
$(Z(A),+)$ este subgrup al lui $(A,+)$ (Propoziția~\ref{prop:centruinel}). Cu $d=(n,p)$ și $un+vp=d$ (Bézout), avem $da=u(na)+v(pa)\in Z(A)$, căci un subgrup este închis la combinații întregi ale elementelor sale. Este analogul aditiv al Lemei~\ref{lema:bezoutexp}.
\end{proof}

\begin{definitie}
Un element $a\neq0$ este \emph{divizor al lui zero} dacă există $b\neq0$ cu $ab=0$ sau $ba=0$. Un inel comutativ, unitar, cu $1\neq0$ și \textbf{fără} divizori ai lui zero se numește \emph{inel integru} (domeniu de integritate).
\end{definitie}

\begin{propozitie}\label{prop:integru}
\emph{(i)} Într-un inel integru funcționează simplificarea: $ab=ac$ și $a\neq0$ implică $b=c$. \emph{(ii)} Un element inversabil nu este divizor al lui zero (în orice inel).
\end{propozitie}

\begin{proof}
(i) $a(b-c)=0$ și $a\neq 0$ dau $b-c=0$. (ii) Dacă $ab=0$ și $a$ inversabil, atunci $b=a^{-1}(ab)=0$; analog pe partea cealaltă.
\end{proof}

\begin{propozitie}[binomul lui Newton în inele]\label{prop:binom}
Dacă $x,y\in A$ \textbf{comută} ($xy=yx$), atunci pentru orice $n\in\NN$:
\[
(x+y)^n=\sum_{k=0}^{n}\binom{n}{k}x^k y^{n-k},
\]
unde $\binom nk=C_n^k$. (Demonstrație: inducție după $n$, exact ca în $\RR$; comutativitatea este esențială pentru gruparea termenilor.)
\end{propozitie}

\subsection*{Centrul și morfismele inelelor de matrice}

\begin{propozitie}[centrul inelului de matrice]\label{prop:centrumatrice}
Fie $A$ un inel comutativ (unitar) și $n\ge1$. Atunci centrul inelului de matrice este format din matricele scalare:
\[
Z\big(M_n(A)\big)=\{a I_n\mid a\in A\},
\]
iar $a\mapsto aI_n$ este un izomorfism de inele $A\xrightarrow{\ \sim\ }Z\big(M_n(A)\big)$.
\end{propozitie}
\begin{proof}
Matricele scalare comută cu orice matrice (folosind comutativitatea lui $A$), deci $\{aI_n\}\subseteq Z(M_n(A))$. Reciproc, fie $B=(b_{k\ell})\in Z(M_n(A))$ și $E_{ij}$ matricea unitate (cu $1$ pe poziția $(i,j)$, $0$ în rest). Din $BE_{ij}=E_{ij}B$, comparând intrarea $(k,\ell)$:
\[
(BE_{ij})_{k\ell}=b_{ki}\,\delta_{j\ell},\qquad (E_{ij}B)_{k\ell}=\delta_{ik}\,b_{j\ell}.
\]
Pentru $k\neq i$, $\ell=j$: $b_{ki}=0$ --- deci $B$ e diagonală. Pentru $k=i$, $\ell=j$: $b_{ii}=b_{jj}$ --- deci toate intrările diagonale sunt egale, $B=aI_n$ cu $a=b_{11}$. Aplicația $a\mapsto aI_n$ e morfism injectiv de inele (căci $(aI_n)(bI_n)=abI_n$, $I_n\mapsto$ unitatea) și, prin cele de mai sus, surjectiv pe centru.
\end{proof}

\begin{definitie}[ideal la stânga, la dreapta, bilateral]\label{def:ideal}
Fie $A$ un inel. O submulțime $I\subseteq A$ se numește:
\begin{itemize}[nosep,leftmargin=1.4em]
\item \emph{ideal la stânga} dacă $(I,+)$ este subgrup al lui $(A,+)$ și $a x\in I$ pentru orice $a\in A$, $x\in I$;
\item \emph{ideal la dreapta} dacă $(I,+)$ este subgrup al lui $(A,+)$ și $x a\in I$ pentru orice $a\in A$, $x\in I$;
\item \emph{ideal bilateral} (pe scurt: \emph{ideal}) dacă este simultan ideal la stânga și la dreapta, adică $(I,+)\le (A,+)$ și
\[
a x\in I \quad\text{și}\quad x a\in I \qquad \text{pentru orice } a\in A,\ x\in I .
\]
\end{itemize}
Într-un inel \emph{comutativ} cele trei noțiuni coincid. Idealele $\{0\}$ și $A$ se numesc \emph{improprii}; un inel nenul fără alte ideale bilaterale se numește \emph{simplu}.
\end{definitie}

\begin{observatie}
Un ideal este mai mult decât un subinel: cere absorbție la înmulțirea cu \emph{orice} element al inelului, nu doar stabilitate internă. De pildă $\ZZ\subseteq\QQ$ este subinel, dar nu ideal ($\frac12\cdot 1\notin\ZZ$); în schimb $n\ZZ\subseteq\ZZ$ este ideal. Reciproc, un ideal propriu nu conține $1$ (altfel $a=a\cdot 1\in I$ pentru orice $a$, deci $I=A$), deci nu e subinel unitar. Atenție și la centru: $Z(A)$ este subinel (Propoziția~\ref{prop:centruinel}), dar în general \textbf{nu} este ideal.
\end{observatie}

\begin{observatie}[nucleul, câtul, idealul maximal]\label{obs:idealcat}
Idealele bilaterale sunt exact nucleele morfismelor de inele: dacă $f:A\to B$ e morfism, atunci $\Ker f$ e ideal bilateral (din $f(x)=0$ rezultă $f(ax)=f(a)f(x)=0$ și $f(xa)=0$); reciproc, pentru orice ideal bilateral $I$ mulțimea claselor $A/I=\{a+I\}$ devine inel cu operațiile
\[
(a+I)+(b+I)=(a+b)+I,\qquad (a+I)(b+I)=ab+I
\]
(bine definite tocmai datorită absorbției), numit \emph{inelul cât}, iar $\pi:A\to A/I$, $\pi(a)=a+I$, e morfism surjectiv cu $\Ker\pi=I$. Un ideal propriu $M\subsetneq A$ se numește \emph{maximal} dacă nu există ideal $J$ cu $M\subsetneq J\subsetneq A$; pentru $A$ comutativ, $M$ este maximal dacă și numai dacă $A/M$ este corp --- faptul folosit la construcția lui Kronecker ($K[X]/(p)$ cu $p$ ireductibil).
\end{observatie}

\begin{lema}[$M_n(K)$ este inel simplu]\label{lema:mnsimplu}
Fie $K$ un corp și $n\ge1$. Singurele ideale bilaterale ale lui $M_n(K)$ sunt $\{O\}$ și $M_n(K)$.
\end{lema}
\begin{proof}
Fie $I\neq\{O\}$ un ideal bilateral și $A\in I$ cu o intrare $a_{pq}\neq0$. Atunci $E_{ip}\,A\,E_{qj}=a_{pq}E_{ij}\in I$, deci $E_{ij}=a_{pq}^{-1}(a_{pq}E_{ij})\in I$ pentru orice $i,j$; sumând, $I_n\in I$, deci $I=M_n(K)$.
\end{proof}

\begin{propozitie}[morfisme între inele de matrice]\label{prop:morfmatrice}
Fie $K,L$ corpuri comutative și $m,n\in\NN^*$ cu $n>m$. Atunci nu există niciun morfism de inele (unitar, conform convenției noastre) $f:M_n(K)\to M_m(L)$. Mai mult, chiar renunțând la condiția $f(I_n)=I_m$, singura aplicație $f:M_n(K)\to M_m(L)$ care păstrează adunarea și înmulțirea este aplicația nulă.
\end{propozitie}
\begin{proof}
Demonstrăm a doua afirmație, care o implică pe prima (un morfism unitar este nenul, căci $f(I_n)=I_m\neq O$). Fie $f$ o aplicație nenulă care păstrează adunarea și înmulțirea. $\Ker f$ e ideal bilateral al lui $M_n(K)$, iar prin Lema~\ref{lema:mnsimplu} $\Ker f\in\{O,\,M_n(K)\}$; cum $f\neq0$, $\Ker f=\{O\}$, deci $f$ e \emph{injectiv}. Fie $e_{ij}=f(E_{ij})$. Atunci $e_{11},\dots,e_{nn}$ sunt idempotente ($e_{ii}^2=f(E_{ii}^2)=e_{ii}$), \emph{nenule} (injectivitate) și \emph{ortogonale} ($e_{ii}e_{jj}=f(E_{ii}E_{jj})=0$ pentru $i\neq j$). Într-un spațiu vectorial peste $L$, imaginea sumei de idempotente ortogonale este suma directă a imaginilor, deci
\[
\rang\Big(\sum_{i=1}^n e_{ii}\Big)=\sum_{i=1}^n\rang(e_{ii})\ge n
\]
(fiecare $e_{ii}\neq0$ are rang $\ge1$). Dar rangul oricărei matrice din $M_m(L)$ este $\le m$, deci $n\le m$ --- contradicție cu $n>m$.
\end{proof}

\begin{corolar}\label{cor:izomatrice}
Pentru corpuri comutative $K,L$ și $m,n\in\NN^*$: inelele $M_n(K)$ și $M_m(L)$ sunt izomorfe dacă și numai dacă $n=m$ și $K\cong L$.
\end{corolar}
\begin{proof}
$(\Leftarrow)$ e clar. $(\Rightarrow)$ Un izomorfism $M_n(K)\to M_m(L)$ e nenul, deci $n\le m$ (Propoziția~\ref{prop:morfmatrice}); inversul dă $m\le n$, deci $n=m$. Un izomorfism de inele duce centrul în centru, deci $Z(M_n(K))\cong Z(M_m(L))$; prin Propoziția~\ref{prop:centrumatrice}, $K\cong Z(M_n(K))\cong Z(M_m(L))\cong L$.
\end{proof}

\section{Inelul $\ZZ_n$. Euler, Fermat, Wilson}\label{sec:Zn}

Reamintim: $\ZZ_n=\{\cl0,\cl1,\dots,\cl{n-1}\}$, cu $\cl a+\cl b=\cl{a+b}$ și $\cl a\cdot\cl b=\cl{ab}$, operații bine definite (independente de reprezentanți, deoarece $n\mid a-a'$ și $n\mid b-b'$ implică $n\mid ab-a'b'$). $(\ZZ_n,+,\cdot)$ este inel comutativ unitar.

\begin{teorema}[structura lui $\ZZ_n$]\label{teo:Zn}
Fie $\cl a\in\ZZ_n$, $\cl a\neq\cl0$. Atunci:
\emph{(i)} $\cl a$ este inversabil $\iff(a,n)=1$;\quad
\emph{(ii)} $\cl a$ este divizor al lui zero $\iff(a,n)>1$.\\
Prin urmare, în $\ZZ_n$ orice element nenul este fie inversabil, fie divizor al lui zero.
\end{teorema}

\begin{proof}
(i) Dacă $(a,n)=1$, relația lui Bézout dă $ax+ny=1$, deci $\cl a\cdot\cl x=\cl1$. Reciproc, $\cl a\cl b=\cl1$ înseamnă $ab-1=kn$, deci orice divizor comun al lui $a$ și $n$ divide $1$.
(ii) Dacă $d=(a,n)>1$, atunci $\cl a\cdot\cl{\tfrac nd}=\cl{\tfrac ad\cdot n}=\cl 0$, cu $\cl{\tfrac nd}\neq\cl0$. Reciproc, dacă $(a,n)=1$ și $\cl a\cl b=\cl0$, atunci $n\mid ab$, deci $n\mid b$ (lema lui Gauss), adică $\cl b=\cl 0$.
\end{proof}

\begin{corolar}\label{cor:Zpcorp}
$\ZZ_n$ este corp $\iff$ $\ZZ_n$ este inel integru $\iff$ $n$ este prim.
\end{corolar}

\begin{proof}
Dacă $n=ab$ cu $1<a,b<n$, atunci $\cl a\cl b=\cl0$ dă divizori ai lui zero. Dacă $n$ este prim, orice $\cl a\neq\cl0$ are $(a,n)=1$, deci este inversabil.
\end{proof}

\begin{definitie}
$\varphi(n)$ (\emph{indicatorul lui Euler}) este numărul întregilor $k\in\{1,\dots,n\}$ cu $(k,n)=1$. Din Teorema~\ref{teo:Zn}, $|U(\ZZ_n)|=\varphi(n)$.
\end{definitie}

\begin{teorema}[Euler; Fermat]\label{teo:euler}
\emph{(i)} Dacă $(a,n)=1$, atunci $a^{\varphi(n)}\equiv1\pmod n$.
\emph{(ii)} \emph{(Mica teoremă a lui Fermat)} Dacă $p$ este prim și $p\nmid a$, atunci $a^{p-1}\equiv1\pmod p$; pentru orice $a\in\ZZ$, $a^p\equiv a\pmod p$.
\end{teorema}

\begin{proof}
(i) $\cl a$ aparține grupului $U(\ZZ_n)$, de ordin $\varphi(n)$; Corolarul~\ref{cor:lagrange}(ii) dă $\cl a^{\,\varphi(n)}=\cl1$. (ii) este cazul $n=p$, cu $\varphi(p)=p-1$; a doua formă rezultă înmulțind cu $a$, respectiv direct când $p\mid a$.
\end{proof}

\begin{teorema}[lema chineză a resturilor]\label{teo:crt}
Fie $(m,n)=1$ și $a,b\in\ZZ$. Sistemul $x\equiv a\pmod m$, $x\equiv b\pmod n$ are soluție, unică modulo $mn$.
\end{teorema}

\begin{proof}
Căutăm $x=a+mt$; condiția a doua devine $mt\equiv b-a\pmod n$, rezolvabilă deoarece $\cl m\in U(\ZZ_n)$. Unicitatea: dacă $x,y$ sunt soluții, atunci $m\mid x-y$ și $n\mid x-y$, deci $mn\mid x-y$ (căci $(m,n)=1$).
\end{proof}

\begin{corolar}[formula lui Euler]\label{cor:phi}
$\varphi$ este multiplicativă: $(m,n)=1\Rightarrow\varphi(mn)=\varphi(m)\varphi(n)$. Cum $\varphi(p^k)=p^k-p^{k-1}$ (dintre $1,\dots,p^k$ eliminăm cei $p^{k-1}$ multipli de $p$), rezultă
$\varphi(n)=n\prod_{p\mid n}\Bigl(1-\frac1p\Bigr)$.
\end{corolar}

\begin{proof}
Prin lema chineză, $x\mapsto(x\bmod m,\;x\bmod n)$ este o bijecție de la resturile modulo $mn$ la perechile de resturi; iar $(x,mn)=1\iff(x,m)=1$ și $(x,n)=1$. Numărând perechile ,,coprime`` obținem $\varphi(mn)=\varphi(m)\varphi(n)$.
\end{proof}

\begin{propozitie}[ecuația liniară în $\ZZ_n$]\label{prop:liniara}
Ecuația $\cl a\,x=\cl b$ în $\ZZ_n$ are soluții dacă și numai dacă $d=(a,n)$ divide $b$; în acest caz are exact $d$ soluții.
\end{propozitie}

\begin{proof}
Ecuația revine la $ax\equiv b\pmod n$, adică $ax+ny=b$: membrul stâng e mereu divizibil cu $d$, deci condiția e necesară. Dacă $d\mid b$, împărțim prin $d$: $a'x\equiv b'\pmod{n'}$ cu $(a',n')=1$, deci $x\equiv x_0\pmod{n'}$; modulo $n$ soluțiile sunt $x_0,\,x_0+n',\dots,x_0+(d-1)n'$, în număr de $d$.
\end{proof}

\begin{teorema}[Wilson]\label{teo:wilson}
Pentru $n\ge2$: $n$ este prim $\iff(n-1)!\equiv-1\pmod n$.
\end{teorema}

\begin{proof}
($\Rightarrow$) În corpul $\ZZ_p$, ecuația $x^2=\cl1$, adică $(x-\cl1)(x+\cl1)=\cl0$, are exact soluțiile $\pm\cl1$ (nu există divizori ai lui zero). Deci în produsul $\cl1\cdot\cl2\cdots\cl{p-1}$ elementele diferite de $\pm\cl1$ se grupează în perechi $\{x,x^{-1}\}$ cu produs $\cl1$, iar produsul total este $\cl1\cdot(-\cl1)=-\cl1$. (Pentru $p=2,3$ verificare directă.)
($\Leftarrow$) Dacă $n=de$ cu $1<d<n$, atunci $d\mid(n-1)!$ și $d\mid(n-1)!+1$, deci $d\mid1$, contradicție.
\end{proof}

\begin{observatie}
Complementar: dacă $n>4$ este compus, atunci $(n-1)!\equiv0\pmod n$. Într-adevăr, $n=ab$ cu $1<a<b<n$ face ca $a$ și $b$ să apară ca factori distincți în $(n-1)!$; iar dacă $n=p^2$ cu $p>2$, factorii $p$ și $2p$ (ambii $<n$) dau $p^2\mid(n-1)!$.
\end{observatie}

\section{Corpuri. Caracteristică. Morfismul lui Frobenius}

\begin{definitie}
Un inel unitar $K$ cu $1\neq0$ în care orice element nenul este inversabil se numește \emph{corp}; dacă înmulțirea este comutativă, \emph{corp comutativ}. Exemple: $\QQ$, $\RR$, $\CC$, $\ZZ_p$ ($p$ prim), precum și
\[
\QQ(\sqrt d\,)=\{a+b\sqrt d\mid a,b\in\QQ\}\subset\RR\quad(d\in\NN\text{ nu este pătrat perfect}),
\]
unde inversul se obține prin raționalizare: $\dfrac1{a+b\sqrt d}=\dfrac{a-b\sqrt d}{a^2-db^2}$, numitorul fiind nenul pentru $(a,b)\neq(0,0)$ deoarece $\sqrt d\notin\QQ$.
\end{definitie}

\begin{propozitie}\label{prop:corpintegru}
Un corp nu are divizori ai lui zero; în particular, orice corp comutativ este inel integru.
\end{propozitie}

\begin{proof}
Elementele nenule sunt inversabile, deci nu pot fi divizori ai lui zero (Propoziția~\ref{prop:integru}(ii)).
\end{proof}

\begin{teorema}[inel integru finit]\label{teo:integrufinit}
Orice inel integru \textbf{finit} este corp. Mai general, într-un inel finit orice element nenul care nu este divizor al lui zero este inversabil.
\end{teorema}

\begin{proof}
Fie $a\neq0$ nedivizor al lui zero. Funcțiile $x\mapsto ax$ și $x\mapsto xa$ sunt injective ($ax=ay\Rightarrow a(x-y)=0\Rightarrow x=y$), deci, $A$ fiind finit, surjective: există $b,c$ cu $ab=1$ și $ca=1$. Atunci $c=c(ab)=(ca)b=b$, deci $a$ este inversabil.
\end{proof}

\begin{propozitie}[când un inel finit este corp]\label{prop:finitcorp}
Un inel finit $A$ (cu $1\neq0$) este corp dacă și numai dacă nu are elemente nilpotente nenule și ecuația $x^2=x$ are numai soluțiile $0$ și $1$.
\end{propozitie}
\begin{proof}
$(\Rightarrow)$ Într-un corp, $a^k=0$ cu $a\neq0$ ar da $1=a^{-k}a^k=0$, exclus; iar $x^2=x\Rightarrow x(x-1)=0\Rightarrow x\in\{0,1\}$ (fără divizori ai lui zero).
$(\Leftarrow)$ Fie $a\neq0$. Cum $A$ e finit, șirul $a,a^2,a^3,\dots$ se repetă, deci există $N\ge1$, multiplu al perioadei și suficient de mare, cu $a^{2N}=a^N$; atunci $e:=a^N$ e idempotent, deci $e\in\{0,1\}$. Dacă $e=0$, $a$ ar fi nilpotent nenul --- exclus; deci $a^N=1$, adică $a\cdot a^{N-1}=a^{N-1}\cdot a=1$ și $a$ e inversabil. Orice element nenul fiind inversabil, $A$ e corp (finit, deci comutativ prin Wedderburn).
\end{proof}

\begin{propozitie}[endomorfisme ale grupului aditiv]\label{prop:endoaditiv}
Fie $A$ un inel cu $1\neq0$.
\emph{(i)} Dacă orice endomorfism nenul al grupului $(A,+)$ este injectiv, atunci $A$ nu are divizori ai lui zero.
\emph{(ii)} Dacă orice endomorfism nenul al grupului $(A,+)$ este surjectiv, atunci $A$ este corp.
\end{propozitie}
\begin{proof}
Pentru $a\in A$, translațiile $L_a(x)=ax$ și $R_a(x)=xa$ sunt endomorfisme ale lui $(A,+)$ (distributivitate).
(i) Fie $a\neq0$; $L_a$ e nenul ($L_a(1)=a\neq0$), deci injectiv. Din $ab=0=L_a(0)$ rezultă $b=0$: $a$ nu e divizor al lui zero la stânga; analog cu $R_a$ la dreapta.
(ii) Fie $a\neq0$; $L_a$ nenul, deci surjectiv: există $b$ cu $ab=1$. Cum $b\neq0$, $L_b$ e nenul, deci surjectiv: există $c$ cu $bc=1$. Atunci $a=a(bc)=(ab)c=c$, deci $ba=bc=1$: $a$ e inversabil (bilateral). Orice element nenul fiind inversabil, $A$ e corp.
\end{proof}

\begin{propozitie}\label{prop:morfcorp}
\emph{(i)} Orice morfism (unitar) de corpuri este injectiv. \emph{(ii)} Unicul morfism unitar de inele $f:\ZZ\to A$ este $f(k)=k\cdot1_A$. \emph{(iii)} Unicul endomorfism (unitar) al corpului $\QQ$ este identitatea.
\end{propozitie}

\begin{proof}
(i) Dacă $f(a)=0$ cu $a\neq0$, atunci $1=f(1)=f(a)f(a^{-1})=0$, contradicție cu $1\neq0$ în codomeniu. (ii) Din $f(1)=1_A$ și aditivitate, $f(k)=k\cdot1_A$ pentru $k\in\ZZ$; se verifică imediat că aceasta este morfism. (iii) Ca la Observația~\ref{obs:endo}(b): $f(1)=1$ forțează $f(k)=k$, apoi $f(p/q)=p/q$.
\end{proof}

\begin{definitie}
\emph{Caracteristica} $\car A$ a unui inel este cel mai mic $n\ge1$ cu $\underbrace{1+\dots+1}_{n}=n\cdot1=0$, dacă există; altfel $\car A=0$. De exemplu $\car\ZZ_n=n$, $\car\QQ=\car\RR=\car\CC=0$.
\end{definitie}

\begin{propozitie}\label{prop:caract}
\emph{(i)} Într-un inel integru (în particular, într-un corp comutativ), caracteristica este $0$ sau un număr prim. \emph{(ii)} Orice inel finit are caracteristică nenulă; deci un corp finit are caracteristică un prim $p$, și atunci $p\cdot x=0$ pentru orice $x$.
\end{propozitie}

\begin{proof}
(i) Fie $\car A=n\ge1$ și $n=ab$ cu $1\le a,b\le n$. Atunci $(a\cdot1)(b\cdot1)=ab\cdot1=0$ (distributivitate), deci $a\cdot1=0$ sau $b\cdot1=0$; minimalitatea lui $n$ dă $a=n$ sau $b=n$, adică $n$ prim (sau $n=1$, imposibil căci $1\neq0$). (ii) Într-un inel finit, șirul $1,2\cdot1,3\cdot1,\dots$ ia de două ori aceeași valoare: $m\cdot1=k\cdot1$ cu $m<k$ dă $(k-m)\cdot1=0$. Ultima afirmație: $p\cdot x=(p\cdot1)x=0$.
\end{proof}

\begin{propozitie}[subinelul prim]\label{prop:subinelprim}
Fie $A$ un inel unitar și $P(A)=\{k\cdot 1_A\mid k\in\ZZ\}$. Atunci $P(A)$ este un subinel al lui $A$, inclus în orice subinel al lui $A$ (deci este cel mai mic subinel --- \emph{subinelul prim}). Mai mult:
\[
\car(A)=0\iff P(A)\cong\ZZ,\qquad \car(A)=n\in\NN^*\iff P(A)\cong\ZZ_n.
\]
\end{propozitie}
\begin{proof}
Aplicația $\varphi:\ZZ\to A$, $\varphi(k)=k\cdot 1_A$, este morfism unitar de inele; imaginea sa este $P(A)$, care e deci subinel. Orice subinel conține $1_A$, deci conține toți $k\cdot 1_A$, adică $P(A)$: minimalitatea. Nucleul este $\Ker\varphi=\{k\in\ZZ\mid k\cdot 1_A=0\}$, subgrup al lui $(\ZZ,+)$, deci $=d\ZZ$ pentru un $d\ge0$; prin definiția caracteristicii, $d=\car(A)$ (cu $d=0$ pentru caracteristică $0$). Teorema fundamentală de izomorfism dă $P(A)=\Ima\varphi\cong\ZZ/\Ker\varphi$, adică $\ZZ$ dacă $\car(A)=0$ și $\ZZ_n$ dacă $\car(A)=n$.
\end{proof}

\begin{propozitie}[caracteristica și morfismele]\label{prop:carmorf}
\emph{(i)} Fie $A,B$ inele cu $\car(A)=n$, $\car(B)=m$. Dacă există un morfism (unitar) de inele $f:A\to B$, atunci $m\mid n$.
\emph{(ii)} Dacă $K,L$ sunt corpuri și există un morfism $f:K\to L$, atunci $\car(K)=\car(L)$.
\end{propozitie}
\begin{proof}
(i) $f(1_A)=1_B$. Din $n\cdot 1_A=0$: $0=f(n\cdot 1_A)=n\cdot 1_B$, deci $\ord(1_B)\mid n$ în $(B,+)$, adică $m\mid n$ (pentru $n=0$ afirmația e vidă). (ii) Un morfism de corpuri e injectiv (Propoziția~\ref{prop:morfcorp}(i)). Dacă $\car(K)=0$, subinelul prim $\ZZ\hookrightarrow K\hookrightarrow L$ arată $\car(L)=0$. Dacă $\car(K)=p$ prim, atunci $p\cdot 1_L=f(p\cdot 1_K)=0$, deci $\car(L)\mid p$; cum $1_L\neq0$, $\car(L)=p$.
\end{proof}

\begin{teorema}[Wedderburn, 1905]\label{teo:wedderburn}
Orice corp finit este comutativ.
\end{teorema}

Demonstrația uzuală (prin polinoame ciclotomice și ecuația claselor) depășește cadrul fișei; teorema se citează. Structural, pentru fiecare $p^{\,n}$ ($p$ prim) există, până la izomorfism, exact un corp cu $p^{\,n}$ elemente, notat $\mathbb F_{p^n}$: \emph{existența} este demonstrată în §\ref{sec:polcorpfinite} (Observația~\ref{obs:existafpn}), iar \emph{unicitatea} până la izomorfism o admitem fără demonstrație.

\begin{teorema}[cardinalul unui inel de caracteristică primă]\label{teo:cardpn}
Fie $A$ un inel finit cu $\car A=p$ prim. Atunci $|A|=p^{\,n}$ pentru un $n\ge1$.
\end{teorema}

\begin{proof}
Din $\car A=p$ avem $p\cdot x=(p\cdot1)x=0$ pentru orice $x\in A$, deci în grupul abelian $(A,+)$ orice element nenul are ordinul $p$. Definind înmulțirea cu scalari $\cl k\cdot x:=k\cdot x$ (bine definită, căci $p\cdot x=0$), $(A,+)$ devine spațiu vectorial peste corpul $\ZZ_p$. Fiind finit, are dimensiune finită $n$, deci $|A|=|\ZZ_p|^{\,n}=p^{\,n}$.
\end{proof}

\begin{corolar}
Orice corp finit $K$ are $|K|=p^{\,n}$, unde $p=\car K$ (combinând Propoziția~\ref{prop:caract}(ii) cu Teorema~\ref{teo:cardpn}). Notăm un astfel de corp $\mathbb F_{p^n}$.
\end{corolar}

\begin{teorema}[funcții polinomiale peste un corp finit]\label{teo:functiipol}
Fie $K$ un corp finit cu $|K|=q$. Atunci \textbf{orice} funcție $f:K\to K$ este polinomială: există $P\in K[X]$ cu $f(a)=P(a)$ pentru orice $a\in K$. Mai mult, există exact un astfel de $P$ de grad $\le q-1$.
\end{teorema}

\begin{proof}
Existența și unicitatea polinomului de grad $\le q-1$ sunt exact interpolarea Lagrange (Teorema~\ref{teo:lagrangeint}) în cele $q$ puncte distincte ale lui $K$, cu valorile impuse $f(a)$. (Numărătoare: funcțiile $K\to K$ sunt în număr de $q^{q}$, la fel ca polinoamele de grad $\le q-1$; corespondența polinom $\mapsto$ funcție este deci o bijecție.)
\end{proof}

\begin{propozitie}[reciproca: funcțiile polinomiale forțează corp finit]\label{prop:functiipolrecip}
Fie $A$ un inel cu $0\neq1$ în care orice funcție $f:A\to A$ este polinomială (adică $f(x)=P(x)$ pentru un $P\in A[X]$). Atunci:
\emph{(i)} $A$ este finit;
\emph{(ii)} $A$ este corp; dacă în plus $A$ este comutativ, $A$ este un corp finit.
\end{propozitie}
\begin{proof}
(i) Dacă $A$ ar fi infinit, atunci $|A[X]|=|A|$, deci mulțimea \emph{funcțiilor polinomiale} (imaginea aplicației $A[X]\to A^A$, $P\mapsto\widetilde P$) are cardinalul cel mult $|A|$. Dar mulțimea \emph{tuturor} funcțiilor $A^A$ are cardinalul $|A|^{|A|}>|A|$, deci nu orice funcție ar putea fi polinomială --- contradicție. Așadar $A$ este finit.

(ii) Considerăm funcția $e:A\to A$, $e(0)=1$ și $e(x)=0$ pentru $x\neq0$. Prin ipoteză există $P=c_0+c_1X+\dots+c_mX^m\in A[X]$ cu $e(x)=P(x)=c_0+c_1x+\dots+c_mx^m$ pentru orice $x\in A$ (coeficienții se scriu la stânga). Pentru $x=0$ obținem $c_0=e(0)=1$. Fie acum $a\neq0$. Atunci
\[
0=e(a)=P(a)=1+t\,a,\qquad\text{unde } t=c_1+c_2a+\dots+c_ma^{m-1},
\]
deci $(-t)\,a=1$: orice element nenul are un invers la stânga. De aici rezultă că orice element nenul este inversabil: dacă $a\neq0$ și $ba=1$, atunci $b\neq0$ (altfel $1=0$), deci există $c$ cu $cb=1$; atunci $c=c(ba)=(cb)a=a$, adică $ab=cb=1$, iar $b$ este inversul (bilateral) al lui $a$. Așadar $A$ este corp; fiind finit (punctul (i)), este comutativ prin teorema lui Wedderburn. Enunțul este reciproca teoremei directe de mai sus.
\end{proof}

\begin{teorema}[corpuri cu puține automorfisme]\label{teo:autcorp}
Un corp (inel cu diviziune) cu un \textbf{număr finit} de automorfisme este comutativ.
\end{teorema}

\begin{proof}
Fie $D$ un corp cu centrul $Z=Z(D)$; notăm $D^*=D\setminus\{0\}$ și $Z^*=Z\setminus\{0\}$. Presupunem, prin absurd, că $D$ nu este comutativ. Pentru $g\in D^*$, aplicația $i_g(x)=gxg^{-1}$ păstrează suma și produsul, $i_g(1)=1$, iar inversa ei este $i_{g^{-1}}$; deci $i_g$ este automorfism al inelului $D$. Ca în Teorema~\ref{teo:inn}, $g\mapsto i_g$ este morfism de grupuri de la $(D^*,\cdot)$ la $(\Aut(D),\circ)$, cu nucleul $\{g\in D^*\mid gx=xg\ \forall x\}=Z^*$; prin teorema fundamentală de izomorfism, $\Aut(D)$ conține un subgrup izomorf cu $D^*/Z^*$. Rămâne să arătăm că $D^*/Z^*$ este infinit. Presupunem contrariul.

\emph{Dacă $Z$ este finit}, atunci $D^*$ este reuniunea unui număr finit de clase $gZ^*$, fiecare cu $|Z^*|$ elemente, deci $D$ este finit și, prin teorema lui Wedderburn, comutativ: contradicție.

\emph{Dacă $Z$ este infinit}, fixăm $x\in D\setminus Z$ (există, căci $D$ nu este comutativ). Pentru orice $z\in Z$ avem $x+z\neq0$ (altfel $x=-z\in Z$), deci $x+z\in D^*$. Elementele $x+z$, $z\in Z$, sunt în număr infinit, iar clasele $gZ^*$ sunt în număr finit; deci există $z_1\neq z_2$ în $Z$ cu $x+z_1$ și $x+z_2$ în aceeași clasă, adică $x+z_1=\lambda(x+z_2)$ pentru un $\lambda\in Z^*$. Rezultă $(1-\lambda)x=\lambda z_2-z_1$. Dacă $\lambda=1$, atunci $z_1=z_2$, exclus. Altfel $1-\lambda\in Z^*$, iar inversul său este tot central (Propoziția~\ref{prop:centruinel}), deci $x=(1-\lambda)^{-1}(\lambda z_2-z_1)\in Z$: contradicție.

Așadar $D^*/Z^*$ este infinit, deci un corp necomutativ are o infinitate de automorfisme (interioare). Prin contrapunere, un corp cu un număr finit de automorfisme este comutativ.
\end{proof}

\begin{teorema}[morfismul lui Frobenius; ,,visul bobocului``]\label{teo:frobenius}
Fie $A$ un inel \textbf{comutativ} de caracteristică $p$ prim. Atunci pentru orice $x,y\in A$:
\[
(x+y)^p=x^p+y^p,\qquad (x-y)^p=x^p-y^p,
\]
iar $F:A\to A$, $F(x)=x^p$, este morfism de inele.
\end{teorema}

\begin{proof}
Cheia: $p\mid\binom pk$ pentru $1\le k\le p-1$. Într-adevăr, din identitatea $k\binom pk=p\binom{p-1}{k-1}$ rezultă $p\mid k\binom pk$, iar $(p,k)=1$ dă $p\mid\binom pk$. În binomul lui Newton (Propoziția~\ref{prop:binom}) toți termenii intermediari au coeficienți multipli de $p$, deci se anulează ($p\cdot z=0$). A doua identitate: pentru $p$ impar, $(-y)^p=-y^p$; pentru $p=2$, $-1=1$. Multiplicativitatea lui $F$ este imediată din comutativitate.
\end{proof}

\begin{aplicatie}[Fermat, a doua demonstrație]
În $\ZZ_p$: $(\cl a+\cl 1)^p=\cl a^p+\cl 1$, deci prin inducție după $a$ obținem $\cl a^p=\cl a$ pentru orice $\cl a$, adică $a^p\equiv a\pmod p$.
\end{aplicatie}

\section{Elemente nilpotente și idempotente}

\begin{definitie}[nilpotent, idempotent, inel redus, idempotent central]
Fie $A$ un inel. Un element $x\in A$ este:
\begin{itemize}[nosep,leftmargin=1.4em]
\item \emph{nilpotent} dacă $x^n=0$ pentru un $n\ge1$;
\item \emph{idempotent} dacă $x^2=x$.
\end{itemize}
Inelul $A$ se numește \emph{redus} dacă nu are elemente nilpotente nenule --- echivalent, dacă $x^2=0\Rightarrow x=0$. Un idempotent $e$ se numește \emph{central} dacă aparține centrului inelului, adică $e\in Z(A)$ (Definiția~\ref{def:centruinel}): $ex=xe$ pentru orice $x\in A$.
\end{definitie}

\begin{propozitie}\label{prop:nilpotent}
Dacă $x^n=0$, atunci $1-x$ și $1+x$ sunt inversabile, cu
\[
\begin{gathered}
(1-x)^{-1}=1+x+x^2+\dots+x^{n-1},\\
(1+x)^{-1}=1-x+x^2-\dots+(-1)^{n-1}x^{n-1}.
\end{gathered}
\]
\end{propozitie}

\begin{proof}
$(1-x)(1+x+\dots+x^{n-1})=1-x^n=1$, la fel în cealaltă ordine (puterile lui $x$ comută între ele și cu $1$). A doua formulă: înlocuim $x$ cu $-x$.
\end{proof}

\begin{propozitie}\label{prop:idempotent}
\emph{(i)} Dacă $x$ este idempotent, atunci $1-x$ este idempotent. \emph{(ii)} Într-un inel integru, singurii idempotenți sunt $0$ și $1$.
\end{propozitie}

\begin{proof}
(i) $(1-x)^2=1-2x+x^2=1-2x+x=1-x$. (ii) $x^2=x\iff x(x-1)=0$, deci $x=0$ sau $x=1$.
\end{proof}

\begin{lema}[idempotenți centrali în inele reduse]\label{teo:machale}
Într-un inel \textbf{redus} $A$, orice idempotent este central: $e^2=e\Rightarrow e\in Z(A)$.
\end{lema}

\begin{proof}
Fie $e^2=e$ și $x\in A$ arbitrar. Punem $u=ex-exe$. Folosind $e^2=e$:
\[
u=e(x-xe),\qquad (x-xe)e=xe-xe^2=xe-xe=0,
\]
deci $u^2=e(x-xe)\cdot e(x-xe)=e\bigl[(x-xe)e\bigr](x-xe)=0$. Inelul fiind redus, $u=0$, adică $ex=exe$. Simetric, pentru $v=xe-exe=(x-ex)e$ avem $e(x-ex)=ex-e^2x=0$, deci $v^2=(x-ex)\bigl[e(x-ex)\bigr]e=0$, de unde $v=0$ și $xe=exe$. Așadar $ex=exe=xe$, adică $e\in Z(A)$.
\end{proof}

\begin{observatie}
Rezultatul rămâne valabil sub ipoteza mai slabă $x^2=0\Rightarrow x\in Z(A)$ (nu neapărat $x=0$): forma generală Buckley--MacHale. Ideea ,,$(ex-exe)^2=0$`` (respectiv identitatea $xy+yx\in Z(A)$) este motorul multor teoreme de comutativitate --- vezi Teorema~\ref{teo:machalecomut} (MacHale) și ierarhia de la finalul secțiunii.
\end{observatie}

\begin{lema}[puteri într-un $n$-inel]\label{lema:ninel}
Fie $n\ge2$ fixat și $A$ un \emph{$n$-inel}, adică $x^n=x$ pentru orice $x\in A$. Atunci pentru orice $a\in A$ și orice $m\ge1$,
\[
a^{\,m(n-1)+1}=a.
\]
\end{lema}

\begin{proof}
Inducție după $m$. Pentru $m=1$: $a^{(n-1)+1}=a^n=a$. Pasul: presupunând $a^{m(n-1)+1}=a$,
\[
a^{(m+1)(n-1)+1}=a^{m(n-1)+1}\cdot a^{n-1}=a\cdot a^{n-1}=a^n=a.\qedhere
\]
\end{proof}

\begin{propozitie}[inele reduse: coborârea exponentului]\label{prop:redus}
Fie $A$ un inel \textbf{redus} și $t\ge1$ fixat. Dacă $x\in A$ satisface $x^{n+t}=x^n$ pentru un $n\ge1$, atunci
\[
x^{t+1}=x.
\]
\end{propozitie}

\begin{proof}
Fie $m\ge1$ cel mai mic exponent cu $x^{m+t}=x^m$. Presupunem, prin absurd, $m\ge2$ și considerăm
\[
y=x^{m+t-1}-x^{m-1}=x^{m-1}\bigl(x^{t}-1\bigr).
\]
Puterile lui $x$ comută între ele, deci $y^2=x^{2m-2}(x^t-1)^2=x^{2m+2t-2}-2x^{2m+t-2}+x^{2m-2}$. Înmulțind relația $x^{m+t}=x^m$ cu $x^{m-2}$ (posibil, $m\ge2$) obținem $x^{2m+t-2}=x^{2m-2}$; înmulțind-o cu $x^{m+t-2}$ obținem $x^{2m+2t-2}=x^{2m+t-2}$. Prin urmare toți cei trei termeni sunt egali cu $x^{2m-2}$ și
\[
y^2=x^{2m-2}-2x^{2m-2}+x^{2m-2}=0.
\]
Inelul fiind redus, $y=0$, adică $x^{(m-1)+t}=x^{m-1}$ --- contrazice minimalitatea lui $m$. Deci $m=1$, adică $x^{1+t}=x^1$.
\end{proof}

\begin{teorema}[condiția $x^{k+1}=x^k$ dă inel boolean]\label{teo:kk1}
Fie $A$ un inel unitar cu proprietatea: pentru orice $x\in A$ există $k=k(x)\ge1$ cu $x^{k+1}=x^k$. Atunci $A$ este redus, $x^2=x$ pentru orice $x$ (deci $A$ este boolean) și, în consecință, $A$ este comutativ.
\end{teorema}

\begin{proof}
\emph{Redus}: fie $t$ nilpotent. Atunci $1+t$ este inversabil (Propoziția~\ref{prop:nilpotent}); dar orice element inversabil $u$ satisface $u^{k+1}=u^k$, de unde, înmulțind cu $(u^{-1})^k$, $u=1$. Deci $1+t=1$, adică $t=0$: $A$ nu are nilpotenți nenuli.
\emph{Boolean}: aplicăm Propoziția~\ref{prop:redus} cu $t=1$; ipoteza $x^{k+1}=x^k$ este exact $x^{n+t}=x^n$ (cu $n=k$), deci $x^{2}=x$ pentru orice $x$.
\emph{Comutativ}: Teorema~\ref{teo:boole}.
\end{proof}



\begin{teorema}[inele booleene]\label{teo:boole}
Dacă $x^2=x$ pentru \textbf{orice} $x\in A$, atunci $1+1=0$ și $A$ este comutativ.
\end{teorema}

\begin{proof}
Pentru orice $x$: $(x+x)=(x+x)^2=x^2+x^2+x^2+x^2=4x$, deci $2x=0$, adică fiecare element este propriul opus. Apoi
$x+y=(x+y)^2=x^2+xy+yx+y^2=x+xy+yx+y$, de unde $xy+yx=0$, deci $xy=-yx=yx$.
\end{proof}

\begin{corolar}\label{cor:boolefarad}
Un inel boolean fără divizori ai lui zero are cel mult două elemente (este $\{0\}$ sau $\simeq\ZZ_2$).
\end{corolar}

\begin{proof}
$x^2=x\iff x(x-1)=0$; fără divizori ai lui zero, $x=0$ sau $x=1$, deci $A\subseteq\{0,1\}$.
\end{proof}

\begin{propozitie}[cardinalul unui inel boolean]\label{prop:boolcard}
Pe o mulțime finită cu $m$ elemente se poate defini o structură de inel boolean (unitar, $x^2=x$ pentru orice $x$) dacă și numai dacă $m=2^k$ pentru un $k\in\NN^*$.
\end{propozitie}
\begin{proof}
$(\Rightarrow)$ Într-un inel boolean $1+1=0$, deci $\car(A)=2$ și $2\cdot x=0$ pentru orice $x$; atunci $(A,+)$ devine spațiu vectorial peste $\ZZ_2$ (cu $\cl k\cdot x:=k\cdot x$, bine definit). Fiind finit de dimensiune $k$, $|A|=2^k$.
$(\Leftarrow)$ Inelul produs $\ZZ_2^{\,k}$ (operații pe componente) este boolean și are $2^k$ elemente.
\end{proof}

\begin{observatie}
Se mai poate arăta (rezultat mai fin, care se citează) că un \emph{$n$-inel} ($x^n=x$ pentru orice $x$) este boolean dacă este îndeplinită una dintre condițiile: \emph{(i)} există $k>m$ naturale și $p$ prim cu $n=2^k+2^m=p+1$; \emph{(ii)} $x+x=0$ pentru orice $x$ și există $k$ cu $n=2^k+1$.
\end{observatie}

\subsection*{Teoreme de comutativitate}

Următoarele sunt \textbf{teoreme clasice de comutativitate}. Demonstrațiile lor generale depășesc programa și le enunțăm fără demonstrație; cazurile mici sunt însă elementare: pentru $x^2=x$ vezi Teorema~\ref{teo:boole}, iar pentru $x^3=x$ vezi demonstrația de după Teorema~\ref{teo:jacobson}. Fac excepție teorema lui MacHale și varianta ei cubică (Propoziția~\ref{prop:machalegen}), demonstrate complet aici, și cazurile deja tratate. Fiecare enunț îl generalizează, în esență, pe cel dinaintea sa.

\begin{teorema}[MacHale]\label{teo:machalecomut}
Fie $A$ un inel și $f:A\to A$ un endomorfism \textbf{surjectiv} cu proprietatea
\[
x^2-f(x)\in Z(A)\qquad\text{pentru orice }x\in A.
\]
Atunci $A$ este comutativ.
\end{teorema}

\begin{proof}
Notăm $g(x)=x^2-f(x)\in Z(A)$. Cum $f$ este aditiv, $f(x+y)=f(x)+f(y)$, iar $(x+y)^2-x^2-y^2=xy+yx$; prin urmare
\[
g(x+y)-g(x)-g(y)=(x+y)^2-x^2-y^2=xy+yx.
\]
Membrul stâng aparține lui $Z(A)$ (grup aditiv), deci
\[
xy+yx\in Z(A)\qquad\text{pentru orice }x,y\in A.\tag{$\star$}
\]
Fie $s=xy+yx\in Z(A)$; fiind central, $s$ comută cu $x$: din $xs=sx$ obținem
\[
x^2y+xyx=xyx+yx^2\ \Longrightarrow\ x^2y=yx^2\quad(\forall y),
\]
deci $x^2\in Z(A)$ pentru orice $x$ (toate pătratele sunt centrale). Atunci
\[
f(x)=x^2-g(x)\in Z(A)\qquad(\text{diferență de elemente centrale}).
\]
În fine, $f$ fiind \textbf{surjectiv}, orice element al lui $A$ este de forma $f(x)$, deci aparține lui $Z(A)$: rezultă $A=Z(A)$, adică $A$ este comutativ.
\end{proof}

\begin{observatie}
Demonstrația folosește doar \emph{aditivitatea} și \emph{surjectivitatea} lui $f$ (nu și multiplicativitatea). Pasul-cheie $(\star)$ --- ,,$xy+yx$ este central`` --- forțează deja \emph{toate pătratele să fie centrale}; de aici, surjectivitatea transferă centralitatea pe tot inelul. Un caz particular des întâlnit la concurs: $f=1_A$ (identitatea, evident surjectivă), adică ipoteza $x^2-x\in Z(A)$ pentru orice $x$ implică $A$ comutativ (Herstein cu $n=2$).
\end{observatie}

\begin{propozitie}[o variantă cubică a teoremei lui MacHale]\label{prop:machalegen}
Fie $A$ un inel cu proprietatea că $x^3+x^2\in Z(A)$ pentru orice $x\in A$. Atunci $A$ este comutativ. Aceeași concluzie are loc dacă $x^3-x^2\in Z(A)$ pentru orice $x\in A$: cele două ipoteze sunt echivalente, căci înlocuind $x$ cu $-x$ obținem $(-x)^3+(-x)^2=-(x^3-x^2)$, iar $Z(A)$ este subgrup aditiv.
\end{propozitie}
\begin{proof}
Notăm $h(x)=x^3+x^2$. Folosim că $Z(A)$ este un subinel care conține $1$ (Propoziția~\ref{prop:centruinel}), deci este închis la sume și diferențe și conține toți multiplii întregi $k\cdot1$. Calculele de mai jos implică doar puteri ale unui singur element și constante întregi, care comută între ele.

\emph{Pasul 1.} $h(-x)=-x^3+x^2\in Z(A)$, deci $h(x)+h(-x)=2x^2\in Z(A)$.

\emph{Pasul 2.} $h(x+1)=(x+1)^3+(x+1)^2=x^3+4x^2+5x+2$, deci $h(x+1)-h(x)-2=3x^2+5x\in Z(A)$. Scăzând $2x^2$ (Pasul 1), obținem
\[
u(x):=x^2+5x\in Z(A)\qquad\text{pentru orice }x\in A.
\]

\emph{Pasul 3.} $u(x+1)-u(x)=(x^2+7x+6)-(x^2+5x)=2x+6\in Z(A)$, deci $2x\in Z(A)$; apoi $x^2+x=u(x)-2\cdot(2x)\in Z(A)$ pentru orice $x\in A$.

\emph{Pasul 4 (mecanismul lui MacHale).} Aplicând Pasul 3 lui $x$, $y$ și $x+y$,
\[
(x+y)^2+(x+y)-(x^2+x)-(y^2+y)=xy+yx\in Z(A).
\]
În particular $xy+yx$ comută cu $x$: $x^2y+xyx=xyx+yx^2$, adică $x^2y=yx^2$ pentru orice $y$, deci $x^2\in Z(A)$. În fine, $x=(x^2+x)-x^2\in Z(A)$ pentru orice $x$, deci $A$ este comutativ.
\end{proof}

\begin{observatie}
O formulare ,,mai generală``, cu $f(x^3)-g(x^2)\in Z(A)$ (sau $f(x)-g(x^3)\in Z(A)$) pentru endomorfisme surjective \emph{arbitrare} $f,g$ ale grupului $(A,+)$, este \textbf{falsă}. În inelul necomutativ $A$ al matricelor $2\times2$ superior triunghiulare peste $\ZZ_2$ (cu centrul $\{O,I_2\}$), un calcul direct pe cele $8$ elemente arată că $x^3-x^2$ și $x-x^3$ aparțin mulțimii $\{O,E_{12}\}$ pentru orice $x$. Luând o bijecție aditivă $f$ a lui $A$ cu $f(E_{12})=I_2$ (există, $(A,+)$ fiind spațiu vectorial de dimensiune $3$ peste $\ZZ_2$) și $g=f$, obținem $f(x^3)-g(x^2)=f(x^3-x^2)\in\{O,I_2\}\subseteq Z(A)$ și, analog, $f(x)-g(x^3)\in Z(A)$, deși $A$ nu este comutativ.
\end{observatie}

\begin{teorema}[Jacobson]\label{teo:jacobson}
Dacă pentru orice $x\in A$ există un întreg $n=n(x)>1$ cu $x^{\,n(x)}=x$, atunci $A$ este comutativ.
\end{teorema}

Inelele cu această proprietate se numesc \emph{potente}. Cazul $n$ fixat (orice inel cu $x^n=x$ pentru un $n\ge2$ fix este comutativ) conține, pentru $n=2$, Teorema~\ref{teo:boole}, iar pentru $n=3$ se demonstrează elementar, astfel. Fie $A$ cu $x^3=x$ pentru orice $x$. \emph{(1)} $A$ este redus: din $x^2=0$ rezultă $x=x^3=x\cdot x^2=0$. \emph{(2)} Pentru orice $x$, $x^2$ este idempotent ($x^4=x\cdot x^3=x^2$), deci central (Lema~\ref{teo:machale}). \emph{(3)} $(x+1)^2=x^2+2x+1$ este central (ca pătrat), iar $x^2$ și $1$ sunt centrale, deci $2x\in Z(A)$. \emph{(4)} Din $(x+1)^3=x+1$ și $x^3=x$ rezultă $3x^2+3x=0$, deci $3x=-3x^2\in Z(A)$. \emph{(5)} $x=3x-2x\in Z(A)$ pentru orice $x$: $A$ este comutativ. Teorema lui Jacobson generalizează teorema lui Wedderburn.

\begin{teorema}[Herstein]\label{teo:herstein}
Dacă pentru orice $x\in A$ există $n=n(x)>1$ astfel încât $x^{\,n(x)}-x$ să fie central (adică în $Z(A)$), atunci $A$ este comutativ.
\end{teorema}

Este o generalizare a teoremei lui Jacobson: dacă $x^{n(x)}=x$, atunci $x^{n(x)}-x=0\in Z(A)$. Reciproca teoremei este banală (într-un inel comutativ orice element este central), deci condiția lui Herstein caracterizează exact inelele comutative.

\begin{teorema}[Kaplansky]\label{teo:kaplansky}
Fie $D$ un corp (inel cu diviziune, nu neapărat comutativ) cu centrul $Z=Z(D)$.
\emph{(i)} Dacă pentru orice $x\in D$ o putere $x^{\,n(x)}$ ($n(x)\ge1$) aparține lui $Z$, atunci $D$ este comutativ.
\emph{(ii)} Dacă grupul cât $D^*/Z^*$ este de torsiune (orice element are ordin finit), atunci $D$ este comutativ.
\end{teorema}

\begin{teorema}[exponenți de parități diferite]\label{teo:parity}
Dacă pentru orice $x\in A$ există întregi $n(x)>m(x)>1$ \textbf{de parități diferite} cu $x^{\,n(x)}=x^{\,m(x)}$, atunci $A$ este comutativ.
\end{teorema}

Aceasta generalizează Propoziția~\ref{prop:redus} (coborârea exponentului în inele reduse), fără a mai cere ca inelul să fie redus. Ipoteza de paritate este esențială: fără ea există contraexemple necomutative.

\begin{teorema}[Nicholson--Yaqub, 1979]\label{teo:nicholson}
Fie $j,k$ două numere naturale \textbf{prime între ele} și $A$ un inel (sau $G$ un grup) în care puterile de exponent $j$, respectiv $k$, comută două câte două:
\[
x^{\,j}y^{\,j}=y^{\,j}x^{\,j}\quad\text{și}\quad x^{\,k}y^{\,k}=y^{\,k}x^{\,k}\qquad\text{pentru orice }x,y.
\]
Atunci $A$ este comutativ (respectiv $G$ este abelian).
\end{teorema}

Rezultatul face parte din familia teoremelor de comutativitate ,,prin aplicația de ridicare la putere`` (W.\,K.~Nicholson, A.~Yaqub, \emph{Canad. Math. Bull.} \textbf{22} (1979); vezi și H.\,E.~Bell, 1978). Dacă unul dintre $j,k$ este $1$, concluzia este banală; interesul este pentru $j,k\ge2$.

\begin{propozitie}[trucul $1-ab$ / $1-ba$]\label{prop:halmos}
Într-un inel $A$, dacă $1-ab$ este inversabil, atunci $1-ba$ este inversabil și
\[
(1-ba)^{-1}=1+b\,(1-ab)^{-1}a.
\]
\end{propozitie}

\begin{proof}
Fie $u=(1-ab)^{-1}$, deci $u-abu=u-uab=1$. Punem $v=1+bua$ și verificăm:
\[
(1-ba)v=1-ba+bua-b(ab\,u)a=1-ba+b(u-abu)a=1-ba+ba=1,
\]
\[
v(1-ba)=1-ba+bua-bu(ab)a=1-ba+b(u-uab)a=1-ba+ba=1.\qedhere
\]
\end{proof}

\begin{observatie}
Formula pare ,,scoasă din pălărie``, dar provine din dezvoltarea formală $(1-ba)^{-1}=1+ba+baba+\dots=1+b(1+ab+abab+\dots)a$; la concurs se prezintă doar verificarea, care este riguroasă.
\end{observatie}

\subsection*{Corpul $U(A)\cup\{0\}$ (S. Rădulescu, M. Piticaru)}

Grupăm aici câteva rezultate de tip olimpiadă despre inele finite în care unitățile plus $0$ formează un corp; ca și celelalte teoreme grele de mai sus (Wedderburn, Jacobson, Kaplansky), se citează.

\begin{propozitie}\label{prop:rp1}
Fie $A$ un inel integru de caracteristică $p\in\NN^*$ cu un număr \emph{finit} de elemente inversabile. Atunci operațiile lui $A$ induc pe $L=U(A)\cup\{0\}$ o structură de corp.
\end{propozitie}

\begin{propozitie}\label{prop:rp2}
Dacă $A$ este un inel finit astfel încât $\big(U(A)\cup\{0\},+,\cdot\big)$ este corp, atunci $A$ nu are elemente nilpotente nenule.
\end{propozitie}

\begin{propozitie}[S. Rădulescu, M. Piticaru]\label{prop:rp3}
Fie $A$ un inel finit de caracteristică $p\ge3$ (număr impar). Dacă $\big(U(A)\cup\{0\},+,\cdot\big)$ este corp, atunci $A$ este corp.
\end{propozitie}

\begin{propozitie}[S. Rădulescu, M. Piticaru]\label{prop:rp4}
Fie $A$ un inel finit de caracteristică $p\ge3$ (număr impar). Dacă pentru orice $x\in U(A)$ avem $1+x\in U(A)\cup\{0\}$, atunci $A$ este corp.
\end{propozitie}

\section{Polinoame peste corpuri finite}\label{sec:polcorpfinite}

Fixăm un prim $p$. Studiem inelul de polinoame $K[X]$ peste un corp comutativ $K$ de caracteristică $p$ (în particular $K=\ZZ_p$), unde apar fenomene specifice caracteristicii, în care morfismul lui Frobenius (Teorema~\ref{teo:frobenius}) joacă rolul principal. \emph{Convenție}: un polinom $f\in K[X]$ de grad $\ge1$ este \emph{ireductibil} dacă nu se scrie ca produs $f=gh$ cu $\deg g,\deg h\ge1$; teoria generală a lui $K[X]$ (grad, împărțire cu rest, factorizare unică în ireductibile, derivată formală) este dezvoltată în Partea~III și o folosim aici ca atare.

\begin{teorema}[Artin--Schreier]\label{teo:artinschreier}
Fie $K$ un corp comutativ de caracteristică $p$ și $a\in K$. Polinomul \emph{Artin--Schreier}
\[
f=X^p-X-a
\]
este ireductibil în $K[X]$ \textbf{dacă și numai dacă} $a$ nu este de forma $\lambda^p-\lambda$ cu $\lambda\in K$.
\end{teorema}

\begin{proof}
Dacă $a=\lambda^p-\lambda$ cu $\lambda\in K$, atunci $f(\lambda)=0$, deci $(X-\lambda)\mid f$ și $f$ este reductibil.

Reciproc, presupunem $f$ reductibil și arătăm că $a=\lambda^p-\lambda$ pentru un $\lambda\in K$. Observăm mai întâi, prin Frobenius pe inelul $K[X]$ (care are tot caracteristica $p$),
\[
f(X+1)=(X+1)^p-(X+1)-a=(X^p+1)-X-1-a=X^p-X-a=f(X).
\]
Cum $f'=pX^{p-1}-1=-1$, un eventual factor pătratic $g^2\mid f$ ar da $g\mid f'=-1$, imposibil pentru $\deg g\ge1$; deci $f$ este \textbf{liber de pătrate}: $f=g_1g_2\cdots g_r$, cu $g_j$ ireductibile monice \emph{distincte} (Teorema~\ref{teo:descompunere}), iar $r\ge2$ prin reductibilitate.

Substituția $\sigma:h(X)\mapsto h(X+1)$ este un automorfism al lui $K[X]$ care fixează $f$, deci \emph{permută} factorii $\{g_1,\dots,g_r\}$ (păstrându-le gradul). Cum $p\cdot1=0$ în $K$, avem $\sigma^p=\mathrm{id}$ ($h(X+p\cdot1)=h(X)$), deci fiecare orbită a lui $\sigma$ pe mulțimea factorilor are cardinalul $1$ sau $p$.

\emph{Nu pot fi toate orbitele de cardinal $1$.} Dacă $\sigma$ ar fixa un factor $g$ de grad $d$, adică $g(X+1)=g(X)$, comparând coeficientul lui $X^{d-1}$ în cei doi membri ($g=X^d+c_{d-1}X^{d-1}+\cdots$ dă, în $g(X+1)$, coeficientul $d+c_{d-1}$) obținem $d=0$ în $K$, adică $p\mid d$; imposibil pentru $1\le d\le\deg f-1<p$.

Rămâne că există o orbită de cardinal $p$: ea conține $p$ factori distincți de același grad $d\ge1$, contribuind cu $pd$ la $\deg f=p$; deci $d=1$ și acei $p$ factori epuizează $f$. Așadar $f$ se descompune în factori liniari peste $K$; în particular are o rădăcină $\lambda\in K$, iar $f(\lambda)=0$ dă $a=\lambda^p-\lambda$.
\end{proof}

\begin{observatie}
Aplicația $\wp:K\to K$, $\wp(\lambda)=\lambda^p-\lambda$, este morfism al grupului aditiv (din nou Frobenius), numită \emph{operatorul Artin--Schreier}; teorema spune că $X^p-X-a$ este ireductibil exact când $a\notin\Ima\wp$. Pe corpul prim $\ZZ_p$ avem $\wp\equiv0$ (Fermat: $\lambda^p=\lambda$), deci $\Ima\wp=\{0\}$: peste $\ZZ_p$, polinomul $X^p-X-a$ este ireductibil pentru \emph{orice} $a\neq0$ (și reductibil pentru $a=0$, căci $X^p-X=\prod_{c\in\ZZ_p}(X-c)$, iarăși prin Fermat).
\end{observatie}

Pentru rezultatele structurale care urmează folosim o construcție standard: dat un polinom monic ireductibil $f\in\ZZ_p[X]$ de grad $d$, \emph{clasele de resturi ale polinoamelor modulo $f$} formează un inel comutativ, notat $\ZZ_p[X]/(f)$ (analog cu $\ZZ_n=\ZZ/n\ZZ$). Fiecare clasă are un unic reprezentant de grad $<d$, deci inelul are $p^{\,d}$ elemente; iar dacă $f$ este ireductibil, orice rest nenul $r$ este inversabil (din $(r,f)=1$ și Bézout polinomial), deci $\ZZ_p[X]/(f)$ este \textbf{corp} --- un corp cu $p^{\,d}$ elemente, în care clasa $\alpha=\widehat X$ este o rădăcină a lui $f$.

\begin{teorema}[ireductibile și $X^{p^n}-X$]\label{teo:xpnx}
Fie $f\in\ZZ_p[X]$ monic ireductibil de grad $d$. Atunci $f\mid X^{p^{\,d}}-X$. Prin urmare, pentru orice polinom ireductibil $f$ există $n$ (anume $n=\deg f$) cu $f\mid X^{p^{\,n}}-X$.
\end{teorema}

\begin{proof}
Fie $L=\ZZ_p[X]/(f)$, corp cu $p^{\,d}$ elemente, și $\alpha=\widehat X\in L$ rădăcină a lui $f$. Grupul multiplicativ $L^*$ are ordinul $p^{\,d}-1$, deci prin Lagrange $\beta^{\,p^{\,d}-1}=1$ pentru orice $\beta\neq0$; adică $\beta^{\,p^{\,d}}=\beta$ pentru \emph{orice} $\beta\in L$ (inclusiv $\beta=0$). În particular $\alpha^{\,p^{\,d}}=\alpha$, deci $\alpha$ este rădăcină a lui $X^{p^{\,d}}-X$. Dar evaluarea $g\mapsto g(\alpha)$ este morfismul de trecere la cât $\ZZ_p[X]\to L$, cu nucleul $(f)$; din $\bigl(X^{p^{\,d}}-X\bigr)(\alpha)=0$ rezultă $X^{p^{\,d}}-X\in(f)$, adică $f\mid X^{p^{\,d}}-X$.
\end{proof}

\begin{propozitie}[subcorpurile unui corp finit]\label{prop:subcorpuri}
Fie $p$ prim, $n\ge1$ și $K=\mathbb F_{p^n}$ un corp finit cu $p^n$ elemente. Pentru fiecare divizor $d\mid n$, $K$ conține \emph{exact un} subcorp cu $p^d$ elemente, anume
\[
K_d=\{x\in K\mid x^{p^d}=x\},
\]
iar acestea sunt \emph{toate} subcorpurile lui $K$.
\end{propozitie}
\begin{proof}
\emph{Orice subcorp are ordinul $p^d$ cu $d\mid n$:} un subcorp $L$ are caracteristica $p$, deci $|L|=p^d$; iar $K$ e spațiu vectorial peste $L$, deci $p^n=(p^d)^{\dim_L K}$, de unde $d\mid n$.

\emph{Existența și forma:} din $d\mid n$ rezultă $p^d-1\mid p^n-1$, deci $X^{p^d-1}-1\mid X^{p^n-1}-1$ și, înmulțind cu $X$, $X^{p^d}-X\mid X^{p^n}-X$. Cele $p^n$ elemente ale lui $K$ sunt exact rădăcinile lui $X^{p^n}-X$ (grupul $K^*$ are ordin $p^n-1$, deci $x^{p^n}=x$ pentru orice $x$). Prin urmare $X^{p^d}-X$ are toate cele $p^d$ rădăcini distincte în $K$, iar mulțimea lor $K_d=\{x:x^{p^d}=x\}$ este un subcorp (mulțimea punctelor fixe ale iteratei $d$-a a morfismului lui Frobenius $x\mapsto x^p$ este închisă la $+,\cdot$ și la inverse). Deci $|K_d|=p^d$.

\emph{Unicitatea:} orice subcorp $L$ cu $p^d$ elemente satisface $x^{p^d}=x$ pentru orice $x\in L$ (grupul $L^*$ are ordin $p^d-1$), deci $L\subseteq K_d$; egalitatea cardinalelor dă $L=K_d$.
\end{proof}

\begin{lema}[factorii ireductibili ai lui $X^{p^n}-X$]\label{lema:ireddivide}
Fie $n\ge1$ și $f\in\ZZ_p[X]$ monic ireductibil de grad $d$. Atunci $f\mid X^{p^{\,n}}-X$ dacă și numai dacă $d\mid n$.
\end{lema}

\begin{proof}
($\Leftarrow$) Prin Teorema~\ref{teo:xpnx}, $f\mid X^{p^{\,d}}-X$. Din $d\mid n$ rezultă $p^d-1\mid p^n-1$ (căci $p^n-1=(p^d)^{n/d}-1$), deci $X^{p^d-1}-1\mid X^{p^n-1}-1$ și, înmulțind cu $X$, $X^{p^{\,d}}-X\mid X^{p^{\,n}}-X$.

($\Rightarrow$) Fie $L=\ZZ_p[X]/(f)$, corp comutativ cu $p^{\,d}$ elemente, și $\alpha=\widehat X$; din $f\mid X^{p^{\,n}}-X$ rezultă $\alpha^{p^n}=\alpha$. Morfismul lui Frobenius $F(\beta)=\beta^p$ este morfism de inele al lui $L$ (Teorema~\ref{teo:frobenius}), deci și $F^n(\beta)=\beta^{p^n}$ este; prin urmare mulțimea $S=\{\beta\in L\mid\beta^{p^n}=\beta\}$ este închisă la adunare și înmulțire. Ea conține clasele constantelor din $\ZZ_p$ (Fermat: $c^p=c$) și pe $\alpha$; cum orice element al lui $L$ este de forma $c_0+c_1\alpha+\dots+c_{d-1}\alpha^{d-1}$ cu $c_i\in\ZZ_p$, rezultă $S=L$, adică $F^n=\mathrm{id}_L$. Pe de altă parte, $\beta^{p^d}=\beta$ pentru orice $\beta\in L$ (grupul $L^*$ are ordinul $p^d-1$), adică $F^d=\mathrm{id}_L$. Deci $F$ este o bijecție a lui $L$, iar ordinul ei în grupul bijecțiilor lui $L$ divide atât $d$, cât și $n$, deci divide $g=(d,n)$: $F^g=\mathrm{id}_L$. Așadar toate cele $p^d$ elemente ale lui $L$ sunt rădăcini ale polinomului $X^{p^g}-X$, care are cel mult $p^g$ rădăcini în corpul $L$ (Teorema~\ref{teo:numarradacini}). Rezultă $p^d\le p^g$, deci $d\le g$, adică $d=g\mid n$.
\end{proof}

\begin{teorema}[existența ireductibilelor de orice grad]\label{teo:existaired}
Pentru $d\ge1$ notăm $I(d)$ numărul polinoamelor monice ireductibile de grad $d$ din $\ZZ_p[X]$. Atunci, pentru orice $n\ge1$,
\[
p^{\,n}=\sum_{d\mid n}d\cdot I(d),
\]
și $I(n)\ge1$: există un polinom ireductibil de grad $n$ în $\ZZ_p[X]$.
\end{teorema}

\begin{proof}
Polinomul $h=X^{p^{\,n}}-X$ este monic și \emph{liber de pătrate}: derivata sa este $p^{\,n}X^{p^n-1}-1=-1$, iar dacă $h=g^2k$ cu $\deg g\ge1$, atunci $h'=g\,(2g'k+gk')$ ar fi divizibil cu $g$, imposibil pentru $h'=-1$. Prin descompunerea în factori ireductibili (Teorema~\ref{teo:descompunere}), $h$ este produsul unor polinoame monice ireductibile distincte, iar prin Lema~\ref{lema:ireddivide} acestea sunt exact polinoamele monice ireductibile de grad $d\mid n$. Comparând gradele, $p^{\,n}=\sum_{d\mid n}d\cdot I(d)$.

Aplicând formula pentru fiecare $d$ obținem $d\cdot I(d)\le p^{\,d}$. Divizorii $d<n$ ai lui $n$ satisfac $d\le n/2$, deci, cu $m=\lfloor n/2\rfloor$,
\[
n\cdot I(n)=p^{\,n}-\sum_{\substack{d\mid n\\ d<n}}d\cdot I(d)\ \ge\ p^{\,n}-\sum_{k=1}^{m}p^{\,k}=p^{\,n}-\frac{p^{\,m+1}-p}{p-1}\ >\ p^{\,n}-p^{\,m+1}\ \ge\ 0,
\]
unde am folosit $\frac{p^{m+1}-p}{p-1}\le p^{m+1}-p<p^{m+1}$ (căci $p-1\ge1$) și $m+1\le n$. Așadar $I(n)\ge1$.
\end{proof}

\begin{observatie}\label{obs:existafpn}
\emph{(a)} Dacă $f\in\ZZ_p[X]$ este ireductibil de grad $n$ (există, prin Teorema~\ref{teo:existaired}), atunci $\ZZ_p[X]/(f)$ este un corp cu $p^{\,n}$ elemente (construcția de mai sus). Aceasta demonstrează \emph{existența} corpului $\mathbb F_{p^n}$ pentru orice $p$ prim și orice $n\ge1$.
\emph{(b)} Aplicând inversiunea Möbius (Propoziția~\ref{prop:mobius}(ii)) relației $p^{\,n}=\sum_{d\mid n}d\cdot I(d)$, cu $M_d=d\cdot I(d)$ și $N_n=p^{\,n}$, obținem \emph{formula lui Gauss}
\[
I(n)=\frac1n\sum_{d\mid n}\mu(d)\,p^{\,n/d}.
\]
De exemplu, peste $\ZZ_2$ există $I(4)=\frac14(2^4-2^2)=3$ polinoame ireductibile de grad $4$, iar peste $\ZZ_3$ există $I(2)=\frac12(3^2-3)=3$ polinoame monice ireductibile de grad $2$.
\end{observatie}

\chapter{Polinoame peste un corp}

\section{Polinoame: grad, împărțire cu rest, divizibilitate}\label{sec:polinoame}

Fie $A$ un inel comutativ (de regulă un corp $K$). $A[X]$ este inelul polinoamelor cu coeficienți în $A$: $f=a_nX^n+\dots+a_1X+a_0$. Dacă $a_n\neq0$, \emph{gradul} este $\deg f=n$, $a_n$ este \emph{coeficientul dominant} ($f$ este \emph{unitar} sau \emph{monic} dacă $a_n=1$), iar $a_0$ este \emph{termenul liber}. Convenim $\deg 0=-\infty$.

\begin{propozitie}[gradele]\label{prop:grade}
\emph{(i)} $\deg(f+g)\le\max(\deg f,\deg g)$, cu egalitate dacă $\deg f\neq\deg g$. \emph{(ii)} Dacă $A$ este integru, $\deg(fg)=\deg f+\deg g$. \emph{(iii)} Dacă $A$ este integru, atunci $A[X]$ este integru și $U(A[X])=U(A)$; în particular, pentru un corp $K$, $U(K[X])=K^*$ (constantele nenule).
\end{propozitie}

\begin{proof}
(ii) Coeficientul lui $X^{\deg f+\deg g}$ în $fg$ este produsul coeficienților dominanți, nenul. (iii) $fg=0$ cu $f,g\neq0$ ar contrazice (ii); iar $fg=1$ forțează $\deg f=\deg g=0$, deci $f,g\in U(A)$.
\end{proof}

\begin{observatie}[{construcția riguroasă a lui $R[X]$}]
Inelul $R[X]$ se poate construi \emph{formal} ca mulțimea șirurilor $(a_0,a_1,a_2,\dots)$ de elemente din $R$ cu doar un număr finit de termeni nenuli, cu adunarea pe componente și înmulțirea dată de produsul Cauchy
\[
(a_i)\cdot(b_j)=(c_k),\qquad c_k=\sum_{i+j=k}a_ib_j.
\]
Notând $X=(0,1,0,0,\dots)$ și identificând $a\in R$ cu $(a,0,0,\dots)$, orice șir se scrie unic $(a_0,\dots,a_n,0,\dots)=\sum_{i=0}^n a_iX^i$; astfel „expresia'' $\sum a_iX^i$ capătă sens riguros ca element al acestui inel.
\end{observatie}

\begin{definitie}[ordinul unui polinom]
Pentru $f=\sum_{i=0}^n a_iX^i\in R[X]$ nenul, \emph{ordinul} $\ord f$ este cel mai mic $j$ cu $a_j\neq0$ (cel mai mic exponent care apare efectiv). Convenim $\ord 0=+\infty$.
\end{definitie}

\begin{propozitie}[ordinul: sumă și produs]\label{prop:ordinpol}
Pentru $f,g\in R[X]$:
\[
\ord(f+g)\ge\min\{\ord f,\ord g\},\qquad \ord(fg)\ge\ord f+\ord g.
\]
Dacă $R$ este domeniu de integritate, atunci $\ord(fg)=\ord f+\ord g$.
\end{propozitie}
\begin{proof}
Prima inegalitate: termenul de exponent $<\min\{\ord f,\ord g\}$ e nul în ambele, deci în sumă. Pentru produs: coeficientul lui $X^k$ în $fg$ este $\sum_{i+j=k}a_ib_j$, care e nul pentru $k<\ord f+\ord g$ (orice pereche are $i<\ord f$ sau $j<\ord g$). Când $R$ e domeniu, coeficientul lui $X^{\ord f+\ord g}$ este exact $a_{\ord f}\,b_{\ord g}\neq0$, deci egalitate. (Rezultat simetric celui pentru grad de mai sus.)
\end{proof}

\begin{observatie}[polinom vs.\ funcție polinomială]\label{obs:functie}
Polinomul $f$ și funcția $\tilde f:A\to A$, $\tilde f(x)=f(x)$, sunt obiecte diferite. Peste un corp \textbf{finit} ele nu se determină reciproc: în $\ZZ_p[X]$, polinomul $X^p-X$ este nenul, dar funcția sa este identic nulă (Fermat). Peste un inel integru \textbf{infinit}, egalitatea funcțiilor implică egalitatea polinoamelor (Corolarul~\ref{cor:identitate}).
\end{observatie}

\begin{teorema}[împărțirea cu rest]\label{teo:impartire}
Fie $K$ un corp comutativ și $f,g\in K[X]$, $g\neq0$. Există și sunt \textbf{unice} $q,r\in K[X]$ cu
\[
f=gq+r,\qquad \deg r<\deg g.
\]
\end{teorema}

\begin{proof}
\emph{Unicitatea}: din $gq_1+r_1=gq_2+r_2$ rezultă $g(q_1-q_2)=r_2-r_1$; dacă $q_1\neq q_2$, membrul stâng are grad $\ge\deg g$, cel drept $<\deg g$ --- contradicție. Deci $q_1=q_2$, apoi $r_1=r_2$.
\emph{Existența}, prin inducție după $\deg f$: dacă $\deg f<\deg g$, luăm $q=0$, $r=f$. Altfel, cu $a,b$ coeficienții dominanți ai lui $f,g$ și $n=\deg f\ge m=\deg g$, polinomul $f_1=f-\frac ab X^{\,n-m} g$ are grad $<n$; aplicăm ipoteza de inducție lui $f_1$ și recompunem.
\end{proof}

\begin{observatie}
Aceeași demonstrație arată că peste \textbf{orice} inel comutativ $A$ împărțirea cu rest funcționează dacă împărțitorul $g$ este \emph{unitar} (sau are coeficient dominant inversabil) --- util pentru $\ZZ[X]$.
\end{observatie}

\begin{definitie}
$g\mid f$ dacă $f=gq$ pentru un $q\in K[X]$. Polinoamele $f$ și $cf$ ($c\in K^*$) se numesc \emph{asociate}. Cel mai mare divizor comun $(f,g)$ este divizorul comun de grad maxim, ales unitar (pentru unicitate).
\end{definitie}

\begin{teorema}[algoritmul lui Euclid; Bézout]\label{teo:bezout}
Pentru $f,g\in K[X]$, nu ambele nule, c.m.m.d.c.\ $d=(f,g)$ există, se obține ca ultimul rest nenul din algoritmul lui Euclid (împărțiri cu rest succesive) și există $u,v\in K[X]$ cu
\[
d=uf+vg.
\]
\end{teorema}

\begin{proof}[Schiță]
Ca în $\ZZ$: dacă $f=gq+r$, atunci divizorii comuni ai perechii $(f,g)$ coincid cu cei ai perechii $(g,r)$; gradele resturilor scad strict, deci algoritmul se oprește. Scriind fiecare rest ca o combinație $u_if+v_ig$ (recurență), ultimul rest nenul are forma cerută.
\end{proof}

\begin{definitie}\label{def:ireductibil}
$f\in K[X]$ cu $\deg f\ge1$ este \emph{ireductibil peste $K$} dacă nu se poate scrie $f=gh$ cu $g,h\in K[X]$, $\deg g,\deg h\ge1$; altfel este \emph{reductibil}.
\end{definitie}

\begin{lema}[Euclid pentru polinoame]\label{lema:euclidpol}
Dacă $p\in K[X]$ este ireductibil și $p\mid fg$, atunci $p\mid f$ sau $p\mid g$.
\end{lema}

\begin{proof}
Dacă $p\nmid f$, atunci $(p,f)=1$ (divizorii lui $p$ sunt constante sau asociați ai lui $p$), deci $1=up+vf$ (Bézout). Atunci $g=upg+v(fg)$ este divizibil cu $p$.
\end{proof}

\begin{teorema}[descompunerea în factori ireductibili]\label{teo:descompunere}
Orice $f\in K[X]$ cu $\deg f\ge1$ se scrie ca produs de polinoame ireductibile peste $K$; scrierea este unică până la ordinea factorilor și înmulțirea cu constante nenule.
\end{teorema}

\begin{proof}[Schiță]
Existența: inducție tare după grad (dacă $f$ nu e ireductibil, se sparge în factori de grade mai mici). Unicitatea: din două descompuneri, Lema~\ref{lema:euclidpol} arată că fiecare factor ireductibil dintr-o parte divide (deci este asociat cu) un factor din cealaltă; simplificăm și repetăm --- exact ca la teorema fundamentală a aritmeticii.
\end{proof}

\section{Rădăcini. Bézout. Derivata formală}

\begin{teorema}[teorema restului și a factorului]\label{teo:rest}
Fie $f\in A[X]$, $A$ comutativ, și $a\in A$. Restul împărțirii lui $f$ la $X-a$ este $f(a)$:
\[
f=(X-a)\,q+f(a).
\]
În particular \emph{(Bézout)}: $a$ este rădăcină a lui $f$ $\iff(X-a)\mid f$.
\end{teorema}

\begin{proof}
Împărțim la polinomul unitar $X-a$: $f=(X-a)q+r$ cu $\deg r\le0$, deci $r$ constant; punând $X=a$ obținem $r=f(a)$.
\end{proof}

\begin{observatie}
\emph{Schema lui Horner} calculează simultan câtul $q$ și valoarea $f(a)$ prin recurența $b_{n-1}=a_n$, $b_{k-1}=a_k+ab_k$; este metoda rapidă de testat rădăcini la concurs.
\end{observatie}

\begin{definitie}
$a$ este rădăcină cu \emph{ordinul de multiplicitate} $m\ge1$ pentru $f$ dacă $(X-a)^m\mid f$ și $(X-a)^{m+1}\nmid f$; echivalent, $f=(X-a)^m g$ cu $g(a)\neq0$.
\end{definitie}

\begin{teorema}[numărul rădăcinilor]\label{teo:numarradacini}
Fie $A$ un inel \textbf{integru} și $f\in A[X]$, $f\neq0$, $\deg f=n$. Atunci $f$ are cel mult $n$ rădăcini în $A$, \emph{numărate cu multiplicități}.
\end{teorema}

\begin{proof}
Inducție după $n$. Dacă $f$ nu are rădăcini, gata. Altfel, fie $a$ o rădăcină de multiplicitate $m$: $f=(X-a)^mg$, $g(a)\neq0$, $\deg g=n-m$. Orice altă rădăcină $b\neq a$ dă $0=(b-a)^m g(b)$; cum $A$ este integru și $b-a\neq0$, rezultă $g(b)=0$, cu aceeași multiplicitate în $g$ ca în $f$ (același argument de integritate). Prin ipoteza de inducție, $g$ are cel mult $n-m$ rădăcini cu multiplicități, deci $f$ are cel mult $m+(n-m)=n$.
\end{proof}

\begin{exemplu}[integritatea este esențială]
În $\ZZ_8[X]$, polinomul $X^2-\cl1$ de gradul $2$ are \textbf{patru} rădăcini: $\cl1,\cl3,\cl5,\cl7$.
\end{exemplu}

\begin{corolar}[principiul identității]\label{cor:identitate}
Fie $A$ integru. \emph{(i)} Dacă $\deg f,\deg g\le n$ și $f,g$ iau valori egale în $n+1$ puncte distincte, atunci $f=g$. \emph{(ii)} Dacă $A$ este infinit și $\tilde f=\tilde g$ ca funcții, atunci $f=g$ ca polinoame.
\end{corolar}

\begin{proof}
$f-g$ are grad $\le n$ și cel puțin $n+1$ rădăcini (respectiv o infinitate), deci este polinomul nul.
\end{proof}

\begin{definitie}
\emph{Derivata formală} a lui $f=\sum_{k=0}^{n}a_kX^k$ este $f'=\sum_{k=1}^{n}ka_kX^{k-1}$ --- definiție pur algebrică, fără limite.
\end{definitie}

\begin{propozitie}\label{prop:derivata}
$(f+g)'=f'+g'$, $(cf)'=cf'$, $(fg)'=f'g+fg'$ și $\bigl((X-a)^m\bigr)'=m(X-a)^{m-1}$.
\end{propozitie}

\begin{proof}[Schiță]
Primele două sunt imediate. Regula produsului se verifică pe monoame $f=X^i$, $g=X^j$ (ambii membri dau $(i+j)X^{i+j-1}$) și se extinde prin biliniaritate; ultima rezultă prin inducție din regula produsului.
\end{proof}

\begin{teorema}[criteriul rădăcinii multiple]\label{teo:multipla}
Fie $K$ corp comutativ, $f\in K[X]$, $a\in K$. Atunci $a$ este rădăcină multiplă (de ordin $\ge2$) a lui $f$ $\iff f(a)=f'(a)=0$.
\end{teorema}

\begin{proof}
Dacă $f(a)=0$, scriem $f=(X-a)g$; atunci $f'=g+(X-a)g'$, deci $f'(a)=g(a)$. Rădăcina $a$ este multiplă $\iff g(a)=0\iff f'(a)=0$.
\end{proof}

\begin{teorema}[multiplicitatea prin derivate succesive]\label{teo:derivatesucc}
Fie $f\in K[X]$ cu $K\in\{\QQ,\RR,\CC\}$. Atunci $a$ are ordinul de multiplicitate \textbf{exact} $m$ $\iff$
\[
f(a)=f'(a)=\dots=f^{(m-1)}(a)=0\quad\text{și}\quad f^{(m)}(a)\neq0.
\]
\end{teorema}

\begin{proof}
Cheia: dacă $f=(X-a)^mg$ cu $g(a)\neq0$ și $m\ge1$, atunci
$f'=(X-a)^{m-1}\bigl[mg+(X-a)g'\bigr]$, iar paranteza ia în $a$ valoarea $mg(a)\neq0$ (aici folosim $\car K=0$). Deci multiplicitatea lui $a$ în $f'$ este exact $m-1$; o inducție încheie demonstrația în ambele sensuri.
\end{proof}

\begin{observatie}
\textbf{Atenție în caracteristică $p$}: în $\ZZ_p[X]$, $f=X^p$ are $f'=pX^{p-1}=0$, deci criteriul cu derivate succesive eșuează; Teorema~\ref{teo:multipla} (cu $f$ și $f'$) rămâne însă valabilă peste orice corp.
\end{observatie}

\section{Relațiile lui Viète și sumele lui Newton}

\begin{teorema}[Viète]\label{teo:viete}
Fie $f=a_nX^n+\dots+a_0\in\CC[X]$, $a_n\neq0$, cu rădăcinile $x_1,\dots,x_n$ (cu multiplicități). Notând sumele simetrice fundamentale $e_k=\sum_{i_1<\dots<i_k}x_{i_1}\cdots x_{i_k}$, avem
\[
e_k=(-1)^k\,\frac{a_{n-k}}{a_n},\qquad k=1,\dots,n.
\]
\end{teorema}

\begin{proof}
Din descompunerea $f=a_n(X-x_1)\cdots(X-x_n)$ (Teorema~\ref{teo:tfa}), dezvoltăm produsul: coeficientul lui $X^{n-k}$ este $a_n(-1)^ke_k$; identificăm cu $a_{n-k}$.
\end{proof}

\begin{exemplu}
$n=2$: $x_1+x_2=-\frac{a_1}{a_2}$, $x_1x_2=\frac{a_0}{a_2}$.\quad
$n=3$: pentru $aX^3+bX^2+cX+d$,
\[
x_1+x_2+x_3=-\frac ba,\qquad x_1x_2+x_2x_3+x_3x_1=\frac ca,\qquad x_1x_2x_3=-\frac da.
\]
\end{exemplu}

\begin{observatie}[reciproca]
Numerele $x_1,\dots,x_n$ cu sumele simetrice $e_1,\dots,e_n$ date sunt exact rădăcinile polinomului
$X^n-e_1X^{n-1}+e_2X^{n-2}-\dots+(-1)^ne_n$. Aceasta transformă sisteme simetrice în ecuații polinomiale --- tipar frecvent la concurs.
\end{observatie}

\begin{teorema}[sumele de puteri]\label{teo:newton}
Fie $f=X^n+c_{n-1}X^{n-1}+\dots+c_0$ unitar, cu rădăcinile $x_1,\dots,x_n$, și $s_k=x_1^k+\dots+x_n^k$ (cu $s_0=n$). Atunci:
\emph{(i)} pentru orice $k\ge n$,
\[
s_k+c_{n-1}s_{k-1}+c_{n-2}s_{k-2}+\dots+c_0\,s_{k-n}=0;
\]
\emph{(ii)} \emph{(identitățile lui Newton, admise)} pentru $1\le k\le n$,
\[
s_k+c_{n-1}s_{k-1}+\dots+c_{n-k+1}s_1+k\,c_{n-k}=0.
\]
\end{teorema}

\begin{proof}[Demonstrație pentru (i)]
Fiecare rădăcină satisface $x_i^n+c_{n-1}x_i^{n-1}+\dots+c_0=0$. Înmulțim cu $x_i^{\,k-n}$ și sumăm după $i$. (Identitățile (ii) se verifică direct pentru $n$ mic; de exemplu $k=1$: $s_1+c_{n-1}=0$, adică Viète.)
\end{proof}

\begin{aplicatie}[identitatea celor trei cuburi]\label{apl:treicuburi}
Pentru $n=k=3$, cu $c_2=-e_1$, $c_1=e_2$, $c_0=-e_3$, relația (ii) dă $s_3=e_1s_2-e_2s_1+3e_3$. Cum $s_2=e_1^2-2e_2$, rezultă
\[
x^3+y^3+z^3-3xyz=(x+y+z)\bigl(x^2+y^2+z^2-xy-yz-zx\bigr).
\]
\end{aplicatie}

\begin{definitie}[inelul de polinoame în mai multe variabile]
Inelul polinoamelor în $n$ nedeterminate cu coeficienți în $R$ se definește inductiv prin
\[
R[X_1,\dots,X_{n+1}]:=\big(R[X_1,\dots,X_n]\big)[X_{n+1}].
\]
Elementele sale se scriu în mod unic ca sume finite $\sum a_{i_1\dots i_n}X_1^{i_1}\cdots X_n^{i_n}$, cu $a_{i_1\dots i_n}\in R$.
\end{definitie}

\begin{definitie}[polinom simetric]
$f\in R[X_1,\dots,X_n]$ se numește \emph{simetric} dacă $f(X_{\sigma(1)},\dots,X_{\sigma(n)})=f(X_1,\dots,X_n)$ pentru orice permutare $\sigma\in S_n$. Polinoamele simetrice fundamentale sunt $e_1=\sum X_i$, $e_2=\sum_{i<j}X_iX_j$, \dots, $e_n=X_1\cdots X_n$ (notația $s_k$ este rezervată sumelor de puteri, Teorema~\ref{teo:newton}).
\end{definitie}

\begin{propozitie}[funcțiile simetrice de rădăcini rămân în inelul de bază]\label{prop:simroots}
Fie $R\subseteq S$ inele integre, $f\in R[X]$ \emph{monic} de grad $n$, cu toate rădăcinile $x_1,\dots,x_n$ în $S$. Atunci pentru orice polinom simetric $g\in R[X_1,\dots,X_n]$ avem $g(x_1,\dots,x_n)\in R$.
\end{propozitie}
\begin{proof}
Prin Teorema fundamentală a polinoamelor simetrice (de mai jos), $g=G(e_1,\dots,e_n)$ cu $G\in R[X_1,\dots,X_n]$. Cum $f=(X-x_1)\cdots(X-x_n)$ (monic), relațiile lui Viète (Teorema~\ref{teo:viete}) dau $e_k(x_1,\dots,x_n)=(-1)^k a_{n-k}\in R$ (coeficienții lui $f$). Deci $g(x_1,\dots,x_n)=G\big(e_1(x),\dots,e_n(x)\big)$ este o expresie polinomială cu coeficienți în $R$ evaluată în elemente din $R$, deci aparține lui $R$.
\end{proof}

\begin{definitie}[corpul funcțiilor raționale]
Pentru un corp comutativ $K$, corpul de fracții al domeniului $K[X_1,\dots,X_n]$ se notează $K(X_1,\dots,X_n)$ și se numește \emph{corpul funcțiilor raționale} în $X_1,\dots,X_n$ cu coeficienți în $K$.
\end{definitie}

\begin{teorema}[teorema fundamentală a polinoamelor simetrice]\label{teo:polsim}
Orice polinom simetric în $x_1,\dots,x_n$ se exprimă în mod \textbf{unic} ca polinom în polinoamele simetrice fundamentale $e_1,\dots,e_n$.
\end{teorema}

La concurs se citează ca atare; calculul efectiv al unei expresii simetrice date se face de obicei cu relațiile lui Newton de mai sus.

\section{Rădăcini pentru polinoame din $\ZZ[X]$, $\QQ[X]$, $\RR[X]$, $\CC[X]$}

\begin{teorema}[rădăcinile raționale]\label{teo:rationale}
Fie $f=a_nX^n+\dots+a_0\in\ZZ[X]$, $a_n\neq 0$, și $\frac pq\in\QQ$ o rădăcină, cu fracția ireductibilă. Atunci $p\mid a_0$ și $q\mid a_n$. În particular, dacă $f$ este unitar, orice rădăcină rațională este întreagă și divide $a_0$.
\end{teorema}

\begin{proof}
Din $a_np^n+a_{n-1}p^{n-1}q+\dots+a_0q^n=0$: toți termenii în afara ultimului sunt divizibili cu $p$, deci $p\mid a_0q^n$, iar $(p,q)=1$ dă $p\mid a_0$. Simetric pentru $q$.
\end{proof}

\begin{propozitie}\label{prop:aminusb}
Dacă $f\in\ZZ[X]$ și $a,b\in\ZZ$, atunci $(a-b)\mid f(a)-f(b)$.
\end{propozitie}

\begin{proof}
$f(a)-f(b)=\sum_k a_k(a^k-b^k)$ și $(a-b)\mid a^k-b^k=(a-b)(a^{k-1}+\dots+b^{k-1})$.
\end{proof}

\begin{observatie}[testul de paritate]
Dacă $f\in\ZZ[X]$ și $f(0),f(1)$ sunt ambele impare, atunci $f$ nu are rădăcini întregi: pentru $k$ par, $f(k)\equiv f(0)\equiv1\pmod2$; pentru $k$ impar, $f(k)\equiv f(1)\equiv1\pmod 2$.
\end{observatie}

\begin{aplicatie}[ciclu de lungime 3 --- problemă clasică]\label{apl:ciclu}
Nu există $f\in\ZZ[X]$ și întregi \emph{distincți} $a,b,c$ cu $f(a)=b$, $f(b)=c$, $f(c)=a$.

\emph{Soluție.} Din Propoziția~\ref{prop:aminusb}: $(a-b)\mid f(a)-f(b)=b-c$, apoi $(b-c)\mid c-a$ și $(c-a)\mid a-b$. Deci $|a-b|\le|b-c|\le|c-a|\le|a-b|$, adică toate cele trei module sunt egale cu un $d>0$. Dar $(a-b)+(b-c)+(c-a)=0$, iar o sumă de trei termeni egali cu $\pm d$ nu poate fi $0$ (ar trebui $\pm d\pm d\pm d=0$, imposibil pentru $d\neq0$). Contradicție.
\end{aplicatie}

\begin{teorema}[rădăcini complexe conjugate]\label{teo:conjugata}
Fie $f\in\RR[X]$ și $z\in\CC\setminus\RR$ o rădăcină. Atunci $\bar z$ este rădăcină, \textbf{cu același ordin de multiplicitate}.
\end{teorema}

\begin{proof}
Conjugarea este compatibilă cu operațiile ($\overline{u+v}=\bar u+\bar v$, $\overline{uv}=\bar u\,\bar v$) și fixează coeficienții reali, deci $f(\bar z)=\overline{f(z)}=0$. Pentru multiplicitate: $q=(X-z)(X-\bar z)=X^2-2\mathrm{Re}(z)X+|z|^2\in\RR[X]$; împărțind în $\RR[X]$, $f=q\,h+r$ cu $\deg r\le1$, iar $r(z)=0$ cu $z\notin\RR$ forțează $r=0$ (un polinom real de grad $\le1$ cu rădăcină nereală este nul). Deci $f=qh$ cu $h\in\RR[X]$, iar multiplicitățile lui $z$ și $\bar z$ în $f$ sunt cu $1$ mai mari decât în $h$; inducția încheie.
\end{proof}

\begin{corolar}
Orice polinom din $\RR[X]$ de grad impar are cel puțin o rădăcină reală. (Rădăcinile nereale se grupează în perechi conjugate cu multiplicități egale, deci sunt în număr par.)
\end{corolar}

\begin{teorema}[conjugata pătratică]\label{teo:conjpatratica}
Fie $f\in\QQ[X]$, $a,b,d\in\QQ$ cu $\sqrt d\notin\QQ$. Dacă $a+b\sqrt d$ este rădăcină a lui $f$, atunci și $a-b\sqrt d$ este rădăcină (cu același ordin de multiplicitate).
\end{teorema}

\begin{proof}
Prin inducție, $(a+b\sqrt d)^k=A_k+B_k\sqrt d$ și $(a-b\sqrt d)^k=A_k-B_k\sqrt d$, cu $A_k,B_k\in\QQ$ (pasul de inducție: înmulțirea cu $a\pm b\sqrt d$ păstrează simetria semnului). Sumând cu coeficienții lui $f$: $f(a+b\sqrt d)=P+Q\sqrt d$ și $f(a-b\sqrt d)=P-Q\sqrt d$, cu $P,Q\in\QQ$. Din $P+Q\sqrt d=0$: dacă $Q\neq0$, atunci $\sqrt d=-P/Q\in\QQ$, contradicție; deci $Q=0$, apoi $P=0$, adică $f(a-b\sqrt d)=0$. Multiplicitatea: se împarte la $\bigl(X-(a+b\sqrt d)\bigr)\bigl(X-(a-b\sqrt d)\bigr)=X^2-2aX+(a^2-db^2)\in\QQ[X]$ și se raționează ca la Teorema~\ref{teo:conjugata}.
\end{proof}

\begin{teorema}[teorema fundamentală a algebrei; d'Alembert--Gauss]\label{teo:tfa}
\emph{(Admisă.)} Orice polinom din $\CC[X]$ de grad $\ge1$ are cel puțin o rădăcină în $\CC$. Prin urmare, orice $f\in\CC[X]$ de grad $n\ge 1$ se descompune
\[
f=a_n(X-x_1)^{m_1}\cdots(X-x_r)^{m_r},\qquad m_1+\dots+m_r=n.
\]
\end{teorema}

\begin{corolar}[{descompunerea în $\RR[X]$}]\label{cor:descR}
Orice $f\in\RR[X]$ de grad $\ge1$ se descompune în $\RR[X]$ ca produs de factori de gradul întâi și factori de gradul al doilea cu discriminant negativ.
\end{corolar}

\begin{proof}
Descompunem în $\CC[X]$ și grupăm rădăcinile nereale în perechi conjugate (Teorema~\ref{teo:conjugata}): $(X-z)(X-\bar z)\in\RR[X]$ cu $\Delta<0$.
\end{proof}

\begin{teorema}[corp de descompunere]\label{teo:corpdescompunere}
Fie $K$ un corp comutativ și $f\in K[X]$ un polinom neconstant de grad $n$. Atunci există un corp $L\supseteq K$ în care $f$ are $n$ rădăcini (adică $f$ se descompune complet în factori de gradul întâi peste $L$).
\end{teorema}
\begin{proof}
Procedăm prin inducție după $n=\deg f$, arătând că putem \emph{adăuga o rădăcină} și apoi itera. Fie $p$ un factor ireductibil al lui $f$ în $K[X]$. Inelul cât $L_1=K[X]/(p)$ este \emph{corp}, deoarece $(p)$ este ideal maximal ($p$ ireductibil), iar $K$ se scufundă în $L_1$ prin $a\mapsto a+(p)$. Clasa $\alpha=X+(p)$ satisface $p(\alpha)=0$, deci și $f(\alpha)=0$: peste $L_1$, $f$ are rădăcina $\alpha$, deci $f=(X-\alpha)f_1$ cu $f_1\in L_1[X]$, $\deg f_1=n-1$. Prin ipoteza de inducție aplicată lui $f_1$ peste $L_1$, există $L\supseteq L_1\supseteq K$ în care $f_1$ (deci și $f$) are toate rădăcinile. (Este construcția lui Kronecker; pentru $K=\ZZ_p$ ea este exact mecanismul folosit mai jos la corpuri finite, $\ZZ_p[X]/(p)$.)
\end{proof}

\section{Ireductibilitate: criterii practice}

\begin{propozitie}[gradele mici]\label{prop:gradmic}
Peste un corp $K$: \emph{(i)} orice polinom de gradul $1$ este ireductibil; \emph{(ii)} un polinom de gradul $2$ sau $3$ este ireductibil $\iff$ nu are rădăcini în $K$.
\end{propozitie}

\begin{proof}
(ii) O descompunere netrivială a unui polinom de grad $2$ sau $3$ conține obligatoriu un factor de gradul $1$, $aX+b$, care dă rădăcina $-b/a\in K$; reciproc, o rădăcină $\alpha$ dă factorul $X-\alpha$ (Bézout).
\end{proof}

\begin{observatie}[capcană la gradul 4]
Pentru grad $\ge4$, lipsa rădăcinilor \textbf{nu} implică ireductibilitatea:
\[
X^4+X^2+1=(X^2+X+1)(X^2-X+1)
\]
nu are nicio rădăcină reală (este $>0$), dar este reductibil peste $\QQ$. Verificarea: $(X^2+1)^2-X^2$.
\end{observatie}

\begin{corolar}[din TFA]
Peste $\CC$, ireductibile sunt exact polinoamele de gradul $1$; peste $\RR$, cele de gradul $1$ și cele de gradul $2$ cu $\Delta<0$ (Corolarul~\ref{cor:descR}). Chestiunea ireductibilității este deci interesantă peste $\QQ$ (și peste $\ZZ_p$).
\end{corolar}

\begin{definitie}
Pentru $f\in\ZZ[X]$ nenul, \emph{conținutul} $c(f)$ este c.m.m.d.c.\ (pozitiv) al coeficienților; $f$ este \emph{primitiv} dacă $c(f)=1$. Orice $f\in\ZZ[X]$ nenul se scrie $f=c(f)\,f_0$ cu $f_0$ primitiv.
\end{definitie}

\begin{lema}[Gauss]\label{lema:gauss}
Produsul a două polinoame primitive din $\ZZ[X]$ este primitiv.
\end{lema}

\begin{proof}
Fie $g,h$ primitive și presupunem că un prim $p$ divide toți coeficienții lui $gh$. Reducem coeficienții modulo $p$: în $\ZZ_p[X]$ avem $\bar g\cdot\bar h=\overline{gh}=0$. Dar $\ZZ_p[X]$ este inel integru (Propoziția~\ref{prop:grade}(iii), $\ZZ_p$ corp), deci $\bar g=0$ sau $\bar h=0$, adică $p$ divide toți coeficienții lui $g$ sau ai lui $h$ --- contradicție cu primitivitatea.
\end{proof}

\begin{teorema}[teorema lui Gauss]\label{teo:gauss}
Fie $f\in\ZZ[X]$, $\deg f\ge1$. Dacă $f=gh$ cu $g,h\in\QQ[X]$ de grade $\ge1$, atunci există $g_1,h_1\in\ZZ[X]$ cu $f=g_1h_1$, $\deg g_1=\deg g$, $\deg h_1=\deg h$. Echivalent: \emph{ireductibil peste $\ZZ$} (în sensul gradelor) $\Rightarrow$ \emph{ireductibil peste $\QQ$}.
\end{teorema}

\begin{proof}
Alegem $m,n\in\NN^*$ care elimină numitorii: $G=mg\in\ZZ[X]$, $H=nh\in\ZZ[X]$, deci $mn\,f=GH$. Scriem $G=c(G)G_0$, $H=c(H)H_0$ cu $G_0,H_0$ primitive; atunci $mn\,f=c(G)c(H)\,G_0H_0$, iar $G_0H_0$ este primitiv (Lema~\ref{lema:gauss}). Luând conținutul ambilor membri: $mn\,c(f)=c(G)c(H)$. Prin urmare
\[
f=\frac{c(G)c(H)}{mn}\,G_0H_0=c(f)\,G_0H_0=\bigl(c(f)G_0\bigr)\cdot H_0,
\]
o descompunere în $\ZZ[X]$ cu aceleași grade.
\end{proof}

\begin{teorema}[criteriul lui Eisenstein]\label{teo:eisenstein}
Fie $f=a_nX^n+\dots+a_0\in\ZZ[X]$ și $p$ un prim cu:
\[
p\nmid a_n,\qquad p\mid a_i\ (0\le i\le n-1),\qquad p^2\nmid a_0.
\]
Atunci $f$ este ireductibil peste $\QQ$.
\end{teorema}

\begin{proof}
Prin Teorema~\ref{teo:gauss} este suficient să arătăm că $f$ nu se scrie $f=gh$ cu $g,h\in\ZZ[X]$, $\deg g=m\ge1$, $\deg h=k\ge1$. Fie $g=b_mX^m+\dots+b_0$, $h=c_kX^k+\dots+c_0$. Din $a_0=b_0c_0$, $p\mid a_0$ și $p^2\nmid a_0$, exact unul dintre $b_0,c_0$ este divizibil cu $p$; zicem $p\mid b_0$, $p\nmid c_0$. Din $a_n=b_mc_k$ și $p\nmid a_n$ rezultă $p\nmid b_m$; fie deci $r=\min\{i\mid p\nmid b_i\}$, cu $1\le r\le m<n$. Atunci
\[
a_r=b_rc_0+\underbrace{b_{r-1}c_1+\dots+b_0c_r}_{\text{divizibil cu }p},
\]
iar $p\mid a_r$ (deoarece $r<n$), deci $p\mid b_rc_0$ --- contradicție, căci $p\nmid b_r$ și $p\nmid c_0$.
\end{proof}

\begin{exemplu}
$X^n-2$ și, mai general, $X^n-p$ sunt ireductibile peste $\QQ$ pentru orice $n\ge1$ (Eisenstein cu primul $p$). Consecință: există polinoame ireductibile peste $\QQ$ de orice grad.
\end{exemplu}

\begin{aplicatie}[polinomul ciclotomic $\Phi_p$]\label{apl:ciclotomic}
Pentru $p$ prim, $f=X^{p-1}+X^{p-2}+\dots+X+1=\dfrac{X^p-1}{X-1}$ este ireductibil peste $\QQ$.

\emph{Soluție.} Substituția $X\mapsto X+1$ păstrează (i)reductibilitatea (este inversabilă, prin $X\mapsto X-1$). Calculăm
\[
f(X+1)=\frac{(X+1)^p-1}{X}=X^{p-1}+\binom p1X^{p-2}+\dots+\binom{p}{p-2}X+\binom{p}{p-1}.
\]
Coeficientul dominant este $1$, toți ceilalți coeficienți $\binom pk$ ($1\le k\le p-1$) sunt divizibili cu $p$ (demonstrat la Teorema~\ref{teo:frobenius}), iar termenul liber este $\binom p{p-1}=p$, nedivizibil cu $p^2$. Eisenstein încheie.
\end{aplicatie}

\begin{teorema}[reducerea modulo $p$]\label{teo:redmodp}
Fie $f\in\ZZ[X]$ și $p$ un prim cu $p\nmid a_n$ (coeficientul dominant). Dacă polinomul redus $\bar f\in\ZZ_p[X]$ este ireductibil peste $\ZZ_p$, atunci $f$ este ireductibil peste $\QQ$.
\end{teorema}

\begin{proof}
Presupunem $f$ reductibil peste $\QQ$; prin Teorema~\ref{teo:gauss}, $f=gh$ în $\ZZ[X]$ cu $\deg g,\deg h\ge1$. Coeficienții dominanți satisfac $a_n=(\text{dom }g)(\text{dom }h)$, deci $p$ nu divide niciunul; prin urmare $\deg\bar g=\deg g$ și $\deg\bar h=\deg h$, iar $\bar f=\bar g\,\bar h$ este o descompunere netrivială în $\ZZ_p[X]$ --- contradicție.
\end{proof}

\begin{observatie}
(a) Reciproca este falsă: un exemplu clasic este $X^4+1$, ireductibil peste $\QQ$, dar reductibil modulo \emph{orice} prim $p$. Deci eșecul criteriului pentru un $p$ nu spune nimic. (b) Ireductibilitatea \textbf{depinde de corp}: $X^2-2$ este ireductibil peste $\QQ$, reductibil peste $\RR$; $X^2+1$ este ireductibil peste $\RR$ și peste $\ZZ_3$, dar reductibil peste $\CC$ și peste $\ZZ_5$ ($X^2+\cl1=(X-\cl2)(X+\cl2)$ în $\ZZ_5[X]$).
\end{observatie}

\section{Polinoame peste $\ZZ_p$. Rădăcini primitive}

\begin{teorema}\label{teo:XpX}
În $\ZZ_p[X]$ ($p$ prim) are loc descompunerea
\[
X^p-X=\prod_{\cl a\in\ZZ_p}(X-\cl a),\qquad\text{deci}\qquad X^{p-1}-\cl1=\prod_{\cl a\in\ZZ_p^*}(X-\cl a).
\]
\end{teorema}

\begin{proof}
Prin Fermat (Teorema~\ref{teo:euler}), fiecare dintre cele $p$ elemente ale lui $\ZZ_p$ este rădăcină a lui $X^p-X$; prin Bézout, produsul $\prod(X-\cl a)$ divide $X^p-X$ (factorii sunt primi între ei doi câte doi). Ambele polinoame sunt unitare de grad $p$, deci sunt egale. A doua egalitate: simplificăm cu $X$ (adică factorul $X-\cl 0$).
\end{proof}

\begin{corolar}[Wilson, a doua demonstrație]
Identificând termenii liberi în $X^{p-1}-\cl 1=\prod_{\cl a\neq\cl 0}(X-\cl a)$:
$-\cl1=(-1)^{p-1}\cdot\cl1\cdot\cl2\cdots\cl{p-1}$, deci pentru $p$ impar $\overline{(p-1)!}=-\cl1$ (iar $p=2$ direct).
\end{corolar}

\begin{lema}[identitatea lui Gauss pentru $\varphi$]\label{lema:sumaphi}
$\displaystyle\sum_{d\mid n}\varphi(d)=n$ pentru orice $n\ge1$.
\end{lema}

\begin{proof}
Împărțim $\{1,2,\dots,n\}$ după c.m.m.d.c.\ cu $n$: pentru fiecare divizor $e\mid n$, numerele $k$ cu $(k,n)=e$ sunt $k=ej$ cu $1\le j\le n/e$ și $(j,n/e)=1$, în număr de $\varphi(n/e)$. Sumând după $e$ și schimbând $d=n/e$: $n=\sum_{e\mid n}\varphi(n/e)=\sum_{d\mid n}\varphi(d)$.
\end{proof}

\begin{teorema}[ciclicitate; rădăcini primitive]\label{teo:ciclicitate}
Orice subgrup finit $G$ al grupului multiplicativ $(K^*,\cdot)$ al unui corp comutativ $K$ este ciclic. În particular, $(\ZZ_p^*,\cdot)$ este ciclic: există $g$ (numit \emph{rădăcină primitivă modulo $p$}) cu $\ZZ_p^*=\{\cl 1,\cl g,\dots,\cl g^{\,p-2}\}$.
\end{teorema}

\begin{proof}
Fie $|G|=n$ și, pentru $d\mid n$, fie $\psi(d)$ numărul elementelor de ordin $d$ din $G$ (ordinele divid $n$, prin Lagrange). Dacă $\psi(d)>0$, alegem $x$ de ordin $d$. Toate cele $d$ elemente din $\langle x\rangle$ satisfac $y^d=1$, iar în corpul $K$ ecuația polinomială $y^d=1$ are \textbf{cel mult} $d$ soluții (Teorema~\ref{teo:numarradacini}); deci mulțimea soluțiilor este exact $\langle x\rangle$. Orice element de ordin $d$ este o astfel de soluție, deci aparține lui $\langle x\rangle$ și este generator al acestuia: sunt exact $\varphi(d)$ asemenea generatori (Propoziția~\ref{prop:generatori}). Așadar $\psi(d)\in\{0,\varphi(d)\}$ pentru fiecare $d\mid n$. Dar
\[
n=\sum_{d\mid n}\psi(d)\le\sum_{d\mid n}\varphi(d)=n\qquad(\text{Lema~\ref{lema:sumaphi}}),
\]
deci peste tot avem egalitate: $\psi(d)=\varphi(d)$ pentru orice $d\mid n$. În particular $\psi(n)=\varphi(n)\ge1$: există un element de ordin $n$, care generează $G$.
\end{proof}

\begin{observatie}
Demonstrația dă mai mult: în $\ZZ_p^*$ există exact $\varphi(d)$ elemente de ordin $d$, pentru fiecare $d\mid p-1$. Împreună cu Propoziția~\ref{prop:subgrupC}, obținem și structura subgrupurilor finite ale lui $\CC^*$.
\end{observatie}

\begin{observatie}[a doua demonstrație a ciclicității]
Corolarul~\ref{cor:lcm} dă o demonstrație alternativă: fie $m$ ordinul maxim al unui element din $G$, atins de $a$. Pentru orice $x\in G$, grupul abelian $G$ conține un element de ordin $[\,\ord x,\,m\,]$; prin maximalitate, $[\,\ord x,\,m\,]\le m$, deci $\ord x\mid m$. Atunci toate cele $|G|$ elemente sunt rădăcini ale polinomului $X^m-1$, care are cel mult $m$ rădăcini în $K$ (Teorema~\ref{teo:numarradacini}); așadar $|G|\le m=\ord a\le|G|$, deci $a$ generează $G$.
\end{observatie}


\begin{propozitie}[criteriul lui Euler]\label{prop:eulercrit}
Fie $p$ prim impar și $a$ cu $p\nmid a$. Atunci $a^{\frac{p-1}2}\equiv\pm1\pmod p$, iar
\[
a^{\frac{p-1}{2}}\equiv1\pmod p\iff a\text{ este rest pătratic mod }p\ (\exists\,x:\ x^2\equiv a).
\]
\end{propozitie}

\begin{proof}
Fie $g$ o rădăcină primitivă și $\cl a=\cl g^{\,t}$. Elementul $h=\cl g^{\,\frac{p-1}2}$ satisface $h^2=\cl 1$ și $h\neq\cl1$ (ordinul lui $g$ este $p-1$), deci $h=-\cl1$ (singurele soluții ale lui $x^2=\cl1$ în corp). Atunci $\cl a^{\frac{p-1}2}=(-\cl 1)^t$. Dacă $t$ este par, $a=g^t=(g^{t/2})^2$ este pătrat; reciproc, $\cl a=\cl x^2$ cu $\cl x=\cl g^{\,s}$ dă $t\equiv 2s\pmod{p-1}$, deci $t$ par ($p-1$ fiind par).
\end{proof}

\begin{corolar}\label{cor:minus1}
$-1$ este rest pătratic modulo primul impar $p$ $\iff p\equiv1\pmod4$. (Se aplică criteriul: $(-1)^{\frac{p-1}2}=1\iff\frac{p-1}2$ par.)
\end{corolar}

\begin{propozitie}[sumele de puteri modulo $p$]\label{prop:sumeputeri}
Pentru $p$ prim și $k\ge 1$:
\[
1^k+2^k+\dots+(p-1)^k\equiv
\begin{cases}
-1\pmod p,&\text{dacă }(p-1)\mid k,\\
\ \ 0\pmod p,&\text{altfel}.
\end{cases}
\]
\end{propozitie}

\begin{proof}
Dacă $(p-1)\mid k$, fiecare termen este $\equiv1$ (Fermat), iar suma este $p-1\equiv-1$. Altfel, fie $g$ o rădăcină primitivă; suma este $S=\sum_{j=0}^{p-2}\cl g^{\,jk}$, iar
\[
(\cl g^{\,k}-\cl 1)\,S=\cl g^{\,k(p-1)}-\cl1=\cl0.
\]
Cum $(p-1)\nmid k$, avem $\cl g^{\,k}\neq\cl1$, deci $\cl g^{\,k}-\cl1$ este inversabil în corpul $\ZZ_p$ și $S=\cl0$.
\end{proof}

\section{Interpolarea Lagrange și polinoame auxiliare}

\begin{teorema}[interpolarea Lagrange]\label{teo:lagrangeint}
Fie $K$ un corp comutativ, $a_0,\dots,a_n\in K$ distincte și $b_0,\dots,b_n\in K$. Există un \textbf{unic} polinom $L\in K[X]$ cu $\deg L\le n$ și $L(a_i)=b_i$ pentru toți $i$, anume
\[
L(X)=\sum_{i=0}^{n}b_i\prod_{j\neq i}\frac{X-a_j}{a_i-a_j}.
\]
\end{teorema}

\begin{proof}
Fiecare produs $\prod_{j\neq i}\frac{X-a_j}{a_i-a_j}$ ia valoarea $1$ în $a_i$ și $0$ în celelalte noduri, deci $L(a_i)=b_i$; gradul este evident $\le n$. Unicitatea: diferența a două soluții are grad $\le n$ și $n+1$ rădăcini, deci este nulă (Corolarul~\ref{cor:identitate}).
\end{proof}

\begin{aplicatie}[polinomul auxiliar --- problemă clasică]\label{apl:auxiliar}
Fie $f\in\RR[X]$ de grad $n$ cu $f(k)=\dfrac{k}{k+1}$ pentru $k=0,1,\dots,n$. Calculați $f(n+1)$.

\emph{Soluție.} Polinomul $g(X)=(X+1)f(X)-X$ are gradul $n+1$ și se anulează în $0,1,\dots,n$, deci
$g(X)=c\,X(X-1)\cdots(X-n)$. Evaluăm în $X=-1$: $g(-1)=0\cdot f(-1)+1=1$, iar produsul dă $c\,(-1)^{n+1}(n+1)!$, de unde $c=\dfrac{(-1)^{n+1}}{(n+1)!}$. Atunci $g(n+1)=c\,(n+1)!=(-1)^{n+1}$ și
\[
f(n+1)=\frac{g(n+1)+(n+1)}{n+2}=
\begin{cases}
1,&n\text{ impar},\\[2pt]
\dfrac{n}{n+2},&n\text{ par}.
\end{cases}
\]
\emph{Morala:} când se dau valorile lui $f$ în noduri, construiți un polinom auxiliar ale cărui rădăcini sunt nodurile; gradul și un nod suplimentar determină constanta.
\end{aplicatie}

\vspace{6pt}
\begin{center}
\rule{0.5\textwidth}{0.4pt}\\[4pt]
\small\emph{Succes la olimpiadă!}
\end{center}











\clearpage
\part{PROBLEME}
\lhead{\small\textcolor{grey}{Probleme de algebră}}
\addcontentsline{toc}{chapter}{\textcolor{blue!60!black}{Probleme de algebră}}
\sectiunealg{Probleme de algebră}{Structurile algebrice ale clasei a XII-a.}
\nopagebreak
\parteamareenunt{Partea I \;·\; Enunțuri}
\subsectiuneclasa{Clasa a XII-a}{Structuri algebrice: grupuri, inele și corpuri --- comutativitate, idempotenți, involuții.}{Enunțuri · Clasa a XII-a}

\nivel{olm}{Nivel OLM}{Faza locală --- o singură idee, bine aplicată.}
\setcounter{cat}{1}\setcounter{problemaX}{0}

\begin{problema}[Clasică]
Fie $G$ un grup în care $x^{2}=e$ pentru orice $x \in G$. Demonstrați
că $G$ este comutativ.
\end{problema}

\begin{problema}
Determinați toate grupurile $G$ pentru care există numere naturale
$m,n$ prime între ele astfel încât $x^{m}=x^{n}=e$ pentru orice
$x \in G$.
\end{problema}

\begin{problema}[Clasică]
Fie $G$ un grup finit comutativ de ordin \emph{impar}. Demonstrați că
produsul tuturor elementelor sale este $e$.
\end{problema}

\begin{problema}[Clasică]
Fie $G$ un grup finit care are exact un element de ordin $2$, notat $a$.
Demonstrați că $a$ comută cu orice element al lui $G$ (adică
$a$ aparține centrului lui $G$).
\end{problema}

\begin{problema}
Fie $G$ un grup și $a,b \in G$ cu $aba^{-1}=b^{-1}$. Demonstrați
că $a^{2}b = ba^{2}$ (adică $a^{2}$ comută cu $b$).
\end{problema}

\begin{problema}
Un număr natural $n$ se numește \emph{liber de pătrate} dacă
nu este divizibil cu pătratul niciunui număr prim (echivalent, în
descompunerea $n=p_{1}p_{2}\cdots p_{r}$ toți factorii primi sunt
distincți). Fie $n$ liber de pătrate și $G$ un grup de ordin
$n$. Găsiți cel mai mic număr natural $k$ cu proprietatea
că $x^{k}=e$ pentru orice $x \in G$.
\end{problema}

\begin{problema}
Determinați toate numerele naturale $n$ pentru care există un grup
$G$ de ordin $n$ cu proprietatea că, pentru orice divizor $d$ al lui
$n$, grupul $G$ conține exact un element de ordin $d$.
\end{problema}

\begin{problema}
Fie $G$ un grup care are un \emph{unic} subgrup ciclic netrivial (adică
un singur subgrup ciclic diferit de $\{e\}$). Demonstrați că $G$
este ciclic; mai precis, $G\cong\mathbb{Z}/p$ pentru un număr prim $p$.
\end{problema}

\begin{problema}
Fie $G$ un grup cu proprietatea că pentru orice $x\in G$ există
$y\in G$ astfel încât $x^{2}=yxy^{-1}$ (adică $x^{2}$ este un
conjugat al lui $x$). Demonstrați că $G$ nu are elemente de ordin
par (în particular, dacă $G$ e finit, $|G|$ este impar).
\end{problema}

\begin{problema}
Există un grup $G$ care are două subgrupuri $K,L$ (nu neapărat
diferite) astfel încât produsul lor $KL=\{\,kl : k\in K,\ l\in L\,\}$
să fie subgrupul trivial $\{e\}$, cu $K$ sau $L$ netrivial? Justificați.
\end{problema}

\begin{problema}
Determinați toate inelele $(A,+,\cdot)$ cu proprietatea că $0=1$
(elementul neutru la adunare coincide cu elementul neutru la înmulțire).
\end{problema}

\begin{problema}
Un grup $G$ se numește \emph{rezolubil} (solubil) dacă există un
lanț finit de subgrupuri
\[
  \{e\}=G_{0}\trianglelefteq G_{1}\trianglelefteq\cdots\trianglelefteq G_{m}=G
\]
în care, pentru fiecare $i$, $G_{i}$ este normal în $G_{i+1}$ și
câtul $G_{i+1}/G_{i}$ este abelian. Arătați că grupul $S_{3}$
al permutărilor mulțimii $\{1,2,3\}$ (de ordin $6$) este rezolubil,
exhibând un astfel de lanț.
\end{problema}

\begin{problema}
Fie $A$ un inel și $x\in A$ un element pentru care există un
întreg $\ell\ge 1$ cu
\[
  x^{\ell}=0\qquad\text{și}\qquad (x+1)^{\ell}=0 .
\]
Demonstrați că $1=0$ în $A$ (adică $A$ este inelul nul).
\end{problema}

\begin{problema}
Fie $A$ un inel. Un element $x\in A$ se numește \emph{idempotent}
dacă $x^{2}=x$, și \emph{nilpotent} dacă există un
întreg $n\ge 1$ cu $x^{n}=0$. Demonstrați că singurul element
care este simultan idempotent și nilpotent este $0$.
\end{problema}

\begin{problema}
Fie $K$ un corp și $a\in K$, cu $a\neq 1$, având proprietatea
că, pentru orice element idempotent $x$ (adică $x^{2}=x$), elementul
$a-x$ este tot idempotent. Demonstrați că $\operatorname{char}K=2$.
\end{problema}

\begin{problema}
Fie $P \in \mathbb{C}[X]$ un polinom neconstant astfel încât $P(X)$ divide $P(X^2)$ în $\mathbb{C}[X]$. Arătați că orice rădăcină nenulă a lui $P$ este o rădăcină a unității.
\end{problema}

\begin{problema}
Determinați toate polinoamele $P \in \RR[X]$ cu proprietatea
\[
  P\bigl(P(x)\bigr) = P(x)^2 \qquad \text{pentru orice } x \in \RR.
\]
\end{problema}

\begin{problema}[Clasică]
Fie $K$ un corp finit cu cel puțin 3 elemente. Arătați că suma tuturor elementelor lui $K$ este 0. Ce se întâmplă pentru $K = \ZZ_2$?
\end{problema}

\begin{problema}[Clasică]
Fie $K$ un corp finit de caracteristică $p$.
\begin{parti}
\item Arătați că aplicația $F : K \to K$, $F(x) = x^p$, este un automorfism al lui $K$.
\item Arătați că $F(x) = x$ dacă și numai dacă $x$ aparține subcorpului prim $\{0,\, 1,\, 1+1,\, \dots\} \simeq \ZZ_p$.
\end{parti}
\end{problema}

\begin{problema}
Fie $G$ un grup cu proprietatea că $xy \in Z(G)$ pentru orice
$x,y\in G$ cu $x\neq e$ și $y\neq e$. Demonstrați că $G$ este
comutativ.
\end{problema}

\begin{problema}
Fie $G$ un grup finit și $f,g$ două endomorfisme ale lui $G$ cu
proprietatea
\[
  f\bigl(x^{2}y^{4}\bigr) = g\bigl(x^{2}\bigr)\,g\bigl(y^{4}\bigr)
  \qquad \text{pentru orice } x,y\in G.
\]
Determinați sub ce condiție asupra ordinului $|G|$ rezultă
$f=g$.
\end{problema}

\begin{problema}
O \emph{relație de echivalență} pe o mulțime $M$ este o
relație $\rho$ care este reflexivă ($x\,\rho\,x$ pentru orice $x$),
simetrică ($x\,\rho\,y\Rightarrow y\,\rho\,x$) și tranzitivă
($x\,\rho\,y$ și $y\,\rho\,z\Rightarrow x\,\rho\,z$). O relație de
echivalență $\rho$ pe un grup $G$ se numește \emph{congruență}
dacă este compatibilă cu operația: $x\,\rho\,y$ și
$x'\,\rho\,y'$ implică $xx'\,\rho\,yy'$. Determinați toate
congruențele pe grupul $(\mathbb{Z},+)$.
\end{problema}

\begin{problema}
Fie $\sigma(z)=5/\bar z$ inversiunea planului față de cercul de
rază $\sqrt5$ centrat în origine. Grupul simetriilor rețelei
$\mathbb{Z}[i]$ (rotațiile $z\mapsto i^{k}z$ și conjugarea
$z\mapsto\bar z$) permută mulțimea $X$ a punctelor din $\mathbb{Z}[i]$
fixate de $\sigma$. Descrieți $X$ ca mulțime cu acțiune de grup.
\end{problema}

\begin{problema}
Fie $G$ un grup finit de ordin $n$ și $\varphi\in\operatorname{Aut}(G)$
un automorfism fixat. Arătați că, pentru orice $x\in G$ fixat,
numărul soluțiilor $g\in G$ ale ecuației $gx=x\varphi(g)$ divide
$n$.
\end{problema}

\begin{problema}
Fie $A$ un inel cu proprietatea că orice element $x\in A$ se scrie ca
sumă $x=e+f$ a două elemente idempotente ($e^{2}=e$, $f^{2}=f$) care
\emph{anticomută}, adică $ef+fe=0$. Demonstrați că $A$ este
comutativ.
\end{problema}

\begin{problema}
Fie $A$ un inel comutativ și fie $I,J$ ideale ale sale astfel încât $I \cup J$ să fie un ideal al lui $A$. Arătați că $I \subseteq J$ sau $J \subseteq I$.
\end{problema}

\begin{problema}
Determinați toate numerele naturale $n$ pentru care \emph{orice} grup
$G$ de ordin $n$ are proprietatea că orice endomorfism neconstant al
său este automorfism. (Un endomorfism se numește \emph{constant}
dacă este morfismul banal $g \mapsto e$.)
\end{problema}

\begin{problema}
Fie $K$ un corp infinit și $f : K \to K$ o funcție
\emph{polinomială} cu proprietatea
\[
  f\bigl(x+f(y)\bigr) = f(x)\,f(y) \qquad \text{pentru orice } x,y \in K .
\]
Determinați toate funcțiile $f$.
\end{problema}

\begin{problema}
Fie $G$ un grup cu proprietatea că, pentru oricare trei elemente
distincte ale sale, există cel puțin o pereche dintre ele care
comută. Demonstrați că $G$ este comutativ.
\end{problema}

\nivel{ojm}{Nivel OJM}{Faza județeană --- combină două rezultate.}
\setcounter{cat}{2}\setcounter{problemaX}{0}

\begin{problema}
Fie $K$ un corp cu 8 elemente. Arătați că polinomul $X^2+X+1$ nu are nicio rădăcină în $K$, dar că polinomul $X^3+X+1$ are trei rădăcini distincte în $K$.
\end{problema}

\begin{problema}[Clasică]
Fie $G$ un grup finit de ordin \emph{par}. Demonstrați că există
un element $x \in G$, $x \neq e$, cu $x^{2}=e$ (adică $G$ conține
un involutiv netrivial).
\end{problema}

\begin{problema}[Clasică]
Fie $G$ un grup și $a,b \in G$ astfel încât $a^{4}=e$ și
$aba^{-1}=b^{2}$. Demonstrați că $b^{15}=e$.
\end{problema}

\begin{problema}
Fie $G$ un grup de ordin $2n$ și $a \in G$ un element de ordin cel
puțin $n$. Presupunem că există $x \in G$ cu $x \neq a^{k}$
pentru orice $k$ (adică $x \notin \langle a\rangle$) astfel încât
$ax=xa$. Demonstrați că $a$ comută cu orice element al lui $G$
(adică $a$ aparține centrului lui $G$).
\end{problema}

\begin{problema}
Fie $G$ un grup și $a,b \in G$ astfel încât orice element al
lui $G$ comută cu $a$ sau cu $b$. Notând $C(x)=\{\,g\in G : gx=xg\,\}$
centralizatorul lui $x$, ipoteza spune că $C(a)\cup C(b)=G$.
Demonstrați că $a$ sau $b$ aparține centrului $Z(G)$ (cel
puțin unul dintre ele este central).
\end{problema}

\begin{problema}
Fie $G$ un grup de ordin $2n$, cu $n$ \emph{impar}, astfel încât
$x^{n+2}\in Z(G)$ pentru orice $x\in G$. Demonstrați că $G$ este
comutativ.
\end{problema}

\begin{problema}
Determinați toate numerele naturale $n$ pentru care există un grup
$G$ de ordin $n$ astfel încât, pentru orice întreg $k$ cu
$1\le k\le \tfrac{n}{3}$, aplicația $f_{k}:G\to G$, $f_{k}(x)=x^{k}$,
este injectivă.
\end{problema}

\begin{problema}[Clasică]
Fie $a<b$ numere reale.
\begin{enumerate}[label=(\alph*),nosep]
  \item Înzestrați mulțimea $G=(a,b)$ (intervalul deschis) cu o
        structură de grup comutativ.
  \item Aceeași cerință pentru intervalul \emph{închis}
        $G=[a,b]$.
\end{enumerate}
\end{problema}

\begin{problema}[Clasică]
Fie $f:G\to H$ un morfism de grupuri finite, cu
$\operatorname{Ker}f=\{x\in G:f(x)=e_{H}\}$ și $\operatorname{Im}f=f(G)$.
Arătați că toate fibrele nevide ale lui $f$ (mulțimile
$f^{-1}(h)$, $h\in\operatorname{Im}f$) au același număr de elemente
și că
\[
  |G|=|\operatorname{Ker}f|\cdot|\operatorname{Im}f| .
\]
\end{problema}

\begin{problema}
Fie $\pi=2+i\in\mathbb{Z}[i]$ și $\pi\mathbb{Z}[i]=\{\pi z : z\in\mathbb{Z}[i]\}$
subgrupul multiplilor săi în grupul aditiv $(\mathbb{Z}[i],+)$.
Arătați că grupul cât $\mathbb{Z}[i]/\pi\mathbb{Z}[i]$ este
izomorf cu $(\mathbb{Z}/5,+)$.
\end{problema}

\begin{problema}[Clasică]
Arătați că nu există niciun corp $K$ pentru care grupul
aditiv $(K,+)$ să fie izomorf cu grupul multiplicativ
$(K^{*},\cdot)$, unde $K^{*}=K\setminus\{0\}$.
\end{problema}

\begin{problema}
Fie $A$ un inel finit al cărui grup aditiv $(A,+)$ este izomorf cu grupul
multiplicativ $(K^{*},\cdot)$ al unui corp finit $K$. Demonstrați că
$A$ este comutativ.
\end{problema}

\begin{problema}
Fie $A$ un inel comutativ de caracteristică $p$, unde $p$ este un
număr prim \emph{impar}, cu proprietatea că $x^{p}=1$ pentru orice
element inversabil $x\in A$. Demonstrați că suma a două elemente
inversabile este tot un element inversabil.
\end{problema}

\begin{problema}
Fie $A$ un inel \emph{finit}, cu $1\neq 0$, având proprietatea că
pentru orice $n\ge 1$ ecuația $x^{n+1}=x^{n}$ are ca singure soluții
pe $0$ și pe $1$. Demonstrați că $A$ este corp.
\end{problema}

\begin{problema}
Fie $A$ un inel cu proprietatea că $x^{3}+x^{2}\in Z(A)$ pentru orice
$x\in A$ (unde $Z(A)=\{z\in A: zy=yz\ \forall y\in A\}$ este centrul
inelului). Demonstrați că $A$ este comutativ.
\end{problema}


\begin{problema}[Clasică]
Fie $K \subseteq L$ două corpuri finite, operațiile lui $K$ fiind cele induse de $L$. Arătați că există un întreg $m \geq 1$ cu
\[
  |L| = |K|^m.
\]
\end{problema}

\begin{problema}
Fie $G$ un grup finit de ordin $n$ și $k$ un număr întreg
pozitiv astfel încât $f(x)=x^{k}$ este morfism al lui $G$.
Presupunem că mulțimea
\[
  A = \{\, a \in G : a^{k}b = ba^{k} \ \text{pentru orice } b \in G \,\}
\]
are cel puțin $\tfrac{n}{2}+1$ elemente. Demonstrați că
$(ab)^{k}=(ba)^{k}$ pentru orice $a,b \in G$.
\end{problema}

\begin{problema}
Fie $G$ un grup abelian finit cu $|G|$ care are cel puțin doi factori
primi distincți. Demonstrați că există două subgrupuri
\emph{proprii și netriviale} $K,L\le G$ (adică diferite de $\{e\}$ și de $G$) cu $KL=G$.
\end{problema}

\begin{problema}
Fie $G$ un grup finit de ordin $n$ și $k$ un întreg astfel încât
aplicația $x\mapsto x^{k}$ este morfism de grupuri (adică
$(xy)^{k}=x^{k}y^{k}$ pentru orice $x,y\in G$). Arătați că, pentru
orice $x\in G$ fixat, numărul soluțiilor $a\in G$ ale ecuației
$xa^{k}=a^{k}x$ divide $n$.
\end{problema}

\begin{problema}
Fie $D$ un inel de diviziune (orice element nenul este inversabil) de
caracteristică $n$, fie $a,c$ numere întregi și $b,d\ge 1$
exponenți, astfel încât, pentru orice $x\in D$,
\[
  a\,x^{b}\in Z(D)\qquad\text{și}\qquad c\,x^{d}\in Z(D),
\]
unde $m\,y$ notează multiplul aditiv $\underbrace{y+\cdots+y}_{m}$.
Demonstrați că dacă
\[
  \gcd(a,n)=\gcd(c,n)=1\quad\text{și}\quad \gcd(b,d)=1,
\]
atunci $D$ este comutativ.
\end{problema}

\nivel{onm}{Nivel ONM}{Faza națională --- mai mulți pași și o idee-cheie.}
\setcounter{cat}{3}\setcounter{problemaX}{0}

\begin{problema}
Fie $G$ un grup finit de ordin $2n$. Câte soluții poate avea
ecuația $x^{n}=e$? (Determinați toate valorile posibile ale
numărului de soluții și arătați că fiecare apare.)
\end{problema}

\begin{problema}
Fie $G$ un grup finit cu un număr \emph{impar} de elemente și fie
$H \subseteq G$ o submulțime nevidă cu proprietatea
\[
  ghg \in H \qquad \text{pentru orice } g \in G,\ h \in H.
\]
Demonstrați că $H = G$.
\end{problema}

\begin{problema}
Fie $k \in \mathbb{N}^{*}$ fixat. Determinați toate polinoamele
$P \in \mathbb{C}[X]$ cu proprietatea
\[
  P(z^{k}) = P(z)\,P(z+1)\cdots P(z+k-1) \qquad \text{pentru orice } z \in \mathbb{C}.
\]
\end{problema}

\begin{problema}
Fie $R$ un inel cu unitate, $1\neq 0$, astfel încât $x^{3}=x^{2}$
pentru orice $x\in R$. Demonstrați că $x^{2}=x$ pentru orice
$x\in R$ (deci $R$ este inel Boolean) și deduceți că $R$ este
comutativ.
\end{problema}

\begin{problema}
Fie $G$ un grup finit de ordin $n$ cu proprietatea că pentru orice
$g\in G$, $g\neq e$, există un element $x\in G$ cu $gx\neq xg$ (altfel
spus, niciun element diferit de $e$ nu comută cu toate elementele lui
$G$). Demonstrați că $n$ divide $|\operatorname{Aut}(G)|$,
numărul automorfismelor lui $G$.
\end{problema}

\begin{problema}
Fie $G$ un grup finit cu $|Z(G)|=2$ care nu admite niciun morfism surjectiv
pe $\mathbb{Z}/2$ (echivalent: $G$ nu are subgrup de indice $2$, adică
$\operatorname{Hom}(G,\mathbb{Z}/2)=\{0\}$). Demonstrați că
$Z(\operatorname{Aut}G)=\{\operatorname{id}\}$.
\end{problema}

\begin{problema}
Fie $G$ un grup și $a,b,c\in\mathbb{N}^{*}$ numere prime între ele
două câte două. Presupunem că
\[
  (xy)^{n}=x^{n}y^{n}\qquad\text{pentru toți } x,y\in G,
\]
ori de câte ori $n$ este multiplu al lui $a$, al lui $b$ sau al lui $c$.
Demonstrați că $G$ este comutativ.
\end{problema}

\begin{problema}
Fie $G$ un grup finit de ordin $p^{\alpha}q^{\beta}$, cu $p\neq q$ numere
prime și $\alpha,\beta\ge 1$. Demonstrați că există
subgrupuri \emph{proprii} $H,K$ ale lui $G$ (adică $H\neq G$, $K\neq G$) astfel încât $HK=G$ (unde
$HK=\{\,hk : h\in H,\ k\in K\,\}$).
\end{problema}

\begin{problema}
Fie $n\ge 5$. Determinați toate relațiile de congruență pe
grupul simetric $S_{n}$. (O congruență este o relație de
echivalență compatibilă cu compunerea permutărilor.)
\end{problema}

\begin{problema}[Clasică]
Fie $G$ un grup finit \emph{neabelian} de ordin $n$. Demonstrați că
numărul perechilor ordonate $(a,b)\in G\times G$ cu $ab=ba$ este cel mult
$\dfrac{5}{8}\,n^{2}$.
\end{problema}

\begin{problema}
Fie $p\neq q$ numere prime, $\alpha\ge 1$, și $G$ un grup de ordin
$p\,q^{\alpha}$ cu proprietatea că
\[
  p\not\equiv 1\pmod q\qquad\text{și}\qquad q^{j}\not\equiv 1\pmod p
  \ \ \text{pentru }1\le j\le\alpha .
\]
Arătați că $G\cong P\times Q$ (produsul subgrupurilor sale Sylow)
și că $pq\mid|Z(G)|$.
\end{problema}

\begin{problema}
Fie $G$ un grup de ordin $p\,q^{\alpha}$ ($p\neq q$ prime) al cărui
$q$-subgrup Sylow $Q$ este normal \emph{și abelian}, iar $P$ un
$p$-subgrup Sylow. Arătați că
\[
  |Z(G)\cap Q|\equiv q^{\alpha}\pmod p,
\]
și deduceți că dacă $q^{\alpha}\not\equiv 1\pmod p$, atunci
$q\mid|Z(G)|$.
\end{problema}

\begin{problema}
Fie $G$ un grup finit și $M\subseteq G$ o submulțime cu $m$ elemente,
cu proprietatea că $gMg^{-1}=M$ pentru orice $g\in G$ (adică $M$ are
un singur conjugat, pe el însuși). Arătați că
$C_{G}(M)=\{g:gx=xg\ \forall x\in M\}$ este subgrup normal și că
$[G:C_{G}(M)]$ divide $m!$. Deduceți că dacă
$\gcd\bigl(|G|,m!\bigr)=1$, atunci $M\subseteq Z(G)$.
\end{problema}

\begin{problema}
Fie $p_{1}<p_{2}<\dots<p_{n}<\dots$ șirul numerelor prime. Demonstrați
că, pentru orice $n\ge 2$, orice grup de ordin $p_{n}\,p_{n+1}$ este
ciclic.
\end{problema}

\begin{problema}
Determinați toate numerele naturale $m$ pentru care există un grup
abelian finit având exact $m$ elemente de ordin $4$.
\end{problema}

\begin{problema}
Fie $G$ un grup abelian finit. Arătați că există un grup $H$
cu $G\cong H\times H$ dacă și numai dacă, pentru orice $n\ge 1$,
numărul soluțiilor ecuației $x^{n}=e$ în $G$ este pătrat
perfect.
\end{problema}

\begin{problema}
Arătați că orice inel comutativ $R$ este izomorf cu centrul unui
inel --- mai precis, există un inel $S$ (necomutativ dacă $R\neq\{0\}$)
al cărui centru $Z(S)=\{s\in S : sx=xs\ \forall x\in S\}$ este izomorf cu
$R$.
\end{problema}

\begin{problema}
Fie $A$ un inel finit de ordin impar, care nu este corp, cu proprietatea
că orice element neinversabil se scrie ca sumă a doi idempotenți care
comută. Arătați că $A$ este comutativ și că $x^{3}=x$ pentru orice
$x\in A$.
\end{problema}

\begin{problema}
Fie $A$ un inel cu proprietatea că $x^{4}=x$ pentru orice $x\in A$.
Demonstrați că $A$ este comutativ.
\end{problema}

\begin{problema}
Fie $K$ un corp al cărui grup multiplicativ $(K^*,\cdot)$ este ciclic. Arătați că $K$ este finit.
\end{problema}

\begin{problema}[Clasică]
Căsuțele unui caroiaj $N \times N$ se colorează cu $n \ge 1$ culori (nu neapărat distincte). Două colorări sunt considerate echivalente dacă una se obține din cealaltă printr-o rotație a caroiajului.
\begin{parti}
\item Determinați numărul colorărilor neechivalente.
\item Deduceți că numărătorul expresiei obținute este divizibil cu 4 pentru orice $n \ge 1$ și orice $N \ge 1$.
\end{parti}
\end{problema}

\begin{problema}
Fie $p$ un număr prim, $b \ge 2$ o bază cu $\gcd(b,p) = 1$ și $f = \ord_p(b)$.
Scrierea lui $\frac{1}{p}$ în baza $b$ este periodică simplă, cu perioada de lungime
exact $f$; fie $N = \frac{b^f-1}{p}$ numărul format din cele $f$ cifre ale unei
perioade (pentru $b = 10$, $p = 7$: $N = 142857$).
\begin{parti}
  \item Arătați că, pentru $1 \le k \le p-1$, numărul $kN$ (scris cu $f$ cifre
    în baza $b$) este o permutare circulară a lui $N$ dacă și numai dacă
    $k \in \langle b \rangle = \{1, b, \dots, b^{f-1}\}$ în $(\ZZ_p^*, \cdot)$;
    deduceți că numerele $N, 2N, \dots, (p-1)N$ se împart în exact $(p-1)/f$
    familii, fiecare formată din cele $f$ rotații ale unui ,,număr ciclic''.
  \item (Midy) Dacă $f = 2g$ este par, cele două jumătăți ale unei perioade, de
    câte $g$ cifre fiecare, au suma $b^g - 1$.
\end{parti}
\end{problema}

\begin{problema}
Fie $p$ un număr prim, $D \mid p-1$ și $m \ge 1$ un întreg. Arătați că suma puterilor a $m$-a ale elementelor de ordin $D$ din $(\ZZ_p^*,\cdot)$ este congruentă modulo $p$ cu \emph{suma lui Ramanujan}
\[
  \sum_{\ord(x)=D} x^m \equiv c_D(m) = \mu\!\left(\tfrac{D}{g}\right)\frac{\varphi(D)}{\varphi(D/g)},
  \qquad g = \gcd(D,m),
\]
unde $\mu$ este funcția lui Möbius, iar $\varphi$ este indicatorul lui Euler. (În particular, ori de câte ori $\gcd(m,D)=1$, deci și pentru $m=1$, suma este egală cu $\mu(D)$.)
\end{problema}

\begin{problema}
Fie $p$ un număr prim impar și $a \in \ZZ$ cu $p \nmid a$. Determinați restul împărțirii numărului
\[
  \prod_{k=1}^{p-1}\bigl(k^2 + a^2\bigr)
\]
la $p$.
\end{problema}

\begin{problema}
Fie $A$ un inel finit, cu $1\neq 0$, cu proprietatea că
\[
  (1+u)^{9}=1+u\qquad\text{pentru orice element inversabil } u\in A .
\]
\begin{parti}
\item Arătați că $x^{9}=x$ pentru orice $x\in A$.
\item Arătați că $A$ este comutativ.
\item Pentru fiecare număr natural $N\ge 2$, determinați numărul inelelor cu $N$ elemente, neizomorfe două câte două, care au proprietatea din enunț.
\end{parti}
\end{problema}

\begin{problema}
Fie $A$ un inel comutativ finit și
\[
  E(A)=\{\,e\in A : e^{2}=e\,\}
\]
mulțimea idempotenților săi. (Pentru $y\in A$ scriem $2y=y+y$ și $4y=y+y+y+y$.)
\begin{parti}
\item Fie $a\in A$ astfel încât elementul $1+4a$ este inversabil. Arătați că
  ecuația $x^{2}-x=a$ fie nu are nicio soluție în $A$, fie are exact $|E(A)|$
  soluții în $A$.
\item Presupunem că $A$ are $2^{m}$ elemente, unde $m\ge1$. Arătați că mulțimea
  \[
    M=\{\,x^{2}-x : x\in A\,\}
  \]
  este un subgrup al grupului $(A,+)$ și că $|M|\cdot|E(A)|=2^{m}$.
\end{parti}
\end{problema}

\begin{problema}
Fie $p$ un număr prim impar. Determinați numărul polinoamelor $f\in\ZZ_p[X]$ unitare, ireductibile peste $\ZZ_p$, de grad $2p$, care verifică egalitatea de polinoame
\[
  f(X+\cl 1)=f(X).
\]
\end{problema}

\begin{problema}
Fie $p$ un număr prim impar.
\begin{parti}
\item Determinați numărul soluțiilor $(x,y,z,t) \in \ZZ_p^4$ ale sistemului
\[
  xy = z + t, \qquad zt = x + y .
\]
\item Presupunem, în plus, că $3 \nmid p-1$. Determinați numărul perechilor $(a,b) \in \ZZ_p \times \ZZ_p$ pentru care fiecare dintre polinoamele $X^2 + aX + b$ și $X^2 + bX + a$ are cel puțin o rădăcină în $\ZZ_p$.
\end{parti}
\end{problema}

\begin{problema}
Fie $p$ un număr prim și $f,g\in\ZZ_p[X]$ două polinoame, de grade $d$, respectiv $e$, unde $2\le d\le p-1$, astfel încât
\[
  g\bigl(f(x)\bigr)=x\qquad\text{pentru orice } x\in\ZZ_p
\]
(așadar $f$ induce o bijecție a lui $\ZZ_p$, iar $g$ induce bijecția inversă). Notăm cu $q$ și $r$ câtul și restul împărțirii lui $p-1$ la $d$, adică $p-1=qd+r$, cu $0\le r\le d-1$.
\begin{parti}
\item Arătați că
\[
  r\,(e-1)\ \ge\ q\,(d-1).
\]
În particular, $d$ nu divide $p-1$.
\item Arătați că, dacă $r\ge 1$ și $r$ divide $d-1$, inegalitatea de la punctul (a) nu poate fi îmbunătățită: există polinoame $f,g$ ca mai sus, cu $\deg f=d$, pentru care $r(e-1)=q(d-1)$.
\item Determinați numerele prime $p\ge 5$ pentru care există un polinom $f\in\ZZ_p[X]$ de grad $3$ cu $f\bigl(f(x)\bigr)=x$ pentru orice $x\in\ZZ_p$.
\end{parti}
\end{problema}

\parteamare{Partea II \;·\; Soluții}{Fiecare enunț este reluat, apoi rezolvat complet.}
\subsectiuneclasa{Clasa a XII-a}{Structuri algebrice --- soluții complete.}{Soluții · Clasa a XII-a}

\nivel{olm}{Nivel OLM}{Faza locală --- o singură idee, bine aplicată.}
\setcounter{cat}{1}\setcounter{problemaX}{0}

\begin{problema}[Clasică]
Fie $G$ un grup în care $x^{2}=e$ pentru orice $x \in G$. Demonstrați
că $G$ este comutativ.
\end{problema}

\begin{solutie}
Fie $x,y \in G$ arbitrare. Din ipoteză, orice element este egal cu
propriul său invers: $x^{-1}=x$, $y^{-1}=y$, iar tot din ipoteză
$(xy)^{2}=e$, deci $(xy)^{-1}=xy$. Pe de altă parte, în orice grup
\[
  (xy)^{-1} = y^{-1}x^{-1} = yx.
\]
Comparând cele două expresii ale lui $(xy)^{-1}$ obținem
$xy=yx$. Cum $x,y$ au fost arbitrare, $G$ este comutativ.
\end{solutie}

\begin{problema}
Determinați toate grupurile $G$ pentru care există numere naturale
$m,n$ prime între ele astfel încât $x^{m}=x^{n}=e$ pentru orice
$x \in G$.
\end{problema}

\noindent\textbf{Răspuns:} doar grupul trivial $G=\{e\}$.

\begin{solutie}
Presupunem că $G$ satisface condiția, cu $\gcd(m,n)=1$. Din lema lui
Bézout există numere întregi $a,b$ cu
\[
  am+bn = 1 .
\]
Atunci, pentru orice $x \in G$,
\[
  x = x^{1} = x^{\,am+bn} = \left(x^{m}\right)^{a}\left(x^{n}\right)^{b}
    = e^{a}\,e^{b} = e .
\]
Așadar orice element al lui $G$ este egal cu $e$, deci $G=\{e\}$.
Reciproc, grupul trivial satisface evident condiția (pentru orice
alegere de $m,n$ prime între ele). Prin urmare singurul grup căutat
este cel trivial.
\end{solutie}

\begin{obs}
Se poate raționa și prin ordine: $x^{m}=e$ dă $\operatorname{ord}(x)\mid m$,
iar $x^{n}=e$ dă $\operatorname{ord}(x)\mid n$, deci
$\operatorname{ord}(x)\mid \gcd(m,n)=1$, adică $x=e$. Argumentul cu
Bézout este însă mai direct și nu folosește noțiunea
de ordin. Ipoteza de coprimalitate este esențială: fără ea
(de exemplu $m=n=2$) condiția devine $x^{2}=e$ și este satisfăcută
de multe grupuri netriviale.
\end{obs}

\begin{problema}[Clasică]
Fie $G$ un grup finit comutativ de ordin \emph{impar}. Demonstrați că
produsul tuturor elementelor sale este $e$.
\end{problema}

\begin{solutie}
Fie $P = \prod_{x \in G} x$ produsul tuturor elementelor; acesta are sens
deoarece $G$ este comutativ, deci ordinea factorilor nu contează.

Ideea este de a împerechea fiecare element cu inversul său. Un
element coincide cu propriul invers exact când $x^{2}=e$; notăm
$A=\{\, x \in G : x^{2}=e \,\}$.

\smallskip
\textit{Pasul 1: $A=\{e\}$.} Fie $x \in A$. Atunci $\operatorname{ord}(x) \mid 2$,
iar din teorema lui Lagrange (ordinul oricărui subgrup al unui grup finit
divide ordinul grupului; în particular, ordinul oricărui element
divide $|G|$) avem $\operatorname{ord}(x) \mid |G|$; cum $|G|$
este impar, rezultă $\operatorname{ord}(x) \mid \gcd(2,|G|)=1$, deci
$\operatorname{ord}(x)=1$ și $x=e$.

\smallskip
\textit{Pasul 2: restul elementelor se anulează în perechi.}
Elementele lui $G \setminus \{e\}$ satisfac $x \neq x^{-1}$, deci se
grupează în perechi disjuncte $\{x, x^{-1}\}$. Fiecare astfel de
pereche contribuie în produs cu $x \cdot x^{-1} = e$.

\smallskip
\textit{Concluzie.} Separând în produs contribuția lui $e$ de a
perechilor,
\[
  P = e \cdot \prod_{\{x,\,x^{-1}\}} \bigl(x\,x^{-1}\bigr)
    = e \cdot e \cdots e = e .
\]
\end{solutie}

\begin{obs}
Ipoteza de ordin impar este cea care garantează răspunsul curat.
La un grup comutativ de ordin \emph{par}, produsul tuturor elementelor este
egal cu produsul involutivelor și nu mai este neapărat $e$: de pildă
în $\mathbb{Z}/4$ produsul este $0+1+2+3=2 \neq 0$, iar în
$(\mathbb{Z}/p)^{\times}$ produsul tuturor elementelor este $-1$ --- exact
teorema lui Wilson.
\end{obs}

\begin{problema}[Clasică]
Fie $G$ un grup finit care are exact un element de ordin $2$, notat $a$.
Demonstrați că $a$ comută cu orice element al lui $G$ (adică
$a$ aparține centrului lui $G$).
\end{problema}

\begin{solutie}
Fie $g \in G$ oarecare și considerăm conjugatul $gag^{-1}$.
Conjugarea păstrează ordinul: $\bigl(gag^{-1}\bigr)^{2}
= ga^{2}g^{-1} = geg^{-1} = e$, iar $gag^{-1} \neq e$ (deoarece
$a \neq e$). Așadar $gag^{-1}$ este și el un element de ordin $2$.
Cum $a$ este \emph{unicul} element de ordin $2$, rezultă
$gag^{-1} = a$, adică $ga = ag$. Prin urmare $a$ comută cu orice
$g \in G$.
\end{solutie}

\begin{problema}
Fie $G$ un grup și $a,b \in G$ cu $aba^{-1}=b^{-1}$. Demonstrați
că $a^{2}b = ba^{2}$ (adică $a^{2}$ comută cu $b$).
\end{problema}

\begin{solutie}
Conjugăm încă o dată cu $a$:
\[
  a^{2}\,b\,a^{-2}
  = a\bigl(aba^{-1}\bigr)a^{-1}
  = a\,b^{-1}a^{-1}
  = \bigl(aba^{-1}\bigr)^{-1}
  = \bigl(b^{-1}\bigr)^{-1}
  = b .
\]
Așadar $a^{2}ba^{-2}=b$, adică $a^{2}b = ba^{2}$.
\end{solutie}

\begin{obs}
Conjugarea cu $a$ inversează pe $b$; aplicată de două ori,
readuce pe $b$, deci $a^{2}$ îi este ,,neutru''. Nu se folosește
nicio ipoteză despre ordinul lui $a$: concluzia ține pentru orice
$a,b$ cu $aba^{-1}=b^{-1}$.
\end{obs}

\begin{problema}
Un număr natural $n$ se numește \emph{liber de pătrate} dacă
nu este divizibil cu pătratul niciunui număr prim (echivalent, în
descompunerea $n=p_{1}p_{2}\cdots p_{r}$ toți factorii primi sunt
distincți). Fie $n$ liber de pătrate și $G$ un grup de ordin
$n$. Găsiți cel mai mic număr natural $k$ cu proprietatea
că $x^{k}=e$ pentru orice $x \in G$.
\end{problema}

\noindent\textbf{Răspuns:} $k=n$.

\begin{solutie}
Fie $k$ cel mai mic număr cu $x^{k}=e$ pentru orice $x$ (exponentul
grupului).

\smallskip
\textit{$k \mid n$.} Din teorema lui Lagrange $\operatorname{ord}(x) \mid n$
pentru orice $x$, deci $x^{n}=e$ pentru orice $x$. Așadar $n$ este un
exponent, iar cel mai mic exponent $k$ îl divide; în particular
$k \le n$.

\smallskip
\textit{$n \mid k$.} Fie $p$ un divizor prim al lui $n$. Din teorema lui
Cauchy (dacă un prim $p$ divide ordinul unui grup finit, atunci grupul
are un element de ordin $p$), $G$ conține un element $a$ de ordin $p$.
Cum $a^{k}=e$, rezultă $p=\operatorname{ord}(a) \mid k$. Prin urmare orice
divizor prim al lui $n$ divide $k$.

\smallskip
\textit{Concluzie.} Scriem $n=p_{1}p_{2}\cdots p_{r}$ cu primii distincți
(aici intervine ipoteza de a fi liber de pătrate). Fiecare $p_{i}$
divide $k$, iar fiind distincți, produsul lor divide $k$:
$n=p_{1}\cdots p_{r} \mid k$. Combinând cu $k \mid n$, obținem
$k=n$.
\end{solutie}

\medskip
\noindent\textbf{Soluție alternativă (exponentul ca cmmmc al ordinelor).}\enspace
Condiția $x^{k}=e$ pentru orice $x$ revine la $\operatorname{ord}(x)\mid k$
pentru orice $x$, adică la faptul că $k$ este un multiplu comun al
tuturor ordinelor. Cel mai mic astfel de $k$ este deci
\[
  k = \operatorname{lcm}\{\,\operatorname{ord}(x) : x\in G\,\}.
\]
Fiecare $\operatorname{ord}(x)$ divide $n$ (Lagrange), deci $k\mid n$. Pe de
altă parte, pentru fiecare prim $p\mid n$ teorema lui Cauchy dă un
element de ordin $p$, deci $p$ apare în acest cmmmc: $p\mid k$. Cum $n$
este liber de pătrate (produs de primi distincți), rezultă
$n\mid k$, așadar $k=n$. De remarcat că niciun element nu are
neapărat ordinul $n$ (de pildă în $S_{3}$, de ordin $6$, ordinele
sunt doar $1,2,3$), însă cmmmc-ul lor este exact $n$.
\hfill$\square$

\begin{obs}
Ipoteza de a fi liber de pătrate este esențială. Dacă $n$
ar avea un factor $p^{2}$, din primii distincți care divid $k$ nu s-ar
mai putea deduce $n \mid k$: de exemplu pentru $G=\mathbb{Z}/2\times\mathbb{Z}/2$
(ordin $4$) orice element satisface $x^{2}=e$, deci $k=2 \neq 4$.
\end{obs}

\begin{problema}
Determinați toate numerele naturale $n$ pentru care există un grup
$G$ de ordin $n$ cu proprietatea că, pentru orice divizor $d$ al lui
$n$, grupul $G$ conține exact un element de ordin $d$.
\end{problema}

\noindent\textbf{Răspuns:} $n \in \{1,2\}$.

\begin{solutie}
Deoarece $d=n$ divide $n$, ipoteza cere un element $g$ de ordin $n$. Atunci
$\langle g\rangle$ are $n$ elemente, deci $\langle g\rangle=G$: grupul este
\emph{ciclic}, $G\cong\mathbb{Z}/n$. Într-un grup ciclic de ordin $n$,
elementele de ordin $n$ sunt exact generatorii, în număr de
$\varphi(n)$. Unicitatea impune $\varphi(n)=1$, adică $n\in\{1,2\}$.
Reciproc, grupul trivial ($n=1$) și $\mathbb{Z}/2$ ($n=2$) verifică
evident condiția.
\end{solutie}

\medskip
\noindent\textbf{Soluție alternativă (numărare).}\enspace
Fie $N_{d}$ numărul elementelor de ordin $d$. Orice element are ordinul
un divizor al lui $n$ (Lagrange), deci
\[
  n = |G| = \sum_{d \mid n} N_{d} .
\]
Ipoteza dă $N_{d}=1$ pentru fiecare $d\mid n$, așadar
$n=\sum_{d\mid n}1=\tau(n)$, numărul divizorilor lui $n$. Însă
$\tau(n)=n$ doar pentru $n\in\{1,2\}$: pentru $n\ge 3$, numărul $n-1$
(care este $\ge 2$) nu divide $n$ (căci $n\equiv 1 \pmod{n-1}$), deci nu
toate numerele din $\{1,\dots,n\}$ sunt divizori și $\tau(n)<n$.
\hfill$\square$

\begin{obs}
Cerința forțează de la început existența unui element de
ordin $n$, deci ciclicitatea; de aici încolo totul se reduce la
$\varphi(n)=1$ (varianta întâi) sau la $\tau(n)=n$ (varianta a doua),
două egalități care se întâmplă doar pentru
$n\in\{1,2\}$.
\end{obs}

\begin{problema}
Fie $G$ un grup care are un \emph{unic} subgrup ciclic netrivial (adică
un singur subgrup ciclic diferit de $\{e\}$). Demonstrați că $G$
este ciclic; mai precis, $G\cong\mathbb{Z}/p$ pentru un număr prim $p$.
\end{problema}

\begin{solutie}
Fie $H$ unicul subgrup ciclic netrivial al lui $G$.

\smallskip
\textit{$G=H$, deci $G$ ciclic.} Pentru orice $x\in G$, $x\neq e$, subgrupul
$\langle x\rangle$ este ciclic și netrivial, deci $\langle x\rangle=H$
(prin unicitate); în particular $x\in H$. Cum și $e\in H$, rezultă
$G=H$, așadar $G$ este ciclic.

\smallskip
\textit{$|G|$ este prim.} Dacă $G$ ar fi ciclic infinit
($G\cong\mathbb{Z}$), ar avea subgrupurile ciclice netriviale
$2\mathbb{Z},3\mathbb{Z},\dots$, distincte --- contrazicând unicitatea.
Deci $G$ este finit, $|G|=n$. Pentru orice divizor $d$ al lui $n$, grupul
ciclic $G$ are un subgrup de ordin $d$, ciclic. Dacă $n$ ar avea un
divizor $d$ cu $1<d<n$, acel subgrup ar fi ciclic, netrivial și diferit
de $H=G$ --- din nou o contradicție. Prin urmare singurii divizori ai
lui $n$ sunt $1$ și $n$, deci $n$ este prim și $G\cong\mathbb{Z}/p$.
\end{solutie}

\begin{obs}
Cuvântul ,,netrivial'' este esențial: subgrupul
$\{e\}=\langle e\rangle$ este și el ciclic, deci dacă am număra
și subgrupul trivial, singurul grup cu un \emph{unic} subgrup ciclic ar
fi $G=\{e\}$ (orice $x\neq e$ ar produce $\langle x\rangle\neq\{e\}$, un al
doilea subgrup ciclic). De aceea condiția cu conținut privește
subgrupurile ciclice netriviale.
\end{obs}

\begin{problema}
Fie $G$ un grup cu proprietatea că pentru orice $x\in G$ există
$y\in G$ astfel încât $x^{2}=yxy^{-1}$ (adică $x^{2}$ este un
conjugat al lui $x$). Demonstrați că $G$ nu are elemente de ordin
par (în particular, dacă $G$ e finit, $|G|$ este impar).
\end{problema}

\begin{solutie}
Fie $x\in G$ un element de ordin finit $d=\operatorname{ord}(x)$. Din
ipoteză, $x^{2}=yxy^{-1}$, deci $x$ și $x^{2}$ sunt conjugate. Dar
conjugarea păstrează ordinul: $\operatorname{ord}\bigl(yxy^{-1}\bigr)=\operatorname{ord}(x)$,
așadar
\[
  \operatorname{ord}(x^{2})=\operatorname{ord}(x)=d .
\]
Pe de altă parte, într-un grup $\operatorname{ord}(x^{2})=\dfrac{d}{\gcd(d,2)}$.
Egalitatea $\dfrac{d}{\gcd(d,2)}=d$ forțează $\gcd(d,2)=1$, adică
$d$ \emph{impar}. Prin urmare niciun element (de ordin finit) nu are ordin
par; în particular, dacă $G$ este finit, prin teorema lui Lagrange
un element de ordin par ar da $2\mid|G|$, deci $|G|$ este impar.
\end{solutie}

\begin{obs}
Reciproca este falsă: ordinul impar \emph{nu} garantează că
$x$ și $x^{2}$ sunt conjugate. De pildă în orice grup abelian
netrivial conjugarea e trivială ($yxy^{-1}=x$), deci ipoteza ar cere
$x^{2}=x$, imposibil pentru $x\neq e$; astfel $\mathbb{Z}/3$ are ordin impar,
dar nu satisface proprietatea. Așadar condiția este strict mai tare
decât ,,ordin impar'' --- ea implică ordinul impar, fără a fi
echivalentă cu el.
\end{obs}

\begin{problema}
Există un grup $G$ care are două subgrupuri $K,L$ (nu neapărat
diferite) astfel încât produsul lor $KL=\{\,kl : k\in K,\ l\in L\,\}$
să fie subgrupul trivial $\{e\}$, cu $K$ sau $L$ netrivial? Justificați.
\end{problema}

\noindent\textbf{Răspuns:} nu; singura posibilitate este $K=L=\{e\}$.

\begin{solutie}
Orice subgrup conține elementul neutru $e$. Prin urmare, luând
$l=e\in L$, avem $k=k\cdot e\in KL$ pentru orice $k\in K$, deci
$K\subseteq KL$; simetric, luând $k=e\in K$, obținem $L\subseteq KL$.
Așadar
\[
  K\cup L\subseteq KL\qquad\text{întotdeauna.}
\]
Dacă $KL=\{e\}$, atunci $K\cup L\subseteq\{e\}$, deci $K=L=\{e\}$. Prin
urmare produsul a două subgrupuri nu poate fi trivial decât dacă
ambele sunt triviale.
\end{solutie}

\begin{obs}
Cauza este că $KL$ e mereu ,,cel puțin cât cel mai mare dintre
$K,L$'': mai precis, pentru grupuri finite
$|KL|=\dfrac{|K|\,|L|}{|K\cap L|}\ge\max(|K|,|L|)$, deci $KL=\{e\}$ forțează
$|K|=|L|=1$. O întrebare cu conținut privește nu \emph{produsul}, ci
\emph{intersecția} a două subgrupuri netriviale. De pildă, în grupul lui
Klein $\{e,a,b,ab\}$ subgrupurile $K=\{e,a\}$ și $L=\{e,b\}$ au
$K\cap L=\{e\}$, iar produsul lor este tot grupul; în schimb, în
$\mathbb{Z}/p^{k}$ ($p$ prim) oricare două subgrupuri netriviale au
intersecția netrivială, căci subgrupurile formează un lanț
(Teorema~\ref{teo:subgrupciclic}).
\end{obs}

\begin{problema}
Determinați toate inelele $(A,+,\cdot)$ cu proprietatea că $0=1$
(elementul neutru la adunare coincide cu elementul neutru la înmulțire).
\end{problema}

\noindent\textbf{Răspuns:} exact unul --- \emph{inelul nul} $A=\{0\}$.

\begin{solutie}
Arătăm mai întâi lema de calcul:
\emph{în orice inel, $x\cdot 0=0$ pentru orice $x$.} Într-adevăr,
din distributivitate,
\[
  x\cdot 0=x\cdot(0+0)=x\cdot 0+x\cdot 0 ,
\]
și scăzând $x\cdot 0$ (adică adunând opusul său în
grupul aditiv $(A,+)$) obținem $x\cdot 0=0$.

Fie acum $A$ un inel cu $0=1$. Pentru orice $x\in A$,
\[
  x=x\cdot 1=x\cdot 0=0 ,
\]
deci $A=\{0\}$: singurul element este $0$. Reciproc, mulțimea cu un
singur element $\{0\}$, cu operațiile $0+0=0$ și $0\cdot 0=0$, este
într-adevăr un inel (toate axiomele se verifică trivial, fiecare
egalitate având ambii membri egali cu $0$), iar în el $0$ joacă
simultan rolul lui $0$ și al lui $1$. Așadar există exact un
asemenea inel: inelul nul.
\end{solutie}

\begin{obs}
Răspunsul depinde de convenția din definiția inelului. Dacă
definiția cere doar existența elementului neutru $1$ (fără
condiția $1\neq 0$), răspunsul este cel de mai sus: inelul nul, unic.
Dacă însă definiția include axioma $1\neq 0$ (convenție
folosită în unele manuale, obligatorie la corpuri), atunci răspunsul
devine ,,niciun inel'' --- un enunț de tip ,,demonstrați că nu
există'', ca la problema cu $KL=\{e\}$ de mai sus. Morala e aceeași
în ambele variante: o singură egalitate între constantele
structurii prăbușește întreg inelul la un punct.
\end{obs}

\begin{problema}
Un grup $G$ se numește \emph{rezolubil} (solubil) dacă există un
lanț finit de subgrupuri
\[
  \{e\}=G_{0}\trianglelefteq G_{1}\trianglelefteq\cdots\trianglelefteq G_{m}=G
\]
în care, pentru fiecare $i$, $G_{i}$ este normal în $G_{i+1}$ și
câtul $G_{i+1}/G_{i}$ este abelian. Arătați că grupul $S_{3}$
al permutărilor mulțimii $\{1,2,3\}$ (de ordin $6$) este rezolubil,
exhibând un astfel de lanț.
\end{problema}

\begin{solutie}
Fie $A_{3}=\{e,(1\,2\,3),(1\,3\,2)\}$ subgrupul permutărilor pare
(subgrupul altern). Propunem lanțul
\[
  \{e\}\ \trianglelefteq\ A_{3}\ \trianglelefteq\ S_{3},
\]
și verificăm cele trei cerințe din definiție.

\smallskip
\textit{(1) $A_{3}\trianglelefteq S_{3}$.} Indicele este
$[S_{3}:A_{3}]=6/3=2$. Orice subgrup de indice $2$ este normal: pentru
$g\notin A_{3}$, clasa la stânga $gA_{3}$ și clasa la dreapta
$A_{3}g$ sunt ambele egale cu complementara $S_{3}\setminus A_{3}$ (cele
două clase, la stânga sau la dreapta, partiționează $S_{3}$
în $A_{3}$ și restul), deci $gA_{3}=A_{3}g$, adică
$gA_{3}g^{-1}=A_{3}$. (Pentru $g\in A_{3}$ egalitatea e clară.)

\smallskip
\textit{(2) $S_{3}/A_{3}$ este abelian.} Are ordinul $[S_{3}:A_{3}]=2$, deci
este ciclic ($\cong\mathbb{Z}/2$); orice grup de ordin $2$ e abelian.

\smallskip
\textit{(3) $A_{3}/\{e\}\cong A_{3}$ este abelian, și
$\{e\}\trianglelefteq A_{3}$.} Subgrupul trivial e normal în orice grup.
Iar $A_{3}$ are ordinul $3$, prim, deci este ciclic (generat de $(1\,2\,3)$,
căci $(1\,2\,3)^{2}=(1\,3\,2)$ și $(1\,2\,3)^{3}=e$), în
particular abelian.

Toate cele trei condiții din definiție sunt îndeplinite, deci
$S_{3}$ este rezolubil.
\end{solutie}

\begin{obs}
$S_{3}$ este cel mai mic grup \emph{neabelian}, și totuși e
rezolubil: rezolubilitatea nu cere comutativitate, ci doar ca structura
să se ,,desfacă'' în trepte cu câturi abeliene. Contrastul
apare la $S_{n}$ cu $n\ge 5$: acolo un asemenea lanț nu mai există,
fiindcă $A_{n}$ este simplu și neabelian --- deci $S_{n}$ nu este
rezolubil, fapt legat de imposibilitatea rezolvării prin radicali a
ecuației generale de grad $\ge 5$.
\end{obs}

\begin{problema}
Fie $A$ un inel și $x\in A$ un element pentru care există un
întreg $\ell\ge 1$ cu
\[
  x^{\ell}=0\qquad\text{și}\qquad (x+1)^{\ell}=0 .
\]
Demonstrați că $1=0$ în $A$ (adică $A$ este inelul nul).
\end{problema}

\begin{solutie}
Observăm întâi că din $x^{\ell}=0$ rezultă $x^{m}=0$
pentru orice $m\ge\ell$ (căci $x^{m}=x^{\ell}\,x^{m-\ell}=0$).

Elementul $1$ comută cu orice element al inelului (în particular
$x\cdot 1=1\cdot x=x$), deci $x$ și $1$ comută între ele. Aceasta
ne permite să aplicăm \emph{binomul lui Newton} --- care este valabil
într-un inel oarecare doar pentru elemente care comută:
\[
  0=(x+1)^{\ell}=\sum_{i=0}^{\ell}\binom{\ell}{i}x^{i}
   =x^{\ell}+\binom{\ell}{\ell-1}x^{\ell-1}+\cdots+\binom{\ell}{1}x+1 .
\]
Termenul $x^{\ell}$ este nul, deci
\[
  \binom{\ell}{\ell-1}x^{\ell-1}+\cdots+\binom{\ell}{1}x+1=0 . \tag{$\star$}
\]

\smallskip
\textit{Coborâre: din $x^{k}=0$ rezultă $x^{k-1}=0$ (pentru $1\le k\le\ell$).}
Înmulțim relația $(\star)$ cu $x^{k-1}$:
\[
  0=x^{k-1}\cdot 0=\binom{\ell}{\ell-1}x^{\ell-1+k-1}+\cdots+\binom{\ell}{1}x^{1+k-1}+x^{k-1}.
\]
Fiecare termen provenit dintr-un $x^{i}$ cu $i\ge 1$ devine $x^{i+k-1}$, cu
exponentul $i+k-1\ge k$; cum $x^{k}=0$, toți acești termeni sunt nuli.
Rămâne doar ultimul, deci
\[
  x^{k-1}=0 .
\]

\smallskip
\textit{Concluzia.} Pornim de la $x^{\ell}=0$ (ipoteză) și
aplicăm coborârea de $\ell$ ori:
\[
  x^{\ell}=0\ \Longrightarrow\ x^{\ell-1}=0\ \Longrightarrow\ \cdots
  \ \Longrightarrow\ x^{1}=0\ \Longrightarrow\ x^{0}=0 .
\]
Dar $x^{0}=1$, deci $1=0$. (Echivalent: ultimul pas se poate citi direct ---
odată ce $x=0$, ipoteza $(x+1)^{\ell}=0$ devine $1^{\ell}=1=0$.) Așadar
$A$ este inelul nul.
\end{solutie}

\begin{obs}
Ipoteza cere ca \emph{atât} $x$, cât și $x+1$ să fie
nilpotente, iar asta e imposibil într-un inel netrivial. O explicație
conceptuală: dacă $x$ e nilpotent, atunci $x+1$ este \emph{inversabil},
cu inversul $1-x+x^{2}-\cdots+(-1)^{\ell-1}x^{\ell-1}$ (produsul telescopează
la $1+(-1)^{\ell-1}x^{\ell}=1$). Un element nu poate fi însă
simultan inversabil și nilpotent: din $y$ inversabil și $y^{\ell}=0$
ar rezulta $1=(y^{-1})^{\ell}y^{\ell}=0$. Soluția de mai sus evită
această discuție, folosind doar binomul și înmulțirea cu
$x^{k-1}$.

Aceeași concluzie se obține și din ipoteza aparent mult mai tare
$(x+a)^{\ell}=0$ pentru \emph{orice} $a\in A$: e destul să luăm
$a=0$ și $a=1$.
\end{obs}

\begin{problema}
Fie $A$ un inel. Un element $x\in A$ se numește \emph{idempotent}
dacă $x^{2}=x$, și \emph{nilpotent} dacă există un
întreg $n\ge 1$ cu $x^{n}=0$. Demonstrați că singurul element
care este simultan idempotent și nilpotent este $0$.
\end{problema}

\begin{solutie}
Fie $x$ simultan idempotent și nilpotent: $x^{2}=x$ și $x^{n}=0$
pentru un $n\ge 1$.

\smallskip
\textit{Pas 1: toate puterile lui $x$ sunt egale cu $x$.} Arătăm prin
inducție că $x^{k}=x$ pentru orice $k\ge 1$. Pentru $k=1$ e evident.
Presupunând $x^{k}=x$, avem
\[
  x^{k+1}=x^{k}\cdot x=x\cdot x=x^{2}=x ,
\]
folosind ipoteza de inducție și idempotența. Deci $x^{k}=x$
pentru orice $k\ge 1$.

\smallskip
\textit{Pas 2: concluzia.} Aplicăm Pasul 1 pentru $k=n$:
\[
  x=x^{n}=0 ,
\]
ultima egalitate fiind chiar nilpotența. Așadar $x=0$.

Reciproc, $0$ este într-adevăr idempotent ($0^{2}=0$) și nilpotent
($0^{1}=0$), deci este singurul element cu ambele proprietăți.
\end{solutie}

\begin{obs}
Idempotența ,,îngheață'' toate puterile lui $x$ la $x$
însuși, iar nilpotența spune că una dintre acele puteri este
$0$ --- de unde $x=0$. Cele două proprietăți sunt, în rest,
foarte diferite: într-un inel pot exista idempotenți netriviali (de
pildă $(1,0)$ în $\mathbb{Z}\times\mathbb{Z}$) și nilpotenți
netriviali (de pildă $t$ în $\mathbb{Z}[t]/(t^{2})$), dar niciun
element nenul nu poate fi de ambele feluri.
\end{obs}

\begin{problema}
Fie $K$ un corp și $a\in K$, cu $a\neq 1$, având proprietatea
că, pentru orice element idempotent $x$ (adică $x^{2}=x$), elementul
$a-x$ este tot idempotent. Demonstrați că $\operatorname{char}K=2$.
\end{problema}

\begin{solutie}
\textit{Pas 1: într-un corp, singurii idempotenți sunt $0$ și $1$.}
Dacă $x^{2}=x$, atunci $x^{2}-x=0$, adică $x(x-1)=0$. Într-un corp
nu există divizori ai lui zero, deci $x=0$ sau $x-1=0$. Așadar
mulțimea idempotenților este exact $\{0,1\}$.

\smallskip
\textit{Pas 2: luăm $x=0$ și deducem $a=0$.} Elementul $0$ este
idempotent ($0^{2}=0$), deci, prin ipoteză, $a-0=a$ este idempotent. Din
Pasul 1, $a\in\{0,1\}$; cum $a\neq 1$ prin ipoteză, rămâne
\[
  a=0 .
\]

\smallskip
\textit{Pas 3: luăm $x=1$ și deducem caracteristica.} Elementul $1$
este idempotent ($1^{2}=1$), deci, prin ipoteză, $a-1$ este idempotent.
Cu $a=0$ (Pasul 2), aceasta înseamnă că $-1$ este idempotent:
\[
  (-1)^{2}=-1 .
\]
Dar $(-1)^{2}=1$, deci $1=-1$. Adunând $1$ în ambii membri,
\[
  2=0 .
\]
Prin urmare caracteristica lui $K$ divide $2$; cum $K$ este corp, avem
$1\neq 0$, deci caracteristica nu este $1$. Așadar
$\operatorname{char}K=2$.
\end{solutie}

\begin{obs}
Reciproc, în orice corp de caracteristică $2$ ipoteza chiar se
realizează, cu $a=0$: idempotenții sunt $0$ și $1$, iar
$a-0=0$ și $a-1=-1=1$ sunt într-adevăr idempotenți. Așadar
condiția din enunț caracterizează exact corpurile de
caracteristică $2$ (și forțează $a=0$).

Esențială este lipsa divizorilor lui zero: într-un inel oarecare
pot exista mulți idempotenți (de pildă toate cele patru elemente
ale lui $\mathbb{Z}/2\times\mathbb{Z}/2$), iar concluzia nu se mai
păstrează în aceeași formă.
\end{obs}


\begin{problema}
Fie $P \in \mathbb{C}[X]$ un polinom neconstant astfel încât $P(X)$ divide $P(X^2)$ în $\mathbb{C}[X]$. Arătați că orice rădăcină nenulă a lui $P$ este o rădăcină a unității.
\end{problema}

\begin{solutie}
Scriem $P(X^2) = P(X)\,Q(X)$, cu $Q \in \mathbb{C}[X]$. Dacă $r$ este o rădăcină a lui $P$, evaluăm în $X = r$:
\[
P(r^2) = P(r)\,Q(r) = 0,
\]
deci și $r^2$ este o rădăcină a lui $P$. Iterând, toate numerele
\[
r,\ r^2,\ r^4,\ r^8,\ \ldots,\ r^{2^k},\ \ldots
\]
sunt rădăcini ale lui $P$. Dar $P$ este nenul, deci are un număr finit de rădăcini: două dintre puterile de mai sus coincid, $r^{2^a} = r^{2^b}$ cu $a < b$. Pentru $r \neq 0$ putem simplifica prin $r^{2^a}$ și obținem
\[
r^{\,2^b - 2^a} = 1,
\]
cu exponentul $2^b - 2^a \geq 1$: așadar $r$ este o rădăcină a unității.

\medskip
\noindent\textbf{Observație.} Enunțul are conținut: $P(X) = X^2 - 1$ satisface ipoteza, deoarece $P(X^2) = X^4 - 1 = (X^2 - 1)(X^2 + 1)$, iar rădăcinile sale $\pm 1$ sunt într-adevăr rădăcini ale unității. În plus, excluderea rădăcinii nule este necesară: $P(X) = X^k$ divide $P(X^2) = X^{2k}$, dar rădăcina sa $0$ nu este o rădăcină a unității.
\end{solutie}

\begin{problema}
Determinați toate polinoamele $P \in \RR[X]$ cu proprietatea
\[
  P\bigl(P(x)\bigr) = P(x)^2 \qquad \text{pentru orice } x \in \RR.
\]
\end{problema}

\begin{solutie}
\textbf{Răspuns:} $P \equiv 0$, $P \equiv 1$ și $P(X) = X^2$.

Dacă $P \equiv c$ este constant, condiția devine $c = c^2$, deci $c \in \{0,1\}$: primele două soluții.

Fie acum $P$ neconstant și fie $y$ un element arbitrar al imaginii lui $P$: există $x_0$ cu $y = P(x_0)$. Înlocuind $x = x_0$ în ecuație:
\[
  P(y) = P\bigl(P(x_0)\bigr) = P(x_0)^2 = y^2.
\]
Așadar polinomul $P(Y) - Y^2$ se anulează în fiecare punct al imaginii lui $P$. Dar imaginea unui polinom neconstant este infinită: fiecare valoare este atinsă de cel mult $\deg P$ ori, iar $\RR$ este o mulțime infinită. Un polinom care se anulează pe o mulțime infinită este identic nul, deci
\[
  P(Y) = Y^2.
\]
Verificare: $P(P(x)) = (x^2)^2 = P(x)^2$. $\checkmark$

\medskip
\noindent\textbf{Observație.}\quad Schema merită reținută: o ecuație funcțională cu polinoame se transformă, printr-o substituție, într-o identitate polinomială valabilă pe o mulțime infinită, iar acolo polinoamele sunt complet rigide. Pasul ,,imaginea este infinită'' este cel care înlocuiește orice calcul de coeficienți.
\end{solutie}

\begin{problema}[Clasică]
Fie $K$ un corp finit cu cel puțin 3 elemente. Arătați că suma tuturor elementelor lui $K$ este 0. Ce se întâmplă pentru $K = \ZZ_2$?
\end{problema}

\begin{solutie}
Cum $|K| \geq 3$, există un element $a \in K \setminus \{0,1\}$. Aplicația $x \mapsto ax$ este o bijecție a lui $K$ (inversa ei este $x \mapsto a^{-1}x$), deci parcurge toate elementele corpului, doar că într-o altă ordine. Notând cu $S$ suma tuturor elementelor,
\[
S = \sum_{x \in K} x = \sum_{x \in K} ax = a \sum_{x \in K} x = aS,
\]
adică $(a-1)S = 0$. Dar $a-1 \neq 0$, iar într-un corp nu există divizori ai lui zero: $S = 0$.

Pentru $K = \ZZ_2$ suma este $0+1 = 1 \neq 0$: argumentul cade tocmai pentru că nu există niciun $a \in K \setminus \{0,1\}$.

\medskip
\noindent\textbf{Observație.} Trucul ,,înmulțirea cu un element fixat permută corpul'' merită reținut: el transformă o sumă fără nicio structură într-o ecuație pentru acea sumă. O idee înrudită (o aplicație injectivă a unei mulțimi finite în ea însăși este o permutare a ei) reapare în problema despre automorfismul lui Frobenius, de mai jos.
\end{solutie}

\begin{problema}
Fie $K$ un corp cu 8 elemente. Arătați că polinomul $X^2+X+1$ nu are nicio rădăcină în $K$, dar că polinomul $X^3+X+1$ are trei rădăcini distincte în $K$.
\end{problema}

\begin{solutie}
\emph{Pasul 1: caracteristica.} Subcorpul prim al lui $K$ este $\ZZ_p$ pentru un anumit număr prim $p$, iar $K$ este spațiu vectorial peste el, deci $8=|K|=p^m$: singura posibilitate este $p=2$. Așadar $1+1=0$ în $K$.

\emph{Pasul 2: $X^2+X+1$ nu are rădăcini.} Presupunem că $x\in K$ satisface $x^2+x+1=0$. Atunci $x\neq 0$ (altfel $1=0$) și, înmulțind cu $x-1$, din identitatea $(x-1)(x^2+x+1)=x^3-1$ obținem $x^3=1$: ordinul lui $x$ în grupul $(K^*,\cdot)$ divide $3$. Dar, din teorema lui Lagrange, ordinul divide și $|K^*|=7$; cum $(3,7)=1$, rezultă că ordinul este $1$, adică $x=1$. Însă atunci $1^2+1+1=1\neq 0$ (caracteristica fiind $2$): contradicție.

\emph{Pasul 3: $X^3+X+1$ are trei rădăcini.} Orice $x\in K^*$ satisface $x^7=1$ (consecință a teoremei lui Lagrange), deci cele $8$ elemente ale lui $K$ sunt toate rădăcini ale polinomului $X^8-X$, care are gradul $8$: elementele lui $K$ sunt \emph{exact} rădăcinile sale, fiecare dintre ele fiind simplă. Căutăm descompunerea lui $X^8-X$ în $\ZZ_2[X]$: efectuând direct înmulțirea, se verifică identitatea (în caracteristică $2$)
\[
(X^3+X+1)(X^3+X^2+1)=X^6+X^5+X^4+X^3+X^2+X+1,
\]
iar membrul drept este $\dfrac{X^7-1}{X-1}$, deci
\[
X^8-X=X\,(X-1)\,(X^3+X+1)\,(X^3+X^2+1).
\]
Cele $8$ rădăcini se repartizează între factori: $X$ are rădăcina $0$, $X-1$ are rădăcina $1$, iar celelalte $6$ elemente sunt rădăcini ale produsului celor două polinoame de gradul $3$. Fiecare factor de gradul trei are cel mult $3$ rădăcini, iar împreună trebuie să acopere $6$ elemente distincte: prin urmare fiecare dintre ei are \emph{exact} $3$. În particular $X^3+X+1$ are trei rădăcini distincte în $K$.

\medskip
\noindent\textbf{Observație.} Numărarea de la Pasul 3 este tipică pentru corpurile finite: identitatea $x^{|K|}=x$, valabilă pentru toate elementele, transformă întrebările despre rădăcini în problema repartizării celor $|K|$ rădăcini ale lui $X^{|K|}-X$ între factorii ireductibili. Observați și că Pasul 2 este în acord cu Pasul 3: polinomul $X^2+X+1$ nu apare în descompunere tocmai pentru că nu are rădăcini în $K$.
\end{solutie}

\begin{problema}[Clasică]
Fie $K$ un corp finit de caracteristică $p$.
\begin{parti}
\item Arătați că aplicația $F : K \to K$, $F(x) = x^p$, este un automorfism al lui $K$.
\item Arătați că $F(x) = x$ dacă și numai dacă $x$ aparține subcorpului prim $\{0,\, 1,\, 1+1,\, \dots\} \simeq \ZZ_p$.
\end{parti}
\end{problema}

\begin{solutie}
\emph{(a)} Multiplicativitatea este imediată: $F(xy) = (xy)^p = x^p y^p$ (elementele comută), iar $F(1) = 1$. Pentru aditivitate folosim binomul lui Newton, valabil pentru elemente care comută:
\[
(x+y)^p = \sum_{k=0}^{p} \binom{p}{k} x^k y^{p-k}.
\]
Pentru $1 \leq k \leq p-1$, numărul $\binom{p}{k} = \frac{p\,(p-1)\cdots(p-k+1)}{k!}$ este divizibil cu $p$: numărătorul conține factorul prim $p$, în timp ce $k!$ nu îl conține (toți factorii săi sunt $< p$). În caracteristică $p$ acești termeni se anulează, deci
\[
F(x+y) = (x+y)^p = x^p + y^p = F(x) + F(y).
\]
$F$ este injectivă: dacă $x^p = y^p$, atunci $(x-y)^p = x^p - y^p = 0$ (aditivitatea, aplicată lui $x-y$; pentru $p = 2$ semnele coincid oricum), iar într-un corp o putere nulă obligă baza să fie nulă: $x = y$. O aplicație injectivă a unei mulțimi finite în ea însăși este bijectivă. Așadar $F$ este un automorfism.

\emph{(b)} Punctele fixe ale lui $F$ sunt rădăcinile polinomului $X^p - X$, care are cel mult $p$ rădăcini în corpul $K$. Pe de altă parte, subcorpul prim are exact $p$ elemente, iar fiecare dintre ele este punct fix: pentru clasa lui $a \in \{0, 1, \dots, p-1\}$, egalitatea $a^p \equiv a \pmod p$ este mica teoremă a lui Fermat. Am găsit deci $p$ puncte fixe într-o mulțime cu cel mult $p$ elemente: punctele fixe sunt \emph{exact} elementele subcorpului prim.

\medskip
\noindent\textbf{Observație.} \ Automorfismul lui Frobenius este instrumentul central al corpurilor finite: de pildă, într-un corp finit de caracteristică $2$ ridicarea la pătrat este chiar $F$, deci o bijecție, așa că orice element este un pătrat; tot el stă în spatele descompunerilor de felul celei din problema precedentă.
\end{solutie}

\nivel{ojm}{Nivel OJM}{Faza județeană --- combină două rezultate.}
\setcounter{cat}{2}\setcounter{problemaX}{0}

\begin{problema}[Clasică]
Fie $G$ un grup și fie
\[
  H = \{\, x \in G \ :\ x^{2}=e \,\}.
\]
Presupunem că $H$ este subgrup al lui $G$. Demonstrați că $H$
este comutativ.
\end{problema}

\begin{solutie}
Fie $x,y \in H$. Din definiția lui $H$ avem $x^{2}=e$ și $y^{2}=e$,
adică $x^{-1}=x$ și $y^{-1}=y$.

Cum $H$ este subgrup, este închis la operația grupului, deci
$xy \in H$; prin urmare $(xy)^{2}=e$, ceea ce înseamnă
$(xy)^{-1}=xy$. Dar, în orice grup,
\[
  (xy)^{-1} = y^{-1}x^{-1} = yx.
\]
Așadar $xy=yx$. Elementele $x,y \in H$ fiind arbitrare, $H$ este
comutativ.
\end{solutie}

\begin{obs}
Ipoteza ,,$H$ este subgrup'' este exact ceea ce garantează
închiderea la produs, iar aceasta este esențială: în
general mulțimea involuțiilor \emph{nu} este închisă. De
exemplu, în grupul $S_{3}$ al permutărilor a trei elemente,
transpozițiile $(1\,2)$ și $(1\,3)$ satisfac $x^{2}=e$, dar
produsul lor $(1\,2)(1\,3)=(1\,3\,2)$ este un ciclu de lungime $3$, deci
nu mai are pătratul egal cu $e$. Astfel
$\{x : x^{2}=e\}=\{e,(1\,2),(1\,3),(2\,3)\}$ are $4$ elemente, număr
care nu divide $6=|S_{3}|$, deci nici măcar nu poate fi subgrup.
\end{obs}

\begin{problema}[Clasică]
Fie $G$ un grup finit de ordin \emph{par}. Demonstrați că există
un element $x \in G$, $x \neq e$, cu $x^{2}=e$ (adică $G$ conține
un involutiv netrivial).
\end{problema}

\begin{solutie}
Considerăm aplicația $\varphi : G \to G$, $\varphi(x)=x^{-1}$.
Ea este o bijecție și satisface $\varphi(\varphi(x))=x$, deci este
o involuție a mulțimii $G$.

Punctele fixe ale lui $\varphi$ sunt exact elementele cu $x^{-1}=x$, adică
\[
  A = \{\, x \in G \ :\ x^{2}=e \,\}.
\]
Un element $x \notin A$ satisface $x \neq x^{-1}$, deci se grupează
împreună cu $x^{-1}$ într-o pereche $\{x, x^{-1}\}$ cu două
elemente distincte. Aceste perechi sunt disjuncte și acoperă
întreaga mulțime $G \setminus A$, prin urmare $|G \setminus A|$
este par. De aici
\[
  |A| = |G| - |G \setminus A|
\]
are aceeași paritate ca $|G|$. Cum $|G|$ este par, rezultă că
$|A|$ este par.

Dar $e \in A$ (căci $e^{2}=e$), deci $|A| \geq 1$; fiind par,
obținem $|A| \geq 2$. Așadar există un element $x \in A$ cu
$x \neq e$, adică $x^{2}=e$ și $x \neq e$.
\end{solutie}

\begin{obs}
Rezultatul este cazul particular $p=2$ al teoremei lui Cauchy: orice grup
finit al cărui ordin este divizibil cu un număr prim $p$ conține
un element de ordin $p$. Argumentul de mai sus îl obține însă
complet elementar, doar prin împărțirea lui $G$ în perechi.
\end{obs}

\begin{problema}[Clasică]
Fie $G$ un grup și $a,b \in G$ astfel încât $a^{4}=e$ și
$aba^{-1}=b^{2}$. Demonstrați că $b^{15}=e$.
\end{problema}

\begin{solutie}
Arătăm prin inducție după $k \geq 1$ că
\[
  a^{k}\,b\,a^{-k} = b^{\,2^{k}} .
\]
Pentru $k=1$ este chiar ipoteza. Presupunând egalitatea pentru $k$,
folosim că $a\,b^{n}\,a^{-1}=\bigl(aba^{-1}\bigr)^{n}$:
\[
  a^{k+1}\,b\,a^{-(k+1)}
  = a\bigl(a^{k}ba^{-k}\bigr)a^{-1}
  = a\,b^{\,2^{k}}a^{-1}
  = \bigl(aba^{-1}\bigr)^{2^{k}}
  = \bigl(b^{2}\bigr)^{2^{k}}
  = b^{\,2^{k+1}} .
\]
Pentru $k=4$, cum $a^{4}=e$, membrul stâng este $b$, deci
$b = b^{\,2^{4}} = b^{16}$, adică $b^{15}=e$.
\end{solutie}

\begin{obs}
Rezultă că $\operatorname{ord}(b) \mid 15$, deci ordinul lui $b$
poate fi doar $1$, $3$, $5$ sau $15$.
\end{obs}

\begin{problema}
Fie $G$ un grup de ordin $2n$ și $a \in G$ un element de ordin cel
puțin $n$. Presupunem că există $x \in G$ cu $x \neq a^{k}$
pentru orice $k$ (adică $x \notin \langle a\rangle$) astfel încât
$ax=xa$. Demonstrați că $a$ comută cu orice element al lui $G$
(adică $a$ aparține centrului lui $G$).
\end{problema}

\begin{solutie}
Fie $H=\langle a\rangle$ subgrupul generat de $a$, deci $|H|=\operatorname{ord}(a)$.

\smallskip
\textit{Pasul 1: $\operatorname{ord}(a)=n$, iar $H$ are indice $2$.}
Existența unui $x \notin H$ arată că $H \neq G$, deci
$\operatorname{ord}(a) < 2n$. Pe de altă parte $\operatorname{ord}(a) \geq n$,
iar din teorema lui Lagrange $\operatorname{ord}(a) \mid 2n$. Singurul
divizor al lui $2n$ situat în intervalul $[\,n,\,2n\,)$ este însă
$n$ (căci un divizor $d \geq n$ al lui $2n$ dă $2n/d \leq 2$, deci
$d=n$ sau $d=2n$). Așadar $\operatorname{ord}(a)=n$, deci $|H|=n$ și
$[G:H]=2$.

\smallskip
\textit{Pasul 2: descompunerea în două clase.} Cum $[G:H]=2$ și
$x \notin H$, clasa $xH$ este diferită de $H$, iar
\[
  G = H \,\cup\, xH .
\]
Prin urmare orice element $g \in G$ este fie de forma $a^{k}$ (dacă
$g \in H$), fie de forma $x\,a^{k}$ (dacă $g \in xH$).

\smallskip
\textit{Pasul 3: $a$ comută cu ambele tipuri.} Cu $a^{k}$ elementul $a$
comută evident (puteri ale lui $a$). Pentru $x\,a^{k}$, folosind
$ax=xa$,
\[
  a\,(x a^{k}) = (ax)\,a^{k} = (xa)\,a^{k} = x\,a^{k+1} = (x a^{k})\,a .
\]
Așadar $a$ comută cu orice element al lui $G$, deci aparține
centrului.
\end{solutie}

\medskip
\noindent\textbf{Soluție alternativă (centralizatorul lui $a$).}\enspace
Fie
\[
  C = \{\, g \in G : ga=ag \,\}
\]
centralizatorul lui $a$; se verifică imediat că este subgrup al lui
$G$. El conține toate puterile lui $a$, deci $\langle a\rangle \subseteq C$
și $|C| \geq \operatorname{ord}(a) \geq n$. În plus $x \in C$
(deoarece $ax=xa$) și $x \notin \langle a\rangle$, deci $C$ are cel
puțin $n+1$ elemente. Dar $|C| \mid 2n$ (Lagrange), iar singurul divizor
al lui $2n$ strict mai mare decât $n$ este $2n$ însuși; așadar
$|C|=2n$, deci $C=G$. Prin urmare $a$ comută cu orice element al lui $G$.
\hfill$\square$

\begin{obs}
Ipoteza $x \notin \langle a\rangle$ joacă un rol dublu: exclude cazul
$\operatorname{ord}(a)=2n$ (în care $\langle a\rangle=G$ ar fi ciclic,
dar nu ar exista niciun $x$ în afară) și, forțând astfel
$\operatorname{ord}(a)=n$, furnizează totodată elementul din a doua
clasă cu care $a$ comută.
\end{obs}

\begin{problema}
Fie $G$ un grup și $a,b \in G$ astfel încât orice element al
lui $G$ comută cu $a$ sau cu $b$. Notând $C(x)=\{\,g\in G : gx=xg\,\}$
centralizatorul lui $x$, ipoteza spune că $C(a)\cup C(b)=G$.
Demonstrați că $a$ sau $b$ aparține centrului $Z(G)$ (cel
puțin unul dintre ele este central).
\end{problema}

Motorul demonstrației este următoarea lemă, de sine
stătătoare.

\begin{lemat}
Un grup nu poate fi reuniunea a două subgrupuri proprii ale sale: dacă
$H,K \leq G$ și $H\cup K=G$, atunci $H=G$ sau $K=G$.
\end{lemat}

\begin{proof}
Presupunem, prin absurd, că $H\neq G$ și $K\neq G$. Cum $H\cup K=G$
dar $K\neq G$, există $h\in H\setminus K$; simetric, există
$k\in K\setminus H$. Elementul $hk$ aparține lui $G=H\cup K$, deci
$hk\in H$ sau $hk\in K$.
\begin{itemize}[nosep,leftmargin=1.4em]
  \item Dacă $hk\in H$, atunci $k=h^{-1}(hk)\in H$, contrazicând
        $k\notin H$.
  \item Dacă $hk\in K$, atunci $h=(hk)k^{-1}\in K$, contrazicând
        $h\notin K$.
\end{itemize}
Ambele cazuri sunt imposibile, deci presupunerea a fost falsă: $H=G$ sau
$K=G$.
\end{proof}

\begin{solutie}
Pentru orice $x\in G$, centralizatorul $C(x)$ este subgrup al lui $G$:
conține pe $e$; dacă $g,h$ comută cu $x$, atunci și $gh$
comută cu $x$; iar din $gx=xg$ rezultă $xg^{-1}=g^{-1}x$, deci
$g^{-1}\in C(x)$.

Aplicăm lema subgrupurilor $H=C(a)$ și $K=C(b)$: din
$C(a)\cup C(b)=G$ rezultă $C(a)=G$ sau $C(b)=G$. Dar $C(x)=G$ înseamnă
exact că $x$ comută cu orice element al lui $G$, adică $x\in Z(G)$.
Așadar $a\in Z(G)$ sau $b\in Z(G)$.
\end{solutie}

\begin{obs}
Lema este specifică numărului \emph{doi}: un grup poate fi reuniunea
a \emph{trei} subgrupuri proprii --- de pildă grupul lui Klein
$\mathbb{Z}/2\times\mathbb{Z}/2$ este reuniunea celor trei subgrupuri ale
sale de ordin $2$. Prin urmare, varianta cu trei condiții
($C(a)\cup C(b)\cup C(c)=G$) nu mai forțează neapărat
centralitatea vreunuia dintre $a,b,c$.
\end{obs}

\begin{problema}
Fie $G$ un grup cu proprietatea că $xy \in Z(G)$ pentru orice
$x,y\in G$ cu $x\neq e$ și $y\neq e$. Demonstrați că $G$ este
comutativ.
\end{problema}

\begin{solutie}
Fie $x,y\in G$ două elemente diferite de $e$. Din ipoteză
$xy\in Z(G)$, deci $xy$ comută cu orice element al lui $G$, în
particular cu $x$:
\[
  (xy)\,x = x\,(xy) \quad\Longrightarrow\quad xyx = x^{2}y .
\]
Scriind membrul stâng ca $x\,(yx)$ și pe cel drept ca $x\,(xy)$,
simplificăm la stânga cu $x$ și obținem $yx = xy$. Așadar
orice două elemente diferite de $e$ comută; cum $e$ comută
oricum cu tot, $G$ este comutativ.
\end{solutie}

\medskip
\noindent\textbf{Soluție alternativă (prin puteri, arătând $Z(G)=G$).}\enspace
Pentru orice $x\neq e$, luând $y=x$ în ipoteză obținem
$x^{2}\in Z(G)$; deci toate pătratele sunt centrale. Fie acum $x\neq e$.
Dacă $x^{2}\neq e$, atunci $x^{3}=x^{2}\cdot x$ este produs de două
elemente diferite de $e$, deci $x^{3}\in Z(G)$; cum și $x^{2}\in Z(G)$
și $Z(G)$ este subgrup,
\[
  x = x^{3}\,\bigl(x^{2}\bigr)^{-1} \in Z(G).
\]
Rămân elementele cu $x^{2}=e$. Dacă există vreun $z$ cu
$z^{2}\neq e$, atunci $z^{2}\in Z(G)\setminus\{e\}$, iar pentru un asemenea
$x$ avem $x = (x\cdot z^{2})\,(z^{2})^{-1}\in Z(G)$ (produs de elemente
centrale, folosind $x, z^{2}\neq e$). Dacă însă $z^{2}=e$ pentru
orice $z$, atunci $G$ este comutativ prin rezultatul clasic al involuțiilor
(problema OLM.1). În toate cazurile fiecare element este central, deci
$Z(G)=G$. \hfill$\square$

\begin{obs}
Concluzia se poate întări: ipoteza dă de fapt $Z(G)=G$, adică
un grup comutativ. Prima soluție o obține dintr-o singură
observație (,,$xy$ comută cu $x$''), pe când a doua urmează
ideea de a exprima $x$ prin puteri centrale, $x=x^{3}x^{-2}$.
\end{obs}

\begin{problema}
Fie $G$ un grup finit și $f,g$ două endomorfisme ale lui $G$ cu
proprietatea
\[
  f\bigl(x^{2}y^{4}\bigr) = g\bigl(x^{2}\bigr)\,g\bigl(y^{4}\bigr)
  \qquad \text{pentru orice } x,y\in G.
\]
Determinați sub ce condiție asupra ordinului $|G|$ rezultă
$f=g$.
\end{problema}

\noindent\textbf{Răspuns:} dacă și numai dacă $|G|$ este
impar.

\begin{solutie}
\textit{Reducerea la pătrate.} Cum $f$ este morfism,
$f(x^{2}y^{4})=f(x^{2})f(y^{4})$, deci ipoteza devine
\[
  f(x^{2})\,f(y^{4}) = g(x^{2})\,g(y^{4}) \qquad \text{pentru orice } x,y.
\]
Luând $y=e$ obținem $f(x^{2})=g(x^{2})$ pentru orice $x$: așadar
$f$ și $g$ coincid pe toate pătratele. (Condiția pe puterile a
patra este redundantă, căci $f(x^{4})=f(x^{2})^{2}=g(x^{2})^{2}=g(x^{4})$.)

\smallskip
\textit{Suficiența ($|G|$ impar).} Dacă $|G|$ este impar, ridicarea
la pătrat este surjectivă: pentru orice $a\in G$, fie
$d=\operatorname{ord}(a)\mid|G|$, deci $d$ impar, și
\[
  a = a^{d+1} = \bigl(a^{(d+1)/2}\bigr)^{2},
\]
unde $(d+1)/2$ este întreg. Prin urmare orice element este un pătrat;
cum $f$ și $g$ coincid pe pătrate, rezultă $f(a)=g(a)$ pentru
orice $a$, adică $f=g$.

\smallskip
\textit{Necesitatea ($|G|$ par).} Dacă $|G|$ este par, concluzia poate
cădea. De pildă, în $G=\mathbb{Z}/2$ orice pătrat este $0$,
deci condiția ,,$f=g$ pe pătrate'' este vidă; endomorfismele
lui $\mathbb{Z}/2$ fiind $\mathrm{id}$ și $0$, alegerea $f=\mathrm{id}$,
$g=0$ satisface ipoteza, dar $f\neq g$.

Același fenomen apare pentru \emph{orice} ordin par $2m$. Fie
$G=(\mathbb{Z}/2m,+)$, $f(x)=mx$ și $g=0$. Aplicația $f$ este
endomorfism ($m(x+y)=mx+my$), iar în notație aditivă ipoteza cere
$f(2x+4y)=g(2x)+g(4y)$, adică $m(2x+4y)=0$; aceasta are loc, căci
$m(2x+4y)=2m\,(x+2y)$ este multiplu de $2m$. Totuși
$f(\widehat 1)=\widehat m\neq\widehat 0=g(\widehat 1)$, deci $f\neq g$.
Așadar, pentru fiecare ordin par există un grup și două endomorfisme
care satisfac ipoteza fără a fi egale.
\end{solutie}

\begin{obs}
Totul depinde de faptul că exponenții $2$ și $4$ sunt
amândoi pari, deci condiția se reduce la pătrate --- de aici
nevoia unei ipoteze pe ordin. Dacă exponenții ar fi \emph{primi
între ei} (de pildă $x^{2}$ și $x^{3}$), atunci din
$f(x)=f(x^{3})\,f(x^{2})^{-1}=g(x^{3})\,g(x^{2})^{-1}=g(x)$ ar rezulta
$f=g$ necondiționat.
\end{obs}

\begin{problema}
Fie $G$ un grup de ordin $2n$, cu $n$ \emph{impar}, astfel încât
$x^{n+2}\in Z(G)$ pentru orice $x\in G$. Demonstrați că $G$ este
comutativ.
\end{problema}

Demonstrația se sprijină pe următoarea lemă (Bézout pentru
puteri centrale).

\begin{lemat}[Bézout pentru puteri centrale]\label{lema:bezoutZ}
Fie $x\in G$ și $a,b\in\mathbb{N}^{*}$. Dacă $x^{a}\in Z(G)$ și
$x^{b}\in Z(G)$, atunci $x^{\gcd(a,b)}\in Z(G)$.
\end{lemat}

\begin{proof}
Fie $d=\gcd(a,b)$. Din relația lui Bézout există $u,v\in\mathbb{Z}$
cu $d=ua+vb$, deci
\[
  x^{d}=x^{ua+vb}=\bigl(x^{a}\bigr)^{u}\,\bigl(x^{b}\bigr)^{v}.
\]
Cum $Z(G)$ este subgrup (închis la puteri întregi --- inclusiv
negative --- și la produse) și $x^{a},x^{b}\in Z(G)$, ambii factori
aparțin lui $Z(G)$, deci și produsul lor $x^{d}$.
\end{proof}

\begin{solutie}
Din teorema lui Lagrange $\operatorname{ord}(x)\mid 2n$, deci $x^{2n}=e\in Z(G)$
pentru orice $x$. Astfel, pentru fiecare $x$, atât $x^{n+2}$ (din
ipoteză) cât și $x^{2n}$ aparțin centrului; aplicând lema de mai sus cu $a=n+2$ și $b=2n$,
\[
  x^{\gcd(n+2,\,2n)}\in Z(G).
\]
Calculăm $\gcd(n+2,2n)$: din $2n=2(n+2)-4$ rezultă
$\gcd(n+2,2n)=\gcd(n+2,4)$, iar $n$ impar face $n+2$ impar, deci
$\gcd(n+2,4)=1$. Așadar $\gcd(n+2,2n)=1$ și
\[
  x=x^{1}\in Z(G)\qquad\text{pentru orice }x\in G.
\]
Prin urmare $Z(G)=G$, adică $G$ este comutativ.
\end{solutie}

\begin{obs}
Paritatea impară a lui $n$ este esențială și ascuțită:
ea dă $\gcd(n+2,2n)=1$. Pentru $n$ par concluzia poate cădea --- de
exemplu $G=Q_{8}$ (ordin $8$, $n=4$, $n+2=6$): orice element satisface
$x^{6}\in\{\pm1\}=Z(Q_{8})$ (căci $x^{6}=x^{2}=-1$ pentru $x$ de ordin
$4$, și $x^{6}=1$ altfel), deci ipoteza e verificată, însă
$Q_{8}$ nu e comutativ; acolo lema dă doar $x^{2}$ central. Lema
însăși este unealta recurentă a ,,exponenților primi
între ei'': cazul $\gcd(a,b)=1$ dă direct $x\in Z(G)$ (compară
cu problema OLM.2, unde apare $H=\{e\}$ în locul lui $Z(G)$).
\end{obs}

\begin{problema}
Determinați toate numerele naturale $n$ pentru care există un grup
$G$ de ordin $n$ astfel încât, pentru orice întreg $k$ cu
$1\le k\le \tfrac{n}{3}$, aplicația $f_{k}:G\to G$, $f_{k}(x)=x^{k}$,
este injectivă.
\end{problema}

\noindent\textbf{Răspuns:} $n\in\{1,4\}$ sau $n$ prim (echivalent:
$n\le 5$ sau $n$ prim).

Soluția aplică următorul rezultat (Corolarul~\ref{lema:puterep} din cartea
de teorie, ,,ridicarea la puterea $p$''), pe care îl redemonstrăm.

\begin{lemat}[ridicarea la puterea $p$; Corolarul~\ref{lema:puterep}]
Fie $G$ un grup finit cu $n$ elemente, $p\ge 2$ un întreg, și
$H=\{x^{p}\mid x\in G\}$ mulțimea puterilor de ordin $p$. Atunci
\[
  (p,n)=1 \iff H=G.
\]
(Aici $H=G$ înseamnă că aplicația $x\mapsto x^{p}$ este
surjectivă; pe un grup finit aceasta echivalează cu injectivitatea.)
\end{lemat}

\begin{proof}
$(\Rightarrow)$ Fie $(p,n)=1$; există $a,b\in\mathbb{Z}$ cu $ap+bn=1$
(Bézout). Pentru orice $g\in G$, ordinul lui $g$ divide $n$ (Lagrange),
deci $g^{n}=e$ și astfel $g^{bn}=e$. Atunci
\[
  g=g^{ap+bn}=\bigl(g^{a}\bigr)^{p}\,g^{bn}=\bigl(g^{a}\bigr)^{p}\in H,
\]
deci $H=G$.

$(\Leftarrow)$ Prin contrapoziție: presupunem $(p,n)=d>1$ și fie $q$
un divizor prim al lui $d$. Cum $q\mid n$, teorema lui Cauchy dă un
element $x$ de ordin $q$. Cum $q\mid p$, avem $x^{p}=e=e^{p}$; deci
$x\neq e$ și $e$ au aceeași imagine, aplicația $x\mapsto x^{p}$
nu e injectivă, deci (pe mulțime finită) nici surjectivă:
$H\neq G$.
\end{proof}

\begin{solutie}
Din lemă, pentru orice grup de ordin $n$ și orice întreg
$k\ge 2$: $f_{k}$ este injectivă dacă și numai dacă
$(k,n)=1$; iar $f_{1}$ este identitatea, injectivă întotdeauna.
Cum echivalența are loc în \emph{orice} grup de ordin $n$,
proprietatea din enunț nu depinde de structura grupului, ci doar de $n$:
ea revine la
\[
  (k,n)=1 \qquad\text{pentru orice întreg } 2\le k\le \tfrac{n}{3},
\]
adică la: \emph{niciun divizor prim $p$ al lui $n$ nu satisface
$p\le \tfrac{n}{3}$}. (Dacă toți divizorii primi depășesc
$\tfrac n3$, orice $k\le\tfrac n3$ e prim cu $n$; reciproc, un divizor prim
$p\le\tfrac n3$ dă exponentul rău $k=p$.)

\smallskip
\textit{Cazul $n$ prim sau $n\le 5$.} Pentru $n$ prim, unicul divizor prim
este $n>\tfrac{n}{3}$, deci condiția e îndeplinită. Pentru
$n\le 5$ nu există întregi $k$ cu $2\le k\le\tfrac{n}{3}<2$, deci
condiția e vidă și orice grup convine; dintre acestea, valorile
neacoperite de primalitate sunt $n=1$ și $n=4$.

\smallskip
\textit{Cazul $n\ge 6$ compus.} Fie $p$ cel mai mic divizor prim al lui
$n$; cum $n$ este compus, $p\le\sqrt{n}$. Condiția cere
$p>\tfrac{n}{3}$, deci $\tfrac{n}{3}<\sqrt{n}$, adică $n<9$. Rămân
de verificat $n=6$ și $n=8$: pentru $n=6$, divizorul prim $2\le 2=\tfrac63$
eșuează (exponentul $k=2$ strică injectivitatea, prin lemă);
pentru $n=8$, la fel $2\le\tfrac83$. Deci niciun $n\ge 6$ compus nu convine.

\smallskip
\textit{Concluzie.} Numerele căutate sunt $n\in\{1,4\}$ și $n$ prim
--- echivalent, $n\le 5$ sau $n$ prim.
\end{solutie}

\begin{obs}
Două remarci. Întâi, deoarece lema este valabilă în
\emph{orice} grup de ordin $n$, formulările ,,există un grup'' și
,,orice grup'' sunt aici echivalente: proprietatea depinde numai de $n$.
Apoi, granița este ascuțită exact la $n=6$: exponentul care
strică este $k=2=\tfrac{n}{3}$, chiar capătul intervalului --- cu
inegalitatea strictă $k<\tfrac{n}{3}$, valorile $n=6$ și $n=9$ ar
intra și ele în răspuns.
\end{obs}

\begin{problema}[Clasică]
Fie $a<b$ numere reale.
\begin{enumerate}[label=(\alph*),nosep]
  \item Înzestrați mulțimea $G=(a,b)$ (intervalul deschis) cu o
        structură de grup comutativ.
  \item Aceeași cerință pentru intervalul \emph{închis}
        $G=[a,b]$.
\end{enumerate}
\end{problema}

\begin{solutie}
Instrumentul comun este \emph{transportul de structură}: dacă
$(H,\cdot)$ este grup și $\varphi:G\to H$ este o \emph{bijecție}
oarecare, atunci operația
\[
  x*y:=\varphi^{-1}\bigl(\varphi(x)\cdot\varphi(y)\bigr)
\]
face din $(G,*)$ un grup, iar $\varphi$ devine izomorfism: asociativitatea,
comutativitatea, neutrul $\varphi^{-1}(e_{H})$ și inversele
$\varphi^{-1}\bigl(\varphi(x)^{-1}\bigr)$ se moștenesc direct prin
bijecție. Nu se cere nicăieri continuitate.

\smallskip
\textit{(a) Intervalul deschis.} Aplicația
\[
  \varphi:(a,b)\to\mathbb{R},\qquad
  \varphi(x)=\tan\!\Bigl(\pi\,\frac{x-a}{b-a}-\frac{\pi}{2}\Bigr)
\]
este bijecție (strict crescătoare, cu limitele $\mp\infty$ la
capete). Transportând $(\mathbb{R},+)$, obținem grupul comutativ
$\bigl((a,b),*\bigr)$ cu $x*y=\varphi^{-1}(\varphi(x)+\varphi(y))$, izomorf
cu $(\mathbb{R},+)$. Concret, elementul neutru este mijlocul
$\tfrac{a+b}{2}$, iar inversul lui $x$ este simetricul $a+b-x$ (se
verifică $\varphi(a+b-x)=-\varphi(x)$).

\smallskip
\textit{(b) Intervalul închis --- unde continuitatea moare.}
Atenție: nu există nicio bijecție \emph{continuă}
$[a,b]\to\mathbb{R}$ (imaginea continuă a unui compact este compactă),
deci rețeta de la (a) nu se poate repeta cu o formulă continuă.
Transportul însă cere doar o bijecție, pe care o construim
,,cu mâna''.

\textit{Pas 1: reducere la $[0,1]$.} Aplicația afină
$x\mapsto a+(b-a)x$ este bijecție $[0,1]\to[a,b]$, deci e suficient
să punem structura pe $[0,1]$ și s-o transportăm afin.

\textit{Pas 2: o bijecție $\psi:[0,1]\to[0,1)$.} Problema e că
$[0,1]$ are ,,un punct în plus'' față de $[0,1)$ --- capătul
$1$ --- și vrem să-l eliminăm fără să pierdem sau
să dublăm vreun punct. Aici intervine imaginea \emph{hotelului lui
Hilbert}: un hotel cu o infinitate de camere $1,\tfrac12,\tfrac13,\dots$,
toate ocupate; ca să eliberăm camera $1$ fără să dăm
pe nimeni afară, rugăm fiecare oaspete să treacă în
camera următoare din șir --- cel din $\tfrac1n$ se mută în
$\tfrac{1}{n+1}$. Toți oaspeții rămân cazați, dar
camera $1$ e acum liberă. Punctele care nu sunt de forma $\tfrac1n$ nu
sunt ,,oaspeți'' și rămân neatinse. Formal:
\[
  \psi(x)=\begin{cases}\dfrac{1}{n+1}, & x=\dfrac{1}{n}\ (n\ge 1),\\[4pt]
  x, & \text{altfel.}\end{cases}
\]
Ea mută $1\mapsto\tfrac12\mapsto\tfrac13\mapsto\cdots$ de-a lungul
șirului și lasă restul pe loc. Verificăm că e
bijecție: \emph{injectivă} --- puncte distincte au imagini
distincte (pe șir, $\tfrac1n\mapsto\tfrac{1}{n+1}$ e strict
descrescător în $n$, iar punctele fixate rămân distincte
între ele și de imaginile șirului, care sunt tot de forma
$\tfrac1m$); \emph{surjectivă pe $[0,1)$} --- capătul $1$ nu mai e
atins (bine, e eliminat), iar orice $y\in[0,1)$ este atins: dacă
$y=\tfrac1m$ cu $m\ge 2$, atunci $y=\psi\!\big(\tfrac{1}{m-1}\big)$, iar
dacă $y$ nu e de forma $\tfrac1m$, atunci $y=\psi(y)$. Așadar
$\psi:[0,1]\to[0,1)$ e bijecție.

\textit{Pas 3: structura pe $[0,1)$ și transportul.} Pe $[0,1)$ avem
grupul comutativ al adunării modulo $1$: $s\oplus t=\{s+t\}$ (partea
fracționară) --- acesta este $\mathbb{R}/\mathbb{Z}$, grupul
cercului, cu neutrul $0$ și inversul lui $t\neq 0$ egal cu $1-t$.
Transportăm prin $\psi$:
\[
  x*y:=\psi^{-1}\bigl(\psi(x)\oplus\psi(y)\bigr),
\]
obținând un grup comutativ pe $[0,1]$ (izomorf cu
$\mathbb{R}/\mathbb{Z}$), apoi conjugăm cu bijecția afină din
Pasul 1 pentru a-l muta pe $[a,b]$.
\end{solutie}

\begin{obs}
Cele două puncte au obstacole complet diferite. La (a) transportul este
continuu și grupul obținut este ,,dreapta'' $(\mathbb{R},+)$; la (b)
obstrucția este topologică --- niciun grup nu poate fi pus pe
$[a,b]$ printr-o formulă continuă în ambele variabile cu inversare
continuă (compactitatea interzice bijecția continuă cu
$\mathbb{R}$, iar segmentul nu e cerc) --- așa că orice soluție
este obligatoriu discontinuă, iar grupul natural obținut este
,,cercul'' $\mathbb{R}/\mathbb{Z}$. Morala: transportul de structură
funcționează prin \emph{orice} bijecție de mulțimi;
continuitatea este un bonus, nu o cerință.
\end{obs}

\begin{problema}
Fie $G$ un grup finit de ordin $2n$. Câte soluții poate avea
ecuația $x^{n}=e$? (Determinați toate valorile posibile ale
numărului de soluții și arătați că fiecare apare.)
\end{problema}

\noindent\textbf{Răspuns:} exact $n$ sau exact $2n$; ambele valori se
realizează.

\begin{solutie}
Fie $N$ numărul de soluții ale ecuației $x^{n}=e$ în $G$.

\smallskip
\textit{Pas 1 (Frobenius).} Teorema lui Frobenius spune că, într-un
grup finit $G$, dacă $d$ divide $|G|$, atunci numărul soluțiilor
ecuației $x^{d}=e$ este multiplu de $d$. Aici $d=n$ divide $|G|=2n$, deci
$n\mid N$. Cum $e$ este soluție ($N\ge 1$) și $N\le|G|=2n$, singurii
multipli de $n$ din intervalul $[1,2n]$ sunt $n$ și $2n$. Așadar
$N\in\{n,2n\}$.

\smallskip
\textit{Pas 2 (ambele valori apar).}
\begin{itemize}[nosep,leftmargin=1.4em]
  \item \emph{$N=n$:} luăm $G=\mathbb{Z}/2n$ (ciclic). Scris aditiv,
        $x^{n}=e$ devine $nx\equiv 0\pmod{2n}$, adică $2\mid x$; soluțiile
        sunt $\{0,2,4,\dots,2n-2\}$ --- exact $n$.
  \item \emph{$N=2n$ (pentru $n$ par):} luăm
        $G=\mathbb{Z}/2\times\mathbb{Z}/n$, de ordin $2n$. Exponentul său
        este $\operatorname{lcm}(2,n)=n$ (căci $n$ e par), deci $x^{n}=e$
        pentru \emph{orice} $x$: toate cele $2n$ elemente sunt soluții.
\end{itemize}
\end{solutie}

\begin{obs}
Valoarea $2n$ apare \emph{doar} pentru $n$ par. Într-adevăr, dacă
$n$ este impar, cum $2\mid 2n=|G|$, teorema lui Cauchy dă un element $t$
de ordin $2$; atunci $t^{n}=t$ ($n$ impar), deci $t\neq e$ nu este
soluție, $N<2n$, și rămâne $N=n$. Prin urmare pentru $n$
impar ecuația are \emph{întotdeauna} exact $n$ soluții, iar
alternativa $n$ / $2n$ apare doar pentru $n$ par.
\end{obs}

\begin{problema}
O \emph{relație de echivalență} pe o mulțime $M$ este o
relație $\rho$ care este reflexivă ($x\,\rho\,x$ pentru orice $x$),
simetrică ($x\,\rho\,y\Rightarrow y\,\rho\,x$) și tranzitivă
($x\,\rho\,y$ și $y\,\rho\,z\Rightarrow x\,\rho\,z$). O relație de
echivalență $\rho$ pe un grup $G$ se numește \emph{congruență}
dacă este compatibilă cu operația: $x\,\rho\,y$ și
$x'\,\rho\,y'$ implică $xx'\,\rho\,yy'$. Determinați toate
congruențele pe grupul $(\mathbb{Z},+)$.
\end{problema}

\noindent\textbf{Răspuns:} exact relațiile de congruență
modulo $n$ ($x\,\rho\,y\iff n\mid x-y$), pentru $n\ge 0$ --- și niciuna
în plus.

\begin{solutie}
Folosim teorema din cartea de teorie (Teorema~\ref{teo:congruenta}, ,,congruențele
provin din subgrupuri normale''): \emph{o relație de echivalență
$\rho$ pe un grup $G$ este congruență dacă și numai dacă
există un subgrup normal $N\trianglelefteq G$ cu
\[
  x\,\rho\,y \iff x^{-1}y\in N ,
\]
iar atunci $N=\widehat e$ (clasa neutrului)}. Prin urmare congruențele pe
$G$ sunt în corespondență bijectivă cu subgrupurile normale
ale lui $G$.

Aplicăm pentru $G=(\mathbb{Z},+)$. Cum $\mathbb{Z}$ este abelian, orice
subgrup este normal, deci congruențele corespund \emph{tuturor}
subgrupurilor lui $\mathbb{Z}$. Dar subgrupurile lui $(\mathbb{Z},+)$ sunt
exact $n\mathbb{Z}=\{nk : k\in\mathbb{Z}\}$, pentru $n\ge 0$: dacă
$H\le\mathbb{Z}$ este netrivial, fie $n$ cel mai mic element strict pozitiv
al său; pentru orice $h\in H$, împărțirea cu rest dă
$h=qn+r$ cu $0\le r<n$, iar $r=h-qn\in H$, deci $r=0$ prin minimalitate ---
așadar $H=n\mathbb{Z}$ (iar $H=\{0\}=0\cdot\mathbb{Z}$ în cazul
trivial).

Subgrupului $n\mathbb{Z}$ îi corespunde congruența (scris aditiv)
\[
  x\,\rho_{n\mathbb{Z}}\,y \iff y-x\in n\mathbb{Z} \iff n\mid x-y \iff x\equiv y \pmod n .
\]
Deci congruențele pe $\mathbb{Z}$ sunt exact relațiile $\equiv\!\!\pmod n$,
$n\ge 0$. Extremele: $n=0$ dă egalitatea ($x\equiv y\pmod 0\iff x=y$), iar
$n=1$ dă relația totală (orice două numere sunt congruente).
\end{solutie}

\begin{obs}
Aceasta luminează și \emph{terminologia}: ,,congruența modulo $n$''
din aritmetică este chiar o congruență în sensul algebric (o
echivalență compatibilă cu $+$), și teorema arată că
pe $\mathbb{Z}$ \emph{nu există altele}. Frumusețea aplicației: nu
verificăm ,,de mână'' compatibilitatea fiecărei relații
candidat, ci reducem totul la clasificarea subgrupurilor lui $\mathbb{Z}$,
via corespondența congruențe $\leftrightarrow$ subgrupuri normale.
\end{obs}

\begin{problema}[Clasică]
Fie $f:G\to H$ un morfism de grupuri finite, cu
$\operatorname{Ker}f=\{x\in G:f(x)=e_{H}\}$ și $\operatorname{Im}f=f(G)$.
Arătați că toate fibrele nevide ale lui $f$ (mulțimile
$f^{-1}(h)$, $h\in\operatorname{Im}f$) au același număr de elemente
și că
\[
  |G|=|\operatorname{Ker}f|\cdot|\operatorname{Im}f| .
\]
\end{problema}

\begin{solutie}
Considerăm pe $G$ relația $x\sim y\iff f(x)=f(y)$. Ea este
reflexivă, simetrică și tranzitivă (fiind definită prin
egalitate în $H$), și compatibilă cu operația: dacă
$f(x)=f(y)$ și $f(x')=f(y')$, atunci
\[
  f(xx')=f(x)f(x')=f(y)f(y')=f(yy'),
\]
deci $xx'\sim yy'$. Așadar $\sim$ este o \emph{congruență}
(Propoziția~\ref{prop:congmorf}, ,,congruența asociată unui morfism''), iar
clasele ei sunt exact fibrele lui $f$. Clasa unui element:
\[
  f(x)=f(y)\iff f(x^{-1}y)=e_{H}\iff x^{-1}y\in\operatorname{Ker}f
  \iff y\in x\operatorname{Ker}f,
\]
deci fibra care îl conține pe $x$ este exact clasa la stânga
$x\operatorname{Ker}f$.

Toate clasele la stânga ale unui subgrup au același cardinal ca
subgrupul însuși (aplicația $k\mapsto xk$ este bijecție
$\operatorname{Ker}f\to x\operatorname{Ker}f$), deci fiecare fibră nevidă
are exact $|\operatorname{Ker}f|$ elemente. Cum $x\sim y\iff f(x)=f(y)$,
clasele sunt în corespondență bijectivă cu valorile lui $f$,
adică cu $\operatorname{Im}f$; numărul lor este $|\operatorname{Im}f|$.
Clasele partajează $G$, deci
\[
  |G|=(\text{număr de clase})\cdot|\operatorname{Ker}f|
     =|\operatorname{Im}f|\cdot|\operatorname{Ker}f| .
\]
\end{solutie}

\begin{obs}
Aceasta este forma cardinală a primei teoreme de izomorfism
($G/\operatorname{Ker}f\cong\operatorname{Im}f$): numărând, ea devine
$|G|=|\operatorname{Ker}f|\cdot|\operatorname{Im}f|$. Ideea-cheie e că
,,a avea aceeași imagine'' este o \emph{congruență}, deci clasele
ei sunt automat clasele unui subgrup (normal), fără vreun calcul
direct de normalitate --- iar de aici numărarea e imediată. În
particular, ordinul lui $\operatorname{Im}f$ divide $|G|$.
\end{obs}

\begin{problema}
Fie $\pi=2+i\in\mathbb{Z}[i]$ și $\pi\mathbb{Z}[i]=\{\pi z : z\in\mathbb{Z}[i]\}$
subgrupul multiplilor săi în grupul aditiv $(\mathbb{Z}[i],+)$.
Arătați că grupul cât $\mathbb{Z}[i]/\pi\mathbb{Z}[i]$ este
izomorf cu $(\mathbb{Z}/5,+)$.
\end{problema}

\begin{solutie}
Definim
\[
  f:\mathbb{Z}[i]\to\mathbb{Z}/5,\qquad f(a+bi)=(a+3b)\bmod 5 .
\]

\textit{$f$ este morfism surjectiv.} Aditivitatea e imediată:
$f\bigl((a+bi)+(c+di)\bigr)=(a+c)+3(b+d)\equiv(a+3b)+(c+3d)=f(a+bi)+f(c+di)$.
Iar $f(1)=1$, deci $f$ e surjectiv și $\operatorname{Im}f=\mathbb{Z}/5$.

\textit{Nucleul este $\pi\mathbb{Z}[i]$.} Pentru $z=(2+i)(c+di)=(2c-d)+(c+2d)i$,
\[
  f(z)=(2c-d)+3(c+2d)=5c+5d\equiv 0\pmod 5 ,
\]
deci $\pi\mathbb{Z}[i]\subseteq\operatorname{Ker}f$. Reciproc, ambele subgrupuri
au indice $5$ în $\mathbb{Z}[i]$: pentru $\operatorname{Ker}f$ pentru
că $f$ e surjectiv pe $\mathbb{Z}/5$; pentru $\pi\mathbb{Z}[i]$ pentru
că $\{0,1,2,3,4\}$ este un sistem complet de reprezentanți --- orice
$a+bi$ verifică $a+bi\equiv a+3b\pmod{\pi}$ (căci $i-3=(2+i)(i-1)\in\pi\mathbb{Z}[i]$),
iar $a+3b$ e întreg, redus modulo $5$ deoarece $5=(2+i)(2-i)\in\pi\mathbb{Z}[i]$.
Un subgrup conținut într-altul, ambele de același indice finit,
coincid: $\operatorname{Ker}f=\pi\mathbb{Z}[i]$.

\textit{Prima teoremă de izomorfism.} Așadar
\[
  \boxed{\;\mathbb{Z}[i]/\pi\mathbb{Z}[i]=\mathbb{Z}[i]/\operatorname{Ker}f\;\cong\;\operatorname{Im}f=\mathbb{Z}/5\;}
\]
(izomorfismul indus fiind $a+bi+\pi\mathbb{Z}[i]\mapsto(a+3b)\bmod 5$).
\end{solutie}

\begin{obs}
Numărul $5=|2+i|^{2}$ e chiar norma lui $\pi=2+i$, iar $\pi$ este un
\emph{prim} al lui $\mathbb{Z}[i]$ (normă primă); câtul
$\mathbb{Z}[i]/\pi\cong\mathbb{Z}/5$ este de fapt corpul rezidual
$\mathbb{F}_{5}$ la $\pi$. Reprezentantul cheie e $i\equiv 3\pmod{\pi}$
(fiindcă $2+i\equiv 0$ dă $i\equiv-2\equiv 3$), de unde
$a+bi\mapsto a+3b$. Este exact pasul care arată că primul $5$ se
\emph{descompune} în $\mathbb{Z}[i]$ ($5=(2+i)(2-i)$, cei doi factori
neasociați), spre deosebire de $2$ (ramificat) sau de un prim
$\equiv 3\pmod 4$ (inert).
\end{obs}

\begin{problema}
Fie $\sigma(z)=5/\bar z$ inversiunea planului față de cercul de
rază $\sqrt5$ centrat în origine. Grupul simetriilor rețelei
$\mathbb{Z}[i]$ (rotațiile $z\mapsto i^{k}z$ și conjugarea
$z\mapsto\bar z$) permută mulțimea $X$ a punctelor din $\mathbb{Z}[i]$
fixate de $\sigma$. Descrieți $X$ ca mulțime cu acțiune de grup.
\end{problema}

\begin{solutie}
\textit{Punctele fixe.} Cum $5$ e real, $\sigma(z)=z\iff 5/\bar z=z\iff
|z|^{2}=5$; punctele fixe formează cercul de rază $\sqrt5$ (iar
$\sigma$ e involuție: $\sigma^{2}=\operatorname{id}$). Printre întregii
Gauss, $|z|^{2}=5$ înseamnă $a^{2}+b^{2}=5$, cu soluțiile
$(a,b)\in\{(\pm1,\pm2),(\pm2,\pm1)\}$; deci
\[
  X=\{\pm1\pm2i,\ \pm2\pm i\},\qquad |X|=8 .
\]

\textit{Grupul.} Simetriile rețelei --- rotațiile $z\mapsto i^{k}z$
($k=0,1,2,3$) și conjugarea $z\mapsto\bar z$ --- formează grupul
diedral $D_{4}$ de ordin $8$; ele păstrează norma, deci permută
$X$.

\textit{Orbită--stabilizator.} Aplicația orbită
$\varphi:D_{4}\to X$, $\varphi(g)=g\cdot(2+i)$, evaluată pe cele opt
simetrii,
\[
\begin{array}{llll}
  \operatorname{id}:2+i, & i\cdot:-1+2i, & i^{2}\cdot:-2-i, & i^{3}\cdot:1-2i,\\[2pt]
  \ \overline{\phantom{x}}:2-i, & i\,\overline{\phantom{x}}:1+2i,
  & i^{2}\,\overline{\phantom{x}}:-2+i, & i^{3}\,\overline{\phantom{x}}:-1-2i,
\end{array}
\]
ia opt valori distincte --- exact $X$. Deci $X$ este o singură orbită
și $\operatorname{Stab}(2+i)=\{\operatorname{id}\}$. Reamintim
\textbf{teorema orbită--stabilizator}: dacă un grup finit $H$
acționează pe o mulțime $X$ și $x\in X$, atunci aplicația
orbită $h\mapsto h\cdot x$ induce o bijecție
$H/\!\operatorname{Stab}(x)\to\operatorname{Orb}(x)$ între clasele la
stânga după stabilizator și orbită; în particular
$|\!\operatorname{Orb}(x)|=[H:\operatorname{Stab}(x)]$. Aplicată aici,
ea induce o bijecție de $D_{4}$-mulțimi
\[
  \boxed{\;X\;\cong\;D_{4}/\!\operatorname{Stab}(2+i)=D_{4}/\{\operatorname{id}\}\;\cong\;D_{4}\;}
\]
Așadar $X$ este un \emph{$D_{4}$-torsor}: cele opt puncte fixe ale
involuției $\sigma$ sunt în bijecție canonică cu simetriile
rețelei, fiecare obținându-se dintr-o unică simetrie
aplicată lui $2+i$.
\end{solutie}

\begin{obs}
Atenție: aceasta \emph{nu} este prima teoremă de izomorfism --- $X$
nu e grup, iar $\operatorname{Stab}(2+i)$ nu e neapărat normal, deci
$D_{4}/\!\operatorname{Stab}$ e doar o mulțime de clase (aici, tot
$D_{4}$, fiindcă stabilizatorul e trivial). Rezultatul e \emph{teorema
orbită--stabilizator}, ruda ei pentru acțiuni. Aritmetic, cele opt
puncte sunt asociații lui $2+i$ și ai lui $2-i$ (factorii lui
$5=(2+i)(2-i)$), câte patru, legați prin conjugare.
\end{obs}

\begin{problema}
Fie $G$ un grup finit de ordin $n$ și $\varphi\in\operatorname{Aut}(G)$
un automorfism fixat. Arătați că, pentru orice $x\in G$ fixat,
numărul soluțiilor $g\in G$ ale ecuației $gx=x\varphi(g)$ divide
$n$.
\end{problema}

\begin{solutie}
Cheia este că $g\cdot x:=g\,x\,\varphi(g)^{-1}$ definește o
\emph{acțiune} a lui $G$ pe mulțimea $G$ (conjugarea răsucită
prin $\varphi$; pentru $\varphi=\operatorname{id}$, conjugarea obișnuită).
Într-adevăr, $e\cdot x=e\,x\,\varphi(e)^{-1}=x$, iar folosind că
$\varphi$ e morfism ($\varphi(gh)^{-1}=\varphi(h)^{-1}\varphi(g)^{-1}$),
\[
  (gh)\cdot x=gh\,x\,\varphi(h)^{-1}\varphi(g)^{-1}
  =g\bigl(h\,x\,\varphi(h)^{-1}\bigr)\varphi(g)^{-1}=g\cdot(h\cdot x).
\]

Ecuația $gx=x\varphi(g)$ se rescrie $g\,x\,\varphi(g)^{-1}=x$, adică
$g\cdot x=x$: mulțimea soluțiilor în $g$ este exact
\emph{stabilizatorul} lui $x$,
\[
  \operatorname{Stab}(x)=\{g\in G : g\cdot x=x\}.
\]
Stabilizatorul unui punct este întotdeauna subgrup al lui $G$, deci
(Lagrange) numărul soluțiilor $|\!\operatorname{Stab}(x)|$ divide
$|G|=n$.
\end{solutie}

\begin{obs}
Prin orbită--stabilizator, și orbitele (numite \emph{clase de
conjugare răsucită}) au mărimi $[G:\operatorname{Stab}(x)]$ care
divid $n$ --- exact ca la clasele de conjugare obișnuite (cazul
$\varphi=\operatorname{id}$). Atenție însă la varianta
\emph{duală}: dacă fixăm $g$ și căutăm $x$-urile cu
$g\cdot x=x$, obținem mulțimea punctelor fixe ale permutării
$x\mapsto g\cdot x$, care \emph{nu} e subgrup și poate fi chiar vidă
--- acolo nu mai există divizibilitate.
\end{obs}

\begin{problema}[Clasică]
Arătați că nu există niciun corp $K$ pentru care grupul
aditiv $(K,+)$ să fie izomorf cu grupul multiplicativ
$(K^{*},\cdot)$, unde $K^{*}=K\setminus\{0\}$.
\end{problema}

\begin{solutie}
Un \emph{izomorfism de grupuri} $f:(K,+)\to(K^{*},\cdot)$ este o
\emph{bijecție} $f:K\to K^{*}$ care duce adunarea în înmulțire:
\[
  f(x+y)=f(x)\cdot f(y)\qquad\text{pentru orice }x,y\in K.
\]
Presupunem că un asemenea $f$ există și ajungem la o
contradicție.

\smallskip
\textit{Pas 1: $f$ duce neutrul în neutru, $f(0)=1$.} Din
$f(0)=f(0+0)=f(0)\cdot f(0)$ și din faptul că $f(0)\in K^{*}$ e
inversabil, împărțind cu $f(0)$ obținem $f(0)=1$.

\smallskip
\textit{Pas 2: ce înseamnă ,,dublul e neutru'' de fiecare parte.}
Aplicând regula lui $f$ la $y=x$:
\[
  f(x+x)=f(x)\cdot f(x)=f(x)^{2}.
\]
Deci, folosind și $f(0)=1$,
\[
  x+x=0\iff f(x+x)=f(0)\iff f(x)^{2}=1 .
\]
Cum $f$ e bijecție, ea pune în corespondență unu-la-unu
mulțimile
\[
  S_{+}=\{x\in K : x+x=0\}\qquad\text{și}\qquad
  S_{\times}=\{y\in K^{*} : y^{2}=1\},
\]
deci acestea au \emph{același număr de elemente}.

\smallskip
\textit{Pas 3: câte elemente are fiecare.} Într-un corp nu există
divizori ai lui zero, deci
\[
  y^{2}=1\iff(y-1)(y+1)=0\iff y=1\ \text{sau}\ y=-1 ,
\]
adică $S_{\times}=\{1,-1\}$: are $2$ elemente dacă $-1\neq 1$
(adică $\operatorname{char}K\neq 2$) și $1$ element dacă $-1=1$
($\operatorname{char}K=2$).

Pe de altă parte, $x+x=0$ înseamnă $2x=0$. Aici $2x$ înseamnă
$x+x$; notând $2_{K}:=1+1\in K$ (elementul ,,doi'' al corpului), avem
$x+x=(1+1)\cdot x=2_{K}\cdot x$. Într-un corp nu există divizori ai lui
zero, deci egalitatea $2_{K}\cdot x=0$ implică $2_{K}=0$ sau $x=0$; iar
$2_{K}=0$ înseamnă exact $1+1=0$, adică $\operatorname{char}K=2$.
Prin urmare:
\begin{itemize}[nosep,leftmargin=1.4em]
  \item dacă $\operatorname{char}K\neq 2$, atunci $2_{K}\neq 0$, deci din
        $2_{K}\cdot x=0$ rămâne $x=0$; așadar $S_{+}=\{0\}$ are
        $1$ element;
  \item dacă $\operatorname{char}K=2$, atunci $2_{K}=0$, deci
        $2_{K}\cdot x=0$ pentru orice $x$, adică $S_{+}=K$.
\end{itemize}

\smallskip
\textit{Pas 4: contradicția.} Din Pasul 2, $|S_{+}|=|S_{\times}|$.
\begin{itemize}[nosep,leftmargin=1.4em]
  \item Dacă $\operatorname{char}K\neq 2$: $|S_{+}|=1$ dar
        $|S_{\times}|=2$ --- imposibil.
  \item Dacă $\operatorname{char}K=2$: $|S_{\times}|=1$, deci ar trebui
        $|S_{+}|=1$; dar $S_{+}=K$, deci $|K|=1$ --- imposibil, căci un
        corp are cel puțin două elemente ($0\neq 1$).
\end{itemize}
În ambele cazuri ajungem la o contradicție, deci nu există niciun
corp $K$ cu $(K,+)\cong(K^{*},\cdot)$.
\end{solutie}

\begin{obs}
Ideea e că un izomorfism păstrează numărul de soluții ale
oricărei ecuații scrise cu operația grupului; aici, ecuația
,,dublul unui element e neutrul'' devine $x+x=0$ de partea aditivă și
$y^{2}=1$ de partea multiplicativă, iar numărul lor diferă
(obstrucția e mereu elementul $-1$). La \emph{inele} întrebarea
analogă degenerează din alt motiv: $(A,\cdot)$ nici măcar nu e
grup, fiindcă $0$ nu are invers; singura posibilitate
$(A,+)\cong(A,\cdot)$ ar forța $0$ să fie inversabil, deci $0=1$,
adică inelul nul $\{0\}$. La corpuri, $0$ e scos din discuție
($K^{*}=K\setminus\{0\}$), așa că obstrucția se mută pe
numărul de soluții ale lui $y^{2}=1$.

Notăm și că discuția are exact două ramuri
($\operatorname{char}K=2$ sau $\neq 2$), fără vreo caracteristică
,,putere a lui $2$'' de tipul $4$ sau $8$: \emph{caracteristica unui corp
este întotdeauna $0$ sau un număr prim}. Într-adevăr, dacă
ea ar fi compusă, $c=ab$ cu $1<a,b<c$, atunci
$0=\underbrace{1+\cdots+1}_{c}=\bigl(\underbrace{1+\cdots+1}_{a}\bigr)\cdot
\bigl(\underbrace{1+\cdots+1}_{b}\bigr)$, iar într-un corp (fără
divizori ai lui zero) unul dintre factori ar fi $0$, dând un ,,reset la
$0$'' mai devreme decât $c$ --- contrazicând minimalitatea lui $c$.
Deci singura putere a lui $2$ care apare drept caracteristică este $2$
însuși. (În schimb, într-un \emph{inel} cu divizori ai lui
zero, precum $\mathbb{Z}/4$, chiar avem $\operatorname{char}=4$ și
$2\cdot 2=0$ cu $2\neq 0$ --- tocmai de aceea argumentul cere corp.)
\end{obs}

\begin{problema}
Fie $A$ un inel finit al cărui grup aditiv $(A,+)$ este izomorf cu grupul
multiplicativ $(K^{*},\cdot)$ al unui corp finit $K$. Demonstrați că
$A$ este comutativ.
\end{problema}

\begin{solutie}
\textit{Pas 1: grupul aditiv al lui $A$ este ciclic.} Folosim următoarea
lemă.

\begin{lemat}
Orice subgrup finit $G$ al grupului multiplicativ $(K^{*},\cdot)$ al unui
corp $K$ este ciclic.
\end{lemat}

\begin{proof}
Fie $m=|G|$. Pentru un divizor $d\mid m$, notăm cu $\psi(d)$ numărul
elementelor din $G$ de ordin \emph{exact} $d$. Ne sprijinim pe două
fapte:

(i) \emph{Un polinom de grad $d$ peste un corp are cel mult $d$
rădăcini} (Teorema~\ref{teo:numarradacini} din cartea de teorie). În particular,
ecuația $x^{d}=1$ are cel mult $d$ soluții în $K$.

(ii) \emph{Identitatea lui Gauss}: $\sum_{d\mid m}\varphi(d)=m$, unde
$\varphi$ e indicatorul lui Euler (Lema~\ref{lema:sumaphi}).

Arătăm că $\psi(d)\in\{0,\varphi(d)\}$ pentru fiecare $d\mid m$.
Presupunem $\psi(d)>0$, deci există $a\in G$ de ordin $d$. Atunci
subgrupul ciclic $\langle a\rangle$ are $d$ elemente, toate soluții ale
lui $x^{d}=1$; cum această ecuație are cel mult $d$ soluții în
$K$ (fapt (i)), rezultă că $\langle a\rangle$ le conține
\emph{pe toate}. În particular, orice element din $G$ de ordin $d$
verifică $x^{d}=1$, deci aparține lui $\langle a\rangle$; iar
elementele de ordin exact $d$ dintr-un grup ciclic de ordin $d$ sunt cei
$\varphi(d)$ generatori ai săi. Așadar $\psi(d)=\varphi(d)$.

Fiecare element al lui $G$ are un ordin care divide $m$, deci
$\sum_{d\mid m}\psi(d)=m$. Combinând cu $\psi(d)\le\varphi(d)$ și cu
$\sum_{d\mid m}\varphi(d)=m$ (fapt (ii)):
\[
  m=\sum_{d\mid m}\psi(d)\le\sum_{d\mid m}\varphi(d)=m,
\]
deci egalitate peste tot, adică $\psi(d)=\varphi(d)$ pentru orice
$d\mid m$. În special $\psi(m)=\varphi(m)\ge 1$: există un element de
ordin $m$, deci $G=\langle\text{acel element}\rangle$ este ciclic.
\end{proof}

Aplicăm lema lui $G=K^{*}$ (subgrup finit al lui $K^{*}$, chiar el
însuși), de ordin $|K|-1$: deci $(K^{*},\cdot)$ este ciclic. Cum
$(A,+)\cong(K^{*},\cdot)$ și izomorfismele păstrează
ciclicitatea, $(A,+)$ este ciclic: există $g\in A$ care generează
grupul aditiv, adică orice element al lui $A$ este un multiplu
întreg al lui $g$,
\[
  A=\{k\cdot g : k\in\mathbb{Z}\},
\]
unde $k\cdot g$ înseamnă $\underbrace{g+\cdots+g}_{k}$ (și
$(-k)\cdot g=-(k\cdot g)$, $0\cdot g=0$).

\smallskip
\textit{Pas 2: înmulțirea se reduce la $g\cdot g$.} Folosim regulile
de distributivitate, care dau, pentru multipli întregi,
$(k\cdot g)\cdot(l\cdot g)=(kl)\cdot(g\cdot g)$ (înmulțirea distribuie
peste sumele repetate). Atunci, pentru $a=k\cdot g$ și $b=l\cdot g$
oarecare,
\[
  a\cdot b=(k\cdot g)(l\cdot g)=(kl)\cdot(g\cdot g),\qquad
  b\cdot a=(l\cdot g)(k\cdot g)=(lk)\cdot(g\cdot g).
\]

\smallskip
\textit{Pas 3: comutativitatea.} Coeficienții $k,l$ sunt întregi, deci
$kl=lk$; prin urmare $a\cdot b=(kl)\cdot(g\cdot g)=(lk)\cdot(g\cdot g)=b\cdot a$.
Cum $a,b$ au fost arbitrare, $A$ este comutativ.
\end{solutie}

\begin{obs}
Toată forța stă în faptul că $(A,+)$ e \emph{ciclic}: un
inel cu grup aditiv ciclic e automat comutativ, fiindcă înmulțirea
e determinată de un singur produs $g\cdot g$, iar coeficienții
întregi comută. Ipoteza $(A,+)\cong(K^{*},\cdot)$ nu e folosită
prin structura multiplicativă a lui $K$, ci doar prin faptul că
împrumută lui $(A,+)$ ciclicitatea lui $K^{*}$. De pildă, un inel
necomutativ ca $M_{2}(\mathbb{F}_{2})$ (matrice $2\times 2$, grup aditiv
$(\mathbb{Z}/2)^{4}$, \emph{neciclic}) nu poate satisface ipoteza --- exact
pentru că $(\mathbb{Z}/2)^{4}$ nu e izomorf cu niciun $K^{*}$ ciclic.
\end{obs}

\begin{problema}
Fie $A$ un inel comutativ de caracteristică $p$, unde $p$ este un
număr prim \emph{impar}, cu proprietatea că $x^{p}=1$ pentru orice
element inversabil $x\in A$. Demonstrați că suma a două elemente
inversabile este tot un element inversabil.
\end{problema}

\begin{solutie}
Fie $x,y\in A$ inversabile.

\smallskip
\textit{Pas 1: ,,visul bobocului'' (morfismul lui Frobenius).} Arătăm
că într-un inel comutativ de caracteristică $p$ (prim) are loc
\[
  (x+y)^{p}=x^{p}+y^{p}.
\]
Într-adevăr, $A$ fiind comutativ, putem aplica binomul lui Newton:
\[
  (x+y)^{p}=\sum_{i=0}^{p}\binom{p}{i}x^{i}y^{p-i}.
\]
Pentru $0<i<p$, coeficientul $\binom{p}{i}=\dfrac{p!}{i!\,(p-i)!}$ este
divizibil cu $p$: numărătorul $p!$ conține factorul prim $p$, pe
când numitorul $i!\,(p-i)!$ nu îl conține (căci $i<p$ și
$p-i<p$, iar $p$ e prim). Cum caracteristica lui $A$ este $p$, avem
$p\cdot a=0$ pentru orice $a\in A$, deci toți termenii intermediari se
anulează, rămânând doar $i=0$ și $i=p$.

\smallskip
\textit{Pas 2: calculul.} Aplicăm Pasul 1 și ipoteza $x^{p}=y^{p}=1$:
\[
  (x+y)^{p}=x^{p}+y^{p}=1+1=2 .
\]

\smallskip
\textit{Pas 3: $2$ este inversabil.} Caracteristica fiind $p$ impar, avem
$p\nmid 2$, deci $2\neq 0$ în $A$. Mai mult, multiplii lui $1$ formează
în $A$ o copie a lui $\mathbb{Z}/p$, care este corp (căci $p$ e prim);
elementul $2$ este nenul în această copie, deci este inversabil
în ea (și, prin urmare, în $A$). Explicit, din $\gcd(2,p)=1$,
relația lui Bézout dă întregi $u,v$ cu $2u+pv=1$, deci
$2\cdot(u\cdot 1)=1-v\,(p\cdot 1)=1$: inversul lui $2$ este $u\cdot 1$.

\smallskip
\textit{Pas 4: concluzia.} Din Pașii 2 și 3,
\[
  (x+y)\cdot\Bigl(2^{-1}(x+y)^{p-1}\Bigr)=2^{-1}(x+y)^{p}=2^{-1}\cdot 2=1 ,
\]
și analog în cealaltă ordine ($A$ e comutativ). Așadar $x+y$
este inversabil, având inversul explicit $2^{-1}(x+y)^{p-1}$.
\end{solutie}

\begin{obs}
Ipoteza că $p$ este \emph{impar} este esențială: pentru $p=2$
concluzia se inversează complet. Într-adevăr, în
caracteristică $2$ avem $2=0$, deci calculul de la Pasul 2 dă
$(x+y)^{2}=1+1=0$: suma a două elemente inversabile este \emph{nilpotentă},
deci (dacă e nenulă) niciodată inversabilă. Contraexemplu
concret: în $A=\mathbb{F}_{2}[t]/(t^{2})$ (caracteristică $2$),
elementele inversabile sunt $1$ și $1+t$, iar amândouă verifică
$x^{2}=1$; totuși $1+(1+t)=t$, care nu e inversabil, căci $t^{2}=0$.

Comutativitatea lui $A$ e la fel de necesară: binomul lui Newton --- deci
și identitatea $(x+y)^{p}=x^{p}+y^{p}$ --- este valabilă numai pentru
elemente care comută.
\end{obs}

\begin{problema}
Fie $A$ un inel cu proprietatea că orice element $x\in A$ se scrie ca
sumă $x=e+f$ a două elemente idempotente ($e^{2}=e$, $f^{2}=f$) care
\emph{anticomută}, adică $ef+fe=0$. Demonstrați că $A$ este
comutativ.
\end{problema}

\begin{solutie}
\textit{Pas 1: orice element este idempotent.} Fie $x\in A$ și fie
$x=e+f$ o scriere ca în ipoteză: $e^{2}=e$, $f^{2}=f$ și
$ef+fe=0$. Atunci
\[
  x^{2}=(e+f)^{2}=e^{2}+ef+fe+f^{2}
       =e+\underbrace{(ef+fe)}_{=0}+f=e+f=x .
\]
(Atenție: nu am folosit comutativitatea --- dezvoltarea păstrează
ordinea factorilor, iar cei doi termeni încrucișați se
anulează tocmai prin anticomutare.) Așadar
\[
  x^{2}=x\qquad\text{pentru orice }x\in A .
\]

\smallskip
Am obținut deci un \emph{inel Boolean}. Pașii următori
demonstrează că orice inel Boolean e comutativ (acesta este exact
conținutul Teoremei~\ref{teo:boole} din cartea de teorie: dacă $x^{2}=x$ pentru
orice $x$, atunci $1+1=0$ și $A$ e comutativ); îl reluăm aici,
pentru ca soluția să fie autonomă.

\smallskip
\textit{Pas 2: relația fundamentală.} Aplicăm Pasul 1 sumei
$a+b$, cu $a,b\in A$ arbitrare:
\[
  (a+b)^{2}=a+b .
\]
Dezvoltând membrul stâng (fără a schimba ordinea factorilor)
și folosind $a^{2}=a$, $b^{2}=b$:
\[
  a^{2}+ab+ba+b^{2}=a+ab+ba+b=a+b ,
\]
de unde
\[
  ab+ba=0 . \tag{$\star$}
\]

\smallskip
\textit{Pas 3: caracteristica este $2$.} Punem $b=a$ în $(\star)$:
\[
  a\cdot a+a\cdot a=0\ \Longrightarrow\ a^{2}+a^{2}=0\ \Longrightarrow\ a+a=0 ,
\]
folosind $a^{2}=a$. Așadar $2a=0$ pentru orice $a\in A$, ceea ce
înseamnă $-z=z$ pentru orice $z\in A$.

\smallskip
\textit{Pas 4: comutativitatea.} Din $(\star)$ avem $ab=-ba$; dar, prin Pasul
3, $-ba=ba$. Prin urmare
\[
  ab=ba\qquad\text{pentru orice }a,b\in A ,
\]
adică $A$ este comutativ.
\end{solutie}

\begin{obs}
Ipoteza este de fapt \emph{echivalentă} cu ,,$x^{2}=x$ pentru orice $x$''
(inel Boolean). Un sens e Pasul 1; reciproc, dacă $x^{2}=x$ pentru orice
$x$, atunci orice $x$ se scrie $x=e+f$ cu $e=x$ și $f=0$: ambele sunt
idempotente și anticomută ($ef+fe=0+0=0$).

Anticomutarea este esențială: suma a doi idempotenți oarecare nu e
idempotentă. De pildă, în $\mathbb{Z}$ elementele $e=f=1$ sunt
idempotente, dar $ef+fe=2\neq 0$, iar suma $e+f=2$ nu e idempotentă
($2^{2}=4\neq 2$).
\end{obs}

\begin{problema}
Fie $A$ un inel \emph{finit}, cu $1\neq 0$, având proprietatea că
pentru orice $n\ge 1$ ecuația $x^{n+1}=x^{n}$ are ca singure soluții
pe $0$ și pe $1$. Demonstrați că $A$ este corp.
\end{problema}

\begin{solutie}
Vom folosi următorul criteriu (Propoziția~\ref{prop:finitcorp} din cartea de teorie,
,,când un inel finit este corp''), pe care îl enunțăm
și îl demonstrăm.

\begin{lemat}[când un inel finit este corp; Propoziția~\ref{prop:finitcorp}]
Un inel finit $A$ (cu $1\neq 0$) este corp dacă și numai dacă nu
are elemente nilpotente nenule și ecuația $x^{2}=x$ are numai
soluțiile $0$ și $1$.
\end{lemat}

\begin{proof}
($\Rightarrow$) Într-un corp, dacă $a\neq 0$ și $a^{k}=0$, atunci
$1=a^{-k}a^{k}=0$ --- exclus; deci nu există nilpotenți nenuli. Iar
$x^{2}=x$ dă $x(x-1)=0$, deci $x\in\{0,1\}$ (un corp nu are divizori ai
lui zero).

($\Leftarrow$) Fie $a\neq 0$. Cum $A$ e finit, șirul $a,a^{2},a^{3},\dots$
nu poate avea toate elementele distincte: există $1\le i<j$ cu
$a^{i}=a^{j}$. Fie $d=j-i\ge 1$. Înmulțind cu $a^{k-i}$ obținem
$a^{k+d}=a^{k}$ pentru orice $k\ge i$, deci relația se poate itera:
$a^{k+ld}=a^{k}$ pentru orice $l\ge 0$ și $k\ge i$. Alegem $N=i\,d$
(așadar $N\ge i$ și $d\mid N$); aplicând iterarea de $N/d$ ori,
\[
  a^{2N}=a^{N+N}=a^{N},
\]
deci $e:=a^{N}$ este idempotent ($e^{2}=e$). Prin ipoteză,
$e\in\{0,1\}$. Dacă $e=0$, atunci $a^{N}=0$, adică $a$ ar fi nilpotent
nenul --- exclus. Rămâne $e=1$, adică $a^{N}=1$, deci
\[
  a\cdot a^{N-1}=a^{N-1}\cdot a=1 :
\]
$a$ este inversabil. Orice element nenul fiind inversabil, $A$ este corp.
\end{proof}

Verificăm acum cele două condiții din lemă pentru inelul
nostru.

\smallskip
\textit{(i) $A$ nu are nilpotenți nenuli.} Fie $t$ cu $t^{k}=0$,
$k\ge 1$. Atunci și $t^{k+1}=0=t^{k}$, deci $t$ este soluție a
ecuației $x^{k+1}=x^{k}$. Prin ipoteză, $t\in\{0,1\}$. Dar $t=1$ ar da
$1=1^{k}=t^{k}=0$, contrazicând $1\neq 0$. Deci $t=0$.

\smallskip
\textit{(ii) $x^{2}=x$ are doar soluțiile $0$ și $1$.} Aceasta este
exact ipoteza pentru $n=1$.

\smallskip
Prin lemă, $A$ este corp.
\end{solutie}

\begin{obs}
Ipoteza de \emph{finitudine} este esențială, deși enunțul
pare să nu aibă nevoie de ea. Într-adevăr, în inelul
$\mathbb{Z}$ ecuația $x^{n+1}=x^{n}$ se scrie $x^{n}(x-1)=0$ și, cum nu
există divizori ai lui zero, are exact soluțiile $0$ și $1$ ---
deci ipoteza e satisfăcută, dar $\mathbb{Z}$ nu e corp. Mai general,
\emph{orice} domeniu de integritate satisface ipoteza; fără
finitudine, concluzia cade.

În plus, $A$ fiind corp finit, el este chiar \emph{comutativ}, prin
teorema lui Wedderburn (Teorema~\ref{teo:wedderburn} din carte: orice corp finit este
comutativ).
\end{obs}

\begin{problema}
Fie $A$ un inel cu proprietatea că $x^{3}+x^{2}\in Z(A)$ pentru orice
$x\in A$ (unde $Z(A)=\{z\in A: zy=yz\ \forall y\in A\}$ este centrul
inelului). Demonstrați că $A$ este comutativ.
\end{problema}

\begin{solutie}
Notăm $h(x)=x^{3}+x^{2}$, deci $h(x)\in Z(A)$ pentru orice $x$.
Reamintim că centrul $Z(A)$ este un subinel care conține $1$: e
închis la sume, diferențe și produse și conține toți multiplii
întregi $k\cdot 1$. Toate calculele de mai jos implică doar puteri ale
unui singur element și constante întregi, care comută între ele, deci
putem dezvolta ca în $\mathbb{Z}[X]$.

\smallskip
\textit{Pasul 1: $2x^{2}$ este central.} Aplicând ipoteza lui $-x$:
$h(-x)=-x^{3}+x^{2}\in Z(A)$. Adunând cu $h(x)$,
\[
  h(x)+h(-x)=2x^{2}\in Z(A).
\]

\smallskip
\textit{Pasul 2: $x^{2}+5x$ este central.} Aplicând ipoteza lui $x+1$:
\[
  h(x+1)=(x+1)^{3}+(x+1)^{2}=x^{3}+4x^{2}+5x+2\in Z(A),
\]
deci $h(x+1)-h(x)-2=3x^{2}+5x\in Z(A)$. Scăzând $2x^{2}$ (Pasul 1),
\[
  u(x):=x^{2}+5x\in Z(A)\qquad\text{pentru orice }x\in A.
\]

\smallskip
\textit{Pasul 3: $x^{2}+x$ este central.} Aplicând Pasul 2 lui $x+1$:
\[
  u(x+1)-u(x)=(x^{2}+7x+6)-(x^{2}+5x)=2x+6\in Z(A),
\]
deci $2x\in Z(A)$. Atunci $x^{2}+x=u(x)-2\cdot(2x)\in Z(A)$ pentru
orice $x\in A$.

\smallskip
\textit{Pasul 4: anticomutatorii sunt centrali.} Aplicând Pasul 3
elementelor $x$, $y$ și $x+y$ (dezvoltarea lui $(x+y)^{2}$ păstrează
ordinea factorilor):
\[
  (x+y)^{2}+(x+y)-(x^{2}+x)-(y^{2}+y)=xy+yx\in Z(A).
\]

\smallskip
\textit{Pasul 5: pătratele sunt centrale.} Elementul central $xy+yx$
comută în particular cu $x$:
\[
  x(xy+yx)=(xy+yx)x\ \Longrightarrow\ x^{2}y+xyx=xyx+yx^{2}
  \ \Longrightarrow\ x^{2}y=yx^{2}.
\]
Cum $y$ este arbitrar, $x^{2}\in Z(A)$ pentru orice $x$.

\smallskip
\textit{Concluzie.} Din Pașii 3 și 5,
$x=(x^{2}+x)-x^{2}\in Z(A)$ pentru orice $x\in A$, deci $Z(A)=A$:
inelul $A$ este comutativ.
\end{solutie}

\begin{obs}
Pașii 4 și 5 sunt exact mecanismul din teorema lui MacHale
(Teorema~\ref{teo:machalecomut}): odată ce $x^{2}+x$ este central,
polarizarea face anticomutatorii centrali, apoi pătratele, apoi
totul. Rolul Pașilor 1, 2 și 3 este să ,,coboare gradul'': substituțiile
$x\mapsto -x$ și $x\mapsto x+1$ elimină pe rând termenii de grad $3$,
apoi pe cei de grad $2$ cu coeficient nepotrivit. Rezultatul apare în
partea de teorie ca Propoziția~\ref{prop:machalegen}.

Ipoteza este echivalentă cu $x^{3}-x^{2}\in Z(A)$ pentru orice $x$:
înlocuind $x$ cu $-x$, obținem $(-x)^{3}+(-x)^{2}=-(x^{3}-x^{2})$, iar
$Z(A)$ este subgrup aditiv. Așadar și această variantă ,,cu semn
schimbat'' implică comutativitatea.

Atenție însă la generalizările ,,prea ambițioase'': dacă înlocuim
$x^{3}$ și $x^{2}$ prin $f(x^{3})$ și $g(x^{2})$, cu $f,g$ bijecții
aditive arbitrare ale lui $A$, concluzia devine falsă (vezi
observația de după Propoziția~\ref{prop:machalegen}, cu un
contraexemplu în inelul matricelor $2\times 2$ superior triunghiulare
peste $\mathbb{Z}_{2}$).
\end{obs}


\begin{problema}[Clasică]
Fie $K \subseteq L$ două corpuri finite, operațiile lui $K$ fiind cele induse de $L$. Arătați că există un întreg $m \geq 1$ cu
\[
  |L| = |K|^m.
\]
\end{problema}

\begin{solutie}
Să privim $L$ ca \emph{spațiu vectorial peste corpul $K$}: ,,vectorii'' sunt elementele lui $L$, adunarea este cea din $L$, iar înmulțirea cu ,,scalarii'' din $K$ este înmulțirea din $L$. Axiomele de spațiu vectorial sunt exact axiome ale corpului $L$ (distributivitate, asociativitate, $1 \cdot x = x$).

Cum $L$ este finit, el este generat de o mulțime finită (de exemplu, de toate elementele sale), deci din orice sistem de generatori se poate extrage o bază: fie $e_1, \dots, e_m$ o bază a lui $L$ peste $K$. Orice element al lui $L$ se scrie atunci \emph{în mod unic}
\[
  x = \lambda_1 e_1 + \cdots + \lambda_m e_m, \qquad \lambda_1, \dots, \lambda_m \in K,
\]
iar corespondența $x \mapsto (\lambda_1, \dots, \lambda_m)$ este o bijecție între $L$ și $K^m$. Numărând: $|L| = |K|^m$.

\medskip
\noindent\textbf{Observație} \quad Două consecințe imediate. Luând drept $K$ subcorpul prim: orice corp finit are $p^m$ elemente, cu $p$ prim. Și o obstrucție concretă: un corp cu 8 elemente nu poate conține un subcorp cu 4 elemente, deoarece 8 nu este o putere a lui 4; subcorpurile lui $\mathbb{F}_8$ au 2 sau 8 elemente. Întreaga aritmetică a extinderilor de corpuri finite este, în fond, o numărare de dimensiuni: încă o dată, algebra liniară lucrează acolo unde enunțul nici nu o pomenește.
\end{solutie}

\begin{problema}
Fie $A$ un inel comutativ și fie $I,J$ ideale ale sale astfel încât $I \cup J$ să fie un ideal al lui $A$. Arătați că $I \subseteq J$ sau $J \subseteq I$.
\end{problema}

\begin{solutie}
Presupunem contrariul: există $u \in I \setminus J$ și $w \in J \setminus I$. Elementele $u$ și $w$ aparțin lui $I \cup J$, care este un ideal, deci este închis la adunare:
\[
u + w \in I \cup J.
\]
Dacă $u + w \in I$, atunci $w = (u + w) - u \in I$ (un ideal este un subgrup aditiv, deci este închis la diferențe): contradicție cu $w \notin I$. Dacă $u + w \in J$, atunci $u = (u + w) - w \in J$: contradicție cu $u \notin J$. Ambele cazuri sunt imposibile, deci unul dintre ideale îl conține pe celălalt.

\medskip
\noindent\textbf{Observație.} Aceasta este a treia parte a unei trilogii: exact aceeași demonstrație de trei rânduri arată că un grup nu este reuniunea a două subgrupuri proprii (Lema~\ref{lema:reuniune}) și că un spațiu vectorial nu este reuniunea a două subspații proprii (problema OJM.28 din volumul de algebră pentru clasa a XI-a). De fapt, singura structură folosită este cea de \emph{subgrup aditiv}: comutativitatea inelului și înmulțirea cu elemente din $A$ nu joacă nicăieri niciun rol.
\end{solutie}

\nivel{onm}{Nivel ONM}{Faza națională --- mai mulți pași și o idee-cheie.}
\setcounter{cat}{3}\setcounter{problemaX}{0}

\begin{problema}
Determinați toate numerele naturale $n$ pentru care \emph{orice} grup
$G$ de ordin $n$ are proprietatea că orice endomorfism neconstant al
său este automorfism. (Un endomorfism se numește \emph{constant}
dacă este morfismul banal $g \mapsto e$.)
\end{problema}

\noindent\textbf{Răspuns:} $n=1$ sau $n$ prim.

\begin{solutie}
\textit{Reformulare.} Într-un grup finit, un endomorfism $f$ este
automorfism dacă și numai dacă este injectiv, adică
$\ker f=\{e\}$; iar $f$ este constant dacă și numai dacă
$\ker f=G$. Condiția din enunț revine deci la: pentru orice
endomorfism $f$ al lui $G$, $\ker f \in \{\{e\},\,G\}$.

\smallskip
\textit{Partea 1: $n=1$ și $n$ prim funcționează.}
Pentru $n=1$, grupul trivial are ca unic endomorfism identitatea. Pentru
$n=p$ prim, singurul grup de ordin $p$ este $\mathbb{Z}/p$ (din teorema lui
Lagrange orice element $\neq e$ are ordinul $p$, deci grupul este ciclic).
Endomorfismele lui $\mathbb{Z}/p$ sunt aplicațiile $x \mapsto kx$; pentru
$k=0$ se obține morfismul banal, iar pentru $k \neq 0$ avem
$\gcd(k,p)=1$, deci $x \mapsto kx$ este bijecție, adică automorfism.
Așadar orice grup de ordin $1$ sau prim satisface condiția.

\smallskip
\textit{Partea 2: dacă $n$ este compus, condiția cade.}
Fie $n$ compus și scriem $n=ab$ cu $1<a,b<n$. Considerăm grupul
ciclic $G=\mathbb{Z}/n$ și aplicația
\[
  \varphi : \mathbb{Z}/n \to \mathbb{Z}/n, \qquad \varphi(x)=a\,x \pmod n.
\]
Ea este endomorfism, căci $\varphi(x+y)=a(x+y)=\varphi(x)+\varphi(y)$, și
este neconstantă, deoarece $\varphi(1)=a \neq 0$ (avem $0<a<n$).
Nucleul său este
\[
  \ker\varphi = \{\, x : n \mid ax \,\}
              = \{\, x : b \mid x \,\}
              = \{\, 0,\, b,\, 2b,\, \dots,\, (a-1)b \,\},
\]
mulțime cu exact $a$ elemente. Cum $a>1$, $\varphi$ nu este injectivă,
deci nu este automorfism. Prin urmare grupul $\mathbb{Z}/n$ posedă un
endomorfism neconstant care nu este automorfism, deci nu orice grup de
ordin $n$ satisface condiția.

\smallskip
\textit{Concluzie.} Afirmația ,,orice grup de ordin $n$ satisface
condiția'' este adevărată exact pentru $n=1$ și $n$ prim.
\end{solutie}

\begin{obs}
Este esențial că enunțul cere ,,\emph{orice} grup de ordin
$n$'', nu ,,\emph{există} un grup''. Cu ,,există'' răspunsul
încetează să mai fie curat: intră în joc și ordine
compuse (de pildă $60$, prin grupul simplu $A_{5}$), iar mulțimea
valorilor $n$ nu mai are o formă închisă elementară.
\end{obs}

\begin{problema}
Fie $G$ un grup finit cu un număr \emph{impar} de elemente și fie
$H \subseteq G$ o submulțime nevidă cu proprietatea
\[
  ghg \in H \qquad \text{pentru orice } g \in G,\ h \in H.
\]
Demonstrați că $H = G$.
\end{problema}

\begin{solutie}
Notăm $e$ elementul neutru și $n=|G|$, cu $n$ impar.

\smallskip
\textit{Pasul 1: $e \in H$.} Cum $H \neq \varnothing$, fixăm
$h \in H$ și fie $d=\operatorname{ord}(h)$. Din teorema lui Lagrange
$d \mid n$, deci $d$ este impar și $\gcd(2,d)=1$. Pentru orice
întreg $m \geq 0$ alegem în ipoteză $g=h^{m}$; puterile lui $h$
comutând,
\[
  g\,h\,g = h^{m}\cdot h \cdot h^{m} = h^{\,2m+1} \in H.
\]
Cum $2$ este inversabil modulo $d$, exponenții $2m+1$ parcurg toate
resturile modulo $d$; în particular există $m$ cu
$2m+1 \equiv 0 \pmod d$, deci $h^{\,2m+1}=e \in H$.

\smallskip
\textit{Pasul 2: orice pătrat este în $H$.} Având $e \in H$,
aplicăm ipoteza cu $h=e$: pentru orice $g \in G$,
\[
  g\,e\,g = g^{2} \in H .
\]

\smallskip
\textit{Pasul 3: ridicarea la pătrat este surjectivă.} Pentru orice
$x \in G$ avem $x^{n}=e$, deci
\[
  x = x^{\,n+1} = \left( x^{\frac{n+1}{2}} \right)^{\!2},
\]
unde $\frac{n+1}{2}$ este întreg tocmai pentru că $n$ este impar.
Așadar orice $x \in G$ este un pătrat.

\smallskip
\textit{Concluzie.} Fie $x \in G$. Din Pasul 3 există $g$ cu $x=g^{2}$,
iar din Pasul 2, $g^{2} \in H$; deci $x \in H$. Prin urmare $G \subseteq H$,
și cum $H \subseteq G$, rezultă $H=G$.
\end{solutie}

\begin{obs}
Concluzia este mai tare decât ar sugera enunțul: nu doar că $H$
este subgrup, ci este chiar întregul grup. Paritatea impară a lui
$|G|$ acționează în două locuri diferite --- forțează
ordinele elementelor să fie impare și face ca ridicarea la pătrat
să fie surjectivă pe $G$.

Ipoteza de paritate este esențială: dacă $|G|$ este par,
ridicarea la pătrat nu mai este neapărat surjectivă (în
$\mathbb{Z}/2\mathbb{Z}$ singurul pătrat este $0$), iar argumentul din
Pasul 3 cade.
\end{obs}

\begin{problema}
Fie $K$ un corp infinit și $f : K \to K$ o funcție
\emph{polinomială} cu proprietatea
\[
  f\bigl(x+f(y)\bigr) = f(x)\,f(y) \qquad \text{pentru orice } x,y \in K .
\]
Determinați toate funcțiile $f$.
\end{problema}

\noindent\textbf{Răspuns:} $f \equiv 0$ sau $f \equiv 1$.

\begin{solutie}
Fie $d=\deg f$. Dacă $d=0$, atunci $f\equiv c$ constantă, iar
relația devine $c=c^{2}$, deci $c(c-1)=0$; corpul neavând divizori
ai lui zero, $c=0$ sau $c=1$. Ambele verifică enunțul.

Presupunem acum $d\ge 1$ și ajungem la o contradicție. Fixăm
$y$; deoarece $K$ este infinit, o egalitate între două funcții
polinomiale care coincid în toate punctele lui $K$ înseamnă
egalitatea polinoamelor înseși. Privim atunci
\[
  f\bigl(x+f(y)\bigr) = f(x)\,f(y)
\]
ca identitate în nedeterminata $x$, cu coeficienți în $K[y]$.
Fie $\alpha \neq 0$ coeficientul dominant al lui $f$ (deci $f(x)=\alpha x^{d}+\cdots$).

\smallskip
\textit{Coeficientul lui $x^{d}$ în cei doi membri.} În membrul
stâng, $f\bigl(x+f(y)\bigr)=\alpha\bigl(x+f(y)\bigr)^{d}+\cdots$, iar
termenul de grad $d$ în $x$ are coeficientul $\alpha$ (translația nu
schimbă coeficientul dominant). În membrul drept,
$f(x)\,f(y)=\alpha f(y)\,x^{d}+\cdots$, deci coeficientul lui $x^{d}$ este
$\alpha\,f(y)$. Egalând,
\[
  \alpha = \alpha\,f(y),
\]
adică $\alpha\bigl(f(y)-1\bigr)=0$ pentru orice $y$.

\smallskip
Cum $K$ nu are divizori ai lui zero și $\alpha\neq0$, rezultă
$f(y)=1$ pentru orice $y$, deci $f$ ar fi constantă $1$ --- în
contradicție cu $d\ge1$. Așadar cazul $d\ge1$ este imposibil,
și singurele soluții sunt $f\equiv0$ și $f\equiv1$.
\end{solutie}

\begin{obs}
Cele două ipoteze intervin în locuri diferite. \emph{Infinitatea}
lui $K$ permite trecerea de la ,,egale în toate punctele'' la
egalitatea coeficienților; peste un corp finit acest pas cade.
\emph{Absența divizorilor lui zero} (proprietate a oricărui corp)
permite simplificarea lui $\alpha$ și faptul că singurii
idempotenți sunt $0$ și $1$. De aceea rezultatul rămâne
valabil, cu aceeași demonstrație, pentru orice \emph{domeniu de
integritate infinit}; însă peste un inel infinit cu divizori ai
lui zero pot apărea soluții suplimentare (idempotenți
netriviali), iar enunțul își pierde forma curată.
\end{obs}

\begin{problema}
Fie $k \in \mathbb{N}^{*}$ fixat. Determinați toate polinoamele
$P \in \mathbb{C}[X]$ cu proprietatea
\[
  P(z^{k}) = P(z)\,P(z+1)\cdots P(z+k-1) \qquad \text{pentru orice } z \in \mathbb{C}.
\]
\end{problema}

\noindent\textbf{Răspuns:}
\[
\begin{cases}
  \text{orice } P, & k=1;\\[2pt]
  P\equiv 0 \ \text{ sau }\ P(z)=(z^{2}-z)^{m},\ m\ge 0, & k=2;\\[2pt]
  P\equiv 0 \ \text{ sau }\ P\equiv c \ \text{ cu } c^{\,k-1}=1, & k\ge 3.
\end{cases}
\]

\begin{solutie}
Pentru $k=1$ relația devine $P(z)=P(z)$, deci orice polinom convine.
Presupunem în continuare $k\ge 2$.

\smallskip
\textit{Soluțiile constante.} Dacă $P\equiv c$, relația dă
$c=c^{k}$, deci $c=0$ sau $c^{\,k-1}=1$; toate acestea verifică.

\smallskip
\textit{Coeficientul dominant.} Fie $P$ neconstant, $\deg P=n\ge1$, cu
coeficient dominant $a$. Termenul dominant în membrul stâng este
$a\,z^{nk}$, iar în dreapta $a^{k}z^{nk}$, deci $a^{k}=a$, adică
$a^{\,k-1}=1$.

\smallskip
\textit{Mărginirea rădăcinilor.} Cele două polinoame fiind
egale, au aceleași rădăcini (cu multiplicități). Fie
$r_{1},\dots,r_{n}$ rădăcinile lui $P$ și $M=\max_{i}|r_{i}|$.
Rădăcinile lui $P(z^{k})$ sunt rădăcinile de ordin $k$ ale
lui $r_{i}$, deci de modul maxim $M^{1/k}$. Rădăcinile membrului
drept sunt numerele $r_{i}-j$ cu $0\le j\le k-1$, de modul maxim
$M'=\max_{i,j}|r_{i}-j|$. Luând $j=0$ obținem $M'\ge M$, deci
\[
  M^{1/k}=M'\ge M .
\]
Dacă toate rădăcinile ar fi nule, atunci $P(z)=az^{n}$, iar
membrul drept ar avea rădăcina $z=-1$ (din factorul cu $j=1$), pe
care membrul stâng nu o are; contradicție. Deci $M>0$, iar din
$M^{1/k}\ge M$ rezultă $M\le 1$. Atunci $M'=M^{1/k}\le 1$, deci
\[
  |r_{i}-j|\le 1 \qquad \text{pentru orice } i \text{ și orice } j\in\{0,1,\dots,k-1\}.
\]

\smallskip
\textit{Mărginirea lui $k$.} Fie $r$ o rădăcină. Din
$|r|\le 1$ și $|r-(k-1)|\le 1$ și inegalitatea triunghiului,
\[
  k-1 = |(k-1)-0| \le |r| + |r-(k-1)| \le 2,
\]
deci $k\le 3$. În particular, pentru $k\ge 4$ un polinom neconstant nu
poate avea rădăcini: rămân doar constantele.

\smallskip
\textit{Cazul $k=3$.} Rădăcinile satisfac $|r|\le 1$ și
$|r-2|\le 1$. Discurile închise de centre $0$ și $2$ și rază
$1$ sunt tangente exterior în punctul $1$, deci singura rădăcină
posibilă este $r=1$: $P(z)=a(z-1)^{n}$. Atunci
\[
  P(z^{3})=a(z^{3}-1)^{n}, \qquad
  P(z)P(z{+}1)P(z{+}2)=a^{3}\bigl(z(z-1)(z+1)\bigr)^{n}=a^{3}(z^{3}-z)^{n}.
\]
Termenul liber al primului este $a(-1)^{n}\neq 0$, iar al celui de-al doilea
este $0$ (factorul $z^{n}$); egalitatea e imposibilă pentru $n\ge1$.
Rămân deci doar constantele.

\smallskip
\textit{Cazul $k=2$.} Aici discurile $|r|\le 1$, $|r-1|\le 1$ se
intersectează după o întreagă regiune, deci e nevoie de un
argument suplimentar. Punem $z=r$ (rădăcină) în
$P(z^{2})=P(z)P(z+1)$: cum $P(r)=0$, obținem $P(r^{2})=0$, deci
\emph{mulțimea rădăcinilor este închisă la $r\mapsto r^{2}$}.

Fie $r\neq 0$ o rădăcină. Toate puterile $r^{2^{m}}$ sunt tot
rădăcini, deci formează o mulțime finită; dacă
$|r|\neq 1$, modulele $|r|^{2^{m}}$ ar fi însă distincte două
câte două --- contradicție. Așadar $|r|=1$, deci
$r=e^{i\varphi}$. Condiția $|r-1|\le 1$ revine la
$2-2\cos\varphi\le 1$, adică $\cos\varphi\ge \tfrac12$, deci
$|\varphi|\le \tfrac{\pi}{3}$.

Rădăcinile $r^{2^{m}}=e^{i\,2^{m}\varphi}$ trebuie să
satisfacă aceeași condiție, deci $2^{m}\varphi$ (redus în
$(-\pi,\pi]$) are modul $\le \tfrac{\pi}{3}$ pentru orice $m$. Dar dublarea
unui unghi din $[-\tfrac{\pi}{3},\tfrac{\pi}{3}]$ dă un unghi din
$[-\tfrac{2\pi}{3},\tfrac{2\pi}{3}]\subset(-\pi,\pi]$, deci nu e nevoie de
reducere: prin inducție unghiul redus este chiar $2^{m}\varphi$, iar
$|2^{m}\varphi|\le\tfrac{\pi}{3}$ pentru orice $m$ forțează
$\varphi=0$. Așadar singura rădăcină nenulă posibilă
este $r=1$, deci rădăcinile lui $P$ sunt în $\{0,1\}$ și
$P(z)=z^{p}(z-1)^{q}$ (cu $a=1$, căci $a^{\,k-1}=a$).

Înlocuind,
\[
  z^{2p}(z-1)^{q}(z+1)^{q}=z^{p+q}(z-1)^{q}(z+1)^{p},
\]
de unde, comparând exponenții lui $z$ și ai lui $z+1$,
$p=q$. Prin urmare $P(z)=\bigl(z(z-1)\bigr)^{m}=(z^{2}-z)^{m}$.
\end{solutie}

\begin{obs}
Verificare pentru $k=2$, $m=1$: cu $P=z^{2}-z$ avem $P(z^{2})=z^{4}-z^{2}$,
iar $P(z)P(z+1)=z(z-1)\cdot z(z+1)=z^{2}(z^{2}-1)=z^{4}-z^{2}$. Structura
răspunsului arată cât de sensibilă este problema la valoarea
lui $k$: pentru $k=2$ apare o familie infinită, iar de la $k=3$ în
sus supraviețuiesc doar constantele.
\end{obs}

\begin{problema}
Fie $R$ un inel cu unitate, $1\neq 0$, astfel încât $x^{3}=x^{2}$
pentru orice $x\in R$. Demonstrați că $x^{2}=x$ pentru orice
$x\in R$ (deci $R$ este inel Boolean) și deduceți că $R$ este
comutativ.
\end{problema}

\begin{solutie}
(Rezultatul apare în cartea de teorie ca Teorema~\ref{teo:kk1}, ,,condiția
$x^{k+1}=x^{k}$ dă inel Boolean'', împreună cu Teorema~\ref{teo:boole},
,,inele Booleene'': dacă $x^{2}=x$ pentru orice $x$, atunci $1+1=0$
și inelul e comutativ. Îl redemonstrăm aici complet, pentru ca
soluția să fie autonomă.)

\smallskip
\textit{Pasul 1: o relație liniarizată.} Aplicăm ipoteza
elementului $x+1$ (care comută cu $1$):
\[
  (x+1)^{3}=(x+1)^{2}
  \;\Longrightarrow\;
  x^{3}+3x^{2}+3x+1 = x^{2}+2x+1 .
\]
Înlocuind $x^{3}=x^{2}$ în membrul stâng și reducând,
\begin{equation*}
  3x^{2}+x = 0 \qquad \text{pentru orice } x . \tag{A}
\end{equation*}

\textit{Pasul 2: caracteristica este $2$.} Aplicăm $(\mathrm{A})$
elementului $x+1$:
\begin{equation*}
  3(x+1)^{2}+(x+1) = 3x^{2}+7x+4 = 0 . \tag{A$'$}
\end{equation*}
Scăzând $(\mathrm{A})$ din $(\mathrm{A}')$,
\begin{equation*}
  6x+4 = 0 \qquad \text{pentru orice } x . \tag{B}
\end{equation*}
Punând $x=0$ în $(\mathrm{B})$ obținem $4\cdot 1=0$, iar $x=1$
dă $10\cdot 1=0$. Cum $2 = 10-2\cdot 4$, rezultă $2\cdot 1 = 0$,
adică $x+x=0$ pentru orice $x$.

\smallskip
\textit{Pasul 3: idempotență.} În caracteristică $2$ avem
$3\cdot 1 = 1$, deci $3x^{2}=x^{2}$, iar $(\mathrm{A})$ devine $x^{2}+x=0$,
adică
\[
  x^{2} = -x = x \qquad \text{pentru orice } x .
\]

\textit{Pasul 4: comutativitate.} Pentru orice $x,y$, din Pasul 3,
$(x+y)^{2}=x+y$. Dezvoltând,
\[
  (x+y)^{2} = x^{2}+xy+yx+y^{2} = x+xy+yx+y ,
\]
iar egalitatea cu $x+y$ dă $xy+yx=0$, deci $xy=-yx=yx$. Așadar $R$
este comutativ.
\end{solutie}

\begin{obs}
Reciproca este imediată: dacă $x^{2}=x$ atunci
$x^{3}=x\cdot x^{2}=x^{2}$. Prin urmare inelele cu $x^{3}=x^{2}$ sunt
\emph{exact} inelele Boolean, iar ipoteza aparent mai slabă
,,$x^{3}=x^{2}$'' este echivalentă cu ,,$x^{2}=x$''. Caracteristica
$2$, care nu apare în enunț, se naște singură din două
substituții $x\mapsto x+1$.
\end{obs}

\begin{problema}
Fie $G$ un grup finit de ordin $n$ și $k$ un număr întreg
pozitiv astfel încât $f(x)=x^{k}$ este morfism al lui $G$.
Presupunem că mulțimea
\[
  A = \{\, a \in G : a^{k}b = ba^{k} \ \text{pentru orice } b \in G \,\}
\]
are cel puțin $\tfrac{n}{2}+1$ elemente. Demonstrați că
$(ab)^{k}=(ba)^{k}$ pentru orice $a,b \in G$.
\end{problema}

\begin{solutie}
Observăm întâi că $A=\{\,a\in G : a^{k}\in Z(G)\,\}$, unde
$Z(G)$ este centrul lui $G$ (elementele care comută cu tot grupul).

\smallskip
\textit{Pasul 1: reducerea prin morfism.} Cum $f$ este morfism,
\[
  (ab)^{k}=a^{k}b^{k}, \qquad (ba)^{k}=b^{k}a^{k}.
\]
Prin urmare concluzia $(ab)^{k}=(ba)^{k}$ este echivalentă cu
$a^{k}b^{k}=b^{k}a^{k}$, adică ,,$a^{k}$ comută cu $b^{k}$''.

\smallskip
\textit{Pasul 2: $A$ este subgrup.} Folosim că $Z(G)$ este subgrup
și că $f$ este morfism:
\begin{itemize}[nosep,leftmargin=1.4em]
  \item $e \in A$, căci $e^{k}=e\in Z(G)$;
  \item dacă $a,b \in A$, atunci $(ab)^{k}=a^{k}b^{k}\in Z(G)$ (produs de
        elemente centrale), deci $ab\in A$;
  \item dacă $a \in A$, atunci $\bigl(a^{-1}\bigr)^{k}=\bigl(a^{k}\bigr)^{-1}\in Z(G)$
        (inversul unui element central), deci $a^{-1}\in A$.
\end{itemize}
Așadar $A \leq G$.

\smallskip
\textit{Pasul 3: $A=G$.} Avem $|A|\ge \tfrac{n}{2}+1 > \tfrac{n}{2}$, iar
$|A|\mid n$ (Lagrange). Singurul divizor al lui $n$ strict mai mare decât
$\tfrac{n}{2}$ este $n$ însuși, deci $|A|=n$, adică $A=G$.

\smallskip
\textit{Pasul 4: concluzie.} Din $A=G$, pentru orice $a$ avem
$a^{k}\in Z(G)$. În particular $a^{k}$ comută cu $b^{k}$, deci
$a^{k}b^{k}=b^{k}a^{k}$; combinând cu Pasul 1,
\[
  (ab)^{k}=a^{k}b^{k}=b^{k}a^{k}=(ba)^{k}.
\]
\end{solutie}

\begin{obs}
Rolul ipotezei de morfism este dublu: pe de o parte transformă
comparația $(ab)^{k}$ cu $(ba)^{k}$ în comparația lui $a^{k}$ cu
$b^{k}$ (Pasul 1), pe de altă parte face ca $A$ să fie subgrup
(închiderea la produs din Pasul 2 folosește $(ab)^{k}=a^{k}b^{k}$).
Fără ea, $A$ nu ar fi neapărat subgrup, iar argumentul de
divizor ar cădea.
\end{obs}

\begin{problema}
Fie $G$ un grup finit de ordin $n$ cu proprietatea că pentru orice
$g\in G$, $g\neq e$, există un element $x\in G$ cu $gx\neq xg$ (altfel
spus, niciun element diferit de $e$ nu comută cu toate elementele lui
$G$). Demonstrați că $n$ divide $|\operatorname{Aut}(G)|$,
numărul automorfismelor lui $G$.
\end{problema}

\begin{solutie}
Pentru fiecare $g\in G$, conjugarea $\alpha_{g}:G\to G$,
$\alpha_{g}(x)=gxg^{-1}$, este automorfism al lui $G$: e morfism, căci
\[
  \alpha_{g}(xy)=g\,xy\,g^{-1}=(gxg^{-1})(gyg^{-1})=\alpha_{g}(x)\,\alpha_{g}(y),
\]
și bijectiv, având inversul $\alpha_{g^{-1}}$.

Considerăm aplicația $\Phi:G\to\operatorname{Aut}(G)$,
$\Phi(g)=\alpha_{g}$. Ea este morfism de grupuri:
\[
  \alpha_{g}\circ\alpha_{h}(x)=g\,(hxh^{-1})\,g^{-1}=(gh)\,x\,(gh)^{-1}=\alpha_{gh}(x),
\]
deci $\Phi(gh)=\Phi(g)\,\Phi(h)$. Nucleul său este mulțimea
elementelor care comută cu întreg $G$ (adică centrul $Z(G)$):
\[
  \ker\Phi=\{\,g\in G : gxg^{-1}=x \ \forall x\in G\,\}
          =\{\,g\in G : gx=xg \ \forall x\in G\,\},
\]
care, conform ipotezei, este $\{e\}$. Deci $\Phi$ este injectiv. Prin urmare imaginea sa --- subgrupul
automorfismelor interioare $\operatorname{Inn}(G)=\Phi(G)$ --- este
izomorfă cu $G$, deci $|\operatorname{Inn}(G)|=|G|=n$.

Cum $\operatorname{Inn}(G)$ este subgrup al lui $\operatorname{Aut}(G)$, din
teorema lui Lagrange $|\operatorname{Inn}(G)|$ divide $|\operatorname{Aut}(G)|$.
Așadar $n\mid|\operatorname{Aut}(G)|$.
\end{solutie}

\begin{obs}
Ipoteza trebuie în forma ei tare: nu e suficient ca $G$ să fie doar
neabelian (adică să \emph{existe} $x,y$ cu $xy\neq yx$); se cere ca
\emph{fiecare} element $\neq e$ să nu comute cu tot grupul, ceea ce
înseamnă exact centru trivial, $Z(G)=\{e\}$. Aceasta este exact ce
face $\Phi$ injectiv, deci $\operatorname{Inn}(G)\cong G$. În general
$\operatorname{Inn}(G)\cong G/Z(G)$, deci $|\operatorname{Inn}(G)|=[G:Z(G)]$
divide întotdeauna $|\operatorname{Aut}(G)|$ --- însă este egal
cu $n$ doar când centrul e trivial. De exemplu $S_{3}$ are
$Z(S_{3})=\{e\}$, deci $6\mid|\operatorname{Aut}(S_{3})|$; de fapt
$\operatorname{Aut}(S_{3})\cong S_{3}$, cu exact $6$ automorfisme.
\end{obs}

\begin{problema}
Fie $G$ un grup finit cu $|Z(G)|=2$ care nu admite niciun morfism surjectiv
pe $\mathbb{Z}/2$ (echivalent: $G$ nu are subgrup de indice $2$, adică
$\operatorname{Hom}(G,\mathbb{Z}/2)=\{0\}$). Demonstrați că
$Z(\operatorname{Aut}G)=\{\operatorname{id}\}$.
\end{problema}

\begin{lemat}[conjugarea unui automorfism interior]\label{lema:conjinn}
Pentru orice $\varphi\in\operatorname{Aut}(G)$ și $g\in G$, notând
$\alpha_{g}(x)=gxg^{-1}$, are loc
$\varphi\circ\alpha_{g}\circ\varphi^{-1}=\alpha_{\varphi(g)}$.
\end{lemat}

\begin{proof}
$(\varphi\alpha_{g}\varphi^{-1})(x)=\varphi\bigl(g\,\varphi^{-1}(x)\,g^{-1}\bigr)
=\varphi(g)\,x\,\varphi(g)^{-1}=\alpha_{\varphi(g)}(x)$.
\end{proof}

\begin{solutie}
Fie $Z(G)=\{e,t\}$ (cu $t^{2}=e$). Folosim că
\[
  \alpha_{a}=\alpha_{b}\iff a^{-1}b\in Z(G),
\]
căci $\alpha_{a}=\alpha_{b}\iff (a^{-1}b)\,x\,(a^{-1}b)^{-1}=x\ \forall x
\iff a^{-1}b\in Z(G)$.

Fie $\varphi\in Z(\operatorname{Aut}G)$. Cum $\varphi$ comută cu orice
automorfism, comută în particular cu fiecare $\alpha_{g}$, deci
$\varphi\alpha_{g}\varphi^{-1}=\alpha_{g}$; din lema de mai sus,
$\alpha_{\varphi(g)}=\alpha_{g}$, adică
\[
  g^{-1}\varphi(g)\in Z(G)=\{e,t\}\qquad\text{pentru orice }g\in G.
\]
Definim $\chi:G\to\{e,t\}$ prin $\chi(g)=g^{-1}\varphi(g)$, deci
$\varphi(g)=g\,\chi(g)$. Cum $\varphi$ este morfism și $t$ este central,
\[
  gh\,\chi(gh)=\varphi(gh)=\varphi(g)\varphi(h)=g\chi(g)\,h\chi(h)=gh\,\chi(g)\chi(h),
\]
de unde $\chi(gh)=\chi(g)\chi(h)$: așadar $\chi$ este morfism
$G\to\{e,t\}\cong\mathbb{Z}/2$.

Prin ipoteză nu există morfisme nenule $G\to\mathbb{Z}/2$, deci
$\chi\equiv e$; atunci $\varphi(g)=g$ pentru orice $g$, adică
$\varphi=\operatorname{id}$. Prin urmare $Z(\operatorname{Aut}G)=\{\operatorname{id}\}$.
\end{solutie}

\begin{obs}
Comparație cu cazul centrului trivial (Lema~\ref{lema:centrutrivial}
din partea de teorie): dacă $Z(G)=\{e\}$, atunci
$Z(\operatorname{Aut}G)=\{\operatorname{id}\}$ \emph{necondiționat},
fiindcă $g\mapsto\alpha_{g}$ e injectivă (exact argumentul din
problema ONM.7), deci $\alpha_{\varphi(g)}=\alpha_{g}$ forțează direct
$\varphi(g)=g$. La
$|Z(G)|=2$ apare libertatea $\varphi(g)\in\{g,\,tg\}$, codificată de un
morfism $\chi:G\to\mathbb{Z}/2$; ipoteza ,,fără subgrup de indice
$2$'' îl omoară pe $\chi$. De exemplu $\operatorname{SL}(2,3)$ are
$|Z|=2$, iar abelianizatul său este $\cong\mathbb{Z}/3$ (deci
$\operatorname{Hom}(\cdot,\mathbb{Z}/2)=\{0\}$), așadar
$Z(\operatorname{Aut}\operatorname{SL}(2,3))=\{\operatorname{id}\}$.
\end{obs}

\begin{problema}
Fie $G$ un grup și $a,b,c\in\mathbb{N}^{*}$ numere prime între ele
două câte două. Presupunem că
\[
  (xy)^{n}=x^{n}y^{n}\qquad\text{pentru toți } x,y\in G,
\]
ori de câte ori $n$ este multiplu al lui $a$, al lui $b$ sau al lui $c$.
Demonstrați că $G$ este comutativ.
\end{problema}

Numim un exponent $m$ \emph{bun} dacă $(xy)^{m}=x^{m}y^{m}$ pentru orice
$x,y\in G$. Ipoteza spune că toți multiplii lui $a$, ai lui $b$
și ai lui $c$ sunt buni.

\begin{lemat}[trei exponenți consecutivi]\label{lema:treiexp}
Dacă trei valori consecutive $n,\ n+1,\ n+2$ sunt exponenți buni,
atunci $G$ este comutativ.
\end{lemat}

\begin{proof}
\textit{Pas 1: dacă $m$ și $m+1$ sunt buni, atunci $y^{m}\in Z(G)$
pentru orice $y$.} Folosind bunătatea lui $m+1$ și apoi a lui $m$,
\[
  x^{m}\,x\,y^{m}\,y = x^{m+1}y^{m+1}=(xy)^{m+1}=(xy)^{m}(xy)=x^{m}y^{m}\,x\,y.
\]
Simplificăm la stânga cu $x^{m}$ și la dreapta cu $y$:
$x\,y^{m}=y^{m}\,x$. Cum $x$ e arbitrar, $y^{m}\in Z(G)$.

\textit{Pas 2.} Aplicând Pasul 1 perechilor $(n,n+1)$ și $(n+1,n+2)$,
obținem $y^{n}\in Z(G)$ și $y^{n+1}\in Z(G)$ pentru orice $y$. Atunci
\[
  y = y^{n+1}\bigl(y^{n}\bigr)^{-1}\in Z(G)
\]
(produs de două elemente centrale), deci orice $y$ este central și
$G$ este comutativ.
\end{proof}

\begin{solutie}
Deoarece $a,b,c$ sunt prime între ele două câte două,
sistemul
\[
  m\equiv 0\pmod a,\qquad m\equiv -1\pmod b,\qquad m\equiv -2\pmod c
\]
are soluție prin teorema chineză a resturilor; alegem pentru $m$ o
soluție întreagă pozitivă. Atunci $m$ este multiplu al lui
$a$, $m+1$ este multiplu al lui $b$, iar $m+2$ este multiplu al lui $c$;
prin ipoteză, exponenții consecutivi $m,\ m+1,\ m+2$ sunt toți
buni. Din lema celor trei exponenți consecutivi, $G$ este comutativ.
\end{solutie}

\begin{obs}
Sunt necesari \emph{trei} exponenți consecutivi: doi nu ajung. De pildă,
în $S_{3}$ exponenții $6$ și $7$ sunt buni (căci $x^{6}=e$
și $x^{7}=x$ pentru orice $x$, deci $(xy)^{6}=e=x^{6}y^{6}$ și
$(xy)^{7}=xy=x^{7}y^{7}$), însă $S_{3}$ nu este comutativ; cu doi
exponenți consecutivi se obține doar că $n$-puterile sunt
centrale. Rolul numerelor prime între ele este ca, prin CRT, să
garanteze existența a trei exponenți buni \emph{consecutivi}
în reuniunea multiplilor lor.
\end{obs}

\begin{problema}
Fie $G$ un grup finit de ordin $p^{\alpha}q^{\beta}$, cu $p\neq q$ numere
prime și $\alpha,\beta\ge 1$. Demonstrați că există
subgrupuri \emph{proprii} $H,K$ ale lui $G$ (adică $H\neq G$, $K\neq G$) astfel încât $HK=G$ (unde
$HK=\{\,hk : h\in H,\ k\in K\,\}$).
\end{problema}

\begin{lemat}[cardinalul unui produs de subgrupuri]
Pentru orice subgrupuri $H,K$ ale unui grup finit,
\[
  |HK|=\frac{|H|\,|K|}{|H\cap K|}.
\]
\end{lemat}

\begin{proof}
Considerăm aplicația surjectivă $H\times K\to HK$,
$(h,k)\mapsto hk$. Fixăm $g=hk\in HK$; căutăm toate perechile
$(h',k')$ cu $h'k'=hk$. Egalitatea revine la $h^{-1}h'=k(k')^{-1}$; membrul
stâng e în $H$, cel drept în $K$, deci valoarea comună
$d:=h^{-1}h'=k(k')^{-1}$ aparține lui $H\cap K$. Reciproc, orice
$d\in H\cap K$ dă o pereche $(h',k')=(hd,\,d^{-1}k)$ cu $h'k'=hk$. Prin
urmare fibra deasupra lui $g$ are exact $|H\cap K|$ elemente, iar cum toate
fibrele au aceeași mărime,
$|H|\,|K|=|HK|\cdot|H\cap K|$.
\end{proof}

\begin{solutie}
Reamintim \emph{teoremele lui Sylow} pentru un grup finit $G$ de ordin
$p^{a}m$ cu $p$ prim și $p\nmid m$: (I) $G$ are un subgrup de ordin
$p^{a}$ (un \emph{$p$-subgrup Sylow}); (II) oricare două $p$-subgrupuri
Sylow sunt conjugate; (III) numărul $n_{p}$ al $p$-subgrupurilor Sylow
satisface $n_{p}\equiv 1\pmod p$ și $n_{p}\mid m$; în particular
$n_{p}=1$ dacă și numai dacă $p$-subgrupul Sylow e normal.

Prin Sylow (I), $G$ are un $p$-subgrup Sylow $P$ (de ordin
$p^{\alpha}$) și un $q$-subgrup Sylow $Q$ (de ordin $q^{\beta}$). Luăm
$H=P$ și $K=Q$. Ambele sunt proprii: $|P|=p^{\alpha}<p^{\alpha}q^{\beta}=|G|$,
căci $\beta\ge1$, și analog $|Q|<|G|$, căci $\alpha\ge1$.

Intersecția $P\cap Q$ este subgrup atât al lui $P$ cât și al
lui $Q$, deci (Lagrange) $|P\cap Q|$ divide atât $p^{\alpha}$ cât
și $q^{\beta}$; cum $p\neq q$, acești doi divizori sunt primi
între ei, deci $|P\cap Q|=1$, adică $P\cap Q=\{e\}$.

Din lemă,
\[
  |PQ|=\frac{|P|\,|Q|}{|P\cap Q|}=\frac{p^{\alpha}q^{\beta}}{1}=p^{\alpha}q^{\beta}=|G|.
\]
Cum $PQ\subseteq G$ și $|PQ|=|G|$, rezultă $PQ=G$. Așadar
$HK=PQ=G$.
\end{solutie}

\begin{obs}
Nu se afirmă că $HK$ este subgrup --- în general $PQ$ e subgrup
doar dacă unul dintre $P,Q$ îl normalizează pe celălalt; aici
se cere doar egalitatea de \emph{mulțimi} $HK=G$, care rezultă
pur din numărare (ordinele se înmulțesc, intersecția e
trivială). Esențial este pasul de la Cauchy la Sylow: Cauchy ar da
doar elemente de ordin $p$ și $q$, pe când aici avem nevoie de
subgrupuri \emph{întregi} de ordin $p^{\alpha}$ și $q^{\beta}$, adică
de existența subgrupurilor Sylow.
\end{obs}

\begin{problema}
Fie $G$ un grup abelian finit cu $|G|$ care are cel puțin doi factori
primi distincți. Demonstrați că există două subgrupuri
\emph{proprii și netriviale} $K,L\le G$ (adică diferite de $\{e\}$ și de $G$) cu $KL=G$.
\end{problema}

\begin{solutie}
Fie $|G|=p_{1}^{\alpha_{1}}\cdots p_{r}^{\alpha_{r}}$, cu $r\ge2$ prime
distincte și $\alpha_{i}\ge1$. Prin \emph{teorema lui Sylow} (un grup
finit al cărui ordin e divizibil cu $p^{\alpha}$, dar nu cu $p^{\alpha+1}$,
are un subgrup de ordin $p^{\alpha}$), pentru fiecare $i$ există un
subgrup $P_{i}\le G$ cu $|P_{i}|=p_{i}^{\alpha_{i}}$. Folosim două fapte:
\begin{itemize}[nosep,leftmargin=1.4em]
  \item într-un grup abelian, produsul $AB$ a două subgrupuri este
        subgrup (din $ab\,(a'b')^{-1}=(aa'^{-1})(bb'^{-1})$ și criteriul
        de subgrup);
  \item formula pentru \emph{două} subgrupuri, demonstrată la problema
        ONM.10: $|AB|=|A|\,|B|/|A\cap B|$.
\end{itemize}

\smallskip
\textit{Afirmație.} Pentru $j=1,\dots,r$, mulțimea
$A_{j}=P_{1}P_{2}\cdots P_{j}$ este subgrup de ordin
$p_{1}^{\alpha_{1}}\cdots p_{j}^{\alpha_{j}}$.

Demonstrăm prin inducție după $j$. Pentru $j=1$, $A_{1}=P_{1}$. Fie
$2\le j\le r$ și presupunem afirmația adevărată pentru $j-1$. Atunci
$A_{j}=A_{j-1}P_{j}$ este subgrup (primul fapt). Ordinele
$|A_{j-1}|=p_{1}^{\alpha_{1}}\cdots p_{j-1}^{\alpha_{j-1}}$ și
$|P_{j}|=p_{j}^{\alpha_{j}}$ sunt prime între ele, iar $|A_{j-1}\cap P_{j}|$
le divide pe amândouă (Lagrange), deci $A_{j-1}\cap P_{j}=\{e\}$. Din al
doilea fapt,
\[
  |A_{j}|=\frac{|A_{j-1}|\,|P_{j}|}{|A_{j-1}\cap P_{j}|}
  =p_{1}^{\alpha_{1}}\cdots p_{j}^{\alpha_{j}} .
\]

Pentru $j=r$ obținem $|A_{r}|=|G|$, deci $P_{1}P_{2}\cdots P_{r}=G$.

\smallskip
\textit{Cele două subgrupuri.} Luăm $K=P_{1}$ și
$L=P_{2}\cdots P_{r}$; exact ca mai sus (inducția aplicată
subgrupurilor $P_{2},\dots,P_{r}$), $L$ este subgrup de ordin
$p_{2}^{\alpha_{2}}\cdots p_{r}^{\alpha_{r}}$. Ambele sunt netriviale
($|K|=p_{1}^{\alpha_{1}}>1$ și $|L|>1$, căci $r\ge2$) și proprii
($|K|<|G|$ și $|L|<|G|$, căci fiecăruia îi lipsește cel puțin un factor
prim al lui $|G|$). În fine, $KL=P_{1}(P_{2}\cdots P_{r})=G$, prin
afirmația de mai sus.
\end{solutie}

\begin{obs}
Inducția folosește la fiecare pas doar formula pentru două
subgrupuri. Ea se generalizează la un \emph{lanț} de $k$ subgrupuri:
dacă $\bigl(H_{1}\cdots H_{i}\bigr)\cap H_{i+1}=\{e\}$ pentru orice $i$,
atunci
\[
  |H_{1}H_{2}\cdots H_{k}|=|H_{1}|\,|H_{2}|\cdots|H_{k}| ,
\]
prin aceeași inducție ($\,|A\,H_{i+1}|=|A|\,|H_{i+1}|$ când
$A\cap H_{i+1}=\{e\}$). Atenție însă: pentru \emph{trei sau mai
multe} subgrupuri nu există o formulă generală în funcție
doar de ordinele intersecțiilor perechi --- ipoteza ,,în lanț''
este esențială. Aici, comutativitatea o asigură: se obține
chiar $G\cong P_{1}\times\cdots\times P_{r}$ (descompunerea în
$p$-componente), iar problema ONM.10 este cazul $r=2$.
\end{obs}

\begin{problema}
Fie $n\ge 5$. Determinați toate relațiile de congruență pe
grupul simetric $S_{n}$. (O congruență este o relație de
echivalență compatibilă cu compunerea permutărilor.)
\end{problema}

\noindent\textbf{Răspuns:} exact \emph{trei} congruențe --- egalitatea,
,,același semn'' (congruența modulo $A_{n}$) și relația
totală.

\begin{solutie}
Prin Teorema~\ref{teo:congruenta}, congruențele pe $S_{n}$ sunt în corespondență
bijectivă cu subgrupurile normale ale lui $S_{n}$: fiecărui
$N\trianglelefteq S_{n}$ îi corespunde $\sigma\,\rho\,\tau\iff\sigma^{-1}\tau\in N$.
Rămâne să determinăm subgrupurile normale ale lui $S_{n}$.

\smallskip
\textit{Subgrupurile normale ale lui $S_{n}$ ($n\ge 5$) sunt $\{e\}$,
$A_{n}$, $S_{n}$.} Fie $N\trianglelefteq S_{n}$, $N\neq\{e\}$. Atunci
$N\cap A_{n}\trianglelefteq A_{n}$; cum $A_{n}$ este \emph{simplu} pentru
$n\ge 5$ (teoremă clasică: pentru $n\ge 5$, singurele subgrupuri
normale ale lui $A_{n}$ sunt $\{e\}$ și $A_{n}$), avem
$N\cap A_{n}=\{e\}$ sau $N\cap A_{n}=A_{n}$.
\begin{itemize}[nosep,leftmargin=1.4em]
  \item Dacă $N\cap A_{n}=A_{n}$, atunci $A_{n}\subseteq N$; cum
        $[S_{n}:A_{n}]=2$, rezultă $N=A_{n}$ sau $N=S_{n}$.
  \item Dacă $N\cap A_{n}=\{e\}$, atunci $N$ se scufundă în
        $S_{n}/A_{n}\cong\mathbb{Z}/2$, deci $|N|\le 2$. Un subgrup normal de
        ordin $2$, $N=\{e,\tau\}$, ar fi central (pentru orice $g$,
        $g\tau g^{-1}\in N\setminus\{e\}=\{\tau\}$), dar $Z(S_{n})=\{e\}$
        pentru $n\ge 3$; deci $|N|=1$, contradicție.
\end{itemize}

\smallskip
Celor trei subgrupuri normale le corespund trei congruențe:
\begin{itemize}[nosep,leftmargin=1.4em]
  \item $N=\{e\}$: $\sigma\,\rho\,\tau\iff\sigma=\tau$ --- \emph{egalitatea};
  \item $N=A_{n}$: $\sigma\,\rho\,\tau\iff\sigma^{-1}\tau\in A_{n}\iff
        \operatorname{sgn}(\sigma)=\operatorname{sgn}(\tau)$ --- două
        permutări sunt în relație exact când au \emph{același
        semn} (ambele pare sau ambele impare);
  \item $N=S_{n}$: orice două permutări sunt în relație ---
        \emph{relația totală}.
\end{itemize}
Acestea sunt toate. $\blacksquare$
\end{solutie}

\begin{obs}
Cazurile mici se abat exact acolo unde $A_{n}$ nu mai e simplu: pentru
$n=4$, $S_{4}$ are în plus subgrupul normal $V_{4}$ (grupul lui Klein),
deci \emph{patru} congruențe; pentru $n=3$, $A_{3}$ e simplu (ordin
prim), deci tot trei; pentru $n=2$, $S_{2}\cong\mathbb{Z}/2$ are două.
Frumusețea: o întrebare despre \emph{relații de echivalență}
pe $S_{n}$ se reduce, prin Teorema~\ref{teo:congruenta}, la o întrebare despre
\emph{subgrupuri normale}, rezolvată de simplitatea lui $A_{n}$.
\end{obs}

\begin{problema}[Clasică]
Fie $G$ un grup finit \emph{neabelian} de ordin $n$. Demonstrați că
numărul perechilor ordonate $(a,b)\in G\times G$ cu $ab=ba$ este cel mult
$\dfrac{5}{8}\,n^{2}$.
\end{problema}

\begin{lemat}
Dacă $G$ este un grup finit neabelian de ordin $n$, atunci
$|Z(G)|\le n/4$.
\end{lemat}

\begin{proof}
Fie $Z=Z(G)$. Dacă $G/Z$ ar fi ciclic, generat de $gZ$, orice element ar
fi de forma $g^{k}z$ cu $z\in Z$, iar două astfel de elemente comută
--- deci $G$ ar fi abelian, contradicție. Prin urmare $G/Z$ nu e ciclic,
deci $|G/Z|\ge 4$ (grupurile de ordin $1,2,3$ sunt ciclice), adică
$|Z|\le n/4$.
\end{proof}

\begin{solutie}
Fie $C=\{(a,b)\in G\times G:ab=ba\}$; vrem $|C|\le\tfrac58 n^{2}$. Pentru
$a\in G$ notăm $C_{G}(a)=\{b:ab=ba\}$ centralizatorul, $\operatorname{cl}(a)$
clasa de conjugare, și $k$ numărul claselor de conjugare ale lui $G$.

\smallskip
\textit{Pas 1: $|C|=k\,n$.} Numărăm perechile din $C$ ,,felicând''
după prima coordonată. Pentru un $a\in G$ fixat, perechile din $C$ cu
prima componentă $a$ sunt exact cele $(a,b)$ cu $b\in C_{G}(a)$, deci felia
lui $a$ are $|C_{G}(a)|$ elemente. Feliile corespunzătoare la $a$ diferiți
sunt disjuncte (au prime coordonate diferite) și acoperă tot $C$, deci
\[
  |C|=\sum_{a\in G}|C_{G}(a)|
\]
--- deocamdată o simplă numărare (împărțim o mulțime
în felii disjuncte), fără vreo teoremă de grup.

Acum evaluăm fiecare $|C_{G}(a)|$: grupul $G$ acționează pe sine
prin conjugare, cu orbita lui $a$ egală cu $\operatorname{cl}(a)$ și
stabilizatorul egal cu $C_{G}(a)$. Prin teorema orbită--stabilizator
(pentru o acțiune a unui grup finit, $|\!\operatorname{Orb}(a)|\cdot
|\!\operatorname{Stab}(a)|=|G|$), rezultă
$|C_{G}(a)|=\dfrac{n}{|\operatorname{cl}(a)|}$. Așadar
\[
  |C|=\sum_{a\in G}\frac{n}{|\operatorname{cl}(a)|}
     =n\sum_{a\in G}\frac{1}{|\operatorname{cl}(a)|}.
\]
În ultima sumă grupăm termenii după clasele de conjugare
(exact partiția lui $G$ care stă la baza \emph{ecuației claselor}):
o clasă $\mathcal K$ are $|\mathcal K|$ elemente, iar fiecare $a\in\mathcal K$
dă $\dfrac{1}{|\operatorname{cl}(a)|}=\dfrac{1}{|\mathcal K|}$, deci
contribuția totală a clasei este $|\mathcal K|\cdot\dfrac{1}{|\mathcal K|}=1$.
Cum sunt $k$ clase, $\sum_{a\in G}\dfrac{1}{|\operatorname{cl}(a)|}=k$, deci
$|C|=k\,n$. (Atenție: prima egalitate $|C|=\sum_a|C_G(a)|$ este numărare,
nu ecuația claselor; aceasta din urmă intervine abia la gruparea pe
clase, unde se folosește partiția lui $G$ în clasele de conjugare.)

\smallskip
\textit{Pas 2: majorarea lui $k$.} Un element are clasa de mărime $1$
exact când e central. Deci cele $|Z|$ elemente centrale dau $|Z|$ clase de
mărime $1$, iar celelalte $n-|Z|$ elemente stau în clase de mărime
$\ge 2$, adică în cel mult $\dfrac{n-|Z|}{2}$ clase. Folosind
și lema ($|Z|\le n/4$),
\[
  k\le |Z|+\frac{n-|Z|}{2}=\frac{n}{2}+\frac{|Z|}{2}
   \le \frac{n}{2}+\frac{n}{8}=\frac{5n}{8}.
\]
Împreună cu Pasul 1, $|C|=k\,n\le\dfrac{5n}{8}\cdot n=\dfrac{5}{8}\,n^{2}$.
\end{solutie}

\begin{obs}
Împărțind la $n^{2}$: probabilitatea ca două elemente alese la
întâmplare să comute este $k/n\le 5/8$ pentru \emph{orice} grup
neabelian --- nu există grupuri ,,aproape comutative'' cu probabilitate
între $5/8$ și $1$. Marginea e atinsă: la $D_{4}$ și la $Q_{8}$
(ordin $8$, $|Z|=2$) avem $k=5$ și $|C|=5\cdot 8=40=\tfrac58\cdot 64$.

Identitatea $|C|=k\,n$ este și \emph{lema lui Burnside} pentru
conjugare, deghizată: punctele fixe ale conjugării cu $g$ sunt
$\operatorname{Fix}(g)=\{a:gag^{-1}=a\}=C_{G}(g)$, deci numărul de orbite
(clase) este $\dfrac{1}{n}\sum_{g\in G}|\operatorname{Fix}(g)|
=\dfrac{1}{n}\sum_{g}|C_{G}(g)|=\dfrac{|C|}{n}=k$.
\end{obs}

\begin{problema}
Fie $G$ un grup finit de ordin $n$ și $k$ un întreg astfel încât
aplicația $x\mapsto x^{k}$ este morfism de grupuri (adică
$(xy)^{k}=x^{k}y^{k}$ pentru orice $x,y\in G$). Arătați că, pentru
orice $x\in G$ fixat, numărul soluțiilor $a\in G$ ale ecuației
$xa^{k}=a^{k}x$ divide $n$.
\end{problema}

\begin{solutie}
Fie $f:G\to G$, $f(a)=a^{k}$; din ipoteză $f$ este morfism, iar
$G^{k}:=f(G)=\{a^{k}:a\in G\}$ este subgrup (imaginea unui morfism).

Fixăm $x$ și notăm cu $T$ mulțimea soluțiilor:
\[
  T=\{a\in G : x a^{k}=a^{k}x\}=\{a\in G : a^{k}\in C_{G}(x)\}=f^{-1}\bigl(C_{G}(x)\bigr),
\]
căci $xa^{k}=a^{k}x$ înseamnă exact că $a^{k}$ comută cu
$x$, adică $a^{k}\in C_{G}(x)$. Centralizatorul $C_{G}(x)$ este subgrup,
iar preimaginea unui subgrup printr-un morfism este subgrup; deci $T$ este
subgrup al lui $G$. Prin teorema lui Lagrange, $|T|$ divide $|G|=n$.

\smallskip
\textit{Soluția a doua (prin acțiune).} Definim o \emph{acțiune}
a lui $G$ pe mulțimea $G$ prin
\[
  a\cdot y:=a^{k}\,y\,a^{-k}.
\]
Este într-adevăr o acțiune: $e\cdot y=y$, iar folosind ipoteza
că $a\mapsto a^{k}$ e morfism (deci $(ab)^{k}=a^{k}b^{k}$),
\[
  a\cdot(b\cdot y)=a^{k}\bigl(b^{k}y\,b^{-k}\bigr)a^{-k}
  =(a^{k}b^{k})\,y\,(a^{k}b^{k})^{-1}=(ab)^{k}\,y\,(ab)^{-k}=(ab)\cdot y .
\]
(Tocmai ipoteza de morfism face ca această acțiune să existe ---
fără ea, $(ab)^{k}\neq a^{k}b^{k}$ și compunerea s-ar rupe.)

Pentru $x$ fixat, ecuația $xa^{k}=a^{k}x$ se rescrie
$a^{k}xa^{-k}=x$, adică $a\cdot x=x$: mulțimea soluțiilor în
$a$ este exact \emph{stabilizatorul} lui $x$ pentru această acțiune,
$T=\operatorname{Stab}(x)$. Stabilizatorul unui punct este întotdeauna
subgrup, iar prin teorema \emph{orbită--stabilizator} (pentru o
acțiune a unui grup finit, $|\!\operatorname{Orb}(x)|\cdot
|\!\operatorname{Stab}(x)|=|G|$)
\[
  |T|=|\!\operatorname{Stab}(x)|=\frac{|G|}{|\!\operatorname{Orb}(x)|}
     =\frac{n}{|\{a^{k}xa^{-k}:a\in G\}|},
\]
deci $|T|$ divide $n$. (Orbita lui $x$ este clasa sa de conjugare sub
subgrupul puterilor $G^{k}=\{a^{k}\}$.)
\end{solutie}

\begin{obs}
Ipoteza că $x\mapsto x^{k}$ e morfism este \emph{esențială}:
ea garantează închiderea lui $T$ la produs (dacă
$a^{k},b^{k}\in C_{G}(x)$, atunci $(ab)^{k}=a^{k}b^{k}\in C_{G}(x)$). Fără
ea, $T$ poate să nu fie subgrup: în $S_{3}$ cu $k=2$ (unde
$x\mapsto x^{2}$ nu e morfism), pentru $x=(1\,2)$ mulțimea $T=\{a:a^{2}\in
C_{G}(x)\}=\{e,(1\,2),(1\,3),(2\,3)\}$ are $4$ elemente, iar $4\nmid 6$.

Interpretare prin acțiune: $a\cdot y:=a^{k}ya^{-k}$ este o acțiune a
lui $G$ pe $G$ (tocmai fiindcă $f$ e morfism, deci $(ab)^{k}=a^{k}b^{k}$),
iar $T$ este stabilizatorul lui $x$; concluzia e orbită--stabilizator.
Când $\gcd(k,n)=1$, $f$ este bijecție (,,ridicarea la puterea $p$'',
Corolarul~\ref{lema:puterep}), deci $G^{k}=G$ și $|T|=|C_{G}(x)|$.
\end{obs}

\begin{problema}
Fie $p\neq q$ numere prime, $\alpha\ge 1$, și $G$ un grup de ordin
$p\,q^{\alpha}$ cu proprietatea că
\[
  p\not\equiv 1\pmod q\qquad\text{și}\qquad q^{j}\not\equiv 1\pmod p
  \ \ \text{pentru }1\le j\le\alpha .
\]
Arătați că $G\cong P\times Q$ (produsul subgrupurilor sale Sylow)
și că $pq\mid|Z(G)|$.
\end{problema}

\begin{solutie}
Din Sylow III, $n_{q}\equiv 1\pmod q$ și $n_{q}\mid p$, deci
$n_{q}\in\{1,p\}$; cum $p\not\equiv 1\pmod q$, rămâne $n_{q}=1$, adică
$q$-subgrupul Sylow $Q$ este normal. La fel, $n_{p}\equiv 1\pmod p$ și
$n_{p}\mid q^{\alpha}$, deci $n_{p}=q^{j}$ cu $q^{j}\equiv 1\pmod p$; ipoteza
exclude $1\le j\le\alpha$, deci $n_{p}=q^{0}=1$ și $P$ e normal.

Ambele Sylow fiind normale, cu $P\cap Q=\{e\}$ (ordine coprime) și
$|PQ|=|P||Q|=|G|$, rezultă $G=P\times Q$. Atunci
\[
  Z(G)=Z(P)\times Z(Q)=P\times Z(Q),
\]
căci $P$ (ordin $p$) e ciclic, deci abelian. Iar $Z(Q)\neq\{e\}$: un
$q$-grup are centru netrivial (ecuația claselor pentru $Q$ --- toți
termenii necentrali sunt multipli de $q$, la fel și $|Q|$, deci
$q\mid|Z(Q)|$). Prin urmare $p\mid|Z(G)|$ și $q\mid|Z(G)|$, adică
$pq\mid|Z(G)|$.
\end{solutie}

\begin{obs}
Congruențele interzic orice număr nenormal de Sylow-uri, deci
,,rup'' grupul în produsul direct $P\times Q$ (nilpotent). Ecuația
claselor intervine exact la $Z(Q)\neq\{e\}$.
\end{obs}

\begin{problema}
Fie $G$ un grup de ordin $p\,q^{\alpha}$ ($p\neq q$ prime) al cărui
$q$-subgrup Sylow $Q$ este normal \emph{și abelian}, iar $P$ un
$p$-subgrup Sylow. Arătați că
\[
  |Z(G)\cap Q|\equiv q^{\alpha}\pmod p,
\]
și deduceți că dacă $q^{\alpha}\not\equiv 1\pmod p$, atunci
$q\mid|Z(G)|$.
\end{problema}

\begin{solutie}
Cum $Q\trianglelefteq G$, conjugarea definește o acțiune a lui $P$ pe
$Q$ prin automorfisme ($a\cdot z=aza^{-1}$). Aplicăm \emph{ecuația
claselor a unei acțiuni}: pentru un grup $H$ ce acționează pe o
mulțime finită $X$, orbitele partiționează $X$, deci, alegând
câte un reprezentant din fiecare orbită și folosind
orbită--stabilizator $|\!\operatorname{Orb}(x)|=[H:\operatorname{Stab}(x)]$,
\[
  |X|=|X^{H}|+\sum_{|\operatorname{Orb}(x_i)|>1}[H:\operatorname{Stab}(x_i)],
\]
unde $X^{H}$ e mulțimea punctelor fixe (orbitele de mărime $1$).
(Pentru $H=G$ acționând pe $X=G$ prin conjugare, $X^{H}=Z(G)$,
$\operatorname{Stab}=C_G$, și se regăsește ecuația claselor
uzuală $|G|=|Z(G)|+\sum_i[G:C_G(x_i)]$.)

Aici $H=P$ acționează pe $X=Q$. Fiecare orbită are mărime
$[P:\operatorname{Stab}(z)]$, divizor al lui $|P|=p$ (prim), deci $1$ sau $p$;
punctele fixe sunt $Q^{P}=\{z\in Q: aza^{-1}=z\ \forall a\in P\}=C_{Q}(P)$.
Ecuația claselor a acestei acțiuni devine
\[
  q^{\alpha}=|Q|=|C_{Q}(P)|+p\cdot\#\{\text{orbite de mărime }p\},
\]
și, reducând modulo $p$ (termenii de mărime $p$ dispar),
\[
  q^{\alpha}\equiv|C_{Q}(P)|\pmod p .
\]

Deoarece $Q$ e abelian, un element $z\in Q$ e central în $G$ dacă
și numai dacă comută cu $P$ (cu $Q$ comută automat), deci
$Z(G)\cap Q=C_{Q}(P)$. Așadar
$|Z(G)\cap Q|\equiv q^{\alpha}\pmod p$.

Dacă $q^{\alpha}\not\equiv 1\pmod p$: $|Z(G)\cap Q|$ este o putere a lui
$q$ (subgrup al $q$-grupului $Q$); dacă ar fi $1$, am avea
$1\equiv q^{\alpha}\pmod p$, fals. Deci $|Z(G)\cap Q|\ge q$, adică
$q\mid|Z(G)|$.
\end{solutie}

\begin{obs}
Aici se folosește \emph{doar} normalitatea lui $Q$ (nu și a lui $P$),
prin ecuația claselor a acțiunii $P\curvearrowright Q$; congruența
$q^{\alpha}\bmod p$ (adică ordinul lui $q$ modulo $p$) decide dacă
partea $q$ a centrului e forțată netrivială.
\end{obs}

\begin{problema}
Fie $G$ un grup finit și $M\subseteq G$ o submulțime cu $m$ elemente,
cu proprietatea că $gMg^{-1}=M$ pentru orice $g\in G$ (adică $M$ are
un singur conjugat, pe el însuși). Arătați că
$C_{G}(M)=\{g:gx=xg\ \forall x\in M\}$ este subgrup normal și că
$[G:C_{G}(M)]$ divide $m!$. Deduceți că dacă
$\gcd\bigl(|G|,m!\bigr)=1$, atunci $M\subseteq Z(G)$.
\end{problema}

\begin{solutie}
Deoarece $gMg^{-1}=M$, pentru orice $x\in M$ și $g\in G$ avem
$gxg^{-1}\in gMg^{-1}=M$: conjugarea permută elementele lui $M$
între ele, deci definește o \emph{acțiune} a lui $G$ pe
mulțimea finită $M$ (prin $g\cdot x=gxg^{-1}$).

Prin Propoziția~\ref{prop:actmorf} (,,acțiuni $=$ morfisme în grupul de
permutări''), acestei acțiuni îi corespunde un morfism
\[
  \varphi:G\to S(M)\cong S_{m},\qquad \varphi(g)(x)=gxg^{-1}.
\]
Nucleul său este
\[
  \operatorname{Ker}\varphi=\{g:\varphi(g)=\operatorname{id}\}
  =\{g:gxg^{-1}=x\ \forall x\in M\}=C_{G}(M),
\]
deci $C_{G}(M)$, fiind nucleu de morfism, este subgrup \emph{normal} al lui
$G$. (Aici invarianța lui $M$ la conjugare e esențială: pentru o
submulțime oarecare, $C_{G}(M)$ nu e neapărat normal.) Prin prima
teoremă de izomorfism,
\[
  G/C_{G}(M)\cong\operatorname{Im}\varphi\hookrightarrow S_{m},
\]
deci $[G:C_{G}(M)]=|\operatorname{Im}\varphi|$ divide $|S_{m}|=m!$.

În fine, $[G:C_{G}(M)]$ divide atât $|G|$ (Lagrange) cât și
$m!$; dacă $\gcd(|G|,m!)=1$, atunci $[G:C_{G}(M)]=1$, adică
$C_{G}(M)=G$. Aceasta înseamnă că orice $g\in G$ comută cu
orice $x\in M$, deci $M\subseteq Z(G)$.
\end{solutie}

\begin{obs}
Este generalizarea la \emph{submulțimi} a scufundărilor clasice
$G/Z(G)\hookrightarrow S_{|G|}$ (Cayley $+$ conjugare) și
$G/C_{G}(a)\hookrightarrow S_{|\operatorname{cl}(a)|}$. Cazuri particulare
de submulțimi invariante la conjugare: orice clasă de conjugare și
orice subgrup normal; pentru un subgrup normal $N$, morfismul se rafinează
la $G/C_{G}(N)\hookrightarrow\operatorname{Aut}(N)$, căci conjugarea
păstrează și structura de grup a lui $N$.
\end{obs}

\begin{problema}
Fie $p_{1}<p_{2}<\dots<p_{n}<\dots$ șirul numerelor prime. Demonstrați
că, pentru orice $n\ge 2$, orice grup de ordin $p_{n}\,p_{n+1}$ este
ciclic.
\end{problema}

\begin{solutie}
Notăm $p=p_{n}$, $q=p_{n+1}$ ($p<q$ prime consecutive, iar $p\ge 3$
căci $n\ge 2$). Folosim două rezultate, pe care le enunțăm
explicit.

\emph{Postulatul lui Bertrand:} pentru orice întreg $m\ge 1$ există
un număr prim $r$ cu $m<r<2m$.

\emph{Criteriul numerelor ciclice:} un întreg $N\ge 1$ se numește
\emph{număr ciclic} dacă orice grup de ordin $N$ este ciclic; are loc
echivalența
\[
  N\ \text{număr ciclic}\iff \gcd\bigl(N,\varphi(N)\bigr)=1 ,
\]
unde $\varphi$ este indicatorul lui Euler.

\smallskip
\textit{Pas 1: $pq$ este număr ciclic $\iff p\nmid q-1$.} Avem
$\varphi(pq)=(p-1)(q-1)$. Cum $p,q$ sunt prime, $q>p>p-1$ și $q>q-1$,
deci $q\nmid(p-1)(q-1)$; iar $p\mid(p-1)(q-1)\iff p\mid q-1$ (fiindcă
$p\nmid p-1$). Așadar
\[
  \gcd\bigl(pq,(p-1)(q-1)\bigr)=1\iff p\nmid q-1 ,
\]
și, prin criteriul numerelor ciclice, $pq$ e număr ciclic $\iff
p\nmid q-1$. (Direcția netrivială ,,$p\nmid q-1\Rightarrow$ ciclic''
se poate obține și direct din Sylow: $n_{q}\equiv 1\pmod q$,
$n_{q}\mid p$, $p<q$ dau $n_{q}=1$; iar $p\nmid q-1$ dă $n_{p}=1$, deci
$G=P\times Q\cong\mathbb{Z}/pq$.)

\smallskip
\textit{Pas 2: $p\nmid q-1$, via postulatul lui Bertrand.} Aplicăm
Bertrand cu $m=p$: există un prim în intervalul $(p,2p)$. Cel mai mic
prim mai mare decât $p$ este chiar $q=p_{n+1}$, deci $q<2p$. Dacă am
avea $p\mid q-1$, atunci $q-1=kp$ cu $k\ge 1$ întreg, iar $kp=q-1<2p$
forțează $k=1$, adică $q=p+1$. Dar două numere prime
consecutive care diferă prin $1$ sunt \emph{doar} $2$ și $3$; cum
$p\ge 3$, contradicție. Așadar $p\nmid q-1$, deci (Pas 1) $pq$ este
număr ciclic: orice grup de ordin $p_{n}p_{n+1}$ este ciclic.
\end{solutie}

\begin{obs}
Ipoteza $n\ge 2$ e necesară: pentru $n=1$, $p_{1}p_{2}=6$ și
$2\mid 3-1$, deci $6$ nu e număr ciclic --- grupul $S_{3}$ de ordin $6$ nu
e ciclic. Rolul postulatului lui Bertrand este exact să garanteze
$q<2p$, ceea ce exclude singurul caz problematic $q=p+1$ (posibil doar la
$2,3$).
\end{obs}

\begin{problema}
Determinați toate numerele naturale $m$ pentru care există un grup
abelian finit având exact $m$ elemente de ordin $4$.
\end{problema}

\noindent\textbf{Răspuns:} $m=0$ sau $m=2^{k}\,(2^{v}-1)$ cu întregi
$1\le v\le k$. Echivalent: scriind $m=2^{k}t$ cu $t$ impar, $m\neq 0$ este
realizabil dacă și numai dacă $t+1$ este o putere $2^{v}$ a lui
$2$, cu $v\le k$.

\begin{solutie}
Folosim clasificarea grupurilor abeliene finite (Corolarul~\ref{cor:abelianfinit} din cartea
de teorie, consecință a Teoremei~\ref{teo:structura} de structură a grupurilor
abeliene finit generate): \emph{orice grup abelian finit este produs direct
de grupuri ciclice}; combinând cu lema chineză a resturilor
(dacă $\gcd(m_1,m_2)=1$, atunci
$\mathbb{Z}/m_1m_2\cong\mathbb{Z}/m_1\times\mathbb{Z}/m_2$), este
produs direct de grupuri ciclice de ordin \emph{putere de prim}
(descompunerea primară). În particular,
\[
  A\cong A_{2}\times A_{\text{impar}},\qquad
  A_{2}\cong \mathbb{Z}/2^{a_{1}}\times\cdots\times\mathbb{Z}/2^{a_{k}}
  \quad(a_{i}\ge 1),
\]
unde $A_{2}$ adună factorii de ordin putere a lui $2$, iar
$A_{\text{impar}}$ pe cei de ordin impar.

\smallskip
\textit{Pas 1: contează doar partea 2-primară.} Un element
$(x,y)\in A_{2}\times A_{\text{impar}}$ are ordinul
$\operatorname{lcm}\bigl(\operatorname{ord}x,\operatorname{ord}y\bigr)$;
acesta este $4$ dacă și numai dacă $\operatorname{ord}(y)=1$
(ordinul lui $y$ e impar și divide $4$) și $\operatorname{ord}(x)=4$.
Așadar numărul căutat se calculează în $A_{2}$.

\smallskip
\textit{Pas 2: numărarea în $A_{2}$.} În factorul
$\mathbb{Z}/2^{a_{i}}$, soluțiile ecuației $2^{j}x=0$ sunt în
număr de $\min(2^{j},2^{a_{i}})$. Notând
$v=\#\{i:a_{i}\ge 2\}$, obținem
\[
  \bigl|A[2]\bigr|=\prod_{i=1}^{k}\min(2,2^{a_{i}})=2^{k},\qquad
  \bigl|A[4]\bigr|=\prod_{i=1}^{k}\min(4,2^{a_{i}})=2^{k-v}\,4^{v}=2^{k+v},
\]
deci numărul elementelor de ordin \emph{exact} $4$ este
\[
  \bigl|A[4]\bigr|-\bigl|A[2]\bigr|=2^{k+v}-2^{k}=2^{k}\,(2^{v}-1),
  \qquad 0\le v\le k
\]
(valoarea $v=0$ dă $0$).

\smallskip
\textit{Pas 3: realizabilitate și unicitatea formei.} Pentru orice
$1\le v\le k$, grupul $(\mathbb{Z}/4)^{v}\times(\mathbb{Z}/2)^{k-v}$ atinge
exact $2^{k}(2^{v}-1)$ elemente de ordin $4$; iar $m=0$ se atinge prin orice
grup de ordin impar. Reciproc, dintr-un $m\neq 0$ realizabil, partea impară
$t=2^{v}-1$ determină $v$, iar puterea lui $2$ determină $k$; deci
$m=2^{k}t$ este realizabil exact când $t+1$ e putere a lui $2$, $t+1=2^{v}$,
cu $v\le k$.
\end{solutie}

\begin{obs}
Trei remarci. Întâi, valoarea $m=6$ este \emph{imposibilă} în
abelian ($t=3$ dă $v=2$, dar $k=1<v$) --- însă grupul
cuaternionilor $Q_{8}$ are \emph{exact} $6$ elemente de ordin $4$
($\pm i,\pm j,\pm k$): comutativitatea este esențială. Apoi,
$m=2026$ este imposibil ($2026=2\cdot 1013$, iar $1014$ nu e putere a lui
$2$). În fine, contrastul cu ordinul $2$: numărul elementelor de
ordin $2$ este $2^{k}-1$ --- un fapt de spațiu vectorial ($A[2]$ peste
$\mathbb{F}_{2}$), care nu cere clasificarea; ordinul $4$ este primul care
distinge factorii $\mathbb{Z}/2$ de cei $\mathbb{Z}/2^{a}$, $a\ge 2$, deci
primul invariant de numărare care are nevoie de descompunerea ciclică.
(Verificare de consistență: numărul e mereu par când e nenul,
căci elementele de ordin $4$ vin în perechi $\{x,x^{-1}\}$ cu
$x\neq x^{-1}$.)
\end{obs}

\begin{problema}
Fie $G$ un grup abelian finit. Arătați că există un grup $H$
cu $G\cong H\times H$ dacă și numai dacă, pentru orice $n\ge 1$,
numărul soluțiilor ecuației $x^{n}=e$ în $G$ este pătrat
perfect.
\end{problema}

\begin{solutie}
Notăm $G[n]=\{x\in G:x^{n}=e\}$ (mulțimea soluțiilor lui
$x^{n}=e$; e subgrup, $G$ fiind
abelian).

\smallskip
\textit{($\Rightarrow$).} Dacă $G=H\times H$, atunci $(x,y)^{n}=e\iff
x^{n}=e$ și $y^{n}=e$, deci $G[n]=H[n]\times H[n]$ și
$|G[n]|=|H[n]|^{2}$ --- pătrat perfect pentru orice $n$.

\smallskip
\textit{($\Leftarrow$).} Presupunem $|G[n]|$ pătrat perfect pentru orice
$n$. Folosim clasificarea grupurilor abeliene finite (Corolarul~\ref{cor:abelianfinit}, cu
Teorema~\ref{teo:structura} în spate): \emph{orice grup abelian finit este produs
direct de grupuri ciclice}, iar împreună cu lema chineză (pentru
$\gcd(m_1,m_2)=1$, $\mathbb{Z}/m_1m_2\cong\mathbb{Z}/m_1\times\mathbb{Z}/m_2$)
--- produs direct de grupuri ciclice de ordin \emph{putere de prim}, cu
familia acestor puteri unic determinată. Grupând factorii după
primul $p$: $G\cong\prod_{p}G_{p}$ (descompunerea primară), iar fiecare
componentă se scrie
\[
  G_{p}\cong\prod_{i\ge 1}\bigl(\mathbb{Z}/p^{i}\bigr)^{c_{i}},
\]
cu multiplicitățile $c_{i}\ge 0$ \emph{bine definite} --- aici este
esențială partea de \emph{unicitate} a teoremei de structură,
fără de care $c_{i}$ nici n-ar fi invarianți ai lui $G$.

\textit{Pas 1: reducere la componente.} Pentru $n=p^{j}$, componentele
$G_{q}$ cu $q\neq p$ contribuie trivial la $G[p^{j}]$ (un element de ordin
putere a lui $q$ cu $x^{p^{j}}=e$ e neutru), deci
$|G[p^{j}]|=|G_{p}[p^{j}]|$: fiecare $|G_{p}[p^{j}]|$ e pătrat perfect.

\textit{Pas 2: numărarea.} În $\mathbb{Z}/p^{i}$, ecuația
$p^{j}x=0$ are $\min(p^{j},p^{i})=p^{\min(i,j)}$ soluții. Prin urmare
\[
  |G_{p}[p^{j}]|=p^{e_{j}},\qquad e_{j}=\sum_{i\ge 1}\min(i,j)\,c_{i},
\]
iar ipoteza spune că toți exponenții $e_{j}$ sunt \emph{pari}.

\textit{Pas 3: extragerea parității multiplicităților.} Fie
$T_{j}=\sum_{i\ge j}c_{i}$. Deoarece $\min(i,j)-\min(i,j-1)=1$ exact când
$i\ge j$ (și $0$ altfel),
\[
  e_{1}=T_{1},\qquad e_{j}-e_{j-1}=T_{j}\quad(j\ge 2).
\]
Toți $e_{j}$ pari $\Rightarrow$ toți $T_{j}$ pari $\Rightarrow$
fiecare $c_{i}=T_{i}-T_{i+1}$ este \emph{par}.

\textit{Pas 4: construcția rădăcinii.} Scriem $c_{i}=2b_{i}$
(pentru fiecare prim $p$) și punem
\[
  H=\prod_{p}\prod_{i\ge 1}\bigl(\mathbb{Z}/p^{i}\bigr)^{b_{i}} .
\]
Atunci $H\times H$ are, în fiecare componentă primară, exact
multiplicitățile $2b_{i}=c_{i}$ ale lui $G$; prin unicitatea
descompunerii, $G\cong H\times H$.
\end{solutie}

\begin{obs}
Comutativitatea este esențială. Pentru $p$ impar, fie
$G=\operatorname{Heis}(p)\times\mathbb{Z}/p$, unde $\operatorname{Heis}(p)$
este grupul Heisenberg modulo $p$ (matricele $3\times 3$ superior
unitriunghiulare peste $\mathbb{Z}/p$) --- neabelian, de ordin $p^{3}$ și
\emph{exponent} $p$. Atunci $|G|=p^{4}$ și orice element satisface
$x^{p}=e$, deci numărul soluțiilor lui $x^{n}=e$ este $p^{4}$ dacă
$p\mid n$ și $1$ altfel --- \emph{toate pătrate perfecte}. Totuși
$G$ nu este pătratul niciunui grup: dacă $K\times K\cong G$, atunci
$|K|=p^{2}$, deci $K$ este abelian. (Într-adevăr, orice grup $K$ de
ordin $p^{2}$ e abelian: prin ecuația claselor un $p$-grup are centru
netrivial, deci $|Z(K)|\in\{p,p^{2}\}$; dacă $|Z(K)|=p$, atunci $K/Z(K)$
ar fi ciclic de ordin $p$, iar dacă $K/Z(K)$ e ciclic atunci $K$ e
abelian --- contradicție cu $|Z(K)|=p$; deci $|Z(K)|=p^{2}$, adică
$K$ abelian.) Prin urmare $K\times K$ ar fi abelian --- contradicție.
Echivalența este
așadar un fenomen strict abelian. De remarcat și că nu ajunge un
singur $n$: pentru $n=|G|$ numărul soluțiilor e $|G|$, pătrat
pentru orice grup de ordin pătrat perfect; e nevoie de toate valorile
$n=p^{j}$ cu $p^{j}\mid\exp(G)$.
\end{obs}

\begin{problema}
Arătați că orice inel comutativ $R$ este izomorf cu centrul unui
inel --- mai precis, există un inel $S$ (necomutativ dacă $R\neq\{0\}$)
al cărui centru $Z(S)=\{s\in S : sx=xs\ \forall x\in S\}$ este izomorf cu
$R$.
\end{problema}

\begin{solutie}
Luăm $S=M_{2}(R)$, inelul matricelor $2\times 2$ cu intrări în
$R$ (adunare și înmulțire de matrice uzuale). Vom arăta
că centrul său este format din matricele scalare, deci
$Z(S)\cong R$.

\smallskip
\textit{Matricele scalare sunt centrale.} Pentru $a\in R$, matricea
$aI=\begin{psmallmatrix}a&0\\0&a\end{psmallmatrix}$ comută cu orice
$N\in M_{2}(R)$: cum $a$ e central în $R$ (chiar comutativ),
$(aI)N=aN=N(aI)$.

\smallskip
\textit{Doar ele sunt centrale.} Fie
$M=\begin{psmallmatrix}a&b\\c&d\end{psmallmatrix}\in Z(S)$. E de ajuns ca
$M$ să comute cu matricele elementare $E_{ij}$ (cu $1$ pe poziția
$(i,j)$ și $0$ în rest), fiindcă acestea generează aditiv tot
$M_{2}(R)$. Din $ME_{12}=E_{12}M$:
\[
  ME_{12}=\begin{psmallmatrix}0&a\\0&c\end{psmallmatrix},\qquad
  E_{12}M=\begin{psmallmatrix}c&d\\0&0\end{psmallmatrix},
\]
deci, egalând intrările, $c=0$ și $a=d$. Analog, din
$ME_{21}=E_{21}M$:
\[
  ME_{21}=\begin{psmallmatrix}b&0\\d&0\end{psmallmatrix},\qquad
  E_{21}M=\begin{psmallmatrix}0&0\\a&b\end{psmallmatrix},
\]
de unde $b=0$. Așadar
$M=\begin{psmallmatrix}a&0\\0&a\end{psmallmatrix}=aI$. Prin urmare
\[
  Z(M_{2}(R))=\{aI : a\in R\}.
\]

\smallskip
\textit{Izomorfismul cu $R$.} Aplicația $\iota:R\to Z(M_{2}(R))$,
$\iota(a)=aI$, este bijectivă și morfism de inele:
$\iota(a+b)=(a+b)I=aI+bI$, $\iota(1)=I$, iar
$\iota(a)\iota(b)=(aI)(bI)=(ab)I=\iota(ab)$ --- unde s-a folosit
comutativitatea lui $R$, ca produsul celor două matrice scalare
să fie tot scalar, egal cu $(ab)I$. Deci $Z(M_{2}(R))\cong R$.

Pentru $R\neq\{0\}$, $M_{2}(R)$ este necomutativ (de pildă
$E_{12}E_{21}=E_{11}\neq E_{22}=E_{21}E_{12}$), deci $R$ apare drept centrul
unui inel \emph{strict} mai mare.
\end{solutie}

\begin{obs}
Construcția arată că centrul nu determină inelul: orice inel
comutativ $R$ este centrul lui $M_{2}(R)$ (și, la fel, al lui
$M_{n}(R)$ pentru orice $n\ge 2$), deși $|M_{2}(R)|=|R|^{4}$ când $R$
e finit. Este ,,umflătorul de centru'' din spatele faptului că
$Z(A)\cong Z(B)$ nu implică $A\cong B$: $R$ și $M_{2}(R)$ au același
centru, dar sunt foarte diferite. Singurul caz degenerat este $R=\{0\}$, unde
$M_{2}(R)=\{0\}$ e tot nul.
\end{obs}

\begin{problema}
Fie $D$ un inel de diviziune (orice element nenul este inversabil) de
caracteristică $n$, fie $a,c$ numere întregi și $b,d\ge 1$
exponenți, astfel încât, pentru orice $x\in D$,
\[
  a\,x^{b}\in Z(D)\qquad\text{și}\qquad c\,x^{d}\in Z(D),
\]
unde $m\,y$ notează multiplul aditiv $\underbrace{y+\cdots+y}_{m}$.
Demonstrați că dacă
\[
  \gcd(a,n)=\gcd(c,n)=1\quad\text{și}\quad \gcd(b,d)=1,
\]
atunci $D$ este comutativ.
\end{problema}

\begin{solutie}
Reamintim că centrul $Z(D)=\{z : zy=yz\ \forall y\in D\}$ este
\emph{subinel}: e închis la sume, la opuse (deci e subgrup aditiv) și
la produse. Folosim de două ori \emph{relația lui Bézout}: pentru
întregi $p,q$ există întregi $u,v$ cu $up+vq=\gcd(p,q)$.

\begin{lemat}[Bézout aditiv: elimină coeficienții]
În condițiile problemei, $x^{b}\in Z(D)$ și $x^{d}\in Z(D)$
pentru orice $x\in D$.
\end{lemat}

\begin{proof}
Cum $\gcd(a,n)=1$, Bézout dă întregi $u,v$ cu $ua+vn=1$. Pe de
altă parte, caracteristica fiind $n$, avem $n\,y=0$ pentru orice $y\in D$,
deci pentru orice $y$
\[
  y=1\cdot y=(ua+vn)\,y=u\,(a\,y)+v\,(n\,y)=u\,(a\,y).
\]
Aplicăm aceasta lui $y=x^{b}$: elementul $a\,x^{b}$ este central prin
ipoteză, iar $Z(D)$ e subgrup aditiv, deci și multiplul întreg
$u\,(a\,x^{b})$ este central. Dar tocmai am văzut că
$u\,(a\,x^{b})=x^{b}$. Așadar $x^{b}\in Z(D)$. Identic, folosind
$\gcd(c,n)=1$, se obține $x^{d}\in Z(D)$.
\end{proof}

\begin{lemat}[Bézout multiplicativ: elimină exponenții]
Dacă $x^{b}\in Z(D)$ și $x^{d}\in Z(D)$ pentru orice $x$, iar
$\gcd(b,d)=1$, atunci orice element inversabil al lui $D$ este central.
\end{lemat}

\begin{proof}
Mai întâi, \emph{inversul unui element central inversabil este
central}: dacă $z\in Z(D)$ e inversabil și $y\in D$, atunci din
$zy=yz$, înmulțind la stânga și la dreapta cu $z^{-1}$,
obținem $yz^{-1}=z^{-1}y$; deci $z^{-1}\in Z(D)$. Prin urmare, pentru
$z$ central inversabil, \emph{toate} puterile întregi $z^{k}$
($k\in\mathbb{Z}$) sunt centrale.

Fie acum $x$ inversabil. Din $\gcd(b,d)=1$, Bézout dă întregi
$s,t$ cu $sb+td=1$. Elementele $x^{b}$ și $x^{d}$ sunt centrale (ipoteză)
și inversabile ($x$ fiind inversabil), deci $(x^{b})^{s}$ și
$(x^{d})^{t}$ sunt centrale; cum $Z(D)$ e închis la produs,
\[
  x=x^{1}=x^{sb+td}=(x^{b})^{s}\,(x^{d})^{t}\in Z(D). \qedhere
\]
\end{proof}

\smallskip
\textit{Combinarea celor două leme.} Prin lema aditivă, ipoteza cu
coeficienți se curăță: $x^{b},x^{d}\in Z(D)$ pentru orice $x$.
Prin lema multiplicativă, orice element inversabil este central. Dar $D$
este inel de diviziune, deci \emph{orice} element nenul este inversabil;
așadar $D\setminus\{0\}\subseteq Z(D)$. Cum și $0\in Z(D)$, rezultă
$Z(D)=D$, adică $xy=yx$ pentru orice $x,y$: inelul $D$ este comutativ.
\end{solutie}

\begin{obs}
Ipotezele de coprimalitate sunt cele care fac demonstrația
\emph{elementară}; ele nu sunt însă toate necesare.

Condiția $\gcd(b,d)=1$ (exponenți) nu este necesară: din lema aditivă,
$x^{b}\in Z(D)$ pentru orice $x$, iar teorema lui Kaplansky
(Teorema~\ref{teo:kaplansky}(i) din partea de teorie: dacă pentru orice
$x$ o putere $x^{n(x)}$, $n(x)\ge1$, este centrală, atunci $D$ este
comutativ) dă deja concluzia. Coprimalitatea exponenților înlocuiește
această teoremă adâncă printr-un simplu calcul de tip Bézout. (Atenție:
în corpul cuaternionilor $\mathbb{H}$, cu $Z(\mathbb{H})=\mathbb{R}$,
pătratul unui cuaternion $x=t+v$, cu $t$ real și $v$ pur, este
$x^{2}=t^{2}-|v|^{2}+2tv$, care \emph{nu} este real pentru $t\neq0$ și
$v\neq0$; deci $\mathbb{H}$ nu satisface ipoteza pentru $b=2$.)

Din același motiv ajunge ca \emph{una} dintre condițiile
$\gcd(a,n)=1$, $\gcd(c,n)=1$ să aibă loc. Nu pot lipsi însă amândouă:
pentru $a=c=0$ ipoteza devine vidă ($0\cdot x^{b}=0\in Z(D)$), iar
corpul cuaternionilor (de caracteristică $n=0$, deci cu
$\gcd(0,0)=0\neq1$) nu este comutativ.
\end{obs}

\begin{problema}
Fie $G$ un grup cu proprietatea că, pentru oricare trei elemente
distincte ale sale, există cel puțin o pereche dintre ele care
comută. Demonstrați că $G$ este comutativ.
\end{problema}

\begin{solutie}
Presupunem prin absurd că $G$ nu e comutativ: există $x,y\in G$ cu
$xy\neq yx$. Considerăm tripletul
\[
  \{\,x,\ y,\ xy\,\}
\]
și arătăm că el contrazice ipoteza.

\smallskip
\textit{Pas 1: cele trei elemente sunt distincte.}
\begin{itemize}[nosep,leftmargin=1.4em]
  \item $x\neq y$: dacă $x=y$, atunci $xy=x\cdot x=yx$, contrazicând
        $xy\neq yx$.
  \item $xy\neq x$: dacă $xy=x$, înmulțind la stânga cu
        $x^{-1}$ obținem $y=e$; dar $e$ comută cu orice element, deci
        $xy=yx$ --- contradicție.
  \item $xy\neq y$: dacă $xy=y$, înmulțind la dreapta cu $y^{-1}$
        obținem $x=e$, și la fel $xy=yx$ --- contradicție.
\end{itemize}
Deci ipoteza problemei se aplică tripletului $\{x,y,xy\}$.

\smallskip
\textit{Pas 2: niciuna dintre cele trei perechi nu comută.}
\begin{itemize}[nosep,leftmargin=1.4em]
  \item \emph{$x$ și $y$}: nu comută, prin presupunere.
  \item \emph{$x$ și $xy$}: dacă ar comuta, atunci $x(xy)=(xy)x$,
        adică $x^{2}y=xyx$; înmulțind la stânga cu $x^{-1}$
        rezultă $xy=yx$ --- contradicție. Deci nu comută.
  \item \emph{$y$ și $xy$}: dacă ar comuta, atunci $y(xy)=(xy)y$,
        adică $yxy=xy^{2}$; înmulțind la dreapta cu $y^{-1}$
        rezultă $yx=xy$ --- contradicție. Deci nu comută.
\end{itemize}

\smallskip
\textit{Pas 3: contradicția.} Tripletul $\{x,y,xy\}$ este format din trei
elemente distincte între care \emph{nicio} pereche nu comută, ceea ce
contrazice ipoteza (care cerea cel puțin o pereche comutantă).
Așadar presupunerea e falsă: $xy=yx$ pentru orice $x,y\in G$, deci $G$
este comutativ.
\end{solutie}

\begin{obs}
Soluția directă arată mai mult decât cere enunțul: în
orice grup neabelian există trei elemente distincte între care
\emph{nicio} pereche nu comută (anume $x$, $y$, $xy$ cu $xy\neq yx$).
Altfel spus, ipoteza ,,cel puțin o pereche comută'' este deja
suficientă, deși pare mult mai slabă decât ,,cel puțin
două perechi comută''. Cheia este alegerea celui de-al treilea element:
produsul $xy$ nu comută nici cu $x$, nici cu $y$ --- necomutarea ,,se
propagă'' la produs.

Soluția funcționează și pentru grupuri infinite. O abordare prin
numărare, pentru $G$ finit de ordin $n$, nu ajunge aici: ipoteza spune
doar că ,,graful necomutării'' (vârfurile sunt elementele lui $G$,
muchiile unesc elementele care nu comută) nu are triunghiuri, deci, prin
teorema lui Mantel, are cel mult $n^{2}/4$ muchii. Rezultă cel puțin
$n^{2}/2$ perechi ordonate care comută, ceea ce este compatibil cu
marginea $\tfrac58 n^{2}$ valabilă în orice grup neabelian (problema
ONM.13); nu se obține nicio contradicție. Argumentul algebric direct
este cel care decide.
\end{obs}

\begin{problema}
Fie $A$ un inel finit de ordin impar cu proprietatea că orice element
neinversabil se scrie ca sumă a doi idempotenți care comută.
Arătați că $A$ este comutativ.
\end{problema}

\begin{solutie}
\textit{Pas 1: $2$ este inversabil, cu invers central.} Grupul aditiv
$(A,+)$ are ordin impar, deci aplicația $\varphi(x)=2x=x+x$ este
injectivă: dacă $2x=0$, ordinul aditiv al lui $x$ divide $2$ și,
prin teorema lui Lagrange (ordinul unui element divide ordinul grupului),
divide și $|A|$, care e impar; deci acel ordin este $1$, adică $x=0$.
Fiind injectivă pe o mulțime finită, $\varphi$ e bijectivă,
deci există $u\in A$ cu $2u=1$. Acest $u$ e central: pentru orice $a$,
\[
  2(ua)=(2u)a=a=a(2u)=2(au),
\]
iar din injectivitatea lui $\varphi$ rezultă $ua=au$. Notăm
$u=2^{-1}$.

\smallskip
\textit{Pas 2: orice element neinversabil $x$ satisface $x(x-1)(x-2)=0$.}
Fie $x=e+f$ cu $e^{2}=e$, $f^{2}=f$ și $ef=fe$. Punem $g=ef$; atunci
$g^{2}=(ef)(ef)=e^{2}f^{2}=ef=g$ (am folosit comutarea), iar
$eg=e(ef)=e^{2}f=ef=g$ și, analog, $ge=fg=gf=g$. Calculăm:
\[
  x^{2}=(e+f)^{2}=e^{2}+ef+fe+f^{2}=e+2g+f=x+2g ,
\]
\[
  x^{3}=x\cdot x^{2}=(e+f)(x+2g)=x^{2}+2(eg+fg)=(x+2g)+2(g+g)=x+6g .
\]
Cum $1$ și $2$ sunt centrale, $x$ comută cu $x-1$ și $x-2$, deci
\[
  x(x-1)(x-2)=x^{3}-3x^{2}+2x=(x+6g)-3(x+2g)+2x=0 .
\]

\smallskip
\textit{Pas 3: $A$ este redus} (adică nu are elemente nilpotente nenule).
Fie $t$ nilpotent, $t^{k}=0$ cu $k\ge 1$. Atunci $t$ nu e inversabil (altfel,
înmulțind $t^{k}=0$ cu $(t^{-1})^{k}$, ar rezulta $1=0$), deci Pasul 2
dă $t(t-1)(t-2)=0$. Arătăm că $t-1$ și $t-2$ sunt
inversabile:
\begin{itemize}[nosep,leftmargin=1.4em]
  \item $1-t$ e inversabil, cu inversul $1+t+\cdots+t^{k-1}$ (produsul
        telescopează: $(1-t)(1+t+\cdots+t^{k-1})=1-t^{k}=1$); deci
        $t-1=-(1-t)$ e inversabil;
  \item $t-2=-2\bigl(1-2^{-1}t\bigr)$; elementul $2^{-1}t$ e nilpotent
        ($2^{-1}$ fiind central, $(2^{-1}t)^{k}=(2^{-1})^{k}t^{k}=0$), deci
        $1-2^{-1}t$ e inversabil prin același argument, iar $2$ e
        inversabil prin Pasul 1; produsul lor e inversabil.
\end{itemize}
Elementele $t$, $t-1$, $t-2$ sunt polinoame în $t$ cu coeficienți
centrali, deci comută între ele. Înmulțind relația
$t(t-1)(t-2)=0$ cu $\bigl[(t-1)(t-2)\bigr]^{-1}$ obținem $t=0$. Așadar
$A$ e redus. (Aici s-a folosit \emph{esențial} ipoteza de ordin impar,
prin inversabilitatea lui $2$.)

\smallskip
\textit{Pas 4: orice element satisface o relație $x^{s}=x$ cu $s\ge 2$.}
$A$ fiind finit, șirul puterilor $x,x^{2},x^{3},\dots$ nu poate avea toate
elementele distincte: există $n\ge 1$ și $t\ge 1$ cu
$x^{n+t}=x^{n}$. Cum $A$ e redus (Pasul 3), Propoziția~\ref{prop:redus} din cartea de
teorie (,,inele reduse: coborârea exponentului'': dacă $A$ e redus
și $x^{n+t}=x^{n}$ pentru un $n\ge 1$, atunci $x^{t+1}=x$) dă
\[
  x^{t+1}=x,\qquad\text{cu } t+1\ge 2 .
\]

\smallskip
\textit{Pas 5: concluzia.} Am arătat că pentru orice $x\in A$
există un întreg $n(x)>1$ cu $x^{n(x)}=x$. Prin Teorema~\ref{teo:jacobson}
(Jacobson) din cartea de teorie --- \emph{dacă pentru orice $x\in A$
există $n(x)>1$ cu $x^{n(x)}=x$, atunci $A$ este comutativ} --- rezultă
că $A$ este comutativ.
\end{solutie}

\begin{obs}
Problema conține, ca un caz particular, \emph{teorema lui Wedderburn}
(Teorema~\ref{teo:wedderburn} din carte: orice corp finit este comutativ). Într-adevăr,
într-un corp finit singurul element neinversabil este $0$, iar
$0=0+0$ este sumă a doi idempotenți care comută; așadar orice
corp finit de ordin impar satisface ipoteza \emph{în mod automat}, iar
concluzia problemei revine acolo exact la comutativitatea lui. Prin urmare
nicio soluție nu poate evita o unealtă de forța lui Wedderburn ---
aici, teorema lui Jacobson, care o generalizează.

Ipoteza de ordin impar intervine o singură dată: ca $2$ să fie
inversabil, deci ca $t-2$ să fie inversabil în Pasul 3. Fără
ea, relația $x(x-1)(x-2)=0$ nu mai forțează anularea
nilpotenților.
\end{obs}


\begin{problema}
Fie $A$ un inel cu proprietatea că $x^{4}=x$ pentru orice $x\in A$.
Demonstrați că $A$ este comutativ.
\end{problema}

\begin{solutie}
\textit{Pasul 1: $x+x=0$ pentru orice $x$.} Aplicând ipoteza lui $-1$:
$-1=(-1)^{4}=1$, deci $1+1=0$ și, înmulțind cu $x$, $x+x=0$ pentru
orice $x\in A$. În particular $-y=y$ pentru orice $y$.

\smallskip
\textit{Pasul 2: $A$ este redus.} Dacă $x^{2}=0$, atunci
$x=x^{4}=x^{2}\cdot x^{2}=0$. Așadar singurul element cu pătratul nul
este $0$ (și, prin inducție, singurul nilpotent este $0$).

\smallskip
\textit{Pasul 3: idempotenții sunt centrali.} Fie $e^{2}=e$ și
$x\in A$. Elementul $u=ex-exe$ satisface
\[
  u^{2}=(ex-exe)(ex-exe)=exex-exexe-exeex+exeexe=exex-exexe-exex+exexe=0
\]
(am folosit $e^{2}=e$), deci $u=0$ prin Pasul 2, adică $ex=exe$. Analog,
$v=xe-exe$ satisface $v^{2}=xexe-xexe-exexe+exexe=0$, deci $xe=exe$.
Prin urmare $ex=xe$: orice idempotent este central (este Lema~\ref{teo:machale}
din partea de teorie).

\smallskip
\textit{Pasul 4: $x^{2}+x$ este idempotent, deci central.} Puterile lui
$x$ comută între ele, deci
\[
  (x^{2}+x)^{2}=x^{4}+x^{3}+x^{3}+x^{2}=x^{4}+2x^{3}+x^{2}=x+x^{2},
\]
unde am folosit $x^{4}=x$ și $2x^{3}=0$ (Pasul 1). Din Pasul 3,
\[
  x^{2}+x\in Z(A)\qquad\text{pentru orice }x\in A.
\]

\smallskip
\textit{Pasul 5: mecanismul lui MacHale.} Aplicând Pasul 4 lui $x$, $y$
și $x+y$ (dezvoltarea lui $(x+y)^{2}$ păstrează ordinea factorilor):
\[
  (x+y)^{2}+(x+y)-(x^{2}+x)-(y^{2}+y)=xy+yx\in Z(A).
\]
Elementul central $xy+yx$ comută cu $x$:
$x^{2}y+xyx=xyx+yx^{2}$, deci $x^{2}y=yx^{2}$ pentru orice $y$, adică
$x^{2}\in Z(A)$. În fine,
\[
  x=(x^{2}+x)-x^{2}\in Z(A)
\]
pentru orice $x\in A$, deci $A$ este comutativ.
\end{solutie}

\begin{obs}
Pasul-cheie este Pasul 4: identitatea $x^{4}=x$ face ca $x^{2}+x$ să fie
idempotent, iar într-un inel redus idempotenții sunt centrali; de aici
încolo lucrează mecanismul din teorema lui MacHale
(Teorema~\ref{teo:machalecomut}), exact ca la problema OJM.22.

Problema este un caz particular al teoremei lui Jacobson
(Teorema~\ref{teo:jacobson}), cu $n(x)=4$ pentru toți $x$; aici însă
demonstrația este complet elementară. Exemple de inele cu $x^{4}=x$:
$\mathbb{Z}/2$, corpul cu $4$ elemente $\mathbb{F}_{4}$ (în care
$x^{3}=1$ pentru $x\neq0$, prin Lagrange în grupul $\mathbb{F}_{4}^{*}$ de
ordin $3$) și produsele directe ale acestora. Un inel necomutativ precum
$M_{2}(\mathbb{Z}/2)$ nu satisface ipoteza: pentru
$x=\begin{psmallmatrix}0&1\\0&0\end{psmallmatrix}$ avem $x^{2}=0$, deci
$x^{4}=0\neq x$.

Pasul 1 folosește esențial faptul că exponentul $4$ este par (el dă
$(-1)^{4}=1$). Pentru exponentul impar $3$, inelul $\mathbb{Z}/3$
satisface $x^{3}=x$ fără ca $1+1$ să fie $0$; concluzia de
comutativitate rămâne adevărată, dar se obține prin alt calcul (vezi demonstrația de după
Teorema~\ref{teo:jacobson} din partea de teorie).
\end{obs}

\begin{problema}
Fie $K$ un corp al cărui grup multiplicativ $(K^*,\cdot)$ este ciclic. Arătați că $K$ este finit.
\end{problema}

\begin{solutie}
Fie $K^* = \langle g \rangle$. Dacă grupul $\langle g \rangle$ este finit, atunci $K = K^* \cup \{0\}$ este finit și am terminat. Presupunem deci, prin reducere la absurd, că $K^*$ este ciclic și \emph{infinit}, și căutăm o contradicție.

\emph{Pasul 1: caracteristica este} $2$. Într-un grup ciclic infinit, singurul element de ordin finit este elementul neutru (puterile unui generator sunt toate distincte). Dar $(-1)^2 = 1$, deci $-1$ are ordinul cel mult $2$: rezultă că $-1 = 1$, adică $1 + 1 = 0$.

\emph{Pasul 2: o relație polinomială pentru $g$.} Considerăm elementul $1+g$. Dacă $1+g = 0$, atunci $g = 1$ (caracteristică $2$), deci $K^* = \{1\}$ ar fi finit: exclus. Așadar $1 + g \in K^* = \langle g \rangle$, adică
\[
  1 + g = g^m \qquad \text{pentru un anumit } m \in \mathbb{Z}.
\]
Cazurile $m = 0$ și $m = 1$ dau $g = 0$, respectiv $1 = 0$: imposibil. Dacă $m \geq 2$, relația se rescrie (în caracteristică $2$) sub forma
\[
  g^m + g + 1 = 0;
\]
dacă $m \leq -1$, înmulțim cu $g^{-m}$ și obținem $g^{-m} + g^{1-m} = 1$, adică
\[
  g^{\,|m|+1} + g^{\,|m|} + 1 = 0.
\]
În ambele cazuri, $g$ este rădăcină a unui polinom \emph{monic nenul} cu coeficienți în subcorpul prim $\mathbb{Z}_2$; fie $d \geq 2$ gradul său.

\emph{Pasul 3: puterile lui $g$ încap într-o mulțime finită.} Fie
\[
  E = \bigl\{ a_0 + a_1 g + \cdots + a_{d-1} g^{\,d-1} :\ a_i \in \mathbb{Z}_2 \bigr\},
\]
o mulțime cu cel mult $2^d$ elemente. Relația de la Pasul 2 exprimă $g^d$ ca o combinație a puterilor mai mici, deci $E$ este stabilă la înmulțirea cu $g$: prin inducție, \emph{toate} puterile $g^k$ cu $k \geq 0$ aparțin lui $E$.

Dar în grupul ciclic infinit $\langle g \rangle$ puterile $g^0, g^1, g^2, \ldots$ sunt distincte două câte două: o infinitate de elemente distincte într-o mulțime cu cel mult $2^d$ elemente. Contradicția încheie demonstrația: $K^*$ nu poate fi ciclic infinit, deci $K$ este finit.

\medskip
\noindent\textbf{Observație} Rezultatul se citește ca o reciprocă: se știe că grupul multiplicativ al oricărui corp finit este ciclic; problema arată că ciclicitatea lui $K^*$ \emph{caracterizează} corpurile finite. De remarcat economia de mijloace: singurele ingrediente sunt ,,ordinul finit într-un grup ciclic infinit'' și un argument de numărare.
\end{solutie}

\begin{problema}[Clasică]
Căsuțele unui caroiaj $N \times N$ se colorează cu $n \ge 1$ culori (nu neapărat distincte). Două colorări sunt considerate echivalente dacă una se obține din cealaltă printr-o rotație a caroiajului.
\begin{parti}
\item Determinați numărul colorărilor neechivalente.
\item Deduceți că numărătorul expresiei obținute este divizibil cu 4 pentru orice $n \ge 1$ și orice $N \ge 1$.
\end{parti}
\end{problema}

\begin{solutie}
Fie $X$ mulțimea celor $n^{N^2}$ colorări ale celor $N^2$ căsuțe și fie $G = \{\mathrm{id}, \rho, \rho^2, \rho^3\}$ grupul rotațiilor caroiajului ($\rho$ = rotația de $90^\circ$), care acționează pe $X$ permutând căsuțele. Colorările neechivalente sunt exact \emph{orbitele} acestei acțiuni.

\emph{Lema lui Burnside.} Numărul orbitelor unei acțiuni a unui grup finit $G$ pe o mulțime finită $X$ este
\[
\#\text{orbite} = \frac{1}{|G|} \sum_{\sigma \in G} \bigl|\operatorname{Fix}(\sigma)\bigr|,
\qquad
\operatorname{Fix}(\sigma) = \{x \in X :\ \sigma \cdot x = x\}.
\]

\emph{Demonstrație.} Numărăm în două moduri perechile $(\sigma, x)$ cu $\sigma \cdot x = x$. Pe de o parte, numărul lor este $\sum_{\sigma} |\operatorname{Fix}(\sigma)|$. Pe de altă parte, el este $\sum_{x \in X} |\operatorname{Stab}(x)|$, unde $\operatorname{Stab}(x)$ este stabilizatorul lui $x$. Prin teorema orbită-stabilizator (aplicația $\sigma \mapsto \sigma \cdot x$ induce o bijecție între clasele la stânga din $G$ după $\operatorname{Stab}(x)$ și orbita lui $x$), avem $|\operatorname{Stab}(x)| = |G|/|\operatorname{Orb}(x)|$. Grupând elementele $x$ pe orbite, fiecare orbită $\mathcal{O}$ contribuie cu $\sum_{x \in \mathcal{O}} |G|/|\mathcal{O}| = |G|$. Suma totală este deci $|G| \cdot \#\text{orbite}$, de unde formula.

\emph{(a) Numărăm ciclurile fiecărei rotații pe cele $N^2$ căsuțe.} O colorare este fixată de $\sigma$ dacă și numai dacă este constantă pe fiecare ciclu al lui $\sigma$, deci $|\operatorname{Fix}(\sigma)| = n^{\,c(\sigma)}$, unde $c(\sigma)$ este numărul de cicluri.

\begin{itemize}
\item[---] $\mathrm{id}$: fiecare căsuță formează singură un ciclu, $c = N^2$.
\item[---] $\rho^2$ ($180^\circ$): căsuțele se grupează în perechi $\{x, \rho^2 x\}$, iar dacă $N$ este impar, căsuța centrală este fixă; $c = \frac{N^2}{2}$ ($N$ par), respectiv $\frac{N^2+1}{2}$ ($N$ impar).
\item[---] $\rho, \rho^3$ ($90^\circ$): orbite de câte 4 căsuțe, plus centrul fix când $N$ este impar; $c = \frac{N^2}{4}$ ($N$ par), respectiv $\frac{N^2-1}{4} + 1 = \frac{N^2+3}{4}$ ($N$ impar).
\end{itemize}

Conform lemei lui Burnside, numărul colorărilor neechivalente este
\[
\begin{cases}
\dfrac{n^{N^2} + n^{N^2/2} + 2\,n^{N^2/4}}{4}, & N \text{ par},\\[3ex]
\dfrac{n^{N^2} + n^{(N^2+1)/2} + 2\,n^{(N^2+3)/4}}{4}, & N \text{ impar}.
\end{cases}
\]

\emph{(b)} Membrul stâng numără orbite, deci este un număr întreg: numărătorul de mai sus este divizibil cu 4 pentru orice $n \ge 1$ și orice $N \ge 1$.

\medskip
\noindent\textbf{Observație.} Pentru caroiajul $2 \times 2$ ($N = 2$, par) formula dă $\frac{n^4+n^2+2n}{4}$, iar pentru $n = 2$ dă $\frac{16+4+4}{4} = 6$ colorări neechivalente, lucru ușor de confirmat prin enumerare directă. Poanta punctului (b) merită savurată: un enunț de \emph{teoria numerelor} este demonstrat printr-o \emph{numărare de orbite}, fără nicio congruență. Pe de altă parte, se poate, desigur, proceda și pe cale aritmetică; dar pentru enunțuri mai complicate (coliere cu multe mărgele), calea combinatorie rămâne singura rezonabilă.
\end{solutie}

\begin{problema}
Fie $p$ un număr prim, $b \ge 2$ o bază cu $\gcd(b,p) = 1$ și $f = \ord_p(b)$.
Scrierea lui $\frac{1}{p}$ în baza $b$ este periodică simplă, cu perioada de lungime
exact $f$; fie $N = \frac{b^f-1}{p}$ numărul format din cele $f$ cifre ale unei
perioade (pentru $b = 10$, $p = 7$: $N = 142857$).
\begin{parti}
  \item Arătați că, pentru $1 \le k \le p-1$, numărul $kN$ (scris cu $f$ cifre
    în baza $b$) este o permutare circulară a lui $N$ dacă și numai dacă
    $k \in \langle b \rangle = \{1, b, \dots, b^{f-1}\}$ în $(\ZZ_p^*, \cdot)$;
    deduceți că numerele $N, 2N, \dots, (p-1)N$ se împart în exact $(p-1)/f$
    familii, fiecare formată din cele $f$ rotații ale unui ,,număr ciclic''.
  \item (Midy) Dacă $f = 2g$ este par, cele două jumătăți ale unei perioade, de
    câte $g$ cifre fiecare, au suma $b^g - 1$.
\end{parti}
\end{problema}

\begin{solutie}
\textit{Pasul 0: perioada are lungimea $f$ și $N = \frac{b^f-1}{p}$.}
Din $b^f \equiv 1 \pmod p$ ($f = \ord_p(b)$), numărul $N = \frac{b^f-1}{p}$ este
întreg și (în baza $b$)
\[
  \frac{1}{p} = \frac{N}{b^f-1}
  = N\Bigl(\frac{1}{b^f} + \frac{1}{b^{2f}} + \cdots\Bigr)
  = 0,\overline{a_1 a_2 \dots a_f},
\]
unde $a_1 \dots a_f$ este scrierea cu $f$ cifre a lui $N$ (eventual cu
zerouri la început). Perioada nu poate fi mai scurtă: o perioadă $L$ ar da
$b^L \equiv 1 \pmod p$, iar $\ord_p(b) = f$ dă $f \mid L$.

\smallskip
\textit{Pasul 1: înmulțirea cu $b^j$ rotește perioada.} Fie $r_j$
restul împărțirii lui $b^j$ la $p$, adică $b^j = pC_j + r_j$. Împărțind la $p$,
$b^j \cdot \frac{1}{p} = C_j + \frac{r_j}{p}$; membrul stâng este perioada
cu virgula mutată cu $j$ poziții, deci partea sa fracționară este
perioada rotită cu $j$ poziții:
\[
  \frac{r_j}{p} = 0,\overline{a_{j+1} \dots a_f a_1 \dots a_j}.
\]
Cum $0 < r_j N < p \cdot N = b^f - 1$, scrierea cu $f$ cifre a lui $r_j N$ este
exact rotația lui $N$ cu $j$ poziții.

\smallskip
\textit{(a) Pasul 2: $kN$ este o rotație $\iff k \in \langle b \rangle$.}
Resturile $r_j = b^j \bmod p$ ($j = 0, \dots, f-1$) parcurg exact
subgrupul $\langle b \rangle = \{1, b, \dots, b^{f-1}\}$, iar Pasul 1 arată că
fiecare $r_j N$ este o rotație a lui $N$. Reciproc, dacă $kN$ (cu $f$ cifre)
este o rotație a lui $N$, atunci $\frac{k}{p}$ are aceeași perioadă ca $\frac{1}{p}$,
deplasată cu un anumit $j$, deci $k \equiv b^j \pmod p$. Așadar $kN$ este o rotație
a lui $N$ dacă și numai dacă $k \in \langle b \rangle$.

\smallskip
\textit{Pasul 3: familiile.} Subgrupul $\langle b \rangle$ are $f$ elemente, care dau cele $f$
rotații distincte ale lui $N$; cele $(p-1)/f$ clase la stânga ale lui $\langle b \rangle$
în $(\ZZ_p^*, \cdot)$ dau tot atâtea familii, fiecare cu propriul său
,,număr ciclic''. Când $b$ este rădăcină primitivă ($f = p-1$), toți
multiplii $N, \dots, (p-1)N$ sunt rotații ale lui $N$.

\smallskip
\textit{(b) Midy.} Fie $f = 2g$. Elementul $b^g$ are ordinul $2$ (căci
$(b^g)^2 = b^f \equiv 1$, dar $b^g \not\equiv 1$, deoarece $\ord_p(b) = 2g$), deci
$b^g \equiv -1 \pmod p$, adică $r_g = p - 1$. Conform Pasului 1, $(p-1)N$ scris cu $2g$
cifre este rotația lui $N$ cu $g$ poziții: dacă $N = A\,b^g + B$ (jumătatea superioară $A$ și
jumătatea inferioară $B$, de câte $g$ cifre), această rotație este $B\,b^g + A$. Dar
$(p-1)N = pN - N = (b^{2g} - 1) - N$, deci
\[
  B\,b^g + A = (b^{2g}-1) - (A\,b^g + B)
  \implies (A+B)(b^g+1) = b^{2g} - 1 = (b^g-1)(b^g+1),
\]
de unde $A + B = b^g - 1$.

\medskip
\noindent\textbf{Observație} Pentru $p = 7$: $2N = 285714$, $3N = 428571$,
$4N = 571428$, $5N = 714285$, $6N = 857142$, toate rotații ale lui $142857$; în schimb
$7N = 999999 = 10^6 - 1$, în acord cu $N = \frac{10^6-1}{7}$. Numerele
care au proprietatea (a) în forma ei \textit{completă} (când $b$ este rădăcină primitivă,
$f = p-1$) se numesc \textbf{numere ciclice}, iar numerele prime corespunzătoare se numesc
\textbf{numere prime cu perioadă maximă} (în engleză, full reptend primes); în baza $10$:
$p = 7, 17, 19, 23, 29, 47, \dots$ Este o problemă \textit{deschisă} (legată de
conjectura lui Artin) dacă există o infinitate de astfel de numere prime. Punctul (b) este
\textbf{teorema lui Midy} (E. Midy, 1836); pentru $f = 2g$ oarecare se vorbește de
teorema lui Midy extinsă. Toată ,,magia'' lui $142857$ este așadar un enunț despre
un generator al unui grup ciclic.
\end{solutie}

\begin{problema}
Fie $p$ un număr prim, $D \mid p-1$ și $m \ge 1$ un întreg. Arătați că suma puterilor a $m$-a ale elementelor de ordin $D$ din $(\ZZ_p^*,\cdot)$ este congruentă modulo $p$ cu \emph{suma lui Ramanujan}
\[
  \sum_{\ord(x)=D} x^m \equiv c_D(m) = \mu\!\left(\tfrac{D}{g}\right)\frac{\varphi(D)}{\varphi(D/g)},
  \qquad g = \gcd(D,m),
\]
unde $\mu$ este funcția lui Möbius, iar $\varphi$ este indicatorul lui Euler. (În particular, ori de câte ori $\gcd(m,D)=1$, deci și pentru $m=1$, suma este egală cu $\mu(D)$.)
\end{problema}

\begin{solutie}
Pentru fiecare $d \mid p-1$, notăm cu $S_d(m)$ suma puterilor a $m$-a ale elementelor de ordin exact $d$ din $\ZZ_p^*$; vrem să arătăm că $S_D(m) = c_D(m)$ (toate egalitățile sunt în $\ZZ_p$).

\emph{Pasul 1: suma puterilor pe soluțiile ecuației $x^d = 1$.} Grupul $\ZZ_p^*$ este ciclic de ordin $p-1$ (rezultat din capitolul despre corpuri); pentru $d \mid p-1$, soluțiile ecuației $x^d = 1$ formează unicul subgrup ciclic $C_d$ de ordin $d$ (polinomul $X^d - 1$ are cel mult $d$ rădăcini, iar $C_d$ le furnizează pe toate). Atunci
\[
  T(d) := \sum_{x^d = 1} x^m = \sum_{y \in C_d} y^m =
  \begin{cases}
    d, & d \mid m,\\
    0, & d \nmid m,
  \end{cases}
\]
căci, dacă $\zeta$ generează $C_d$, suma devine suma geometrică $\sum_{j=0}^{d-1} \zeta^{jm}$, egală cu $d$ când $\zeta^m = 1$ (adică $d \mid m$) și cu $0$ în caz contrar.

\emph{Pasul 2: clasificarea după ordin.} Orice soluție a ecuației $x^D = 1$ are ca ordin un divizor al lui $D$, și reciproc; grupând după ordin,
\[
  T(D) = \sum_{d \mid D} S_d(m). \tag{1}
\]

\emph{Pasul 3: inversiunea Möbius.} Relația (1) are loc pentru orice divizor $D$ al lui $p-1$, deci și pentru orice divizor al lui $D$. Prin inversiunea Möbius (Propoziția~\ref{prop:mobius}(ii) din partea de teorie, aplicată în grupul abelian $(\ZZ_p,+)$) și Pasul 1,
\[
  S_D(m) = \sum_{d \mid D} \mu\!\left(\tfrac{D}{d}\right) T(d)
  = \sum_{\substack{d \mid D\\ d \mid m}} \mu\!\left(\tfrac{D}{d}\right) d
  = \sum_{d \mid g} d\,\mu\!\left(\tfrac{D}{d}\right),
  \qquad g=\gcd(D,m).
\]

\emph{Pasul 4: forma închisă.} Arătăm identitatea între numere întregi
\[
  \Sigma:=\sum_{d \mid g} d\,\mu\!\left(\tfrac{D}{d}\right)=\mu(k)\,\frac{\varphi(D)}{\varphi(k)},\qquad k=\frac{D}{g}.
\]
Scriem $d=g/e$, cu $e$ parcurgând divizorii lui $g$; atunci $D/d=ke$ și $\Sigma=\sum_{e\mid g}\frac{g}{e}\,\mu(ke)$. Dacă $k$ și $e$ au un factor prim comun $q$, atunci $q^2\mid ke$ și $\mu(ke)=0$; dacă $(k,e)=1$, atunci $\mu(ke)=\mu(k)\mu(e)$ (factorii primi ai lui $ke$ sunt cei ai lui $k$ reuniți disjunct cu cei ai lui $e$). Deci
\[
  \Sigma=\mu(k)\sum_{\substack{e\mid g\\ (e,k)=1}}\frac{g}{e}\,\mu(e)
  =\mu(k)\,g\prod_{\substack{q\mid g\\ q\nmid k}}\Bigl(1-\frac1q\Bigr),
\]
produsul fiind după numerele prime $q$ (dezvoltând produsul obținem exact termenii $\mu(e)/e$, cu $e$ produs de prime distincte care divid $g$ și nu divid $k$; ceilalți $e$ au $\mu(e)=0$). Dacă $\mu(k)=0$, ambii membri sunt $0$. Altfel, din formula lui Euler $\varphi(N)=N\prod_{q\mid N}(1-\frac1q)$ și $D=gk$,
\[
  \frac{\varphi(D)}{\varphi(k)}=\frac{D}{k}\prod_{\substack{q\mid D\\ q\nmid k}}\Bigl(1-\frac1q\Bigr)=g\prod_{\substack{q\mid g\\ q\nmid k}}\Bigl(1-\frac1q\Bigr),
\]
căci un prim care divide $D=gk$ fără să dividă $k$ divide $g$, iar reciproc orice prim care divide $g$ divide $D$. Așadar $\Sigma=\mu(k)\varphi(D)/\varphi(k)=c_D(m)$, deci $S_D(m)=c_D(m)$. Pentru $m = 1$ avem $g = 1$, $k=D$ și $S_D(1) = \mu(D)$. $\blacksquare$

\medskip
\noindent\textbf{Observație}\quad Verificare pentru $p = 7$, $m = 1$: elementele de ordin $3$ sunt $2$ și $4$, cu suma $6 \equiv -1 = \mu(3)$; cele de ordin $6$ sunt $3$ și $5$, cu suma $8 \equiv 1 = \mu(6)$. $\checkmark$\quad Cazul $m = 1$ este un rezultat clasic al lui \textbf{Gauss}: suma rădăcinilor primitive de ordin $D$ ale unității este $\mu(D)$ (pentru $D = p-1$: suma rădăcinilor primitive modulo $p$). Pentru $m$ oarecare se obțin \textbf{sumele lui Ramanujan} $c_D(m)$, obiect central al analizei Fourier a funcțiilor aritmetice; aici ele sunt transpuse în corpul $\ZZ_p$. Schema: ciclicitatea dă inventarul, suma geometrică adună puterile, iar Möbius separă ordinele.
\end{solutie}

\begin{problema}
Fie $p$ un număr prim impar și $a \in \ZZ$ cu $p \nmid a$. Determinați restul împărțirii numărului
\[
  \prod_{k=1}^{p-1}\bigl(k^2 + a^2\bigr)
\]
la $p$.
\end{problema}

\begin{solutie}
\textbf{Răspuns:} produsul este congruent cu $0$ modulo $p$ dacă $p \equiv 1 \pmod 4$ și cu $4$ modulo $p$ dacă $p \equiv 3 \pmod 4$, \emph{independent} de $a$ (cu $p \nmid a$). Lucrăm în corpul $\ZZ_p$.

\emph{Pasul 0: reducerea la $a = 1$.} Cum $p \nmid a$, elementul $a$ este inversabil în $\ZZ_p$, deci $k \mapsto ak$ permută mulțimea $\{1,\dots,p-1\}$; astfel, prin mica teoremă a lui Fermat ($a^{\,p-1} \equiv 1$),
\[
  \prod_{k=1}^{p-1}(k^2 + a^2)
  = \prod_{k=1}^{p-1}\bigl((ak)^2 + a^2\bigr)
  = a^{\,2(p-1)} \prod_{k=1}^{p-1}(k^2 + 1)
  \equiv \prod_{k=1}^{p-1}(k^2 + 1) \pmod p .
\]
Rămâne deci să calculăm $\prod_{k=1}^{p-1}(k^2+1)$.

\emph{Pasul 1: reducerea la pătrate.} Când $k$ parcurge valorile $1,\dots,p-1$, pătratul $k^2$ parcurge mulțimea $Q$ a pătratelor nenule, fiecare element fiind atins \emph{exact de două ori} ($k$ și $-k$ dau același pătrat, iar $x^2 = y^2$ implică $y = \pm x$). Prin urmare $|Q| = \frac{p-1}{2}$ și
\[
  \prod_{k=1}^{p-1}(k^2+1) = \Bigl( \prod_{r \in Q} (r+1) \Bigr)^{\!2} .
\]

\emph{Pasul 2: pătratele sunt rădăcinile lui $X^{(p-1)/2}-1$.} Pentru $r = s^2 \in Q$, prin mica teoremă a lui Fermat, $r^{(p-1)/2} = s^{\,p-1} = 1$: cele $\frac{p-1}{2}$ elemente ale lui $Q$ sunt rădăcini ale polinomului $X^{(p-1)/2}-1$, care are cel mult $\frac{p-1}{2}$ rădăcini. Așadar $Q$ este exact mulțimea rădăcinilor și
\[
  \prod_{r \in Q} (X - r) = X^{\frac{p-1}{2}} - 1 . \tag{$*$}
\]

\emph{Pasul 3: evaluarea.} Punem $X = -1$ în $(*)$:
\[
  \prod_{r \in Q} (-1 - r) = (-1)^{\frac{p-1}{2}} - 1
  \Longrightarrow
  \prod_{r \in Q} (r + 1) = (-1)^{\frac{p-1}{2}} \Bigl( (-1)^{\frac{p-1}{2}} - 1 \Bigr) .
\]
Dacă $p \equiv 1 \pmod 4$, exponentul $\frac{p-1}{2}$ este par și membrul drept este $1\cdot(1-1) = 0$; produsul cerut este $0^2 = 0$. Dacă $p \equiv 3 \pmod 4$, exponentul este impar și membrul drept este $(-1) \cdot (-1-1) = 2$; produsul cerut este $2^2 = 4$. $\blacksquare$

\emph{Soluție alternativă (numere complexe / Frobenius), direct pentru $a$ oarecare.} Factorizăm $k^2 + a^2 = (k - ai)(k + ai)$, unde $i^2 = -1$. Pornim de la identitatea polinomială $\prod_{k=1}^{p-1}(X-k) = X^{\,p-1}-1$ în $\ZZ_p[X]$ (ambii membri sunt polinoame monice de grad $p-1$, cu aceleași rădăcini $1,\dots,p-1$). Evaluând-o în $X = ai$ și în $X = -ai$ (în $\ZZ_p$ dacă $-a^2$ este pătrat, altfel în corpul $\mathbb{F}_{p^2}$), cum $p-1$ este par și $a^{\,p-1} \equiv 1$ (Fermat), obținem:
\[
  \begin{aligned}
    \prod_{k=1}^{p-1}(k^2+a^2)
    &= \Bigl( \prod_{k=1}^{p-1} (k - ai) \Bigr) \Bigl( \prod_{k=1}^{p-1} (k + ai) \Bigr) \\
    &= \bigl( (ai)^{\,p-1} - 1 \bigr) \bigl( (-ai)^{\,p-1} - 1 \bigr)
     = \bigl( i^{\,p-1} - 1 \bigr) \bigl( (-i)^{\,p-1} - 1 \bigr) .
  \end{aligned}
\]
Factorul $a^{\,p-1}$ s-a simplificat, deci răspunsul nu depinde de $a$. Dacă $p \equiv 1 \pmod 4$, atunci $i \in \ZZ_p$ și $i^{\,p-1} = 1$ (Fermat), deci produsul este $(1-1)(1-1) = 0$. Dacă $p \equiv 3 \pmod 4$, atunci $i \notin \ZZ_p$, iar automorfismul Frobenius $x \mapsto x^p$ al lui $\mathbb{F}_{p^2}$ schimbă între ele cele două rădăcini ale lui $X^2+1$, deci $i^p = -i$; prin urmare $i^{\,p-1} = i^p/i = -1$ și, analog, $(-i)^{\,p-1} = -1$, iar produsul este $(-1-1)(-1-1) = 4$.

\medskip
\noindent\textbf{Observație} Verificare pentru $p = 7$: modulo $7$, factorii $k^2+1$ sunt $2,5,3,3,5,2$, cu produsul $\equiv 4$. $\checkmark$ \quad Anularea în cazul $p \equiv 1 \pmod 4$ are o interpretare directă: identitatea $(*)$, evaluată în $-1$, dă zero exact când $-1$ este pătrat, adică atunci când ecuația $k^2+1 \equiv 0$ are soluții, iar atunci un factor al produsului este nul. Întreaga soluție se sprijină pe o singură identitate: rădăcinile unui polinom monic, numărate toate, îl \emph{determină}, iar după aceea orice evaluare nu mai costă nimic. Este același procedeu ca în $X^p - X = \prod_{k \in \ZZ_p}(X - k)$, din care se deduce, de exemplu, teorema lui Wilson.
\end{solutie}



%% ============================================================
%%  REFERINȚE
%% ============================================================
\clearpage
\thispagestyle{empty}
\addcontentsline{toc}{chapter}{Referințe}
\vspace*{1cm}
{\noindent\color{onm}\rule{\linewidth}{3pt}}\par\vspace{3mm}
{\noindent\fontsize{28}{32}\selectfont\bfseries\color{ink} Referințe}\par\vspace{3mm}
{\noindent\color{onm}\rule{\linewidth}{3pt}}\par\vspace{8mm}

\noindent Volumul de față este autoconținut; demonstrațiile, definițiile și exemplele nu trimit la resurse externe. Lista de mai jos cuprinde lucrările la care autorul s-a raportat în procesul de redactare --- fie pentru a verifica formulări, fie ca surse de inspirație pentru structura prezentării.

\vspace{6mm}
\begin{thebibliography}{9}

\bibitem{exc11vol1}
Marius Perianu, Ioan Balică, Dan Zaharia,
\textit{Matematică de excelență pentru clasa a XI-a, Vol.\ I: Algebră},
Editura Paralela 45, Pitești, 2012.

\bibitem{exc11vol2}
Marius Perianu, Ioan Balică, Dan Zaharia,
\textit{Matematică de excelență pentru clasa a XI-a, Vol.\ II: Analiză},
Editura Paralela 45, Pitești, 2012.

\bibitem{exc12vol1}
Marius Perianu, Ioan Balică, Dan Zaharia,
\textit{Matematică de excelență pentru clasa a XII-a, Vol.\ I: Algebră},
Editura Paralela 45, Pitești, 2013.

\bibitem{exc12vol2}
Marius Perianu, Ioan Balică, Dan Zaharia,
\textit{Matematică de excelență pentru clasa a XII-a, Vol.\ II: Analiză},
Editura Paralela 45, Pitești, 2013.

\bibitem{titu}
Titu Andreescu, Dorin Andrica,
\textit{Algebră liniară, culegere de probleme},
GIL, Zalău, 2004.

\end{thebibliography}


\end{document}
